\documentclass[11pt,a4paper]{article}
% Typography and shared environments for all chapters.
\usepackage[T1]{fontenc}
\usepackage{iftex}
\ifPDFTeX\usepackage[utf8]{inputenc}\fi
\usepackage{lmodern}
\usepackage[margin=25mm,headheight=14pt]{geometry}
\usepackage{amsmath,amssymb,amsthm}
\usepackage{xcolor}
\usepackage{etoolbox}
\usepackage{enumitem}
\usepackage{microtype}
\usepackage{fancyhdr}
\usepackage{hyperref}
\definecolor{solutiontext}{RGB}{25,55,85}
\definecolor{linkcolor}{RGB}{25,75,115}
\hypersetup{colorlinks=true,linkcolor=linkcolor,urlcolor=linkcolor,
  pdftitle={Stochastic Methods: Selected Exercises},
  pdfauthor={Pietro Forni and Riccardo Visconti},bookmarksnumbered=true}
\setlength{\parindent}{0pt}
\setlength{\parskip}{0.35em}
\setlength{\emergencystretch}{2em}
\setlist{itemsep=0.2em,topsep=0.4em}
\allowdisplaybreaks[1]
\pagestyle{fancy}
\fancyhf{}
\fancyhead[L]{\small Stochastic Methods}
\fancyhead[R]{\small\iftoggle{showsolutions}{Exercises and Solutions}{Exercises}}
\fancyfoot[C]{\thepage}
\renewcommand{\headrulewidth}{0.3pt}
\DeclareMathOperator{\E}{\mathbb{E}}
\newtoggle{showsolutions}
\toggletrue{showsolutions}
% Set \togglefalse{showsolutions} above to build the practice edition.
% The optional \hideanswers override is also supported.
\ifdefined\hideanswers\togglefalse{showsolutions}\fi
\newtheoremstyle{exercisestyle}{1em}{0.5em}{\normalfont}{}%
  {\bfseries}{.}{0.5em}{}
\theoremstyle{exercisestyle}
\newtheorem{exercise}{Exercise}
% Preserve the supplied numbers; make PDF destinations unique across topics.
\providecommand{\theHexercise}{}
\renewcommand{\theHexercise}{\thesection.\arabic{exercise}}
\AtBeginEnvironment{exercise}{\setlist[enumerate,1]{label=(\alph*)}}
\newenvironment{solutionbody}{%
  \par\smallskip\begingroup\color{solutiontext}%
  \noindent\textbf{Solution.}\enspace\ignorespaces
}{\par\endgroup\medskip}
\newcommand{\solution}[1]{%
  \iftoggle{showsolutions}{\begin{solutionbody}#1\end{solutionbody}}{}%
}

\begin{document}
\hypersetup{pageanchor=false}
\begin{titlepage}
\centering
\vspace*{12mm}
{\Large\iftoggle{showsolutions}{Selected Exercises and Solutions}{Selected Exercises}\par}
\vspace{8mm}
{\Huge\bfseries Stochastic Methods\par}
\vspace{28mm}
{\Large Pietro Forni\par}
\vspace{3mm}
{\Large Riccardo Visconti\par}
\vfill
{\normalsize Independently written solutions to selected exercises from\par
\emph{Introduction to Probability Models}\par
by Sheldon M. Ross\par}
\vspace{10mm}
{\small Faculty of Informatics\par
Institute of Computational Science and Artificial Intelligence\par
Universit\`a della Svizzera italiana\par
Lugano, Switzerland\par}
\vspace{10mm}
{\large June 2026\par}
\end{titlepage}

\pagenumbering{roman}
\hypersetup{pageanchor=true}
\section*{About these notes}
\addcontentsline{toc}{section}{About these notes}
These notes collect selected exercises and independently written solutions
by Pietro Forni and Riccardo Visconti, prepared while studying stochastic
methods. They are based on \emph{Introduction to Probability Models}
by Sheldon M. Ross and are unofficial study notes. No university or
publisher endorsement is claimed.

Exercise numbers are retained from the original manuscript and restart
within each topic. Cite a topic title and exercise number together.
Some statements have been adapted or expanded with explicit assumptions
so that the exercises can be read without the textbook. In particular,
the two-ball urn and insurance premium models in Markov Chains
Exercises~9 and~53 are fully specified versions chosen for these notes;
they are not asserted to reproduce a particular textbook edition.

The collection has received a first mathematical review. Corrections and
further improvements are welcome.

\section*{Conventions}
Exponential and gamma distributions use rate parameters. Geometric
random variables count trials through the first success and take values
$1,2,\ldots$, unless explicitly defined otherwise. Rates are positive.
Conditioning formulas apply where the conditioning event has positive
probability or the marginal density is positive; conditional expectations
for continuous variables are defined almost everywhere. Variance
identities assume finite second moments. Queueing models use independent
service times and arrival processes unless stated otherwise.
\clearpage

\tableofcontents
\clearpage
\pagenumbering{arabic}
\section{Probability}
\label{sec:probability}

\setcounter{exercise}{0}

\begin{exercise}
\label{ex:probability-1}
    A box contains three marbles: one red, one green, and one blue. Consider an experiment that consists of taking one marble from the box then replacing it in the box and drawing a second marble from the box. What is the sample space? If, at all times, each marble in the box is equally likely to be selected, what is the probability of each point in the sample space?
\end{exercise}

\solution{
    Let the marbles be denoted by R (red), G (green), and B (blue). Since the first marble is replaced before drawing the second, the sample space $S$ consists of all possible ordered pairs of these three colors.
    \[
    S = \{(R,R), (R,G), (R,B), (G,R), (G,G), (G,B), (B,R), (B,G), (B,B)\}
    \]
    There are $3 \times 3 = 9$ possible outcomes in the sample space. Because each marble is equally likely to be selected at all times, every point in the sample space has the same probability:
    \[
    P(\text{each point}) = \frac{1}{9}
    \]
}

\begin{exercise}
\label{ex:probability-2}
    Repeat Exercise 1 when the second marble is drawn without replacing the first marble.
\end{exercise}

\solution{
    Without replacing the first marble, the same marble cannot be drawn twice. The sample space $S$ now consists of all ordered pairs of distinct colors:
    \[
    S = \{(R,G), (R,B), (G,R), (G,B), (B,R), (B,G)\}
    \]
    There are $3 \times 2 = 6$ possible outcomes. Since the selections are random and equally likely, the probability of each point in this new sample space is:
    \[
    P(\text{each point}) = \frac{1}{6}
    \]
}

\begin{exercise}
\label{ex:probability-3}
    A coin is to be tossed until a head appears twice in a row. What is the sample space for this experiment? If the coin is fair, what is the probability that it will be tossed exactly four times?
\end{exercise}

\solution{
    Let H denote a Head and T denote a Tail. The experiment ends as soon as two consecutive Heads (HH) appear. The sample space $S$ consists of all finite sequences of H and T that end in HH and do not contain HH anywhere prior.
    \[
    S = \{HH, THH, HTHH, TTHH, THTHH, HTTHH, TTTHH, \dots\}
    \]
    For the coin to be tossed exactly four times, the sequence must have a length of exactly 4 and end in HH. The only possible sequences that fit this criteria are HTHH and TTHH.
    \noindent
    Since the coin is fair, the probability of any specific sequence of length 4 is $\left(\frac{1}{2}\right)^4 = \frac{1}{16}$. Therefore, the probability that the coin is tossed exactly four times is the sum of the probabilities of these two specific outcomes:
    \[
    P(\text{exactly 4 tosses}) = P(HTHH) + P(TTHH) = \frac{1}{16} + \frac{1}{16} = \frac{2}{16} = \frac{1}{8}
    \]
}

\setcounter{exercise}{4}

\begin{exercise}
\label{ex:probability-5}
    An individual uses the following gambling system at Las Vegas. He bets \$1 that the roulette wheel will come up red. If he wins, he quits. If he loses then he makes the same bet a second time only this time he bets \$2; and then regardless of the outcome, quits. Assuming that he has a probability of $\frac{1}{2}$ of winning each bet, what is the probability that he goes home a winner? Why is this system not used by everyone?
\end{exercise}

\solution{
    Let W represent winning a bet and L represent losing a bet. The possible outcomes of this gambling system are:
    \begin{itemize}
        \item \textbf{W}: Wins the first \$1 bet. Net profit: $+\$1$. Probability: $\frac{1}{2}$
        \item \textbf{LW}: Loses the first \$1 bet, wins the second \$2 bet. Net profit: $-\$1 + \$2 = +\$1$. Probability: $\frac{1}{2} \times \frac{1}{2} = \frac{1}{4}$
        \item \textbf{LL}: Loses the first \$1 bet, loses the second \$2 bet. Net profit: $-\$1 - \$2 = -\$3$. Probability: $\frac{1}{2} \times \frac{1}{2} = \frac{1}{4}$
    \end{itemize}
    The individual goes home a winner in the first two scenarios. Thus, the probability of going home a winner is:
    \[
    P(\text{Winner}) = P(W) + P(LW) = \frac{1}{2} + \frac{1}{4} = \frac{3}{4}
    \]

    \noindent
    Although the probability of going home a winner is high ($75\%$), this system (a simplified Martingale strategy) is not used by everyone for two main reasons:
    \begin{enumerate}
        \item \textbf{Zero or Negative Expected Value:} Even with a high probability of winning, the expected value is zero with fair odds: $E[\text{Profit}] = (+1)\left(\frac{3}{4}\right) + (-3)\left(\frac{1}{4}\right) = 0$. In real Las Vegas roulette, the probability of red is actually less than $\frac{1}{2}$ (usually $\frac{18}{38}$ due to the 0 and 00 green slots), meaning the expected value of this system is mathematically negative. Over time, you will lose money.
        \item \textbf{Risk of Ruin:} To scale this system for a longer streak of losses, a player would need to keep doubling their bet exponentially. This requires unlimited capital and is ultimately prevented by table maximum bet limits set by casinos to stop this exact strategy.
    \end{enumerate}
}

\setcounter{exercise}{10}
\begin{exercise}
\label{ex:probability-11}
    If two fair dice are tossed, what is the probability that the sum is $i$, $i = 2, 3, \dots, 12$?
\end{exercise}

\solution{
    When rolling two fair dice, the sample space contains $6 \times 6 = 36$ equally likely ordered outcomes. Let the random variable $X$ represent the sum of the values on the two dice. We count the number of outcomes that yield each sum $i$:
    \begin{itemize}
        \item For $i = 2$: \{(1,1)\} $\rightarrow 1$ outcome $\rightarrow P(X=2) = \frac{1}{36}$
        \item For $i = 3$: \{(1,2), (2,1)\} $\rightarrow 2$ outcomes $\rightarrow P(X=3) = \frac{2}{36} = \frac{1}{18}$
        \item For $i = 4$: \{(1,3), (2,2), (3,1)\} $\rightarrow 3$ outcomes $\rightarrow P(X=4) = \frac{3}{36} = \frac{1}{12}$
        \item For $i = 5$: \{(1,4), (2,3), (3,2), (4,1)\} $\rightarrow 4$ outcomes $\rightarrow P(X=5) = \frac{4}{36} = \frac{1}{9}$
        \item For $i = 6$: \{(1,5), (2,4), (3,3), (4,2), (5,1)\} $\rightarrow 5$ outcomes $\rightarrow P(X=6) = \frac{5}{36}$
        \item For $i = 7$: \{(1,6), (2,5), (3,4), (4,3), (5,2), (6,1)\} $\rightarrow 6$ outcomes $\rightarrow P(X=7) = \frac{6}{36} = \frac{1}{6}$
        \item For $i = 8$: \{(2,6), (3,5), (4,4), (5,3), (6,2)\} $\rightarrow 5$ outcomes $\rightarrow P(X=8) = \frac{5}{36}$
        \item For $i = 9$: \{(3,6), (4,5), (5,4), (6,3)\} $\rightarrow 4$ outcomes $\rightarrow P(X=9) = \frac{4}{36} = \frac{1}{9}$
        \item For $i = 10$: \{(4,6), (5,5), (6,4)\} $\rightarrow 3$ outcomes $\rightarrow P(X=10) = \frac{3}{36} = \frac{1}{12}$
        \item For $i = 11$: \{(5,6), (6,5)\} $\rightarrow 2$ outcomes $\rightarrow P(X=11) = \frac{2}{36} = \frac{1}{18}$
        \item For $i = 12$: \{(6,6)\} $\rightarrow 1$ outcome $\rightarrow P(X=12) = \frac{1}{36}$
    \end{itemize}
    In general, the symmetric probability distribution can be compactly expressed via the formula:
    \[
    P(X = i) = \frac{6 - |7 - i|}{36}, \quad \text{for } i = 2, 3, \dots, 12
    \]
}

\setcounter{exercise}{12}

\begin{exercise}
\label{ex:probability-13}
    The dice game craps is played as follows. The player throws two dice, and if the sum is seven or eleven, then she wins. If the sum is two, three, or twelve, then she loses. If the sum is anything else, then she continues throwing until she either throws that number again (in which case she wins) or she throws a seven (in which case she loses). Calculate the probability that the player wins.
\end{exercise}

\solution{
    We can compute the total probability of winning by conditioning on the value of the initial roll. Let $P(i)$ represent the probability that the initial sum is $i$. From Exercise 11, we have:
    \begin{itemize}
        \item Instant Win: $P(7) + P(11) = \frac{6}{36} + \frac{2}{36} = \frac{8}{36}$
        \item Instant Loss: $P(2) + P(3) + P(12) = \frac{1}{36} + \frac{2}{36} + \frac{1}{36} = \frac{4}{36}$
    \end{itemize}
    If the initial roll results in a "point" $i \in \{4, 5, 6, 8, 9, 10\}$, all subsequent rolls are ignored unless they result in either $i$ (a win) or $7$ (a loss). Therefore, the conditional probability of winning given that point $i$ was established is:
    \[
    P(\text{Win} \mid \text{point } i) = \frac{P(i)}{P(i) + P(7)}
    \]
    By the Law of Total Probability, the overall probability of winning $P(W)$ is:
    \[
    P(W) = P(\text{instant win}) + \sum_{i \in \{4,5,6,8,9,10\}} P(i) \cdot \frac{P(i)}{P(i) + P(7)}
    \]
    Using symmetry ($P(4)=P(10)$, $P(5)=P(9)$, and $P(6)=P(8)$), we evaluate each distinct pair:
    \begin{itemize}
        \item For $i = 4, 10$: $2 \times \left( \frac{3}{36} \times \frac{3}{3+6} \right) = 2 \times \frac{3}{36} \times \frac{1}{3} = \frac{2}{36}$
        \item For $i = 5, 9$: $2 \times \left( \frac{4}{36} \times \frac{4}{4+6} \right) = 2 \times \frac{4}{36} \times \frac{2}{5} = \frac{16}{180} = \frac{4}{45}$
        \item For $i = 6, 8$: $2 \times \left( \frac{5}{36} \times \frac{5}{5+6} \right) = 2 \times \frac{25}{396} = \frac{50}{396} = \frac{25}{198}$
    \end{itemize}
    Summing all terms together:
    \[
    P(W) = \frac{8}{36} + \frac{2}{36} + \frac{4}{45} + \frac{25}{198} = \frac{10}{36} + \frac{4}{45} + \frac{25}{198} = \frac{5}{18} + \frac{4}{45} + \frac{25}{198}
    \]
    Finding a common denominator ($990$):
    \[
    P(W) = \frac{275}{990} + \frac{88}{990} + \frac{125}{990} = \frac{488}{990} = \frac{244}{495} \approx 0.4929
    \]
    Thus, the player has a $\frac{244}{495}$ (or approximately $49.29\%$) chance of winning.
}

\begin{exercise}
\label{ex:probability-14}
    The probability of winning on a single toss of the dice is $p$. $A$ starts, and if he fails, he passes the dice to $B$, who then attempts to win on her toss. They continue tossing the dice back and forth until one of them wins. What are their respective probabilities of winning?
\end{exercise}

\solution{
    Assume $0<p\le1$ and independent attempts. If $p=0$, neither player ever wins. Let $P(A)$ be the probability that $A$ wins the game and $P(B)$ be the probability that $B$ wins the game.

    We can set up a renewal relation by conditioning on the first toss:
    \begin{enumerate}
        \item With probability $p$, player $A$ wins on the very first toss.
        \item With probability $1-p$, player $A$ fails. At this exact moment, the game resets but roles are effectively reversed: it is now $B$'s turn to toss first, meaning $B$ now inherits the exact same probability of winning from this position that $A$ had at the start of the game ($P(A)$).
    \end{enumerate}
    Therefore, $A$'s remaining probability of winning if the first toss fails is $1 - P(A)$. This yields the equation:
    \[
    P(A) = p + (1-p)(1 - P(A))
    \]
    Expanding and solving for $P(A)$:
    \[
    P(A) = p + 1 - P(A) - p + pP(A) \implies P(A) = 1 - P(A) + pP(A)
    \]
    \[
    2P(A) - pP(A) = 1 \implies P(A)(2-p) = 1 \implies P(A) = \frac{1}{2-p}
    \]
    Since one of the two players must eventually win with probability 1, $B$'s probability of winning is the complement:
    \[
    P(B) = 1 - P(A) = 1 - \frac{1}{2-p} = \frac{2-p-1}{2-p} = \frac{1-p}{2-p}
    \]
    Thus, $P(A) = \frac{1}{2-p}$ and $P(B) = \frac{1-p}{2-p}$.
}

\setcounter{exercise}{16}

\begin{exercise}
\label{ex:probability-17}
    Suppose each of three persons tosses a coin. If the outcome of one of the tosses differs from the other outcomes, then the game ends. If not, then the persons start over and retoss their coins. Assuming fair coins, what is the probability that the game will end with the first round of tosses? If all three coins are biased and have probability $\frac{1}{4}$ of landing heads, what is the probability that the game will end at the first round?
\end{exercise}

\solution{
    The game continues to a new round if and only if all three coins yield identical outcomes (either all Heads or all Tails). Let $C$ be the event that the game continues. The event that the game ends on the first round is the complement, $C^c$.

    \noindent
    \textbf{Case 1: Fair Coins ($p = \frac{1}{2}$ for Heads, $q = \frac{1}{2}$ for Tails)}

    \noindent
    The total number of outcomes for 3 fair coins is $2^3 = 8$ equally likely configurations. The game continues only on the outcomes $\{HHH, TTT\}$. Thus:
    \[
    P(C) = P(HHH) + P(TTT) = \frac{1}{8} + \frac{1}{8} = \frac{2}{8} = \frac{1}{4}
    \]
    The probability that the game ends on the first round is:
    \[
    P(C^c) = 1 - P(C) = 1 - \frac{1}{4} = \frac{3}{4}
    \]

    \noindent
    \textbf{Case 2: Biased Coins ($p = \frac{1}{4}$ for Heads, $q = \frac{3}{4}$ for Tails)}

    \noindent
    Since the tosses are independent, we calculate the joint probabilities for the continuing configurations directly:
    \[
    P(HHH) = p^3 = \left(\frac{1}{4}\right)^3 = \frac{1}{64}
    \]
    \[
    P(TTT) = q^3 = \left(\frac{3}{4}\right)^3 = \frac{27}{64}
    \]
    Therefore, the probability that the game continues is:
    \[
    P(C) = \frac{1}{64} + \frac{27}{64} = \frac{28}{64} = \frac{7}{16}
    \]
    The probability that the game ends on the first round under this bias is:
    \[
    P(C^c) = 1 - P(C) = 1 - \frac{7}{16} = \frac{9}{16}
    \]
}

\begin{exercise}
\label{ex:probability-18}
    Assume that each child who is born is equally likely to be a boy or a girl. If a family has two children, what is the probability that both are girls given that (a) the eldest is a girl, (b) at least one is a girl?
\end{exercise}

\solution{
    Let the sample space for a family with two children be represented by ordered pairs based on age order: $S = \{(B,B), (B,G), (G,B), (G,G)\}$. Since boys and girls are equally likely, each sample point has an equal probability of $\frac{1}{4}$. Let $A$ be the event that both children are girls: $A = \{(G,G)\}$.

    \noindent
    \textbf{(a) Given that the eldest is a girl}
    Let $E$ be the condition that the eldest child is a girl: $E = \{(G,B), (G,G)\}$.
    \noindent
    The intersection of our target event and the condition is $A \cap E = \{(G,G)\}$. Using conditional probability:
    \[
    P(A \mid E) = \frac{P(A \cap E)}{P(E)} = \frac{1/4}{2/4} = \frac{1}{2}
    \]

    \noindent
    \textbf{(b) Given that at least one is a girl}
    Let $F$ be the condition that at least one child is a girl: $F = \{(B,G), (G,B), (G,G)\}$.
    \noindent
    The intersection is $A \cap F = \{(G,G)\}$. Applying the conditional probability definition:
    \[
    P(A \mid F) = \frac{P(A \cap F)}{P(F)} = \frac{1/4}{3/4} = \frac{1}{3}
    \]
}

\setcounter{exercise}{18}

\begin{exercise}
\label{ex:probability-19}
    Two dice are rolled. What is the probability that at least one is a six? If the two faces are different, what is the probability that at least one is a six?
\end{exercise}

\solution{
    Let $X_1, X_2$ denote the outcomes of the two fair dice. Total outcomes = $6 \times 6 = 36$.

    \noindent
    \textbf{Part 1: Probability that at least one is a six}
    The complement event is that neither die lands on a six. There are $5 \times 5 = 25$ such outcomes.
    \[
    P(\text{at least one 6}) = 1 - P(\text{no 6}) = 1 - \frac{25}{36} = \frac{11}{36}
    \]

    \noindent
    \textbf{Part 2: Probability that at least one is a six given faces are different}
    Let $D$ be the event that the two faces are different. Out of 36 outcomes, 6 are duplicates ($(1,1), \dots, (6,6)$), so $|D| = 36 - 6 = 30$.
    The outcomes in $D$ that include at least one six are: $(1,6), (2,6), (3,6), (4,6), (5,6)$ and their reversals $(6,1), (6,2), (6,3), (6,4), (6,5)$. This yields exactly 10 outcomes.
    \[
    P(\text{at least one 6} \mid D) = \frac{10}{30} = \frac{1}{3}
    \]
}

\begin{exercise}
\label{ex:probability-20}
    Three dice are thrown. What is the probability the same number appears on exactly two of the three dice?
\end{exercise}

\solution{
    When throwing three dice, the total size of the sample space is $6^3 = 216$ equally likely outcomes.
    To find the number of favorable combinations where exactly two dice show the matching number:
    \begin{enumerate}
        \item Choose which value is repeated: $\binom{6}{1} = 6$ choices.
        \item Choose which value appears on the remaining distinct die: $\binom{5}{1} = 5$ choices.
        \item Select which two of the three positions host the repeated number: $\binom{3}{2} = 3$ choices.
    \end{enumerate}
    Multiplying these independent choices gives the total number of favorable outcomes:
    \[
    \text{Favorable Outcomes} = 6 \times 5 \times 3 = 90
    \]
    Thus, the calculated probability is:
    \[
    P(\text{exactly two identical}) = \frac{90}{216} = \frac{5}{12}
    \]
}

\begin{exercise}
\label{ex:probability-21}
    Suppose that 5 percent of men and 0.25 percent of women are color-blind. A color-blind person is chosen at random. What is the probability of this person being male? Assume that there are an equal number of males and females.
\end{exercise}

\solution{
    Let $M$ and $F$ represent the events of selecting a male and a female, respectively. We are given $P(M) = 0.5$ and $P(F) = 0.5$. Let $CB$ represent the event that a person is color-blind. The conditional probabilities are $P(CB \mid M) = 0.05$ and $P(CB \mid F) = 0.0025$.
    We calculate the posterior probability $P(M \mid CB)$ using Bayes' Theorem:
    \[
    P(M \mid CB) = \frac{P(CB \mid M)P(M)}{P(CB \mid M)P(M) + P(CB \mid F)P(F)}
    \]
    Substitute the known quantities into the equation:
    \[
    P(M \mid CB) = \frac{0.05 \times 0.5}{(0.05 \times 0.5) + (0.0025 \times 0.5)} = \frac{0.025}{0.025 + 0.00125} = \frac{0.025}{0.02625}
    \]
    Simplifying the fraction reveals:
    \[
    P(M \mid CB) = \frac{2500}{2625} = \frac{20}{21} \approx 0.9524
    \]
    Thus, there is a $\frac{20}{21}$ (or approximately $95.24\%$) chance that the randomly selected color-blind individual is male.
}

\begin{exercise}
\label{ex:probability-22}
    $A$ and $B$ play until one has 2 more points than the other. Assuming that each point is independently won by $A$ with probability $p$, what is the probability they will play a total of $2n$ points? What is the probability that $A$ will win?
\end{exercise}

\solution{
    Because a player must win by 2 points, the match can only end on an even number of total points.

    \noindent
    \textbf{Part 1: Probability that they will play a total of $2n$ points}
    For a match to continue exactly to $2n$ points, the score must remain tied after pairs of rounds $2, 4, \dots, 2n-2$. A tie occurs within any distinct pair of points if and only if the points are split (either $A$ wins then $B$ wins, or $B$ wins then $A$ wins).
    \[
    P(\text{Split pair}) = p(1-p) + (1-p)p = 2p(1-p)
    \]
    For the match to terminate exactly at point $2n$, the game must progress as a tie for the initial $n-1$ pairs of rounds, and then someone must win both points in the $n$-th pair:
    \[
    P(\text{End at pair } n) = 1 - 2p(1-p) = p^2 + (1-p)^2
    \]
    Thus, the probability that they play exactly $2n$ total points is:
    \[
    P(\text{Total } 2n \text{ points}) = [2p(1-p)]^{n-1} [p^2 + (1-p)^2], \quad n = 1, 2, \dots
    \]

    \noindent
    \textbf{Part 2: Probability that $A$ will win}
    Let $P_A$ be the probability that $A$ wins the complete match. We condition on the results of the first two points played:
    \begin{itemize}
        \item $A$ wins both points (probability $p^2$): $A$ wins the match immediately.
        \item $B$ wins both points (probability $(1-p)^2$): $B$ wins the match ($A$ loses).
        \item The points are split (probability $2p(1-p)$): The score returns to a balanced state, renewing the identical starting match probability $P_A$ for player $A$.
    \end{itemize}
    This establishes the following total probability equation:
    \[
    P_A = p^2 \cdot 1 + (1-p)^2 \cdot 0 + 2p(1-p) \cdot P_A
    \]
    Isolating terms to solve for $P_A$:
    \[
    P_A [1 - 2p(1-p)] = p^2 \implies P_A = \frac{p^2}{1 - 2p(1-p)} = \frac{p^2}{p^2 + (1-p)^2}
    \]
}

\setcounter{exercise}{24}

\begin{exercise}
\label{ex:probability-25}
    Two cards are randomly selected from a deck of 52 playing cards.
    \begin{enumerate}
        \item[(a)] What is the probability they constitute a pair (that is, that they are of the same denomination)?
        \item[(b)] What is the conditional probability they constitute a pair given that they are of different suits?
    \end{enumerate}
\end{exercise}

\solution{

    \noindent
    \textbf{(a)}

    \noindent
    \textbf{Method 1 (Combinatorics):} Let $P$ be the event of drawing a pair. The total number of ways to draw 2 cards from 52 is $\binom{52}{2}$. A pair consists of 2 cards of the same denomination. There are 13 denominations, and for each denomination, there are $\binom{4}{2}$ ways to select 2 suits out of 4.
    \[P(P) = \frac{13 \times \binom{4}{2}}{\binom{52}{2}} = \frac{13 \times 6}{1326} = \frac{78}{1326} = \frac{1}{17}\]

    \noindent
    \textbf{Method 2 (Sequential):} The first card drawn can be any card (probability 1). To form a pair, the second card must match the denomination of the first. There are exactly 3 cards remaining of that denomination out of the 51 remaining cards in the deck.
    \[P(P) = \frac{3}{51} = \frac{1}{17}\]

    \noindent
    \textbf{(b)}

    \noindent
    Let $D$ be the event of drawing two cards of different suits. We want to find $P(P \mid D)$. By definition of conditional probability:
    \[P(P \mid D) = \frac{P(P \cap D)}{P(D)}\]
    Since a pair must necessarily be composed of different suits (as there is only one card of each suit per denomination), $P(P \cap D) = P(P) = \frac{1}{17}$.
    To find $P(D)$, consider drawing sequentially: the first card can be anything. The second card must be a different suit. There are 12 remaining cards of the same suit as the first card, leaving $51 - 12 = 39$ cards of different suits.
    \[P(D) = \frac{39}{51} = \frac{13}{17}\]
    Therefore, the conditional probability is:
    \[P(P \mid D) = \frac{1/17}{13/17} = \frac{1}{13}\]
}

\begin{exercise}
\label{ex:probability-26}
    A deck of 52 playing cards, containing all 4 aces, is randomly divided into 4 piles of 13 cards each. Define events $E_1, E_2, E_3$, and $E_4$ as follows:
    \begin{align*}
        E_1 &= \{\text{the first pile has exactly 1 ace}\}, \\
        E_2 &= \{\text{the second pile has exactly 1 ace}\}, \\
        E_3 &= \{\text{the third pile has exactly 1 ace}\}, \\
        E_4 &= \{\text{the fourth pile has exactly 1 ace}\}
    \end{align*}
    Find $P(E_1 E_2 E_3 E_4)$, the probability that each pile has an ace.
\end{exercise}

\solution{
    The total number of ways to divide 52 distinct cards into 4 distinct piles of 13 cards each is given by the multinomial coefficient: $\frac{52!}{(13!)^4}$.

    To find the number of favorable ways where each pile receives exactly 1 ace, we can break the process into two steps:
    \begin{enumerate}
        \item Distribute the 4 aces among the 4 piles. Since the piles are distinct (first, second, third, fourth), there are $4!$ ways to do this.
        \item Distribute the remaining 48 non-ace cards into the 4 piles so that each pile gets exactly 12 more cards. There are $\frac{48!}{(12!)^4}$ ways to do this.
    \end{enumerate}

    The total number of favorable outcomes is $4! \times \frac{48!}{(12!)^4}$.

    The probability that each pile has exactly one ace is the ratio of favorable outcomes to total outcomes:
    \[P(E_1 E_2 E_3 E_4) = \frac{4! \cdot \frac{48!}{(12!)^4}}{\frac{52!}{(13!)^4}} = \frac{4! \cdot 48! \cdot (13!)^4}{52! \cdot (12!)^4}\]
    Since $\frac{13!}{12!} = 13$, this simplifies to:
    \[P(E_1 E_2 E_3 E_4) = \frac{24 \cdot 13^4}{52 \cdot 51 \cdot 50 \cdot 49} = \frac{24 \cdot 28561}{6497400} = \frac{2197}{20825} \approx 0.10549\]
}

\begin{exercise}
\label{ex:probability-27}
    Suppose in Exercise 26 we had defined the events $E_i, i = 1, 2, 3, 4$, by
    \begin{align*}
        E_1 &= \{\text{one of the piles contains the ace of spades}\}, \\
        E_2 &= \{\text{the ace of spades and the ace of hearts are in different piles}\}, \\
        E_3 &= \{\text{the ace of spades, the ace of hearts, and the ace of diamonds are in different piles}\}, \\
        E_4 &= \{\text{all 4 aces are in different piles}\}
    \end{align*}
    Use the general multiplication rule for probabilities to find $P(E_1 E_2 E_3 E_4)$, the probability that each pile has an ace. Compare your answer with the one you obtained in Exercise 26.
\end{exercise}

\solution{
    We use the general multiplication rule for probability:
    \[P(E_1 E_2 E_3 E_4) = P(E_1) P(E_2 \mid E_1) P(E_3 \mid E_1 E_2) P(E_4 \mid E_1 E_2 E_3)\]

    \begin{enumerate}
        \item $P(E_1) = 1$: The ace of spades is guaranteed to end up in exactly one of the piles.
        \item $P(E_2 \mid E_1)$: The ace of spades occupies 1 of the 52 spots in the deck, leaving 51 open spots. To be in a different pile, the ace of hearts must be placed in one of the other 3 piles. Those 3 piles collectively contain $3 \times 13 = 39$ open spots.
        \[P(E_2 \mid E_1) = \frac{39}{51}\]
        \item $P(E_3 \mid E_1 E_2)$: Given the aces of spades and hearts are in different piles, they occupy 2 spots across 2 different piles, leaving 50 open spots. The ace of diamonds must go into one of the 2 remaining empty piles, which collectively have $2 \times 13 = 26$ open spots.
        \[P(E_3 \mid E_1 E_2) = \frac{26}{50}\]
        \item $P(E_4 \mid E_1 E_2 E_3)$: Given three aces are in three different piles, there are 49 open spots remaining. The ace of clubs must go into the single remaining pile that does not yet have an ace. That pile has exactly 13 spots available.
        \[P(E_4 \mid E_1 E_2 E_3) = \frac{13}{49}\]
    \end{enumerate}

    Multiplying these probabilities together yields:
    \[P(E_1 E_2 E_3 E_4) = 1 \cdot \frac{39}{51} \cdot \frac{26}{50} \cdot \frac{13}{49} = \frac{13 \cdot 13 \cdot 13}{17 \cdot 25 \cdot 49} = \frac{2197}{20825} \approx 0.10549\]

    \noindent
    \textbf{Comparison:}
    The final probability $\frac{2197}{20825}$ is identical to the one obtained in Exercise 26. However, utilizing the multiplication rule with this alternative sequence of events makes the calculation vastly simpler and more intuitive to compute by hand compared to using raw multinomial coefficients.
}

\setcounter{exercise}{29}

\begin{exercise}
\label{ex:probability-30}
    Bill and George go target shooting together. Both shoot at a target at the same time. Suppose Bill hits the target with probability $0.7$, whereas George, independently, hits the target with probability $0.4$.
    \begin{enumerate}
        \item[(a)] Given that exactly one shot hit the target, what is the probability that it was George's shot?
        \item[(b)] Given that the target is hit, what is the probability that George hit it?
    \end{enumerate}
\end{exercise}

\solution{
    Let $B$ be the event that Bill hits the target, with $P(B) = 0.7$, and let $G$ be the event that George hits the target, with $P(G) = 0.4$. Since their shots are independent, the probabilities of their complements are $P(B^c) = 0.3$ and $P(G^c) = 0.6$.

    \noindent
    \textbf{(a)}

    \noindent
    Let $E$ be the event that exactly one shot hits the target. This can happen in two disjoint ways: Bill hits and George misses, or Bill misses and George hits.
    \[
    P(E) = P(B \cap G^c) + P(B^c \cap G) = P(B)P(G^c) + P(B^c)P(G)
    \]
    \[
    P(E) = (0.7)(0.6) + (0.3)(0.4) = 0.42 + 0.12 = 0.54
    \]
    We want to find the conditional probability $P(B^c \cap G \mid E)$:
    \[
    P(B^c \cap G \mid E) = \frac{P(B^c \cap G)}{P(E)} = \frac{0.12}{0.54} = \frac{12}{54} = \frac{2}{9}
    \]

    \noindent
    \textbf{(b)}

    \noindent
    Let $H$ be the event that the target is hit by at least one of them.
    \[
    P(H) = 1 - P(B^c \cap G^c) = 1 - P(B^c)P(G^c) = 1 - (0.3)(0.6) = 1 - 0.18 = 0.82
    \]
    We want to find the conditional probability that George hit the target given that it was hit, which is $P(G \mid H)$:
    \[
    P(G \mid H) = \frac{P(G \cap H)}{P(H)} = \frac{P(G)}{P(H)} = \frac{0.4}{0.82} = \frac{40}{82} = \frac{20}{41}
    \]
}

\begin{exercise}
\label{ex:probability-31}
    What is the conditional probability that the first die is six given that the sum of the dice is seven?
\end{exercise}

\solution{
    Let $D_1$ and $D_2$ represent the outcomes of the first and second fair six-sided dice, respectively. The total number of outcomes when rolling two dice is $6 \times 6 = 36$.

    Let $S_7$ be the event that the sum of the two dice is $7$. The favorable outcomes for $S_7$ are:
    \[
    S_7 = \{(1,6), (2,5), (3,4), (4,3), (5,2), (6,1)\}
    \]
    Thus, there are 6 possible outcomes, so $P(S_7) = \frac{6}{36} = \frac{1}{6}$.

    Let $F_6$ be the event that the first die is a $6$. The intersection event $F_6 \cap S_7$ contains exactly one outcome: $\{(6,1)\}$. Thus, $P(F_6 \cap S_7) = \frac{1}{36}$.

    Using the definition of conditional probability:
    \[
    P(F_6 \mid S_7) = \frac{P(F_6 \cap S_7)}{P(S_7)} = \frac{1/36}{6/36} = \frac{1}{6}
    \]
}

\begin{exercise}
\label{ex:probability-32}
    Suppose all $n$ men at a party throw their hats in the center of the room. Each man then randomly selects a hat. Show that the probability that none of the $n$ men selects his own hat is
    \[
    \frac{1}{2!} - \frac{1}{3!} + \frac{1}{4!} - \cdots + \frac{(-1)^n}{n!}
    \]
    Note that as $n \to \infty$ this converges to $e^{-1}$. Is this surprising?
\end{exercise}

\solution{
    This is the classic matching problem, representing the probability of a derangement. Let $E_i$ be the event that the $i$-th man selects his own hat for $i = 1, 2, \dots, n$. We want to find the probability that no man selects his own hat, which is $P(\bigcap_{i=1}^n E_i^c) = 1 - P(\bigcup_{i=1}^n E_i)$.

    By the Inclusion-Exclusion Principle, the probability of the union is:
    \[
    P\left(\bigcup_{i=1}^n E_i\right) = \sum_{i=1}^n P(E_i) - \sum_{1 \le i < j \le n} P(E_i \cap E_j) + \cdots + (-1)^{n-1} P(E_1 \cap E_2 \cap \dots \cap E_n)
    \]
    Since the hats are selected completely at random, the probability that any specific set of $k$ men all select their own hats is:
    \[
    P(E_{i_1} \cap E_{i_2} \cap \dots \cap E_{i_k}) = \frac{(n-k)!}{n!}
    \]
    The number of ways to choose $k$ specific men out of $n$ is $\binom{n}{k}$. Therefore, the $k$-th term in the inclusion-exclusion sum is:
    \[
    \binom{n}{k} \frac{(n-k)!}{n!} = \frac{n!}{k!(n-k)!} \cdot \frac{(n-k)!}{n!} = \frac{1}{k!}
    \]
    Substituting this back into the union formula gives:
    \[
    P\left(\bigcup_{i=1}^n E_i\right) = \sum_{k=1}^n (-1)^{k-1} \frac{1}{k!} = 1 - \frac{1}{2!} + \frac{1}{3!} - \cdots + (-1)^{n-1}\frac{1}{n!}
    \]
    Thus, the probability that no one gets their own hat is:
    \[
    P\left(\bigcap_{i=1}^n E_i^c\right) = 1 - P\left(\bigcup_{i=1}^n E_i\right) = \frac{1}{2!} - \frac{1}{3!} + \frac{1}{4!} - \cdots + \frac{(-1)^n}{n!}
    \]
    As $n \to \infty$, this expression matches the Taylor series expansion for $e^{-1}$:
    \[
    e^{-1} = \sum_{k=0}^{\infty} \frac{(-1)^k}{k!} = 1 - 1 + \frac{1}{2!} - \frac{1}{3!} + \frac{1}{4!} - \cdots
    \]
    \textbf{Is this surprising?} At first glance, it can be intuitive to think that as the number of people $n$ grows larger, the probability of a complete derangement should either drop to $0$ or climb to $1$. However, it rapidly stabilizes to approximately $0.36788$ for any $n \ge 7$, meaning the size of the party has virtually no effect on the probability once $n$ is reasonably large.
}

\begin{exercise}
\label{ex:probability-33}
    In a class there are four freshman boys, six freshman girls, and six sophomore boys. How many sophomore girls must be present if sex and class are to be independent when a student is selected at random?
\end{exercise}

\solution{
    Let $x$ be the number of sophomore girls in the class. We can organize the students into a contingency table:

    \begin{center}
    \begin{tabular}{r|cc|c}
                   & Boys & Girls & Total \\ \hline
        Freshmen   & 4    & 6     & 10    \\
        Sophomores & 6    & $x$   & $6+x$ \\ \hline
        Total      & 10   & $6+x$ & $16+x$
    \end{tabular}
    \end{center}

    For the events ``Sex'' and ``Class'' to be independent when a student is selected at random, the joint probability of selecting a student who is both a Freshman and a Boy must equal the product of their individual probabilities:
    \[
    P(\text{Freshman} \cap \text{Boy}) = P(\text{Freshman}) \times P(\text{Boy})
    \]
    From the table, we have:
    \[
    P(\text{Freshman} \cap \text{Boy}) = \frac{4}{16+x}, \quad P(\text{Freshman}) = \frac{10}{16+x}, \quad P(\text{Boy}) = \frac{10}{16+x}
    \]
    Setting up the independence equation:
    \[
    \frac{4}{16+x} = \left(\frac{10}{16+x}\right) \left(\frac{10}{16+x}\right)
    \]
    Since $16+x \neq 0$, we can multiply both sides by $(16+x)^2$ to simplify:
    \[
    4(16+x) = 100 \implies 64 + 4x = 100 \implies 4x = 36 \implies x = 9
    \]
    Therefore, there must be $9$ sophomore girls present.
}

\begin{exercise}
\label{ex:probability-34}
    Mr. Jones has devised a gambling system for winning at roulette. When he bets, he bets on red, and places a bet only when the ten previous spins of the roulette have landed on a black number. He reasons that his chance of winning is quite large since the probability of eleven consecutive spins resulting in black is quite small. What do you think of this system?
\end{exercise}

\solution{
    Mr. Jones's system is based on a classic misconception known as the \textbf{Gambler's Fallacy}.

    Each spin of a roulette wheel is an independent trial, meaning the outcome of any given spin has no memory of or connection to previous outcomes. Let $R$ be the event of landing on red on the 11th spin, and $B_1, B_2, \dots, B_{10}$ be the events that the first 10 spins land on black. By the definition of independent events, the conditional probability is simply the unconditional probability of hitting red on any single spin:
    \[
    P(R \mid B_1 \cap B_2 \cap \dots \cap B_{10}) = P(R)
    \]
    On a standard American roulette wheel (containing 18 red, 18 black, and 2 green slots), $P(R) = \frac{18}{38} \approx 0.4737$.

    While it is true that the a priori probability of observing 11 consecutive black spins is very small ($\left(\frac{18}{38}\right)^{11} \approx 0.0005$), once the 10 black spins have already occurred, they belong to the past and their historical probability is $1$. The wheel does not "owe" a red outcome to restore balance. Thus, Mr. Jones's system does absolutely nothing to improve his chances of winning.
}

\begin{exercise}
\label{ex:probability-35}
    A fair coin is continually flipped. What is the probability that the first four flips are
    \begin{enumerate}
        \item[(a)] $H, H, H, H$?
        \item[(b)] $T, H, H, H$?
        \item[(c)] What is the probability that the pattern $T, H, H, H$ occurs before the pattern $H, H, H, H$?
    \end{enumerate}
\end{exercise}

\solution{
    Since the coin is fair, each flip lands on Heads ($H$) or Tails ($T$) independently with probability $\frac{1}{2}$.

    \noindent
    \textbf{(a)}

    \noindent
    The probability of any specific sequence of 4 independent flips is:
    \[
    P(HHHH) = \left(\frac{1}{2}\right)^4 = \frac{1}{16}
    \]

    \noindent
    \textbf{(b)}

    \noindent
    Similarly, the probability of this specific ordered sequence is:
    \[
    P(THHH) = \left(\frac{1}{2}\right)^4 = \frac{1}{16}
    \]

    \noindent
    \textbf{(c)}

    \noindent
    Let's analyze how the sequence $HHHH$ can appear before $THHH$.

    The pattern $HHHH$ can only win if it occurs as the absolute first four flips of the experiment. If the first four flips are anything else, $THHH$ will inevitably appear first. To see why, suppose $HHHH$ appears for the first time at some later point in the sequence, say starting at flip index $k$ (where $k > 1$), so that flips $k, k+1, k+2, k+3$ are all $H$. Since this is the \textit{first} time $HHHH$ appears, the flip immediately preceding it (flip $k-1$) must have been a $T$ (if it were an $H$, then a sequence of four heads would have already occurred starting at index $k-1$). Consequently, the four consecutive flips from index $k-1$ to $k+2$ form the pattern $THHH$.

    Thus, $THHH$ will always show up before $HHHH$ unless the game starts directly with $HHHH$. Therefore:
    \[
    P(THHH \text{ before } HHHH) = 1 - P(\text{first 4 flips are } HHHH) = 1 - \frac{1}{16} = \frac{15}{16}
    \]
}

\setcounter{exercise}{35}

\begin{exercise}
\label{ex:probability-36}
    Consider two boxes, one containing one black and one white marble, the other, two black and one white marble. A box is selected at random and a marble is drawn at random from the selected box. What is the probability that the marble is black?
\end{exercise}

\solution{
    Let $B_1$ be the event that the first box is selected, and $B_2$ be the event that the second box is selected. Since a box is selected at random, we have:
    \[
    P(B_1) = \frac{1}{2}, \quad P(B_2) = \frac{1}{2}
    \]
    Let $B$ be the event that a black marble is drawn. From the contents of the boxes, the conditional probabilities are:
    \[
    P(B \mid B_1) = \frac{1}{2}, \quad P(B \mid B_2) = \frac{2}{3}
    \]
    Using the Law of Total Probability, we find the total probability of drawing a black marble:
    \[
    P(B) = P(B \mid B_1)P(B_1) + P(B \mid B_2)P(B_2)
    \]
    \[
    P(B) = \left(\frac{1}{2}\right)\left(\frac{1}{2}\right) + \left(\frac{2}{3}\right)\left(\frac{1}{2}\right) = \frac{1}{4} + \frac{1}{3} = \frac{7}{12}
    \]
}

\begin{exercise}
\label{ex:probability-37}
    In Exercise 36, what is the probability that the first box was the one selected given that the marble is white?
\end{exercise}

\solution{
    Let $W$ be the event that a white marble is drawn. We can find $P(W)$ either as the complement of drawing a black marble, or directly using the Law of Total Probability:
    \[
    P(W) = 1 - P(B) = 1 - \frac{7}{12} = \frac{5}{12}
    \]
    We want to find the conditional probability $P(B_1 \mid W)$. Using Bayes' Theorem:
    \[
    P(B_1 \mid W) = \frac{P(W \mid B_1)P(B_1)}{P(W)}
    \]
    Since the first box contains 1 white marble out of 2, $P(W \mid B_1) = \frac{1}{2}$. Substituting the values in:
    \[
    P(B_1 \mid W) = \frac{\left(\frac{1}{2}\right)\left(\frac{1}{2}\right)}{\frac{5}{12}} = \frac{\frac{1}{4}}{\frac{5}{12}} = \frac{1}{4} \times \frac{12}{5} = \frac{3}{5}
    \]
}

\begin{exercise}
\label{ex:probability-38}
    Urn 1 contains two white balls and one black ball, while urn 2 contains one white ball and five black balls. One ball is drawn at random from urn 1 and placed in urn 2. A ball is then drawn from urn 2. It happens to be white. What is the probability that the transferred ball was white?
\end{exercise}

\solution{
    Let $T_W$ be the event that a white ball is transferred from Urn 1 to Urn 2, and $T_B$ be the event that a black ball is transferred. Based on the initial composition of Urn 1 ($2W, 1B$), we have:
    \[
    P(T_W) = \frac{2}{3}, \quad P(T_B) = \frac{1}{3}
    \]
    Let $W_2$ be the event that a white ball is drawn from Urn 2. We analyze the composition of Urn 2 after the transfer:
    \begin{itemize}
        \item If a white ball was transferred ($T_W$), Urn 2 now contains 2 white balls and 5 black balls (7 total). Thus, $P(W_2 \mid T_W) = \frac{2}{7}$.
        \item If a black ball was transferred ($T_B$), Urn 2 now contains 1 white ball and 6 black balls (7 total). Thus, $P(W_2 \mid T_B) = \frac{1}{7}$.
    \end{itemize}
    We want to find $P(T_W \mid W_2)$. Applying Bayes' Theorem:
    \[
    P(T_W \mid W_2) = \frac{P(W_2 \mid T_W)P(T_W)}{P(W_2 \mid T_W)P(T_W) + P(W_2 \mid T_B)P(T_B)}
    \]
    \[
    P(T_W \mid W_2) = \frac{\left(\frac{2}{7}\right)\left(\frac{2}{3}\right)}{\left(\frac{2}{7}\right)\left(\frac{2}{3}\right) + \left(\frac{1}{7}\right)\left(\frac{1}{3}\right)} = \frac{\frac{4}{21}}{\frac{4}{21} + \frac{1}{21}} = \frac{\frac{4}{21}}{\frac{5}{21}} = \frac{4}{5}
    \]
}

\begin{exercise}
\label{ex:probability-39}
    Stores $A$, $B$, and $C$ have 50, 75, and 100 employees, and, respectively, 50, 60, and 70 percent of these are women. Resignations are equally likely among all employees, regardless of sex. One employee resigns and this is a woman. What is the probability that she works in store C?
\end{exercise}

\solution{
    Let $A, B$, and $C$ denote the events that the resigning employee belongs to Store A, Store B, or Store C, respectively. Since resignations are equally likely among all employees, the prior probabilities are proportional to the size of each store's workforce. The total number of employees is $50 + 75 + 100 = 225$. Thus:
    \[
    P(A) = \frac{50}{225}, \quad P(B) = \frac{75}{225}, \quad P(C) = \frac{100}{225}
    \]
    Let $W$ be the event that the resigning employee is a woman. The given percentages provide the conditional probabilities:
    \[
    P(W \mid A) = 0.50, \quad P(W \mid B) = 0.60, \quad P(W \mid C) = 0.70
    \]
    We want to calculate the posterior probability $P(C \mid W)$ using Bayes' Theorem:
    \[
    P(C \mid W) = \frac{P(W \mid C)P(C)}{P(W \mid A)P(A) + P(W \mid B)P(B) + P(W \mid C)P(C)}
    \]
    Substituting the known values:
    \[
    P(C \mid W) = \frac{(0.70)\left(\frac{100}{225}\right)}{(0.50)\left(\frac{50}{225}\right) + (0.60)\left(\frac{75}{225}\right) + (0.70)\left(\frac{100}{225}\right)}
    \]
    We can cancel out the common denominator of 225 across the expression:
    \[
    P(C \mid W) = \frac{0.70 \times 100}{(0.50 \times 50) + (0.60 \times 75) + (0.70 \times 100)} = \frac{70}{25 + 45 + 70} = \frac{70}{140} = \frac{1}{2}
    \]
    Therefore, the probability that she works in store C is $\frac{1}{2}$ (or $0.5$).
}

\setcounter{exercise}{39}

\begin{exercise}
\label{ex:probability-40}
    \begin{enumerate}
        \item[(a)] A gambler has in his pocket a fair coin and a two-headed coin. He selects one of the coins at random, and when he flips it, it shows heads. What is the probability that it is the fair coin?
        \item[(b)] Suppose that he flips the same coin a second time and again it shows heads. Now what is the probability that it is the fair coin?
        \item[(c)] Suppose that he flips the same coin a third time and it shows tails. Now what is the probability that it is the fair coin?
    \end{enumerate}
\end{exercise}

\solution{
    Let $F$ be the event that the fair coin is chosen, and let $2H$ be the event that the two-headed coin is chosen. Since the choice is random, the prior probabilities are:
    \[
    P(F) = \frac{1}{2}, \quad P(2H) = \frac{1}{2}
    \]

    \noindent
    \textbf{(a)}

    \noindent
    Let $H_1$ be the event that the first flip results in heads. The conditional probabilities are $P(H_1 \mid F) = \frac{1}{2}$ and $P(H_1 \mid 2H) = 1$. Using Bayes' Theorem to find $P(F \mid H_1)$:
    \[
    P(F \mid H_1) = \frac{P(H_1 \mid F)P(F)}{P(H_1 \mid F)P(F) + P(H_1 \mid 2H)P(2H)} = \frac{\left(\frac{1}{2}\right)\left(\frac{1}{2}\right)}{\left(\frac{1}{2}\right)\left(\frac{1}{2}\right) + (1)\left(\frac{1}{2}\right)} = \frac{\frac{1}{4}}{\frac{1}{4} + \frac{1}{2}} = \frac{1}{3}
    \]

    \noindent
    \textbf{(b)}

    \noindent
    Let $H_1 H_2$ be the event that both the first and second flips result in heads. The conditional probabilities for independent tosses of the chosen coin are $P(H_1 H_2 \mid F) = \left(\frac{1}{2}\right)^2 = \frac{1}{4}$ and $P(H_1 H_2 \mid 2H) = 1^2 = 1$. Applying Bayes' Theorem:
    \[
    P(F \mid H_1 H_2) = \frac{P(H_1 H_2 \mid F)P(F)}{P(H_1 H_2 \mid F)P(F) + P(H_1 H_2 \mid 2H)P(2H)} = \frac{\left(\frac{1}{4}\right)\left(\frac{1}{2}\right)}{\left(\frac{1}{4}\right)\left(\frac{1}{2}\right) + (1)\left(\frac{1}{2}\right)} = \frac{\frac{1}{8}}{\frac{1}{8} + \frac{1}{2}} = \frac{1}{5}
    \]

    \noindent
    \textbf{(c)}

    \noindent
    A two-headed coin can never land on tails. Since a tail is observed on the third flip, the probability that the coin is the two-headed coin drops to $0$, meaning it must be the fair coin with absolute certainty:
    \[
    P(F \mid \text{third flip is tails}) = 1
    \]
}

\begin{exercise}
\label{ex:probability-41}
    In a certain species of rats, black dominates over brown. Suppose that a black rat with two black parents has a brown sibling.
    \begin{enumerate}
        \item[(a)] What is the probability that this rat is a pure black rat (as opposed to being a hybrid with one black and one brown gene)?
        \item[(b)] Suppose that when the black rat is mated with a brown rat, all five of their offspring are black. Now, what is the probability that the rat is a pure black rat?
    \end{enumerate}
\end{exercise}

\solution{
    Let $B$ denote the dominant black gene and $b$ denote the recessive brown gene.
    Since the rat has a brown sibling, that sibling must have the genotype $bb$. This implies that both parents must carry the recessive gene, meaning both parents are hybrids with genotype $Bb$.

    A genetic cross between two $Bb$ parents yields offspring genotypes with the following probabilities:
    \[
    P(BB) = \frac{1}{4}, \quad P(Bb) = \frac{1}{2}, \quad P(bb) = \frac{1}{4}
    \]

    \noindent
    \textbf{(a)}

    \noindent
    We are given that the rat in question is black, which rules out the $bb$ genotype. We want to find the conditional probability that it is pure black ($BB$) given that it is black ($BB \cup Bb$):
    \[
    P(BB \mid \text{Black}) = \frac{P(BB)}{P(BB) + P(Bb)} = \frac{\frac{1}{4}}{\frac{1}{4} + \frac{1}{2}} = \frac{1}{3}
    \]
    Thus, the prior probabilities for our target rat are $P(BB) = \frac{1}{3}$ and $P(Bb) = \frac{2}{3}$.

    \noindent
    \textbf{(b)}

    \noindent
    The black rat is mated with a brown rat ($bb$). Let $E$ be the event that all 5 offspring are black.
    \begin{itemize}
        \item If the rat is pure black ($BB$), a $BB \times bb$ cross will always produce $Bb$ (black) offspring. Thus, $P(E \mid BB) = 1$.
        \item If the rat is a hybrid ($Bb$), a $Bb \times bb$ cross gives each offspring a $\frac{1}{2}$ chance of being black ($Bb$). For 5 independent offspring, $P(E \mid Bb) = \left(\frac{1}{2}\right)^5 = \frac{1}{32}$.
    \end{itemize}
    Using Bayes' Theorem to find the updated probability $P(BB \mid E)$:
    \[
    P(BB \mid E) = \frac{P(E \mid BB)P(BB)}{P(E \mid BB)P(BB) + P(E \mid Bb)P(Bb)} = \frac{(1)\left(\frac{1}{3}\right)}{(1)\left(\frac{1}{3}\right) + \left(\frac{1}{32}\right)\left(\frac{2}{3}\right)}
    \]
    \[
    P(BB \mid E) = \frac{\frac{1}{3}}{\frac{1}{3} + \frac{1}{48}} = \frac{\frac{1}{3}}{\frac{16 + 1}{48}} = \frac{1}{3} \times \frac{48}{17} = \frac{16}{17} \approx 0.9412
    \]
}

\begin{exercise}
\label{ex:probability-42}
    There are three coins in a box. One is a two-headed coin, another is a fair coin, and the third is a biased coin that comes up heads 75 percent of the time. When one of the three coins is selected at random and flipped, it shows heads. What is the probability that it was the two-headed coin?
\end{exercise}

\solution{
    Let $C_1, C_2,$ and $C_3$ represent the selection of the two-headed coin, the fair coin, and the biased coin, respectively. Since a coin is selected at random:
    \[
    P(C_1) = P(C_2) = P(C_3) = \frac{1}{3}
    \]
    Let $H$ be the event that the flipped coin shows heads. The conditional probabilities are:
    \[
    P(H \mid C_1) = 1, \quad P(H \mid C_2) = 0.5 = \frac{1}{2}, \quad P(H \mid C_3) = 0.75 = \frac{3}{4}
    \]
    We use Bayes' Theorem to find $P(C_1 \mid H)$:
    \[
    P(C_1 \mid H) = \frac{P(H \mid C_1)P(C_1)}{P(H \mid C_1)P(C_1) + P(H \mid C_2)P(C_2) + P(H \mid C_3)P(C_3)}
    \]
    \[
    P(C_1 \mid H) = \frac{(1)\left(\frac{1}{3}\right)}{(1)\left(\frac{1}{3}\right) + \left(\frac{1}{2}\right)\left(\frac{1}{3}\right) + \left(\frac{3}{4}\right)\left(\frac{1}{3}\right)}
    \]
    Canceling out the common factor of $\frac{1}{3}$ from the numerator and denominator:
    \[
    P(C_1 \mid H) = \frac{1}{1 + \frac{1}{2} + \frac{3}{4}} = \frac{1}{\frac{4 + 2 + 3}{4}} = \frac{1}{\frac{9}{4}} = \frac{4}{9}
    \]
}

\begin{exercise}
\label{ex:probability-43}
    Suppose we have ten coins which are such that if the $i$th one is flipped then heads will appear with probability $i/10, i = 1, 2, \dots, 10$. When one of the coins is randomly selected and flipped, it shows heads. What is the conditional probability that it was the fifth coin?
\end{exercise}

\solution{
    Let $C_i$ be the event that the $i$-th coin is selected, where $i \in \{1, 2, \dots, 10\}$. Since the coin choice is completely random:
    \[
    P(C_i) = \frac{1}{10} \quad \text{for all } i
    \]
    Let $H$ be the event that the flipped coin lands on heads. The conditional probability for the $i$-th coin is given as:
    \[
    P(H \mid C_i) = \frac{i}{10}
    \]
    We want to find $P(C_5 \mid H)$. First, compute the total probability of landing on heads, $P(H)$, using the Law of Total Probability:
    \[
    P(H) = \sum_{i=1}^{10} P(H \mid C_i)P(C_i) = \sum_{i=1}^{10} \left(\frac{i}{10}\right)\left(\frac{1}{10}\right) = \frac{1}{100} \sum_{i=1}^{10} i
    \]
    Using the arithmetic series sum formula $\sum_{i=1}^{n} i = \frac{n(n+1)}{2}$:
    \[
    \sum_{i=1}^{10} i = \frac{10 \times 11}{2} = 55 \implies P(H) = \frac{55}{100}
    \]
    Now, apply Bayes' Theorem for the fifth coin ($C_5$):
    \[
    P(C_5 \mid H) = \frac{P(H \mid C_5)P(C_5)}{P(H)} = \frac{\left(\frac{5}{10}\right)\left(\frac{1}{10}\right)}{\frac{55}{100}} = \frac{\frac{5}{100}}{\frac{55}{100}} = \frac{5}{55} = \frac{1}{11}
    \]
}

\setcounter{exercise}{43}
\begin{exercise}
\label{ex:probability-44}
    Urn 1 has five white and seven black balls. Urn 2 has three white and twelve black balls. We flip a fair coin. If the outcome is heads, then a ball from urn 1 is selected, while if the outcome is tails, then a ball from urn 2 is selected. Suppose that a white ball is selected. What is the probability that the coin landed tails?
\end{exercise}

\solution{
    Let $H$ be the event that the coin lands heads (Urn 1 is chosen), and $T$ be the event that the coin lands tails (Urn 2 is chosen). Since the coin is fair, we have:
    \[
    P(H) = \frac{1}{2}, \quad P(T) = \frac{1}{2}
    \]
    Let $W$ be the event that a white ball is selected. The conditional probabilities of selecting a white ball given the chosen urn are:
    \[
    P(W|H) = \frac{5}{5 + 7} = \frac{5}{12}
    \]
    \[
    P(W|T) = \frac{3}{3 + 12} = \frac{3}{15} = \frac{1}{5}
    \]
    We want to find the probability that the coin landed tails given that a white ball was selected, $P(T|W)$. By Bayes' Theorem:
    \[
    P(T|W) = \frac{P(W|T)P(T)}{P(W|T)P(T) + P(W|H)P(H)}
    \]
    Substituting the known values:
    \[
    P(T|W) = \frac{\left(\frac{1}{5}\right)\left(\frac{1}{2}\right)}{\left(\frac{1}{5}\right)\left(\frac{1}{2}\right) + \left(\frac{5}{12}\right)\left(\frac{1}{2}\right)} = \frac{\frac{1}{10}}{\frac{1}{10} + \frac{5}{24}}
    \]
    Multiply the numerator and the denominator by 120 (the least common multiple of 10 and 24):
    \[
    P(T|W) = \frac{12}{12 + 25} = \frac{12}{37}
    \]
}

\setcounter{exercise}{44}
\begin{exercise}
\label{ex:probability-45}
    An urn contains $b$ black balls and $r$ red balls. One of the balls is drawn at random, but when it is put back in the urn $c$ additional balls of the same color are put in with it. Now suppose that we draw another ball. Show that the probability that the first ball drawn was black given that the second ball drawn was red is $b/(b + r + c)$.
\end{exercise}

\solution{
    Let $B_1$ be the event that the first ball drawn is black, and $R_1$ be the event that the first ball drawn is red. Their initial probabilities are:
    \[
    P(B_1) = \frac{b}{b+r}, \quad P(R_1) = \frac{r}{b+r}
    \]
    Let $R_2$ be the event that the second ball drawn is red. We determine the conditional probabilities for the second draw based on the first draw's outcome:
    \begin{itemize}
        \item If $B_1$ occurs, the urn contains $b+c$ black balls and $r$ red balls. Thus, $P(R_2|B_1) = \frac{r}{b+r+c}$.
        \item If $R_1$ occurs, the urn contains $b$ black balls and $r+c$ red balls. Thus, $P(R_2|R_1) = \frac{r+c}{b+r+c}$.
    \end{itemize}
    We need to find the probability that the first ball was black given the second was red, $P(B_1|R_2)$. Using Bayes' Theorem:
    \[
    P(B_1|R_2) = \frac{P(R_2|B_1)P(B_1)}{P(R_2|B_1)P(B_1) + P(R_2|R_1)P(R_1)}
    \]
    Substitute the probabilities:
    \[
    P(B_1|R_2) = \frac{\left(\frac{r}{b+r+c}\right)\left(\frac{b}{b+r}\right)}{\left(\frac{r}{b+r+c}\right)\left(\frac{b}{b+r}\right) + \left(\frac{r+c}{b+r+c}\right)\left(\frac{r}{b+r}\right)}
    \]
    Since the denominator $(b+r+c)(b+r)$ is common to all terms in the main fraction, we can cancel it out:
    \[
    P(B_1|R_2) = \frac{rb}{rb + (r+c)r} = \frac{rb}{rb + r^2 + cr}
    \]
    Factor out $r$ from the denominator:
    \[
    P(B_1|R_2) = \frac{rb}{r(b + r + c)} = \frac{b}{b+r+c}
    \]
    This completes the proof.
}

\setcounter{exercise}{45}
\begin{exercise}
\label{ex:probability-46}
    Three prisoners are informed by their jailer that one of them has been chosen at random to be executed, and the other two are to be freed. Prisoner $A$ asks the jailer to tell him privately which of his fellow prisoners will be set free, claiming that there would be no harm in divulging this information, since he already knows that at least one will go free. The jailer refuses to answer this question, pointing out that if $A$ knew which of his fellows were to be set free, then his own probability of being executed would rise from $\frac{1}{3}$ to $\frac{1}{2}$, since he would then be one of two prisoners. What do you think of the jailer's reasoning?
\end{exercise}

\solution{
    The jailer's conclusion does not follow without specifying how he chooses which freed prisoner to name. Under the usual symmetric reporting rule, A's probability remains $1/3$. We can prove the jailer wrong by calculating the conditional probability of A being executed using Bayes' Theorem.

    \noindent
    Let $E_A, E_B,$ and $E_C$ be the events that prisoner A, B, or C is executed, respectively. Since the choice is random, the prior probabilities are:
    \[
    P(E_A) = P(E_B) = P(E_C) = \frac{1}{3}
    \]
    Let $J_B$ be the event that the jailer tells A that prisoner B will go free. We evaluate the probability of the jailer naming B under each condition:
    \begin{itemize}
        \item If A is to be executed ($E_A$), both B and C go free. The jailer can randomly choose between B and C. Assuming he picks with equal probability: $P(J_B|E_A) = \frac{1}{2}$.
        \item If B is to be executed ($E_B$), the jailer must name C to go free (he cannot name A to A, nor can he name the condemned). Thus, $P(J_B|E_B) = 0$.
        \item If C is to be executed ($E_C$), the jailer must name B to go free. Thus, $P(J_B|E_C) = 1$.
    \end{itemize}

    \noindent
    We want to find $P(E_A|J_B)$, the probability that A is executed given the jailer names B. Using Bayes' Theorem:
    \[
    P(E_A|J_B) = \frac{P(J_B|E_A)P(E_A)}{P(J_B|E_A)P(E_A) + P(J_B|E_B)P(E_B) + P(J_B|E_C)P(E_C)}
    \]
    \[
    P(E_A|J_B) = \frac{\left(\frac{1}{2}\right)\left(\frac{1}{3}\right)}{\left(\frac{1}{2}\right)\left(\frac{1}{3}\right) + (0)\left(\frac{1}{3}\right) + (1)\left(\frac{1}{3}\right)} = \frac{\frac{1}{6}}{\frac{1}{6} + 0 + \frac{1}{3}} = \frac{\frac{1}{6}}{\frac{3}{6}} = \frac{1}{3}
    \]

    \noindent
    Therefore, even after learning that B goes free, prisoner A's probability of being executed remains exactly $\frac{1}{3}$, not $\frac{1}{2}$. More generally, let $a=P(J_B\mid E_A)$. Then Bayes' rule gives $P(E_A\mid J_B)=a/(a+1)$. The value $1/3$ uses $a=1/2$; if the jailer always names B when he can, the posterior after naming B is $1/2$. Thus his reporting protocol matters.
}

\setcounter{exercise}{47}
\begin{exercise}
\label{ex:probability-48}
    Sixty percent of the families in a certain community own their own car, thirty percent own their own home, and twenty percent own both their own car and their own home. If a family is randomly chosen, what is the probability that this family owns a car or a house but not both?
\end{exercise}

\solution{
    Let $C$ be the event that a randomly chosen family owns a car, and $H$ be the event that they own a home. We are given the following probabilities:
    \[
    P(C) = 0.60
    \]
    \[
    P(H) = 0.30
    \]
    \[
    P(C \cap H) = 0.20
    \]
    We want to find the probability that a family owns a car or a house, but not both. This is mathematically represented by the symmetric difference:
    \[
    P(C \text{ only}) + P(H \text{ only}) = P(C \cap H^c) + P(C^c \cap H)
    \]
    We can calculate these disjoint probabilities by subtracting the intersection from the individual marginal probabilities:
    \[
    P(C \cap H^c) = P(C) - P(C \cap H) = 0.60 - 0.20 = 0.40
    \]
    \[
    P(C^c \cap H) = P(H) - P(C \cap H) = 0.30 - 0.20 = 0.10
    \]
    Adding these together gives the final probability:
    \[
    P(\text{car or house, not both}) = 0.40 + 0.10 = 0.50
    \]

    \noindent
    \textit{Alternative method:} You can also find the probability of the union and subtract the intersection:
    \[
    P(C \cup H) - P(C \cap H) = (P(C) + P(H) - P(C \cap H)) - P(C \cap H)
    \]
    \[
    = (0.60 + 0.30 - 0.20) - 0.20 = 0.70 - 0.20 = 0.50
    \]
}

\section{Random Variables}
\label{sec:random-variables}

\setcounter{exercise}{0}

\begin{exercise}
\label{ex:random-variables-1}
An urn contains five red, three orange, and two blue balls. Two balls are randomly selected. What is the sample space of this experiment? Let $X$ represent the number of orange balls selected. What are the possible values of $X$? Calculate $P\{X=0\}$.
\end{exercise}

\solution{
The experiment consists of selecting two balls from an urn containing red, orange, and blue balls. If we describe outcomes only by color, the sample space is
\[
S=\{RR,RO,RB,OO,OB,BB\}.
\]
Since $X$ represents the number of orange balls selected, the possible values are
\[
X\in\{0,1,2\}.
\]
There are $10$ total balls, of which $7$ are not orange. Thus,
\[
P\{X=0\}
=
\frac{\binom{7}{2}}{\binom{10}{2}}
=
\frac{21}{45}
=
\frac{7}{15}.
\]
}

\begin{exercise}
\label{ex:random-variables-2}
Let $X$ represent the difference between the number of heads and the number of tails obtained when a coin is tossed $n$ times. What are the possible values of $X$?
\end{exercise}

\solution{
Let $H$ be the number of heads and $T$ be the number of tails. Since the coin is tossed $n$ times,
\[
H+T=n.
\]
The random variable is
\[
X=H-T.
\]
Since $T=n-H$, we get
\[
X=H-(n-H)=2H-n.
\]
As $H$ can take the values $0,1,2,\ldots,n$, the possible values of $X$ are
\[
X\in\{-n,-n+2,-n+4,\ldots,n-2,n\}.
\]
}

\begin{exercise}
\label{ex:random-variables-3}
In Exercise 2, if the coin is assumed fair, then, for $n=2$, what are the probabilities associated with the values that $X$ can take on?
\end{exercise}

\solution{
For $n=2$, the sample space is
\[
S=\{HH,HT,TH,TT\}.
\]
Each outcome has probability $\frac14$.

The possible values of $X=H-T$ are
\[
X\in\{-2,0,2\}.
\]
We compute:
\[
P\{X=2\}=P(HH)=\frac14,
\]
\[
P\{X=0\}=P(HT\text{ or }TH)=\frac24=\frac12,
\]
\[
P\{X=-2\}=P(TT)=\frac14.
\]
Thus,
\[
P_X(x)=
\begin{cases}
\frac14, & x=-2,\\
\frac12, & x=0,\\
\frac14, & x=2,\\
0, & \text{otherwise}.
\end{cases}
\]
}

\begin{exercise}
\label{ex:random-variables-4}
Suppose a die is rolled twice. What are the possible values that the following random variables can take on?
\begin{enumerate}
\item The maximum value to appear in the two rolls.
\item The minimum value to appear in the two rolls.
\item The sum of the two rolls.
\item The value of the first roll minus the value of the second roll.
\end{enumerate}
\end{exercise}

\solution{
Let the two rolls be denoted by $(i,j)$, where
\[
i,j\in\{1,2,3,4,5,6\}.
\]

\noindent
\textbf{(a) Maximum value}

The maximum value can be any number from $1$ to $6$, so
\[
\{1,2,3,4,5,6\}.
\]

\noindent
\textbf{(b) Minimum value}

The minimum value can also be any number from $1$ to $6$, so
\[
\{1,2,3,4,5,6\}.
\]

\noindent
\textbf{(c) Sum of the two rolls}

The smallest possible sum is $1+1=2$, and the largest possible sum is $6+6=12$. Hence the possible values are
\[
\{2,3,4,5,6,7,8,9,10,11,12\}.
\]

\noindent
\textbf{(d) First roll minus second roll}

The smallest possible difference is $1-6=-5$, and the largest possible difference is $6-1=5$. Hence the possible values are
\[
\{-5,-4,-3,-2,-1,0,1,2,3,4,5\}.
\]
}

\begin{exercise}
\label{ex:random-variables-5}
If the die in Exercise 4 is assumed fair, calculate the probabilities associated with the random variables in (i)--(iv).
\end{exercise}

\solution{
Since the die is fair and rolled twice, the sample space has $36$ equally likely outcomes.

\noindent
\textbf{(i) Maximum value}

Let $M$ be the maximum of the two rolls. For $m=1,2,\ldots,6$,
\[
P\{M=m\}
=
P\{M\le m\}-P\{M\le m-1\}
=
\left(\frac{m}{6}\right)^2-\left(\frac{m-1}{6}\right)^2.
\]
Thus,
\[
P\{M=m\}=\frac{2m-1}{36}, \qquad m=1,2,\ldots,6.
\]

\noindent
\textbf{(ii) Minimum value}

Let $N$ be the minimum of the two rolls. For $n=1,2,\ldots,6$,
\[
P\{N=n\}
=
P\{N\ge n\}-P\{N\ge n+1\}
=
\left(\frac{7-n}{6}\right)^2-\left(\frac{6-n}{6}\right)^2.
\]
Thus,
\[
P\{N=n\}=\frac{13-2n}{36}, \qquad n=1,2,\ldots,6.
\]

\noindent
\textbf{(iii) Sum of the two rolls}

Let $S$ be the sum. Then
\[
P\{S=s\}
=
\begin{cases}
\dfrac{s-1}{36}, & s=2,3,4,5,6,7,\\[6pt]
\dfrac{13-s}{36}, & s=8,9,10,11,12,\\[6pt]
0, & \text{otherwise}.
\end{cases}
\]

\noindent
\textbf{(iv) First roll minus second roll}

Let $D$ be the first roll minus the second roll. Then $D$ can take values
\[
-5,-4,-3,-2,-1,0,1,2,3,4,5.
\]
For $d=-5,-4,\ldots,5$,
\[
P\{D=d\}=\frac{6-|d|}{36}.
\]
}

\begin{exercise}
\label{ex:random-variables-6}
Suppose five fair coins are tossed. Let $E$ be the event that all coins land heads. Define the random variable $I_E$ by
\[
I_E=
\begin{cases}
1, & \text{if } E \text{ occurs},\\
0, & \text{if } E^c \text{ occurs}.
\end{cases}
\]
For what outcomes in the original sample space does $I_E$ equal $1$? What is $P\{I_E=1\}$?
\end{exercise}

\solution{
The original sample space consists of all sequences of five coin tosses:
\[
S=\{H,T\}^5.
\]
The random variable $I_E$ equals $1$ exactly when event $E$ occurs, that is, when all five coins land heads.

Thus,
\[
I_E=1
\]
only for the outcome
\[
HHHHH.
\]
Since the coins are fair,
\[
P\{I_E=1\}=P(E)=P(HHHHH)=\left(\frac12\right)^5=\frac{1}{32}.
\]
}

\begin{exercise}
\label{ex:random-variables-7}
Suppose a coin having probability $0.7$ of coming up heads is tossed three times. Let $X$ denote the number of heads that appear in the three tosses. Determine the probability mass function of $X$.
\end{exercise}

\solution{
Since the coin is tossed three times independently and the probability of heads is $0.7$, the random variable $X$ has a binomial distribution with parameters
\[
n=3, \qquad p=0.7.
\]
Therefore,
\[
P\{X=k\}
=
\binom{3}{k}(0.7)^k(0.3)^{3-k},
\qquad k=0,1,2,3.
\]
Calculating each value:
\[
P\{X=0\}=\binom{3}{0}(0.7)^0(0.3)^3=0.027,
\]
\[
P\{X=1\}=\binom{3}{1}(0.7)^1(0.3)^2=0.189,
\]
\[
P\{X=2\}=\binom{3}{2}(0.7)^2(0.3)^1=0.441,
\]
\[
P\{X=3\}=\binom{3}{3}(0.7)^3(0.3)^0=0.343.
\]
Thus, the probability mass function is
\[
P_X(k)=
\begin{cases}
\binom{3}{k}(0.7)^k(0.3)^{3-k}, & k=0,1,2,3,\\
0, & \text{otherwise}.
\end{cases}
\]
}

\begin{exercise}
\label{ex:random-variables-8}
Suppose the distribution function of $X$ is given by
\[
F(b)=
\begin{cases}
0, & b<0,\\
\dfrac12, & 0\le b<1,\\
1, & b\ge 1.
\end{cases}
\]
What is the probability mass function of $X$?
\end{exercise}

\solution{
For a discrete random variable, the probability mass at a point equals the jump of the distribution function at that point:
\[
P(X=x)=F(x)-\lim_{t\to x^-}F(t).
\]

The distribution function has jumps at $x=0$ and $x=1$.

At $x=0$:
\[
P(X=0)=F(0)-F(0^-)
=\frac12-0
=\frac12.
\]

At $x=1$:
\[
P(X=1)=F(1)-F(1^-)
=1-\frac12
=\frac12.
\]

Hence the probability mass function is
\[
p_X(x)=
\begin{cases}
\dfrac12, & x=0,\\[6pt]
\dfrac12, & x=1,\\[6pt]
0, & \text{otherwise}.
\end{cases}
\]
}

\begin{exercise}
\label{ex:random-variables-9}
If the distribution function of $X$ is given by
\[
F(b)=
\begin{cases}
0, & b<0,\\[4pt]
\dfrac12, & 0\le b<1,\\[4pt]
\dfrac35, & 1\le b<2,\\[4pt]
\dfrac45, & 2\le b<3,\\[4pt]
\dfrac9{10}, & 3\le b<3.5,\\[4pt]
1, & b\ge 3.5,
\end{cases}
\]
calculate the probability mass function of $X$.
\end{exercise}

\solution{
The probability masses are obtained from the jumps of the distribution function.

\[
P(X=0)=\frac12-0=\frac12.
\]

\[
P(X=1)=\frac35-\frac12=\frac1{10}.
\]

\[
P(X=2)=\frac45-\frac35=\frac15.
\]

\[
P(X=3)=\frac9{10}-\frac45=\frac1{10}.
\]

\[
P(X=3.5)=1-\frac9{10}=\frac1{10}.
\]

Checking:
\[
\frac12+\frac1{10}+\frac15+\frac1{10}+\frac1{10}
=1.
\]

Therefore,
\[
p_X(x)=
\begin{cases}
\dfrac12, & x=0,\\[6pt]
\dfrac1{10}, & x=1,\\[6pt]
\dfrac15, & x=2,\\[6pt]
\dfrac1{10}, & x=3,\\[6pt]
\dfrac1{10}, & x=3.5,\\[6pt]
0, & \text{otherwise}.
\end{cases}
\]
}

\begin{exercise}
\label{ex:random-variables-10}
Suppose three fair dice are rolled. What is the probability that at most one six appears?
\end{exercise}

\solution{
Let $X$ denote the number of sixes obtained in the three rolls.

Since each die is fair,
\[
P(\text{six})=\frac16,
\qquad
P(\text{not six})=\frac56.
\]

We want
\[
P(X\le 1)=P(X=0)+P(X=1).
\]

First,
\[
P(X=0)=\left(\frac56\right)^3=\frac{125}{216}.
\]

Next,
\[
P(X=1)
=
\binom31\left(\frac16\right)\left(\frac56\right)^2
=
3\cdot\frac16\cdot\frac{25}{36}
=
\frac{75}{216}.
\]

Therefore,
\[
P(X\le1)
=
\frac{125}{216}+\frac{75}{216}
=
\frac{200}{216}
=
\frac{25}{27}.
\]
}

\begin{exercise}
\label{ex:random-variables-11}
A ball is drawn from an urn containing three white and three black balls. After the ball is drawn, it is replaced and another ball is drawn. This goes on indefinitely. What is the probability that of the first four balls drawn, exactly two are white?
\end{exercise}

\solution{
Since the ball is replaced after each draw, the draws are independent.

Let $X$ denote the number of white balls among the first four draws.

Since
\[
P(\text{white})=\frac36=\frac12,
\]
we have
\[
X\sim \operatorname{Binomial}(4,\tfrac12).
\]

Therefore,
\[
P(X=2)
=
\binom42
\left(\frac12\right)^2
\left(\frac12\right)^2
=
\binom42\left(\frac12\right)^4.
\]

Since
\[
\binom42=6,
\]
we obtain
\[
P(X=2)
=
6\cdot\frac1{16}
=
\frac38.
\]
}

\begin{exercise}
\label{ex:random-variables-12}
On a multiple-choice exam with three possible answers for each of the five questions, what is the probability that a student would get four or more correct answers just by guessing?
\end{exercise}

\solution{
Let $X$ be the number of correct answers.

Each question is answered correctly with probability
\[
p=\frac13.
\]

Thus
\[
X\sim \operatorname{Binomial}(5,\tfrac13).
\]

We seek
\[
P(X\ge4)=P(X=4)+P(X=5).
\]

Compute:
\[
P(X=4)
=
\binom54
\left(\frac13\right)^4
\left(\frac23\right)
=
5\cdot\frac1{81}\cdot\frac23
=
\frac{10}{243}.
\]

\[
P(X=5)
=
\binom55
\left(\frac13\right)^5
=
\frac1{243}.
\]

Hence
\[
P(X\ge4)
=
\frac{10}{243}+\frac1{243}
=
\frac{11}{243}.
\]
}

\begin{exercise}
\label{ex:random-variables-13}
An individual claims to have extrasensory perception (ESP). As a test, a fair coin is flipped ten times, and he is asked to predict in advance the outcome. Our individual gets seven out of ten correct. What is the probability he would have done at least this well if he had no ESP?
\end{exercise}

\solution{
If the individual has no ESP, each prediction is simply a random guess.

Let $X$ denote the number of correct predictions among the ten flips.

Since each prediction is correct with probability
\[
p=\frac12,
\]
we have
\[
X\sim \operatorname{Binomial}(10,\tfrac12).
\]

The phrase ``at least this well'' means obtaining seven or more correct predictions. Therefore the relevant probability is
\[
P(X\ge7),
\]
not merely $P(X=7)$, because outcomes with $8$, $9$, or $10$ correct predictions would provide at least as much evidence in favor of ESP as getting exactly $7$ correct.

Thus,
\[
P(X\ge7)
=
P(X=7)+P(X=8)+P(X=9)+P(X=10).
\]

Using the binomial formula,
\[
P(X\ge7)
=
\sum_{k=7}^{10}
\binom{10}{k}
\left(\frac12\right)^{10}.
\]

Hence
\[
P(X\ge7)
=
\frac{\binom{10}{7}+\binom{10}{8}+\binom{10}{9}+\binom{10}{10}}
{2^{10}}.
\]

Substituting the binomial coefficients,
\[
P(X\ge7)
=
\frac{120+45+10+1}{1024}
=
\frac{176}{1024}
=
\frac{11}{64}.
\]

Therefore,
\[
{P(X\ge7)=\frac{11}{64}\approx0.1719.}
\]
}

\begin{exercise}
\label{ex:random-variables-14}
Suppose $X$ has a binomial distribution with parameters $6$ and $\frac12$. Show that $X=3$ is the most likely outcome.
\end{exercise}

\solution{
Since $X\sim \operatorname{Binomial}(6,\frac12)$,
\[
P(X=k)=\binom{6}{k}\left(\frac12\right)^6,
\qquad k=0,1,\ldots,6.
\]
Thus the probabilities are proportional to $\binom{6}{k}$.

The binomial coefficients are
\[
\binom60=1,\quad
\binom61=6,\quad
\binom62=15,\quad
\binom63=20,\quad
\binom64=15,\quad
\binom65=6,\quad
\binom66=1.
\]
The largest coefficient is
\[
\binom63=20.
\]
Therefore,
\[
P(X=3)
\]
is the largest probability, so $X=3$ is the most likely outcome.
}

\begin{exercise}
\label{ex:random-variables-15}
Let $X$ be binomially distributed with parameters $n$ and $p$. Show that as $k$ goes from $0$ to $n$, $P(X=k)$ increases monotonically, then decreases monotonically, reaching its largest value
\begin{enumerate}
\item when $(n+1)p$ is an integer, at $k=(n+1)p-1$ and $k=(n+1)p$,
\item when $(n+1)p$ is not an integer, at the unique integer $k$ satisfying
\[
(n+1)p-1<k<(n+1)p.
\]
\end{enumerate}
\end{exercise}

\solution{
Assume $0<p<1$. At $p=0$ or $p=1$ the variable is constant with its unique mode at $0$ or $n$, respectively. For $X\sim \operatorname{Binomial}(n,p)$,
\[
P(X=k)=\binom{n}{k}p^k(1-p)^{n-k}.
\]
Consider the ratio
\[
\frac{P(X=k)}{P(X=k-1)}
=
\frac{\binom{n}{k}p^k(1-p)^{n-k}}
{\binom{n}{k-1}p^{k-1}(1-p)^{n-k+1}}.
\]
Simplifying,
\[
\frac{P(X=k)}{P(X=k-1)}
=
\frac{n-k+1}{k}\cdot \frac{p}{1-p}.
\]
Therefore,
\[
\frac{P(X=k)}{P(X=k-1)}>1
\]
if and only if
\[
\frac{n-k+1}{k}\cdot \frac{p}{1-p}>1.
\]
This is equivalent to
\[
(n-k+1)p>k(1-p).
\]
Expanding,
\[
np-kp+p>k-kp.
\]
Hence
\[
(n+1)p>k.
\]
So
\[
P(X=k)>P(X=k-1)
\]
if and only if
\[
k<(n+1)p.
\]

Similarly,
\[
P(X=k)<P(X=k-1)
\]
if and only if
\[
k>(n+1)p.
\]

Thus the probabilities increase while $k<(n+1)p$ and decrease after $k>(n+1)p$.

If $(n+1)p$ is an integer, then at $k=(n+1)p$,
\[
\frac{P(X=k)}{P(X=k-1)}=1.
\]
Thus the two adjacent values
\[
k=(n+1)p-1
\qquad\text{and}\qquad
k=(n+1)p
\]
have equal largest probability.

If $(n+1)p$ is not an integer, then there is a unique integer $k$ satisfying
\[
(n+1)p-1<k<(n+1)p.
\]
This value of $k$ gives the unique largest probability.
}

\begin{exercise}
\label{ex:random-variables-16}
An airline knows that $5$ percent of the people making reservations on a certain flight will not show up. Consequently, their policy is to sell $52$ tickets for a flight that can hold only $50$ passengers. What is the probability that there will be a seat available for every passenger who shows up?
\end{exercise}

\solution{
Let $X$ be the number of passengers who show up. Since each ticketed passenger shows up with probability
\[
0.95,
\]
we have
\[
X\sim \operatorname{Binomial}(52,0.95).
\]

There will be a seat available for every passenger if and only if at most $50$ passengers show up. Therefore, we want
\[
P(X\le 50).
\]
It is easier to compute the complement:
\[
P(X\le 50)=1-P(X=51)-P(X=52).
\]
Now,
\[
P(X=51)=\binom{52}{51}(0.95)^{51}(0.05),
\]
and
\[
P(X=52)=\binom{52}{52}(0.95)^{52}.
\]
Therefore,
\[
P(X\le 50)
=
1-\binom{52}{51}(0.95)^{51}(0.05)
-(0.95)^{52}.
\]
Since $\binom{52}{51}=52$, this becomes
\[
P(X\le 50)
=
1-52(0.95)^{51}(0.05)-(0.95)^{52}.
\]
Thus,
\[
{
P(X\le 50)
=
1-52(0.95)^{51}(0.05)-(0.95)^{52}
}
\]
and numerically,
\[
P(X\le 50)\approx 0.7405.
\]
}

\begin{exercise}
\label{ex:random-variables-17}
Suppose that an experiment can result in one of $r$ possible outcomes, the $i$th outcome having probability $p_i$, $i=1,\ldots,r$, with
\[
\sum_{i=1}^r p_i=1.
\]
If $n$ of these experiments are performed independently, show that the probability that the first outcome appears $x_1$ times, the second $x_2$ times, and the $r$th $x_r$ times is
\[
\frac{n!}{x_1!x_2!\cdots x_r!}p_1^{x_1}p_2^{x_2}\cdots p_r^{x_r},
\]
where
\[
x_1+x_2+\cdots+x_r=n.
\]
\end{exercise}

\solution{
Consider any particular ordering in which outcome $1$ appears $x_1$ times, outcome $2$ appears $x_2$ times, and so on, up to outcome $r$ appearing $x_r$ times.

By independence, the probability of one such particular ordering is
\[
p_1^{x_1}p_2^{x_2}\cdots p_r^{x_r}.
\]

It remains to count how many such orderings are possible. We must arrange $n$ outcomes, with $x_1$ of type $1$, $x_2$ of type $2$, and so on. The number of such arrangements is the multinomial coefficient
\[
\frac{n!}{x_1!x_2!\cdots x_r!}.
\]

Therefore, multiplying the probability of one ordering by the number of orderings gives
\[
P(X_1=x_1,\ldots,X_r=x_r)
=
\frac{n!}{x_1!x_2!\cdots x_r!}
p_1^{x_1}p_2^{x_2}\cdots p_r^{x_r},
\]
where
\[
x_1+x_2+\cdots+x_r=n.
\]
}

\begin{exercise}
\label{ex:random-variables-18}
Show that when $r=2$ the multinomial reduces to the binomial.
\end{exercise}

\solution{
When $r=2$, there are only two possible outcomes. Let their probabilities be
\[
p_1=p
\qquad\text{and}\qquad
p_2=1-p.
\]
Let
\[
x_1=k
\qquad\text{and}\qquad
x_2=n-k.
\]
The multinomial formula becomes
\[
P(X_1=k,X_2=n-k)
=
\frac{n!}{k!(n-k)!}p^k(1-p)^{n-k}.
\]
Since
\[
\frac{n!}{k!(n-k)!}=\binom{n}{k},
\]
we get
\[
P(X=k)=\binom{n}{k}p^k(1-p)^{n-k},
\]
which is exactly the binomial probability mass function.
}

\begin{exercise}
\label{ex:random-variables-19}
In Exercise 17, let $X_i$ denote the number of times the $i$th outcome appears, $i=1,\ldots,r$. What is the probability mass function of
\[
X_1+X_2+\cdots+X_k?
\]
\end{exercise}

\solution{
Let
\[
Y=X_1+X_2+\cdots+X_k.
\]
This random variable counts the number of experiments whose outcome is one of the first $k$ possible outcomes.

For each experiment, define ``success'' as obtaining one of the first $k$ outcomes. Then
\[
P(\text{success})=p_1+p_2+\cdots+p_k.
\]
Let
\[
p=p_1+p_2+\cdots+p_k.
\]
Since the $n$ experiments are independent,
\[
Y\sim \operatorname{Binomial}(n,p).
\]
Therefore, the probability mass function of $Y$ is
\[
P(Y=j)
=
\binom{n}{j}
\left(p_1+p_2+\cdots+p_k\right)^j
\left(1-p_1-p_2-\cdots-p_k\right)^{n-j},
\]
for
\[
j=0,1,\ldots,n.
\]
Thus,
\[
P(X_1+\cdots+X_k=j)
=
\binom{n}{j}
\left(\sum_{i=1}^k p_i\right)^j
\left(1-\sum_{i=1}^k p_i\right)^{n-j}.
\]
}

\begin{exercise}
\label{ex:random-variables-20}
A television store owner figures that $50$ percent of the customers entering his store will purchase an ordinary television set, $20$ percent will purchase a color television set, and $30$ percent will just be browsing. If five customers enter his store on a certain day, what is the probability that two customers purchase color sets, one customer purchases an ordinary set, and two customers purchase nothing?
\end{exercise}

\solution{
There are three possible customer outcomes:
\[
\text{ordinary set}: 0.5,
\qquad
\text{color set}: 0.2,
\qquad
\text{nothing}: 0.3.
\]
We want exactly $1$ ordinary purchase, $2$ color purchases, and $2$ customers purchasing nothing.

Using the multinomial distribution,
\[
P
=
\frac{5!}{1!2!2!}
(0.5)^1(0.2)^2(0.3)^2.
\]
Now,
\[
\frac{5!}{1!2!2!}=30.
\]
Thus,
\[
P
=
30(0.5)(0.04)(0.09).
\]
Therefore,
\[
P=0.054.
\]
}

\begin{exercise}
\label{ex:random-variables-21}
In Exercise 20, what is the probability that our store owner sells three or more televisions on that day?
\end{exercise}

\solution{
A television is sold if the customer purchases either an ordinary set or a color set. Hence
\[
P(\text{sale})=0.5+0.2=0.7.
\]
Let $X$ be the number of televisions sold among the five customers. Then
\[
X\sim \operatorname{Binomial}(5,0.7).
\]
We want
\[
P(X\ge 3).
\]
Thus,
\[
P(X\ge 3)
=
\sum_{k=3}^{5}
\binom{5}{k}(0.7)^k(0.3)^{5-k}.
\]
Compute:
\[
P(X=3)=\binom53(0.7)^3(0.3)^2
=10(0.343)(0.09)=0.3087,
\]
\[
P(X=4)=\binom54(0.7)^4(0.3)
=5(0.2401)(0.3)=0.36015,
\]
\[
P(X=5)=\binom55(0.7)^5
=0.16807.
\]
Therefore,
\[
P(X\ge3)=0.3087+0.36015+0.16807=0.83692.
\]
}

\begin{exercise}
\label{ex:random-variables-22}
If a fair coin is successively flipped, find the probability that a head first appears on the fifth trial.
\end{exercise}

\solution{
For the first head to appear on the fifth trial, the first four flips must be tails and the fifth flip must be heads.

Thus the required outcome is
\[
TTTTH.
\]
Since the coin is fair and flips are independent,
\[
P(TTTTH)
=
\left(\frac12\right)^4\left(\frac12\right)
=
\left(\frac12\right)^5
=
\frac1{32}.
\]
}

\begin{exercise}
\label{ex:random-variables-23}
A coin having probability $p$ of coming up heads is successively flipped until the $r$th head appears. Argue that $X$, the number of flips required, will be $n$, $n\ge r$, with probability
\[
P(X=n)=\binom{n-1}{r-1}p^r(1-p)^{n-r}.
\]
\end{exercise}

\solution{
For the $r$th head to appear on the $n$th flip, two things must happen:

\begin{enumerate}
\item In the first $n-1$ flips, there must be exactly $r-1$ heads.
\item The $n$th flip must be a head.
\end{enumerate}

The number of ways to arrange exactly $r-1$ heads among the first $n-1$ flips is
\[
\binom{n-1}{r-1}.
\]

For any such arrangement, the probability of having $r-1$ heads and $n-r$ tails in the first $n-1$ flips is
\[
p^{r-1}(1-p)^{n-r}.
\]
The probability that the $n$th flip is a head is
\[
p.
\]
Therefore,
\[
P(X=n)
=
\binom{n-1}{r-1}
p^{r-1}(1-p)^{n-r}p.
\]
Hence
\[
P(X=n)
=
\binom{n-1}{r-1}
p^r(1-p)^{n-r},
\qquad n\ge r.
\]
}

\setcounter{exercise}{23}

\begin{exercise}
\label{ex:random-variables-24}
The probability mass function of $X$ is given by
\[
p(k) = \binom{r+k-1}{r-1} p^r (1-p)^k, \quad k = 0, 1, \ldots
\]
Give a possible interpretation of the random variable $X$.
\end{exercise}

\solution{
From Exercise 23, $X$ can be interpreted as the \textbf{number of failures} (tails) that occur before the $r$-th success (head) when a biased coin with probability $p$ of heads is successively flipped. For the $r$-th head to appear after exactly $k$ tails, the first $r+k-1$ flips must contain exactly $r-1$ heads (in $\binom{r+k-1}{r-1}$ arrangements), and the $(r+k)$-th flip must be a head. This yields:
\[
p(k) = \binom{r+k-1}{r-1}p^{r-1}(1-p)^k \cdot p = \binom{r+k-1}{r-1}p^r(1-p)^k.
\]
This is the \textbf{negative binomial distribution} with parameters $r$ and $p$.
}

\noindent
\textit{In Exercises 25 and 26, suppose that two teams are playing a series of games, each of which is independently won by team $A$ with probability $p$ and by team $B$ with probability $1-p$. The winner of the series is the first team to win $i$ games.}

\begin{exercise}
\label{ex:random-variables-25}
If $i = 4$, find the probability that a total of $7$ games are played. Also show that this probability is maximized when $p = \frac{1}{2}$.
\end{exercise}

\solution{
For the series to last exactly 7 games, the score after 6 games must be tied at $3$--$3$, after which game 7 is decisive. The probability that exactly 3 of the first 6 games are won by team $A$ is $\binom{6}{3}p^3(1-p)^3$, so:
\[
P(\text{7 games played}) = \binom{6}{3}p^3(1-p)^3 = 20\,p^3(1-p)^3.
\]
To maximise, let $f(p) = [p(1-p)]^3$. Differentiating:
\[
f'(p) = 3[p(1-p)]^2(1-2p) = 0 \implies p = \tfrac{1}{2}.
\]
Since $f''(1/2) < 0$, $p = \frac{1}{2}$ is a global maximum (on $(0,1)$), giving
$P = 20 \cdot \left(\frac{1}{2}\right)^6 = \frac{20}{64} = \frac{5}{16}$.
}

\begin{exercise}
\label{ex:random-variables-26}
Find the expected number of games that are played when
\begin{enumerate}
    \item [(a)] $i = 2$;
    \item [(b)] $i = 3$.
\end{enumerate}
In both cases, show that this number is maximized when $p = \frac{1}{2}$.
\end{exercise}

\solution{
\textbf{(a) $i = 2$ (first to win 2 games).} The series lasts 2 games if one team wins both immediately (probability $p^2+(1-p)^2$), and 3 games if the first two split (probability $2p(1-p)$, after which game 3 is decisive). Thus:
\[
E[N] = 2\bigl(p^2+(1-p)^2\bigr) + 3\cdot 2p(1-p) = -2p^2+2p+2.
\]
Differentiating: $\frac{d}{dp}(-2p^2+2p+2) = -4p+2 = 0 \Rightarrow p = \frac{1}{2}$. Since the second derivative is $-4 < 0$, $p = \frac{1}{2}$ is a maximum, giving $E[N] = 5/2$.

\textbf{(b) $i = 3$ (first to win 3 games).} The series can last 3, 4, or 5 games:
\begin{align*}
P(3\text{ games}) &= p^3+(1-p)^3,\\
P(4\text{ games}) &= 3\bigl[p^3(1-p)+p(1-p)^3\bigr] = 3p(1-p)\bigl[p^2+(1-p)^2\bigr],\\
P(5\text{ games}) &= 6p^2(1-p)^2.
\end{align*}
Therefore:
\[
E[N] = 3\bigl[p^3+(1-p)^3\bigr] + 12p(1-p)\bigl[p^2+(1-p)^2\bigr] + 30\,p^2(1-p)^2.
\]
Since $E[N]$ is symmetric in $p$ and $1-p$, it has a critical point at $p = \frac{1}{2}$. Verifying it is a maximum (rather than a minimum) by evaluating: at $p = \frac{1}{2}$,
\[
E[N] = 3\cdot\tfrac{1}{4} + 4\cdot\tfrac{3}{8} + 5\cdot\tfrac{3}{8} = \tfrac{3}{4}+\tfrac{3}{2}+\tfrac{15}{8} = \frac{33}{8}.
\]
To prove global maximality, put $u=p(1-p)\in[0,1/4]$. The expression simplifies to $E[N]=3+3u+6u^2$, which is strictly increasing in $u$. Hence its maximum is attained at $u=1/4$, or $p=1/2$.
}

\begin{exercise}
\label{ex:random-variables-27}
A fair coin is independently flipped $n$ times, $k$ times by $A$ and $n-k$ times by $B$. Show that the probability that $A$ and $B$ flip the same number of heads is equal to the probability that there are a total of $k$ heads.
\end{exercise}

\solution{
Let $H_A \sim \mathrm{Bin}(k,\tfrac{1}{2})$ and $H_B \sim \mathrm{Bin}(n-k,\tfrac{1}{2})$ be independent. Then:
\[
P(H_A = H_B) = \sum_{i=0}^{\min(k,n-k)} P(H_A=i)\,P(H_B=i) = \sum_{i=0}^{\min(k,n-k)} \frac{\binom{k}{i}}{2^k}\cdot\frac{\binom{n-k}{i}}{2^{n-k}} = \frac{1}{2^n}\sum_i \binom{k}{i}\binom{n-k}{i}.
\]
Using the Vandermonde identity $\sum_i \binom{k}{i}\binom{n-k}{n-k-i} = \binom{n}{n-k} = \binom{n}{k}$ (and noting $\binom{n-k}{i} = \binom{n-k}{n-k-i}$):
\[
P(H_A = H_B) = \frac{\binom{n}{k}}{2^n}.
\]
The probability of exactly $k$ heads in $n$ fair flips is $\binom{n}{k}/2^n$. The two quantities are equal. 
}

\begin{exercise}
\label{ex:random-variables-28}
Suppose that we want to generate a random variable $X$ equally likely to be either $0$ or $1$, and all we have at our disposal is a biased coin with (unknown) probability $p$ of landing heads. Consider the following procedure:
\begin{enumerate}
    \item Flip the coin, and let $O_1$ be the result.
    \item Flip the coin again, and let $O_2$ be the result.
    \item If $O_1$ and $O_2$ are the same, return to step 1.
    \item If $O_2$ is heads, set $X = 0$; otherwise set $X = 1$.
\end{enumerate}
\begin{enumerate}
    \item [(a)] Show that $X$ is equally likely to be either $0$ or $1$.
    \item [(b)] Could we use a simpler procedure that continues to flip the coin until the last two flips are different, and then sets $X=0$ if the final flip is a head, and $X=1$ if it is a tail?
\end{enumerate}
\end{exercise}

\solution{
Assume $0<p<1$ and independent flips; for $p=0$ or $p=1$, neither procedure terminates.\par
\textbf{(a)} The procedure terminates only when the two outcomes differ (either HT or TH). Conditional on the outcomes being different, each configuration is equally likely:
\[P(O_1 = T, O_2 = H \mid O_1 \neq O_2) = \frac{P(TH)}{P(TH) + P(HT)} = \frac{p(1-p)}{p(1-p) + p(1-p)} = \frac{p(1-p)}{2p(1-p)} = \frac{1}{2}\]
Since $X=0$ iff $O_2=H$ and $X=1$ iff $O_2=T$, it follows that $P(X=0) = P(X=1) = \frac{1}{2}$. Because the procedure starts over whenever $O_1=O_2$, it terminates in finite time (with probability 1) and the output is unbiased regardless of $p$.

\textbf{(b)} No, except when $p=1/2$. The first change is from tails to heads exactly when the initial flip is tails, so $P(X=0)=1-p$ and $P(X=1)=p$. Overlapping pairs therefore produce a biased output.
}

\begin{exercise}
\label{ex:random-variables-29}
Consider $n$ independent flips of a coin having probability $p$ of landing heads. Say a \emph{changeover} occurs whenever an outcome differs from the one preceding it. If $p = \frac{1}{2}$, what is the probability there are $k$ changeovers?
\end{exercise}

\solution{
Let $Y_i = \mathbf{1}[\text{flip } i \text{ differs from flip } i-1]$ for $i = 2, \ldots, n$. When $p = \frac{1}{2}$, adjacent flips are i.i.d.\ Bernoulli$(\frac{1}{2})$, so $P(Y_i = 1) = 2 \cdot \frac{1}{2} \cdot \frac{1}{2} = \frac{1}{2}$.

Moreover, for any specified vector $(y_2,\ldots,y_n)\in\{0,1\}^{n-1}$, exactly two flip sequences produce that vector: choosing the first flip determines every subsequent flip. Thus each changeover vector has probability $2/2^n=2^{-(n-1)}$. This proves mutual independence of the $Y_i$, not merely pairwise independence.

Therefore the total number of changeovers $K = \sum_{i=2}^n Y_i \sim \mathrm{Bin}(n-1,\tfrac{1}{2})$, and:
\[
P(K = k) = \binom{n-1}{k}\left(\frac{1}{2}\right)^{n-1}, \quad k = 0, 1, \ldots, n-1.
\]
}

\begin{exercise}
\label{ex:random-variables-30}
Let $X$ be a Poisson random variable with parameter $\lambda$. Show that $P\{X = i\}$ increases monotonically and then decreases monotonically as $i$ increases, reaching its maximum when $i$ is the largest integer not exceeding $\lambda$.
\end{exercise}

\solution{
Assume $\lambda>0$; for $\lambda=0$, the unique mode is $0$. Consider the ratio of successive Poisson probabilities:
\[
\frac{P\{X=i\}}{P\{X=i-1\}} = \frac{e^{-\lambda}\lambda^i/i!}{e^{-\lambda}\lambda^{i-1}/(i-1)!} = \frac{\lambda}{i}.
\]
This ratio exceeds $1$ iff $i < \lambda$, equals $1$ iff $i = \lambda$ (integer), and is less than $1$ iff $i > \lambda$. Therefore:
\begin{itemize}
    \item $P\{X=i\}$ is strictly increasing while $i < \lambda$;
    \item $P\{X=i\}$ is strictly decreasing for $i > \lambda$.
\end{itemize}
If $\lambda$ is not an integer, the unique maximum is attained at $k = \lfloor\lambda\rfloor$. If $\lambda$ is an integer, then $P\{X=\lambda\}/P\{X=\lambda-1\} = 1$, so both $k = \lambda-1$ and $k = \lambda$ achieve the same (largest) probability. 
}

\begin{exercise}
\label{ex:random-variables-31}
Compare the Poisson approximation with the correct binomial probability for the following cases:
\begin{enumerate}
    \item [(a)] $P\{X=2\}$ when $n=8$, $p=0.1$.
    \item [(b)] $P\{X=9\}$ when $n=10$, $p=0.95$.
    \item [(c)] $P\{X=0\}$ when $n=10$, $p=0.1$.
    \item [(d)] $P\{X=4\}$ when $n=9$, $p=0.2$.
\end{enumerate}
\end{exercise}

\solution{
For the Poisson approximation we use parameter $\lambda = np$ in each case.

\textbf{(a)} $n=8$, $p=0.1$, $\lambda=0.8$:
\[
\text{Binomial: }\binom{8}{2}(0.1)^2(0.9)^6 \approx 0.14880, \qquad \text{Poisson: } \frac{e^{-0.8}(0.8)^2}{2!}\approx 0.14379.
\]
The approximation is very accurate ($\approx 3\%$ relative error).

\textbf{(b)} $n=10$, $p=0.95$, $\lambda=9.5$:
\[
\text{Binomial: }\binom{10}{9}(0.95)^9(0.05)^1 \approx 0.31512, \qquad \text{Poisson: }\frac{e^{-9.5}(9.5)^9}{9!}\approx 0.13000.
\]
The approximation is very poor, as the Poisson approximation is valid only when $p$ is small and $n$ is large, not when $p \approx 1$.

\textbf{(c)} $n=10$, $p=0.1$, $\lambda=1$:
\[
\text{Binomial: }(0.9)^{10}\approx 0.34868, \qquad \text{Poisson: }e^{-1}\approx 0.36788.
\]
A reasonable approximation.

\textbf{(d)} $n=9$, $p=0.2$, $\lambda=1.8$:
\[
\text{Binomial: }\binom{9}{4}(0.2)^4(0.8)^5\approx 0.06606, \qquad \text{Poisson: }\frac{e^{-1.8}(1.8)^4}{4!}\approx 0.07230.
\]
A fair approximation.
}

\begin{exercise}
\label{ex:random-variables-32}
If you buy a lottery ticket in 50 lotteries, in each of which your chance of winning a prize is $\frac{1}{100}$, what is the (approximate) probability that you will win a prize (a) at least once, (b) exactly once, (c) at least twice?
\end{exercise}

\solution{
The number of prizes won $X \sim \mathrm{Bin}(50, \frac{1}{100})$ is well approximated by $X \sim \mathrm{Poisson}(\lambda)$ with $\lambda = 50 \times \frac{1}{100} = \frac{1}{2}$.

\begin{enumerate}
    \item [(a)] $P(X \geq 1) = 1 - P(X=0) \approx 1 - e^{-1/2} \approx 0.3935$.
    \item [(b)] $P(X = 1) \approx e^{-1/2} \cdot \frac{1}{2} \approx 0.3033$.
    \item [(c)] $P(X \geq 2) = 1 - P(X=0) - P(X=1) \approx 1 - e^{-1/2} - \frac{1}{2}e^{-1/2} = 1 - \frac{3}{2}e^{-1/2} \approx 0.0902$.
\end{enumerate}
}

\begin{exercise}
\label{ex:random-variables-33}
Let $X$ be a random variable with probability density
\[
f(x) = \begin{cases} c(1-x^2), & -1 < x < 1, \\ 0, & \text{otherwise.} \end{cases}
\]
\begin{enumerate}
    \item [(a)] What is the value of $c$?
    \item [(b)] What is the cumulative distribution function of $X$?
\end{enumerate}
\end{exercise}

\solution{
\textbf{(a)} Normalising the density:
\[
1 = \int_{-1}^{1}c(1-x^2)\,dx = c\left[x - \frac{x^3}{3}\right]_{-1}^{1} = c\left[\left(1-\frac{1}{3}\right)-\left(-1+\frac{1}{3}\right)\right] = c\cdot\frac{4}{3}.
\]
Hence $c = \dfrac{3}{4}$.

\textbf{(b)} For $-1 < x < 1$:
\[
F(x) = \int_{-1}^{x}\frac{3}{4}(1-t^2)\,dt = \frac{3}{4}\left[t-\frac{t^3}{3}\right]_{-1}^{x} = \frac{3}{4}\left(x-\frac{x^3}{3}+\frac{2}{3}\right).
\]
The complete CDF is $F(x) = 0$ for $x \leq -1$, the expression above for $-1 < x < 1$, and $F(x) = 1$ for $x \geq 1$.
}

\begin{exercise}
\label{ex:random-variables-34}
Let the probability density of $X$ be given by
\[
f(x) = \begin{cases} c(4x-2x^2), & 0 < x < 2, \\ 0, & \text{otherwise.}\end{cases}
\]
\begin{enumerate}
    \item [(a)] What is the value of $c$?
    \item [(b)] $P\!\left\{\frac{1}{2} < X < \frac{3}{2}\right\} = ?$
\end{enumerate}
\end{exercise}

\solution{
\textbf{(a)} Normalising:
\[
1 = \int_0^2 c(4x-2x^2)\,dx = c\left[2x^2-\frac{2x^3}{3}\right]_0^2 = c\left(8-\frac{16}{3}\right) = \frac{8c}{3}.
\]
Hence $c = \dfrac{3}{8}$.

\textbf{(b)}
\[
P\!\left\{\tfrac{1}{2}<X<\tfrac{3}{2}\right\} = \frac{3}{8}\int_{1/2}^{3/2}(4x-2x^2)\,dx = \frac{3}{8}\left[2x^2-\frac{2x^3}{3}\right]_{1/2}^{3/2}.
\]
Evaluating at $x=\frac{3}{2}$: $2\cdot\frac{9}{4}-\frac{2}{3}\cdot\frac{27}{8} = \frac{9}{2}-\frac{9}{4} = \frac{9}{4}$.
Evaluating at $x=\frac{1}{2}$: $2\cdot\frac{1}{4}-\frac{2}{3}\cdot\frac{1}{8} = \frac{1}{2}-\frac{1}{12} = \frac{5}{12}$.
Therefore:
\[
P\!\left\{\tfrac{1}{2}<X<\tfrac{3}{2}\right\} = \frac{3}{8}\cdot\frac{11}{6} = \frac{11}{16}.
\]
}

\begin{exercise}
\label{ex:random-variables-35}
The density of $X$ is given by
\[
f(x) = \begin{cases} 10/x^2, & x > 10,\\ 0, & x \leq 10.\end{cases}
\]
What is the distribution of $X$? Find $P\{X>20\}$.
\end{exercise}

\solution{
This is a \textbf{Pareto distribution}. The CDF for $x > 10$ is:
\[
F(x) = \int_{10}^{x}\frac{10}{t^2}\,dt = \left[-\frac{10}{t}\right]_{10}^{x} = 1 - \frac{10}{x}.
\]
Therefore:
\[
P\{X > 20\} = 1 - F(20) = \frac{10}{20} = \frac{1}{2}.
\]
}

\begin{exercise}
\label{ex:random-variables-36}
A point is uniformly distributed within the disk of radius $1$. That is, its density is
\[
f(x,y) = C, \quad 0 \leq x^2+y^2 \leq 1.
\]
Find the probability that its distance from the origin is less than $x$, $0 \leq x \leq 1$.
\end{exercise}

\solution{
Since the disk has area $\pi$, the normalising constant is $C = 1/\pi$. Let $R = \sqrt{X^2+Y^2}$ be the distance from the origin. By symmetry, the probability that $R < x$ equals the ratio of the area of a disk of radius $x$ to the area of the unit disk:
\[
P\{R < x\} = \frac{\pi x^2}{\pi} = x^2, \quad 0 \leq x \leq 1.
\]
So the distribution function of $R$ is $F_R(x) = x^2$, and the corresponding density is $f_R(x) = 2x$ for $0 \leq x \leq 1$.
}

\begin{exercise}
\label{ex:random-variables-37}
Let $X_1, X_2, \ldots, X_n$ be independent random variables, each having a uniform distribution over $(0,1)$. Let $M = \max(X_1, \ldots, X_n)$. Show that the distribution function of $M$ is given by $F_M(x) = x^n$ for $0 \leq x \leq 1$. What is the probability density function of $M$?
\end{exercise}

\solution{
Since the $X_i$ are independent and each has CDF $F_{X_i}(x) = x$ on $[0,1]$:
\[
F_M(x) = P(M \leq x) = P(X_1 \leq x, \ldots, X_n \leq x) = \prod_{i=1}^n P(X_i \leq x) = x^n, \quad 0 \leq x \leq 1.
\]
Differentiating gives the probability density function:
\[
f_M(x) = nx^{n-1}, \quad 0 \leq x \leq 1. \qquad
\]
}

\begin{exercise}
\label{ex:random-variables-38}
If the density function of $X$ equals
\[
f(x) = \begin{cases} ce^{-2x}, & 0 < x < \infty,\\ 0, & x < 0,\end{cases}
\]
find $c$. What is $P\{X > 2\}$?
\end{exercise}

\solution{
Normalising: $\int_0^{\infty}ce^{-2x}\,dx = \frac{c}{2} = 1$, so $c = 2$.
Thus $X \sim \mathrm{Exp}(2)$, and:
\[
P\{X > 2\} = \int_2^{\infty}2e^{-2x}\,dx = \left[-e^{-2x}\right]_2^{\infty} = e^{-4}.
\]
}

\begin{exercise}
\label{ex:random-variables-39}
The random variable $X$ has the following probability mass function:
\[
p(1) = \tfrac{1}{2},\quad p(2) = \tfrac{1}{3},\quad p(24) = \tfrac{1}{6}.
\]
Calculate $E[X]$.
\end{exercise}

\solution{
By definition:
\[
E[X] = 1\cdot\frac{1}{2} + 2\cdot\frac{1}{3} + 24\cdot\frac{1}{6} = \frac{1}{2}+\frac{2}{3}+4 = \frac{3}{6}+\frac{4}{6}+\frac{24}{6} = \frac{31}{6}.
\]
}

\begin{exercise}
\label{ex:random-variables-40}
Suppose that two teams are playing a series of games, each of which is independently won by team $A$ with probability $p$ and by team $B$ with probability $1-p$. The winner of the series is the first team to win four games. Find the expected number of games that are played, and evaluate this quantity when $p = \frac{1}{2}$.
\end{exercise}

\solution{
The series can last 4, 5, 6, or 7 games. For team $A$ to win the series in exactly $k$ games, $A$ must win the $k$-th game and exactly 3 of the first $k-1$ games. By symmetry with team $B$:
\begin{align*}
P(4\text{ games}) &= p^4+(1-p)^4,\\
P(5\text{ games}) &= 4\bigl[p^4(1-p)+p(1-p)^4\bigr],\\
P(6\text{ games}) &= 10\bigl[p^4(1-p)^2+p^2(1-p)^4\bigr],\\
P(7\text{ games}) &= 20\,p^3(1-p)^3.
\end{align*}
Therefore:
\[
E[N] = 4P(4) + 5P(5) + 6P(6) + 7P(7).
\]
At $p = \tfrac{1}{2}$, the probabilities simplify to:
\[
P(4)=\tfrac{1}{8},\quad P(5)=\tfrac{1}{4},\quad P(6)=\tfrac{5}{16},\quad P(7)=\tfrac{5}{16},
\]
giving:
\[
E[N] = 4\cdot\frac{1}{8}+5\cdot\frac{1}{4}+6\cdot\frac{5}{16}+7\cdot\frac{5}{16} = \frac{8+20+30+35}{16} = \frac{93}{16} \approx 5.81.
\]
}

\begin{exercise}
\label{ex:random-variables-41}
Consider the case of arbitrary $p$ in Exercise 29. Compute the expected number of changeovers.
\end{exercise}

\solution{
A changeover occurs at position $i$ (for $i=2,\ldots,n$) with probability:
\[
P(\text{changeover at } i) = p(1-p)+(1-p)p = 2p(1-p).
\]
Let $Y_i = \mathbf{1}[\text{changeover at position } i]$. By linearity of expectation, the expected number of changeovers is:
\[
E\!\left[\sum_{i=2}^n Y_i\right] = (n-1)\cdot 2p(1-p).
\]
}

\begin{exercise}
\label{ex:random-variables-42}
Suppose that each coupon obtained is, independent of what has been previously obtained, equally likely to be any of $m$ different types. Find the expected number of coupons one needs to obtain in order to have at least one of each type.
\end{exercise}

\solution{
Let $X_j$ be the additional number of coupons needed to obtain a new type after having collected $j-1$ distinct types. When $j-1$ types have been seen, the probability that the next coupon is new is $\frac{m-(j-1)}{m}$, so $X_j \sim \mathrm{Geom}\!\left(\frac{m-j+1}{m}\right)$ and:
\[
E[X_j] = \frac{m}{m-j+1}.
\]
The total number $X = X_1 + X_2 + \cdots + X_m$ satisfies:
\[
E[X] = \sum_{j=1}^m \frac{m}{m-j+1} = m\sum_{k=1}^m \frac{1}{k} = m\,H_m,
\]
where $H_m = \sum_{k=1}^m \frac{1}{k}$ is the $m$-th harmonic number.
}

\begin{exercise}
\label{ex:random-variables-43}
An urn contains $n+m$ balls, of which $n$ are red and $m$ are black. They are withdrawn one at a time without replacement. Let $X$ be the number of red balls removed before the first black ball is chosen.
\begin{enumerate}
    \item [(a)] Express $X$ in terms of indicator random variables $X_i$, $i = 1,\ldots,n$, where $X_i = 1$ if red ball $i$ is taken before any black ball.
    \item  [(b)] Find $E[X]$.
\end{enumerate}
\end{exercise}

\solution{
\textbf{(a)} Define $X_i = 1$ if red ball $i$ is drawn before all black balls, and $X_i = 0$ otherwise. Then $X = \sum_{i=1}^n X_i$.

\textbf{(b)} Consider the $m+1$ balls consisting of red ball $i$ and the $m$ black balls. In a uniformly random ordering of all $n+m$ balls, each of the $m+1$ positions within this sub-sequence is equally likely to be first. Therefore:
\[
E[X_i] = P(\text{red ball } i \text{ precedes all black balls}) = \frac{1}{m+1}.
\]
By linearity of expectation:
\[
E[X] = \sum_{i=1}^n E[X_i] = \frac{n}{m+1}.
\]
}

\begin{exercise}
\label{ex:random-variables-44}
In Exercise 43, let $Y$ denote the number of red balls chosen after the first but before the second black ball has been chosen.
\begin{enumerate}
    \item [(a)] Express $Y$ as the sum of $n$ indicator random variables.
    \item [(b)] Find $E[Y]$.
    \item [(c)] Compare $E[Y]$ to $E[X]$ obtained in Exercise 43.
    \item [(d)]Can you explain the result obtained in (c)?
\end{enumerate}
\end{exercise}

\solution{
\textbf{(a)} Define $X_i = 1$ if red ball $i$ appears between the first and second black ball (i.e., after the first black but before the second). Then $Y = \sum_{i=1}^n X_i$.

\textbf{(b)} Among the $m+1$ objects consisting of red ball $i$ and the $m$ black balls, each of the $m+1$ relative positions is equally likely. Red ball $i$ is exactly second in this sub-sequence (i.e., immediately after the first black ball, before the second) with probability $\frac{1}{m+1}$. Therefore:
\[
E[Y] = \sum_{i=1}^n \frac{1}{m+1} = \frac{n}{m+1}.
\]

\textbf{(c)} $E[Y] = E[X] = \dfrac{n}{m+1}$.

\textbf{(d)} By symmetry, in a random permutation of $n+m$ balls, for each red ball $i$, the position of red ball $i$ relative to the $m$ black balls (i.e., before ball 1, between ball 1 and 2, \ldots, after ball $m$) is equally likely to be any of the $m+1$ gaps. Hence the expected number of red balls falling in each particular gap is the same: $\frac{n}{m+1}$.
}

\begin{exercise}
\label{ex:random-variables-45}
A total of $r$ keys are to be put, one at a time, in $k$ boxes, with each key independently being put in box $i$ with probability $p_i$, $\sum_{i=1}^k p_i = 1$. Each time a key is put in a nonempty box, a collision occurs. Find the expected number of collisions.
\end{exercise}

\solution{
A collision occurs if and only if a key is placed in an already-occupied box. The expected number of collisions equals $r$ minus the expected number of distinct boxes used:
\[
E[\text{collisions}] = r - E[\text{distinct boxes used}].
\]
Box $i$ is occupied after $r$ keys are placed with probability $1-(1-p_i)^r$ (complementary to all $r$ keys avoiding box $i$). By linearity:
\[
E[\text{distinct boxes used}] = \sum_{i=1}^k \bigl[1-(1-p_i)^r\bigr] = k - \sum_{i=1}^k (1-p_i)^r.
\]
Therefore:
\[
E[\text{collisions}] = r - k + \sum_{i=1}^k (1-p_i)^r.
\]
}

\begin{exercise}
\label{ex:random-variables-46}
If $X$ is a nonnegative integer-valued random variable, show that
\[
E[X] = \sum_{n=1}^{\infty} P\{X \geq n\} = \sum_{n=0}^{\infty} P\{X > n\}.
\]
Also, if $X$ and $Y$ are both nonnegative integer-valued, show that
\[
E[XY] = \sum_{n=1}^{\infty}\sum_{m=1}^{\infty} P(X \geq n,\, Y \geq m).
\]
\end{exercise}

\solution{
\textbf{(a)} Define $I_n = \mathbf{1}[n \leq X]$ for $n \geq 1$. Then $X = \sum_{n=1}^{\infty} I_n$, so by linearity of expectation:
\[
E[X] = \sum_{n=1}^{\infty} E[I_n] = \sum_{n=1}^{\infty} P\{X \geq n\}.
\]
Re-indexing with $n' = n-1$ gives the second formula $\sum_{n=0}^{\infty} P\{X > n\}$.

\textbf{(b)} Write $X = \sum_{n=1}^{\infty} I_n^X$ and $Y = \sum_{m=1}^{\infty} I_m^Y$ with $I_n^X = \mathbf{1}[X \geq n]$ and $I_m^Y = \mathbf{1}[Y \geq m]$. Then:
\[
XY = \Bigl(\sum_{n=1}^{\infty} I_n^X\Bigr)\Bigl(\sum_{m=1}^{\infty} I_m^Y\Bigr) = \sum_{n=1}^{\infty}\sum_{m=1}^{\infty} I_n^X I_m^Y.
\]
Taking expectations and exchanging with the nonnegative sum by Tonelli's theorem (equivalently, monotone convergence):
\[
E[XY] = \sum_{n=1}^{\infty}\sum_{m=1}^{\infty} E[I_n^X I_m^Y] = \sum_{n=1}^{\infty}\sum_{m=1}^{\infty} P(X \geq n,\, Y \geq m). \qquad
\]
}

\begin{exercise}
\label{ex:random-variables-47}
Consider three trials, each of which is either a success or not. Let $X$ denote the number of successes. Suppose that $E[X] = 1.8$.
\begin{enumerate}
    \item [(a)] What is the largest possible value of $P\{X = 3\}$?
    \item [(b)] What is the smallest possible value of $P\{X = 3\}$?
\end{enumerate}
In both cases, construct a probability scenario that results in $P\{X=3\}$ having the desired value.
\end{exercise}

\solution{
Write $X = X_1 + X_2 + X_3$ where $X_i$ are Bernoulli with $p_i = E[X_i]$, and $p_1+p_2+p_3 = 1.8$.

\textbf{(a) Maximum:} Since $P\{X=3\} = P\{X_1=1, X_2=1, X_3=1\} \leq \min(p_1,p_2,p_3)$, to maximise, take all $p_i$ equal: $p_1=p_2=p_3=0.6$. Let $X_1=X_2=X_3$ (perfectly correlated), so $X \in \{0,3\}$ with $P(X=3) = P(X_i=1) = 0.6$ and $E[X] = 1.8$.
\[{P\{X=3\}_{\max} = 0.6.}\]

\textbf{(b) Minimum:} $P\{X=3\} \geq 0$ trivially. Setting $p_3 = 0$ (trial 3 always fails) with $p_1 = 1$, $p_2 = 0.8$ gives $p_1+p_2+p_3 = 1.8$ and $X_3 = 0$ always, so $X \leq 2$ and thus $P\{X=3\} = 0$.
\[{P\{X=3\}_{\min} = 0.}\]
}

\begin{exercise}
\label{ex:random-variables-48}
If $X$ is a nonnegative random variable and $g$ is a differentiable function with $g(0) = 0$, then
\[
E[g(X)] = \int_0^{\infty} P(X > t)\, g'(t)\, dt.
\]
Prove the above when $X$ is a continuous random variable.
\end{exercise}

\solution{
To justify the interchange, assume $g$ is locally absolutely continuous (for example, continuously differentiable), $g(0)=0$, and $\int_0^\infty P(X>t)|g'(t)|\,dt<\infty$. If $g'\ge0$, Tonelli also gives the identity with value $+\infty$. Differentiability alone does not ensure these conditions. Since $X\ge0$ has density $f$:
\begin{align*}
\int_0^{\infty} P(X > t)\, g'(t)\, dt &= \int_0^{\infty} g'(t) \left[\int_t^{\infty} f(x)\, dx\right] dt.
\end{align*}
Switching the order of integration (Fubini's theorem, noting $0 \leq t \leq x$):
\begin{align*}
&= \int_0^{\infty} f(x) \left[\int_0^x g'(t)\, dt\right] dx = \int_0^{\infty} f(x)\, g(x)\, dx = E[g(X)]. \qquad
\end{align*}
}

\begin{exercise}
\label{ex:random-variables-49}
Prove that $E[X^2] \geq (E[X])^2$. When do we have equality?
\end{exercise}

\solution{
The variance $\mathrm{Var}(X) = E[X^2] - (E[X])^2$ satisfies $\mathrm{Var}(X) = E[(X - E[X])^2] \geq 0$ since it is an expectation of a nonnegative quantity. Therefore:
\[
E[X^2] \geq (E[X])^2,
\]
with equality if and only if $(X - E[X])^2 = 0$ with probability 1, i.e., $X$ is a \textbf{degenerate} (constant) random variable. 
}

\begin{exercise}
\label{ex:random-variables-50}
Let $c$ be a constant. Show that
\begin{enumerate}
    \item [(a)] $\mathrm{Var}(cX) = c^2\,\mathrm{Var}(X)$;
    \item [(b)] $\mathrm{Var}(c+X) = \mathrm{Var}(X)$.
\end{enumerate}
\end{exercise}

\solution{
Assume $E[X^2]<\infty$.
\textbf{(a)} Since $cX-E[cX]=c(X-E[X])$,
\[
\operatorname{Var}(cX)=E[c^2(X-E[X])^2]=c^2\operatorname{Var}(X).
\]
\textbf{(b)} Adding $c$ shifts the mean by $c$, so
\[
\operatorname{Var}(c+X)=E[(c+X-(c+E[X]))^2]
=\operatorname{Var}(X).
\]
}

\begin{exercise}
\label{ex:random-variables-51}
A coin having probability $p$ of landing heads is flipped until a head appears for the $r$-th time. Let $N$ denote the number of flips required. Calculate $E[N]$.
\end{exercise}

\solution{
Write $N = G_1 + G_2 + \cdots + G_r$, where $G_j$ is the number of additional flips needed to get the $j$-th head, given $j-1$ heads have already been obtained. Each $G_j$ is independent and has a $\mathrm{Geom}(p)$ distribution (geometric with success probability $p$), so $E[G_j] = 1/p$. By linearity:
\[
E[N] = \sum_{j=1}^r E[G_j] = \frac{r}{p}.
\]
}

\begin{exercise}
\label{ex:random-variables-52}
\begin{enumerate}
    \item [(a)] Calculate $E[X]$ for the maximum random variable $M$ of Exercise 37.
    \item [(b)] Calculate $E[X]$ for $X$ as in Exercise 33.
    \item [(c)] Calculate $E[X]$ for $X$ as in Exercise 34.
\end{enumerate}
\end{exercise}

\solution{
\textbf{(a)} From Exercise 37, $M$ has density $f_M(x) = nx^{n-1}$ on $[0,1]$:
\[
E[M] = \int_0^1 x \cdot nx^{n-1}\,dx = n\int_0^1 x^n\,dx = \frac{n}{n+1}.
\]

\textbf{(b)} From Exercise 33, $f(x) = \frac{3}{4}(1-x^2)$ on $(-1,1)$, which is symmetric about $0$. Since $x \cdot f(x)$ is an odd function:
\[
E[X] = \int_{-1}^1 x\cdot\frac{3}{4}(1-x^2)\,dx = 0.
\]

\textbf{(c)} From Exercise 34, $f(x) = \frac{3}{8}(4x-2x^2)$ on $(0,2)$:
\[
E[X] = \frac{3}{8}\int_0^2 x(4x-2x^2)\,dx = \frac{3}{8}\int_0^2 (4x^2-2x^3)\,dx = \frac{3}{8}\left[\frac{4x^3}{3}-\frac{x^4}{2}\right]_0^2 = \frac{3}{8}\cdot\left(\frac{32}{3}-8\right) = \frac{3}{8}\cdot\frac{8}{3} = 1.
\]
}

\begin{exercise}
\label{ex:random-variables-53}
If $X$ is uniform over $(0,1)$, calculate $E[X^n]$ and $\mathrm{Var}(X^n)$.
\end{exercise}

\solution{
Since $X \sim \mathrm{Uniform}(0,1)$ with density $f(x) = 1$:
\[
E[X^n] = \int_0^1 x^n\,dx = \frac{1}{n+1}, \qquad E[X^{2n}] = \int_0^1 x^{2n}\,dx = \frac{1}{2n+1}.
\]
Therefore:
\[
\mathrm{Var}(X^n) = E[X^{2n}] - \left(E[X^n]\right)^2 = \frac{1}{2n+1} - \frac{1}{(n+1)^2} = \frac{(n+1)^2-(2n+1)}{(n+1)^2(2n+1)} = \frac{n^2}{(n+1)^2(2n+1)}.
\]
}

\begin{exercise}
\label{ex:random-variables-54}
Let $X$ and $Y$ each take on either the value $1$ or $-1$. Let $p(1,1)$, $p(1,-1)$, $p(-1,1)$, $p(-1,-1)$ be the joint probabilities. Suppose $E[X] = E[Y] = 0$. Show that:
\begin{enumerate}
    \item  [(a)] $p(1,1) = p(-1,-1)$;
    \item [(b)] $p(1,-1) = p(-1,1)$.
\end{enumerate}
Let $p = 2p(1,1)$. Find $\mathrm{(c)}\ \mathrm{Var}(X)$; $\mathrm{(d)}\ \mathrm{Var}(Y)$; $\mathrm{(e)}\ \mathrm{Cov}(X,Y)$.
\end{exercise}

\solution{
\textbf{(a) and (b).} $E[X]=0$ implies $P(X=1)=P(X=-1)=\frac{1}{2}$, so $p(1,1)+p(1,-1)=p(-1,1)+p(-1,-1)=\frac{1}{2}$. Similarly, $E[Y]=0$ gives $p(1,1)+p(-1,1)=p(1,-1)+p(-1,-1)=\frac{1}{2}$. Summing the first and second pairs of equations and using total probability:
\[
p(1,-1)=\tfrac{1}{2}-p(1,1), \quad p(-1,1)=\tfrac{1}{2}-p(1,1), \quad p(-1,-1)=p(1,1).
\]
Thus $p(1,1)=p(-1,-1)$ and $p(1,-1)=p(-1,1)$.

With $p = 2p(1,1)$: $p(1,1)=p(-1,-1)=\frac{p}{2}$, $p(1,-1)=p(-1,1)=\frac{1-p}{2}$.

\textbf{(c)} Since $X^2 = 1$ always, $E[X^2]=1$ and $\mathrm{Var}(X) = 1-(E[X])^2 = 1$.

\textbf{(d)} Similarly, $\mathrm{Var}(Y) = 1$.

\textbf{(e)}
\[
E[XY] = (1)(1)\tfrac{p}{2}+(1)(-1)\tfrac{1-p}{2}+(-1)(1)\tfrac{1-p}{2}+(-1)(-1)\tfrac{p}{2} = p - (1-p) = 2p-1.
\]
Since $E[X]E[Y]=0$:
\[
\mathrm{Cov}(X,Y) = E[XY] - E[X]E[Y] = 2p-1.
\]
}
\begin{exercise}
\label{ex:random-variables-55}
Suppose that the joint probability mass function of $X$ and $Y$ is
\[
P(X=i,\,Y=j) = \binom{j}{i}e^{-2\lambda}\frac{\lambda^j}{j!}, \quad 0 \leq i \leq j, \quad j=0,1,2,\ldots
\]
\begin{enumerate}
    \item[(a)] Find the probability mass function of $Y$.
    \item[(b)] Find the probability mass function of $X$.
    \item[(c)] Find the probability mass function of $Y - X$.
\end{enumerate}
\end{exercise}

\solution{
\textbf{(a)} Summing over $i = 0, 1, \ldots, j$:
\[P(Y=j) = \sum_{i=0}^j \binom{j}{i}e^{-2\lambda}\frac{\lambda^j}{j!} = e^{-2\lambda}\frac{\lambda^j}{j!}\cdot 2^j = e^{-2\lambda}\frac{(2\lambda)^j}{j!}.
\]
Thus $Y \sim \mathrm{Poisson}(2\lambda)$.

\textbf{(b)} Summing over $j = i, i+1, \ldots$ and using $\binom{j}{i}/j! = 1/(i!\,(j-i)!)$:
\[
P(X=i) = \sum_{j=i}^{\infty}\binom{j}{i}e^{-2\lambda}\frac{\lambda^j}{j!} = \frac{e^{-2\lambda}\lambda^i}{i!}\sum_{k=0}^{\infty}\frac{\lambda^k}{k!} = \frac{e^{-2\lambda}\lambda^i e^{\lambda}}{i!} = \frac{e^{-\lambda}\lambda^i}{i!}.
\]
Thus $X \sim \mathrm{Poisson}(\lambda)$.

\textbf{(c)} For $k = 0, 1, 2, \ldots$, letting $Z = Y - X \geq 0$:
\[
P(Z=k) = \sum_{i=0}^{\infty}P(X=i,\,Y=i+k) = \sum_{i=0}^{\infty}\binom{i+k}{i}e^{-2\lambda}\frac{\lambda^{i+k}}{(i+k)!} = \frac{e^{-2\lambda}\lambda^k}{k!}\sum_{i=0}^{\infty}\frac{\lambda^i}{i!} = \frac{e^{-\lambda}\lambda^k}{k!}.
\]
Thus $Z = Y-X \sim \mathrm{Poisson}(\lambda)$.
}

\begin{exercise}
\label{ex:random-variables-56}
There are $n$ types of coupons. Each newly obtained coupon is, independently, type $i$ with probability $p_i$, $i=1,\ldots,n$. Find the expected number and the variance of the number of distinct types obtained in a collection of $k$ coupons.
\end{exercise}

\solution{
Let $D$ be the number of distinct types and $X_i = \mathbf{1}[\text{type } i \text{ appears in the collection}]$. Then $D = \sum_{i=1}^n X_i$.

\textbf{Expected value:} $P(X_i=1)=1-(1-p_i)^k$, so:
\[
E[D] = \sum_{i=1}^n \bigl[1-(1-p_i)^k\bigr].
\]

\textbf{Variance:} Since $\mathrm{Var}(D) = \sum_i\mathrm{Var}(X_i) + 2\sum_{i<j}\mathrm{Cov}(X_i,X_j)$, with:
\[
\mathrm{Var}(X_i) = (1-p_i)^k\bigl[1-(1-p_i)^k\bigr],
\]
\[
E[X_iX_j] = P(\text{types } i \text{ and } j \text{ both appear}) = 1-(1-p_i)^k-(1-p_j)^k+(1-p_i-p_j)^k,
\]
\[
\mathrm{Cov}(X_i,X_j) = (1-p_i-p_j)^k - (1-p_i)^k(1-p_j)^k,
\]
we obtain:
\[
\mathrm{Var}(D) = \sum_{i=1}^n(1-p_i)^k\bigl[1-(1-p_i)^k\bigr] + 2\sum_{i<j}\bigl[(1-p_i-p_j)^k-(1-p_i)^k(1-p_j)^k\bigr].
\]
}

\begin{exercise}
\label{ex:random-variables-57}
Suppose that $X$ and $Y$ are independent binomial random variables with parameters $(n,p)$ and $(m,p)$, respectively. Argue probabilistically (no computations necessary) that $X+Y$ is binomial with parameters $(n+m,p)$.
\end{exercise}

\solution{
$X$ counts the number of successes in $n$ independent Bernoulli$(p)$ trials, and $Y$ counts the successes in $m$ additional independent Bernoulli$(p)$ trials, with $X$ and $Y$ independent. Therefore $X+Y$ counts the total number of successes in $n+m$ independent Bernoulli$(p)$ trials, which by definition has a $\mathrm{Bin}(n+m,p)$ distribution. 
}

\begin{exercise}
\label{ex:random-variables-58}
An urn contains $2n$ balls, of which $r$ are red. The balls are randomly removed in $n$ successive pairs. Let $X$ denote the number of pairs in which both balls are red.
\begin{enumerate}
    \item[(a)] Find $E[X]$.
    \item[(b)] Find $\mathrm{Var}(X)$.
\end{enumerate}
\end{exercise}

\solution{
Let $X_i = \mathbf{1}[\text{pair } i \text{ consists of two red balls}]$, so $X = \sum_{i=1}^n X_i$.

\textbf{(a)} Consider the two balls assigned to pair $i$: the first is red with probability $r/(2n)$ and given that, the second is red with probability $(r-1)/(2n-1)$, so:
\[
E[X_i] = \frac{r(r-1)}{2n(2n-1)}, \qquad E[X] = \frac{n\,r(r-1)}{2n(2n-1)} = \frac{r(r-1)}{2(2n-1)}.
\]

\textbf{(b)} For $i \neq j$, the probability that pairs $i$ and $j$ are both all-red is:
\[
E[X_iX_j] = \frac{r(r-1)(r-2)(r-3)}{2n(2n-1)(2n-2)(2n-3)}.
\]
For $n=1$, $X$ is deterministic and its variance is $0$. For $n\ge2$, let $p = E[X_i]$ and $q = E[X_iX_j]$. Then $\mathrm{Var}(X_i) = p(1-p)$ and $\mathrm{Cov}(X_i,X_j) = q - p^2$, giving:
\[
\mathrm{Var}(X) = n\,p(1-p) + n(n-1)(q-p^2).
\]
}

\begin{exercise}
\label{ex:random-variables-59}
Let $X_1, X_2, X_3, X_4$ be independent continuous random variables with a common distribution function $F$ and let
\[
p = P\{X_1 < X_2 > X_3 < X_4\}.
\]
\begin{enumerate}
    \item[(a)] Argue that the value of $p$ is the same for all continuous distribution functions $F$.
    \item[(b)] Find $p$ by integrating the joint density function over the appropriate region.
    \item[(c)] Find $p$ by using the fact that all $4!$ possible orderings of $X_1,\ldots,X_4$ are equally likely.
\end{enumerate}
\end{exercise}

\solution{
\textbf{(a)} Since $F$ is continuous, the probability integral transform $U_i = F(X_i)$ gives i.i.d.\ $\mathrm{Uniform}(0,1)$ random variables, and the event $\{X_1<X_2>X_3<X_4\}$ is equivalent to $\{U_1<U_2>U_3<U_4\}$. The probability of this event depends only on the joint density of $(U_1,\ldots,U_4)$, which is always uniform on $[0,1]^4$ regardless of $F$. Hence $p$ is the same for all continuous $F$.

\textbf{(b)} Taking $F = \mathrm{Uniform}(0,1)$ and integrating over the region $x_1 < x_2$, $x_2 > x_3$, $x_3 < x_4$ (with all variables in $[0,1]$):
\begin{align*}
p &= \int_0^1\!\int_0^{x_2}\!\int_0^{x_2}\!\int_{x_3}^1 dx_4\,dx_1\,dx_3\,dx_2 = \int_0^1 x_2(x_2 - x_2^2/2)\,dx_2 = \int_0^1\!\left(x_2^2 - \frac{x_2^3}{2}\right)\!dx_2 = \frac{1}{3}-\frac{1}{8} = \frac{5}{24}.
\end{align*}

\textbf{(c)} All $4! = 24$ orderings of $(X_1,X_2,X_3,X_4)$ are equally likely. Enumerating permutations of $\{1,2,3,4\}$ satisfying $a < b$, $b > c$, $c < d$:
\[
(1,3,2,4),\;(1,4,2,3),\;(2,3,1,4),\;(2,4,1,3),\;(3,4,1,2).
\]
There are exactly 5 such permutations, so $p = 5/24$.
}

\begin{exercise}
\label{ex:random-variables-60}
Calculate the moment generating function of the uniform distribution on $(0,1)$. Obtain $E[X]$ and $\mathrm{Var}[X]$ by differentiating.
\end{exercise}

\solution{
For $X \sim \mathrm{Uniform}(0,1)$ and $t \neq 0$:
\[
M(t) = E[e^{tX}] = \int_0^1 e^{tx}\,dx = \frac{e^t - 1}{t}.
\]
Using the Taylor expansion $e^t - 1 = t + t^2/2 + t^3/6 + \cdots$, we have $M(t) = 1 + t/2 + t^2/6 + \cdots$.

Differentiating $M(t) = (e^t-1)/t$ using L'Hôpital or the Taylor expansion:
\[
M'(t) = \frac{te^t - (e^t-1)}{t^2} \xrightarrow{t\to 0} \frac{1}{2} = E[X].
\]
\[
M''(t) = \frac{(t^2-2t+2)e^t - 2}{t^3} \xrightarrow{t\to 0} \frac{1}{3} = E[X^2].
\]
Therefore: $\mathrm{Var}(X) = E[X^2] - (E[X])^2 = \dfrac{1}{3} - \dfrac{1}{4} = \dfrac{1}{12}$.
}

\begin{exercise}
\label{ex:random-variables-61}
Let $X$ and $W$ be the working and subsequent repair times of a certain machine. Let $Y = X + W$ and suppose the joint density of $X$ and $Y$ is
\[
f_{X,Y}(x,y) = \lambda^2 e^{-\lambda y}, \quad 0 < x < y < \infty.
\]
\begin{enumerate}
    \item[(a)] Find the density of $X$.
    \item[(b)] Find the density of $Y$.
    \item[(c)] Find the joint density of $X$ and $W$.
    \item[(d)] Find the density of $W$.
\end{enumerate}
\end{exercise}

\solution{
\textbf{(a)} $f_X(x) = \int_x^{\infty}\lambda^2 e^{-\lambda y}\,dy = \lambda e^{-\lambda x}$ for $x > 0$. Thus $X \sim \mathrm{Exp}(\lambda)$.

\textbf{(b)} $f_Y(y) = \int_0^y \lambda^2 e^{-\lambda y}\,dx = \lambda^2 y e^{-\lambda y}$ for $y > 0$. Thus $Y \sim \mathrm{Gamma}(2,\lambda)$.

\textbf{(c)} Since $W = Y - X$, use the change of variables $(x, y) \mapsto (x, w)$ with Jacobian
\[
\left|\frac{\partial(x,y)}{\partial(x,w)}\right|=1.
\]
The transformed density is
\[
f_{X,W}(x,w) = f_{X,Y}(x, x+w) = \lambda^2 e^{-\lambda(x+w)} = \bigl(\lambda e^{-\lambda x}\bigr)\bigl(\lambda e^{-\lambda w}\bigr)
\]
for $x, w > 0$. Since this factors, $X$ and $W$ are independent with each having an $\mathrm{Exp}(\lambda)$ distribution.

\textbf{(d)} From (c), $f_W(w) = \lambda e^{-\lambda w}$ for $w > 0$. Thus $W \sim \mathrm{Exp}(\lambda)$.
}

\begin{exercise}
\label{ex:random-variables-62}
Insurance companies sometimes use the exponential principle: with $X$ the random amount to be paid in claims, the premium is $P = \frac{1}{a}\ln E[e^{aX}]$. Find $P$ when $X$ is exponential with parameter $\lambda$, and $a = \alpha\lambda$ where $0 < \alpha < 1$.
\end{exercise}

\solution{
Here $X\sim\mathrm{Exp}(\lambda)$ and $a=\alpha\lambda<\lambda$. Using its density,
\[
E[e^{aX}] = \int_0^{\infty}e^{ax}\lambda e^{-\lambda x}\,dx = \frac{\lambda}{\lambda - a} = \frac{\lambda}{\lambda - \alpha\lambda} = \frac{1}{1-\alpha}.
\]
Therefore the premium is:
\[
P = \frac{1}{a}\ln\!\left(\frac{1}{1-\alpha}\right) = \frac{-\ln(1-\alpha)}{\alpha\lambda}.
\]
}

\begin{exercise}
\label{ex:random-variables-63}
Calculate the moment generating function of a geometric random variable.
\end{exercise}

\solution{
For $p=1$, $X=1$ almost surely and its MGF is $e^t$ for all real $t$. Otherwise assume $0<p<1$. Let $X \sim \mathrm{Geom}(p)$ with $P(X=n) = p(1-p)^{n-1}$ for $n = 1, 2, \ldots$ Then:
\[
M(t) = E[e^{tX}] = \sum_{n=1}^{\infty} e^{tn}\,p(1-p)^{n-1} = pe^t\sum_{n=1}^{\infty}\bigl[(1-p)e^t\bigr]^{n-1} = \frac{pe^t}{1-(1-p)e^t},
\]
provided $(1-p)e^t < 1$, i.e., $t < -\ln(1-p)$.
}

\begin{exercise}
\label{ex:random-variables-64}
Show that the sum of independent identically distributed exponential random variables has a gamma distribution.
\end{exercise}

\solution{
Let $X_1, \ldots, X_n \stackrel{\text{i.i.d.}}{\sim} \mathrm{Exp}(\lambda)$ with MGF $M_{X_i}(t) = \lambda/(\lambda - t)$ for $t < \lambda$. Since the $X_i$ are independent, the MGF of $S = X_1 + \cdots + X_n$ is:
\[
M_S(t) = \prod_{i=1}^n M_{X_i}(t) = \left(\frac{\lambda}{\lambda - t}\right)^n.
\]
This is precisely the MGF of a $\mathrm{Gamma}(n, \lambda)$ distribution (which has density $f(x) = \lambda e^{-\lambda x}(\lambda x)^{n-1}/(n-1)!$ for $x > 0$). Since MGFs uniquely determine distributions, $S \sim \mathrm{Gamma}(n, \lambda)$. 
}

\begin{exercise}
\label{ex:random-variables-65}
Let $Z_1,\ldots,Z_n$ be independent standard normal random variables.
For fixed real coefficients $a_{rs}$ and constants $\mu_r$, define
$X_r=\mu_r+\sum_{s=1}^n a_{rs}Z_s$. Find
$\operatorname{Cov}(X_i,X_j)$ in terms of the coefficients.
\end{exercise}

\solution{
By the given representation, $X_i = \sum_{s=1}^n a_{is} Z_s + \mu_i$ and $X_j = \sum_{s=1}^n a_{js}Z_s + \mu_j$, where $Z_1, \ldots, Z_n$ are independent standard normal random variables. Since $\mathrm{Cov}(Z_r, Z_s) = \mathbf{1}[r=s]$:
\[
\mathrm{Cov}(X_i, X_j) = \mathrm{Cov}\!\left(\sum_{r=1}^n a_{ir}Z_r,\; \sum_{s=1}^n a_{js}Z_s\right) = \sum_{r=1}^n\sum_{s=1}^n a_{ir}\,a_{js}\,\mathrm{Cov}(Z_r, Z_s) = \sum_{r=1}^n a_{ir}\,a_{jr}.
\]
}

\begin{exercise}
\label{ex:random-variables-66}
Use Chebyshev's inequality to prove the weak law of large numbers: if $X_1, X_2, \ldots$ are i.i.d.\ with mean $\mu$ and variance $\sigma^2$, then for any $\varepsilon > 0$,
\[
P\!\left\{\left|\frac{X_1+\cdots+X_n}{n}-\mu\right| > \varepsilon\right\} \to 0 \quad \text{as } n \to \infty.
\]
\end{exercise}

\solution{
Let $\bar{X}_n = (X_1+\cdots+X_n)/n$. Then $E[\bar{X}_n] = \mu$ and $\mathrm{Var}(\bar{X}_n) = \sigma^2/n$. By Chebyshev's inequality:
\[
P\!\left\{\left|\bar{X}_n - \mu\right| > \varepsilon\right\} \leq \frac{\mathrm{Var}(\bar{X}_n)}{\varepsilon^2} = \frac{\sigma^2}{n\varepsilon^2} \to 0 \quad \text{as } n \to \infty. \qquad
\]
}

\begin{exercise}
\label{ex:random-variables-67}
Suppose that $X$ is a random variable with mean $10$ and variance $15$. What can we say about $P\{5 < X < 15\}$?
\end{exercise}

\solution{
By Chebyshev's inequality, with $E[X] = \mu = 10$ and $\mathrm{Var}(X) = \sigma^2 = 15$:
\[
P\{|X - 10| \geq 5\} \leq \frac{\mathrm{Var}(X)}{5^2} = \frac{15}{25} = \frac{3}{5}.
\]
Therefore:
\[
P\{5 < X < 15\} = P\{|X - 10| < 5\} \geq 1 - \frac{3}{5} = \frac{2}{5}.
\]
}

\begin{exercise}
\label{ex:random-variables-68}
Let $X_1, X_2, \ldots, X_{10}$ be independent Poisson random variables with mean $1$.
\begin{enumerate}
    \item[(a)] Use the Markov inequality to obtain a bound on $P\{X_1+\cdots+X_{10} \geq 15\}$.
    \item[(b)] Use the central limit theorem to approximate $P\{X_1+\cdots+X_{10} \geq 15\}$.
\end{enumerate}
\end{exercise}

\solution{
Put $S=\sum_{i=1}^{10}X_i$. Then $S\sim\operatorname{Poisson}(10)$,
with mean and variance both $10$.
\textbf{(a)} Markov's inequality gives $P(S\ge15)\le10/15=2/3$.

\textbf{(b)} Approximate $S$ by a normal random variable with mean
$10$ and variance $10$. Without a continuity correction,
\[
P(S\ge15)\approx1-\Phi\left(\frac{15-10}{\sqrt{10}}\right)
\approx0.0569.
\]
Using the continuity correction instead gives
\[
P(S\ge15)\approx1-\Phi\left(\frac{14.5-10}{\sqrt{10}}\right)
\approx0.0774.
\]
These are normal approximations; the exact Poisson tail is about $0.0835$.
}

\begin{exercise}
\label{ex:random-variables-69}
If $X$ is normally distributed with mean $1$ and variance $4$, use the tables to find $P\{2 < X < 3\}$.
\end{exercise}

\solution{
With $Z=(X-1)/2\sim N(0,1)$,
\[
P(2<X<3)=P(0.5<Z<1)=\Phi(1)-\Phi(0.5)
\approx0.8413-0.6915=0.1498.
\]
}

\begin{exercise}
\label{ex:random-variables-70}
Show that
\[
\lim_{n\to\infty} e^{-n}\sum_{k=0}^n \frac{n^k}{k!} = \frac{1}{2}.
\]
\end{exercise}

\solution{
Let $X_n \sim \mathrm{Poisson}(n)$. Then:
\[
P\{X_n \leq n\} = e^{-n}\sum_{k=0}^n \frac{n^k}{k!}.
\]
By the central limit theorem, $\frac{X_n - n}{\sqrt{n}} \xrightarrow{d} Z \sim N(0,1)$, so:
\[
P\{X_n \leq n\} = P\!\left\{\frac{X_n - n}{\sqrt{n}} \leq 0\right\} \to P\{Z \leq 0\} = \Phi(0) = \frac{1}{2}. \qquad
\]
}

\begin{exercise}
\label{ex:random-variables-71}
Let $X$ denote the number of white balls selected when $k$ balls are chosen at random from an urn containing $n$ white and $m$ black balls.
\begin{enumerate}
    \item[(a)] Compute $P\{X = i\}$.
    \item[(b)] Let $X_i = \mathbf{1}[\text{$i$-th ball selected is white}]$ and $Y_j = \mathbf{1}[\text{white ball $j$ is selected}]$. Compute $E[X]$ using first $X_i$'s then $Y_j$'s.
\end{enumerate}
\end{exercise}

\solution{
\textbf{(a)} This is the hypergeometric distribution:
\[
P\{X = i\} = \frac{\binom{n}{i}\binom{m}{k-i}}{\binom{n+m}{k}}, \quad \max(0,k-m) \leq i \leq \min(n,k).
\]

\textbf{(b) Via $X_i$'s:} $X = \sum_{i=1}^k X_i$ and $E[X_i] = P(\text{$i$-th selected ball is white}) = \frac{n}{n+m}$ (by symmetry). So:
\[
E[X] = k\cdot\frac{n}{n+m}.
\]
\textbf{Via $Y_j$'s:} $X = \sum_{j=1}^n Y_j$ and $E[Y_j] = P(\text{white ball $j$ is among the $k$ chosen}) = \frac{k}{n+m}$. So:
\[
E[X] = n\cdot\frac{k}{n+m} = \frac{kn}{n+m}.
\]
Both methods give the same result.
}

\begin{exercise}
\label{ex:random-variables-72}
Show that $\mathrm{Var}(X) = 1$ when $X$ is the number of men who select their own hats (as in the hat-matching problem with $n$ men).
\end{exercise}

\solution{
Assume $n\ge2$; for $n=1$ the variance is $0$. Let $X_i = \mathbf{1}[\text{man $i$ gets his own hat}]$, so $X = \sum_{i=1}^n X_i$. By symmetry, $E[X_i] = 1/n$ and $E[X_iX_j] = P(\text{both $i$ and $j$ get own hats}) = 1/[n(n-1)]$ for $i \neq j$. Then:
\[
\mathrm{Var}(X_i) = \frac{1}{n} - \frac{1}{n^2} = \frac{n-1}{n^2}, \qquad \mathrm{Cov}(X_i,X_j) = \frac{1}{n(n-1)}-\frac{1}{n^2} = \frac{1}{n^2(n-1)}.
\]
Therefore:
\[
\mathrm{Var}(X) = n\cdot\frac{n-1}{n^2} + n(n-1)\cdot\frac{1}{n^2(n-1)} = \frac{n-1}{n}+\frac{1}{n} = 1. \qquad
\]
}

\begin{exercise}
\label{ex:random-variables-73}
For the multinomial distribution (Exercise 17), let $N_i$ denote the number of times outcome $i$ occurs. Find:
\begin{enumerate}
    \item[(a)] $E[N_i]$;
    \item[(b)] $\mathrm{Var}(N_i)$;
    \item[(c)] $\mathrm{Cov}(N_i, N_j)$ for $i \neq j$;
    \item[(d)] the expected number of outcomes that do not occur.
\end{enumerate}
\end{exercise}

\solution{
Each of the $n$ experiments independently produces outcome $i$ with probability $p_i$, so $N_i \sim \mathrm{Bin}(n, p_i)$.

\textbf{(a)} $E[N_i] = np_i$.

\textbf{(b)} $\mathrm{Var}(N_i) = np_i(1-p_i)$.

\textbf{(c)} $N_i + N_j \sim \mathrm{Bin}(n, p_i+p_j)$, so $\mathrm{Var}(N_i+N_j) = n(p_i+p_j)(1-p_i-p_j)$. Since $\mathrm{Var}(N_i+N_j) = \mathrm{Var}(N_i)+\mathrm{Var}(N_j)+2\mathrm{Cov}(N_i,N_j)$:
\[
\mathrm{Cov}(N_i,N_j) = \frac{n(p_i+p_j)(1-p_i-p_j)-np_i(1-p_i)-np_j(1-p_j)}{2} = -np_ip_j.
\]

\textbf{(d)} Let $A_i = \mathbf{1}[N_i = 0]$ for $i=1,\ldots,r$. Then $P(A_i=1) = (1-p_i)^n$ and:
\[
E[\text{outcomes that don't occur}] = \sum_{i=1}^r (1-p_i)^n.
\]
}

\section{Conditioning and Expectation}
\label{sec:conditioning-expectation}

\setcounter{exercise}{0}

\begin{exercise}
\label{ex:conditioning-expectation-1}
If $X$ and $Y$ are both discrete, show that $\sum_x p_{X|Y}(x|y) = 1$ for all $y$ such that $p_Y(y) > 0$.
\end{exercise}

\solution{
By definition, the conditional probability mass function of $X$ given $Y=y$ is:
\[
p_{X|Y}(x|y) = \frac{p(x,y)}{p_Y(y)}
\]
Summing this over all possible values of $x$ gives:
\[
\sum_x p_{X|Y}(x|y) = \sum_x \frac{p(x,y)}{p_Y(y)} = \frac{1}{p_Y(y)} \sum_x p(x,y)
\]
Since $\sum_x p(x,y) = p_Y(y)$ by the definition of marginal probability, we have:
\[
\sum_x p_{X|Y}(x|y) = \frac{p_Y(y)}{p_Y(y)} = 1
\]
}

\begin{exercise}
\label{ex:conditioning-expectation-2}
Let $X_1$ and $X_2$ be independent geometric random variables having the same parameter $p$. Guess the value of $P\{X_1 = i \mid X_1 + X_2 = n\}$.

\textbf{Hint:} Suppose a coin having probability $p$ of coming up heads is continually flipped. If the second head occurs on flip number $n$, what is the conditional probability that the first head was on flip number $i$, $i = 1, \ldots, n-1$? Verify your guess analytically.
\end{exercise}

\solution{
\textbf{Guess:} Given that the second head occurs exactly at flip $n$, the first head must have occurred equally likely at any of the preceding $n-1$ flips. Thus, we guess the probability is $\frac{1}{n-1}$ for $i = 1, \dots, n-1$.

\textbf{Analytical Verification:}
Using the definition of conditional probability:
\[
P\{X_1 = i \mid X_1 + X_2 = n\} = \frac{P\{X_1 = i, X_2 = n-i\}}{P\{X_1 + X_2 = n\}}
\]
Since $X_1$ and $X_2$ are independent geometric random variables:
\[
P\{X_1 = i\} = (1-p)^{i-1}p \quad \text{and} \quad P\{X_2 = n-i\} = (1-p)^{n-i-1}p
\]
The numerator becomes:
\[
P\{X_1 = i\} P\{X_2 = n-i\} = (1-p)^{i-1}p \cdot (1-p)^{n-i-1}p = (1-p)^{n-2}p^2
\]
The denominator is the sum over all possible values of $i$ from $1$ to $n-1$:
\[
P\{X_1 + X_2 = n\} = \sum_{j=1}^{n-1} P\{X_1 = j, X_2 = n-j\} = \sum_{j=1}^{n-1} (1-p)^{n-2}p^2 = (n-1)(1-p)^{n-2}p^2
\]
Dividing the numerator by the denominator:
\[
P\{X_1 = i \mid X_1 + X_2 = n\} = \frac{(1-p)^{n-2}p^2}{(n-1)(1-p)^{n-2}p^2} = \frac{1}{n-1}
\]
for $i = 1, \ldots, n-1$.
}

\begin{exercise}
\label{ex:conditioning-expectation-3}
The joint probability mass function of $X$ and $Y$, $p(x,y)$, is given by
\begin{align*}
    p(1,1) &= \frac{1}{9}, & p(2,1) &= \frac{1}{3}, & p(3,1) &= \frac{1}{9}, \\
    p(1,2) &= \frac{1}{9}, & p(2,2) &= 0,           & p(3,2) &= \frac{1}{18}, \\
    p(1,3) &= 0,           & p(2,3) &= \frac{1}{6}, & p(3,3) &= \frac{1}{9}
\end{align*}
Compute $E[X \mid Y = i]$ for $i = 1, 2, 3$.
\end{exercise}

\solution{
First, compute the marginal probabilities $p_Y(i) = \sum_x p(x,i)$:
\begin{itemize}
    \item $p_Y(1) = \frac{1}{9} + \frac{1}{3} + \frac{1}{9} = \frac{5}{9}$
    \item $p_Y(2) = \frac{1}{9} + 0 + \frac{1}{18} = \frac{3}{18} = \frac{1}{6}$
    \item $p_Y(3) = 0 + \frac{1}{6} + \frac{1}{9} = \frac{5}{18}$
\end{itemize}

Now compute $E[X \mid Y = i] = \sum_x x \frac{p(x,i)}{p_Y(i)}$:
\begin{itemize}
    \item For $i = 1$:
    \[
    E[X \mid Y = 1] = 1\left(\frac{1/9}{5/9}\right) + 2\left(\frac{1/3}{5/9}\right) + 3\left(\frac{1/9}{5/9}\right) = 1\left(\frac{1}{5}\right) + 2\left(\frac{3}{5}\right) + 3\left(\frac{1}{5}\right) = \frac{10}{5} = 2
    \]
    \item For $i = 2$:
    \[
    E[X \mid Y = 2] = 1\left(\frac{1/9}{1/6}\right) + 2\left(\frac{0}{1/6}\right) + 3\left(\frac{1/18}{1/6}\right) = 1\left(\frac{2}{3}\right) + 0 + 3\left(\frac{1}{3}\right) = \frac{5}{3}
    \]
    \item For $i = 3$:
    \[
    E[X \mid Y = 3] = 1\left(\frac{0}{5/18}\right) + 2\left(\frac{1/6}{5/18}\right) + 3\left(\frac{1/9}{5/18}\right) = 0 + 2\left(\frac{3}{5}\right) + 3\left(\frac{2}{5}\right) = \frac{12}{5}
    \]
\end{itemize}
}

\begin{exercise}
\label{ex:conditioning-expectation-4}
In Exercise 3, are the random variables $X$ and $Y$ independent?
\end{exercise}

\solution{
No, $X$ and $Y$ are not independent. For independence, we must have $p(x,y) = p_X(x)p_Y(y)$ for all $x, y$.

Let's check $x=1, y=3$:
From the table, $p(1,3) = 0$.
However, the marginal probabilities are:
\[
p_X(1) = p(1,1) + p(1,2) + p(1,3) = \frac{1}{9} + \frac{1}{9} + 0 = \frac{2}{9} > 0
\]
\[
p_Y(3) = \frac{5}{18} > 0
\]
Since $p(1,3) = 0 \neq p_X(1)p_Y(3) = \frac{10}{162}$, $X$ and $Y$ are dependent.
}

\begin{exercise}
\label{ex:conditioning-expectation-5}
An urn contains three white, six red, and five black balls. Six of these balls are randomly selected from the urn. Let $X$ and $Y$ denote respectively the number of white and black balls selected. Compute the conditional probability mass function of $X$ given that $Y = 3$. Also compute $E[X \mid Y = 1]$.
\end{exercise}

\solution{
The total number of balls is $3 + 6 + 5 = 14$.

Given $Y = 3$, we know exactly 3 black balls are selected out of the 5 available. The remaining $6 - 3 = 3$ choices must be selected from the remaining non-black balls ($3 \text{ white} + 6 \text{ red} = 9$ balls total). Thus, $X \mid Y=3$ follows a hypergeometric distribution:
\[
P\{X = x \mid Y = 3\} = \frac{\binom{3}{x}\binom{6}{3-x}}{\binom{9}{3}} \quad \text{for } x = 0, 1, 2, 3
\]

To compute $E[X \mid Y = 1]$, we note that given $Y = 1$, 1 black ball is chosen, so the remaining $6 - 1 = 5$ balls are chosen from the 9 non-black balls (3 white, 6 red). The distribution of $X \mid Y = 1$ is hypergeometric with sample size $n = 5$, from a population of $N = 9$ containing $K = 3$ white balls. The expected value is:
\[
E[X \mid Y = 1] = n \cdot \frac{K}{N} = 5 \cdot \frac{3}{9} = \frac{5}{3}
\]
}

\begin{exercise}
\label{ex:conditioning-expectation-6}
Repeat Exercise 5 but under the assumption that when a ball is selected its color is noted, and it is then replaced in the urn before the next selection is made.
\end{exercise}

\solution{
With replacement, the selections become independent trials. The probability of picking a white, red, or black ball on any trial is:
\[
P(W) = \frac{3}{14}, \quad P(R) = \frac{6}{14} = \frac{3}{7}, \quad P(B) = \frac{5}{14}
\]
Given that $Y = 3$ (exactly 3 black balls are chosen in 6 trials), the remaining $6 - 3 = 3$ trials must result in non-black balls. The conditional probability that a non-black ball is white is:
\[
P(W \mid \text{not } B) = \frac{P(W)}{P(W) + P(R)} = \frac{3/14}{9/14} = \frac{1}{3}
\]
Thus, given $Y=3$, the distribution of $X$ is Binomial$(3, 1/3)$:
\[
P\{X = x \mid Y = 3\} = \binom{3}{x} \left(\frac{1}{3}\right)^x \left(\frac{2}{3}\right)^{3-x} \quad \text{for } x = 0, 1, 2, 3
\]

For $E[X \mid Y = 1]$, given that 1 black ball is selected, the remaining $6 - 1 = 5$ trials must be non-black balls. Thus, $X \mid Y = 1 \sim \text{Binomial}(5, 1/3)$, yielding:
\[
E[X \mid Y = 1] = 5 \cdot \frac{1}{3} = \frac{5}{3}
\]
}

\begin{exercise}
\label{ex:conditioning-expectation-7}
Suppose $p(x,y,z)$, the joint probability mass function of the random variables $X$, $Y$, and $Z$, is given by
\begin{align*}
    p(1,1,1) &= \frac{1}{8},  & p(2,1,1) &= \frac{1}{4}, \\
    p(1,1,2) &= \frac{1}{8},  & p(2,1,2) &= \frac{3}{16}, \\
    p(1,2,1) &= \frac{1}{16}, & p(2,2,1) &= 0, \\
    p(1,2,2) &= 0,           & p(2,2,2) &= \frac{1}{4}
\end{align*}
What is $E[X \mid Y = 2]$? What is $E[X \mid Y = 2, Z = 1]$?
\end{exercise}

\solution{
First, to find $E[X \mid Y = 2]$, we calculate $P\{Y = 2\}$:
\[
P\{Y = 2\} = p(1,2,1) + p(1,2,2) + p(2,2,1) + p(2,2,2) = \frac{1}{16} + 0 + 0 + \frac{1}{4} = \frac{5}{16}
\]
Now compute the conditional probabilities for $X$ given $Y=2$:
\[
P\{X = 1 \mid Y = 2\} = \frac{p(1,2,1) + p(1,2,2)}{P\{Y = 2\}} = \frac{1/16 + 0}{5/16} = \frac{1}{5}
\]
\[
P\{X = 2 \mid Y = 2\} = \frac{p(2,2,1) + p(2,2,2)}{P\{Y = 2\}} = \frac{0 + 1/4}{5/16} = \frac{4}{5}
\]
Thus:
\[
E[X \mid Y = 2] = 1\left(\frac{1}{5}\right) + 2\left(\frac{4}{5}\right) = \frac{9}{5}
\]

Next, to find $E[X \mid Y = 2, Z = 1]$, we check the entries where $Y=2$ and $Z=1$:
\[
P\{Y=2, Z=1\} = p(1,2,1) + p(2,2,1) = \frac{1}{16} + 0 = \frac{1}{16}
\]
The conditional probabilities are:
\[
P\{X = 1 \mid Y = 2, Z = 1\} = \frac{1/16}{1/16} = 1, \quad P\{X = 2 \mid Y = 2, Z = 1\} = 0
\]
Thus:
\[
E[X \mid Y = 2, Z = 1] = 1(1) + 2(0) = 1
\]
}

\begin{exercise}
\label{ex:conditioning-expectation-8}
An unbiased die is successively rolled. Let $X$ and $Y$ denote, respectively, the number of rolls necessary to obtain a six and a five. Find (a) $E[X]$, (b) $E[X \mid Y = 1]$, (c) $E[X \mid Y = 5]$.
\end{exercise}

\solution{
\textbf{(a)} $X$ is a geometric random variable with success probability $p = 1/6$. Therefore:
\[
E[X] = \frac{1}{p} = 6
\]

\textbf{(b)} Given $Y = 1$, the first roll is a 5. Since it cannot be a 6, by the memoryless property of successive independent rolls, we start looking for a 6 from the second roll onwards. Thus:
\[
E[X \mid Y = 1] = 1 + E[X] = 1 + 6 = 7
\]

\textbf{(c)} Given $Y = 5$, the first four rolls are not 5, and the fifth roll is 5. We partition based on whether a 6 occurred in the first four rolls:
\begin{itemize}
    \item The probability that a 6 occurs on roll $x$ (for $x \le 4$) given that no 5 occurs in the first 4 trials is $\frac{1}{5}\left(\frac{4}{5}\right)^{x-1}$.
    \item If no 6 occurs in the first 4 rolls (which happens with probability $\left(\frac{4}{5}\right)^4$), the fifth roll is 5 (not 6), so we must wait an additional expected 6 rolls. The total rolls will be $5 + 6 = 11$.
\end{itemize}
Using the properties of the geometric series:
\[
E[X \mid Y = 5] = \sum_{x=1}^{4} x \frac{1}{5}\left(\frac{4}{5}\right)^{x-1} + 11 \left(\frac{4}{5}\right)^4 = 5 + 2\left(\frac{4}{5}\right)^4 = 5 + 2\left(\frac{256}{625}\right) = \frac{3637}{625} = 5.8192
\]
}

\begin{exercise}
\label{ex:conditioning-expectation-9}
Show in the discrete case that if $X$ and $Y$ are independent, then $E[X \mid Y = y] = E[X]$ for all $y$.
\end{exercise}

\solution{
For $p_Y(y)>0$ and integrable $X$, by definition of conditional expectation:
\[
E[X \mid Y = y] = \sum_x x p_{X|Y}(x|y)
\]
Since $X$ and $Y$ are independent, $p_{X|Y}(x|y) = p_X(x)$ for all $x$ and all $y$ with $p_Y(y)>0$. Substituting this into the sum yields:
\[
E[X \mid Y = y] = \sum_x x p_X(x) = E[X]
\]
}

\begin{exercise}
\label{ex:conditioning-expectation-10}
Suppose $X$ and $Y$ are independent continuous random variables. Show that $E[X \mid Y = y] = E[X]$ for all $y$.
\end{exercise}

\solution{
Where $f_Y(y)>0$ (and almost everywhere for conditional expectations), for integrable $X$:
\[
E[X \mid Y = y] = \int_{-\infty}^{\infty} x f_{X|Y}(x|y) \, dx
\]
Since $X$ and $Y$ are independent, the conditional density satisfies $f_{X|Y}(x|y) = f_X(x)$. Thus:
\[
E[X \mid Y = y] = \int_{-\infty}^{\infty} x f_X(x) \, dx = E[X]
\]
}

\begin{exercise}
\label{ex:conditioning-expectation-11}
The joint density of $X$ and $Y$ is
\[
f(x,y) = \frac{y^2 - x^2}{8} e^{-y}, \quad 0 < y < \infty, \quad -y \le x \le y
\]
Show that $E[X \mid Y = y] = 0$.
\end{exercise}

\solution{
The conditional density is given by $f_{X|Y}(x|y) = \frac{f(x,y)}{f_Y(y)}$. Note that for any fixed $y$, the joint density function $f(x,y)$ is an \textbf{even} function with respect to $x$, because:
\[
f(-x,y) = \frac{y^2 - (-x)^2}{8} e^{-y} = \frac{y^2 - x^2}{8} e^{-y} = f(x,y)
\]
Since the domain of $x$ given $Y=y$ is $[-y, y]$, which is symmetric about 0, $f_{X|Y}(x|y)$ is also an even function of $x$.

When computing the conditional expectation:
\[
E[X \mid Y = y] = \int_{-y}^{y} x f_{X|Y}(x|y) \, dx
\]
The integrand $x f_{X|Y}(x|y)$ is the product of an odd function ($x$) and an even function ($f_{X|Y}$), making the entire integrand an \textbf{odd} function. The integral of an odd function over a symmetric interval $[-y, y]$ is identically $0$. Thus, $E[X \mid Y = y] = 0$.
}

\begin{exercise}
\label{ex:conditioning-expectation-12}
The joint density of $X$ and $Y$ is given by
\[
f(x,y) = \frac{e^{-x/y}e^{-y}}{y}, \quad 0 < x < \infty, \quad 0 < y < \infty
\]
Show $E[X \mid Y = y] = y$.
\end{exercise}

\solution{
First, compute the marginal density $f_Y(y)$ by integrating over $x$:
\[
f_Y(y) = \int_{0}^{\infty} \frac{e^{-x/y}e^{-y}}{y} \, dx = \frac{e^{-y}}{y} \int_{0}^{\infty} e^{-x/y} \, dx = \frac{e^{-y}}{y} \left[ -y e^{-x/y} \right]_{0}^{\infty} = \frac{e^{-y}}{y} (0 - (-y)) = e^{-y}
\]
Next, find the conditional density $f_{X|Y}(x|y)$:
\[
f_{X|Y}(x|y) = \frac{f(x,y)}{f_Y(y)} = \frac{\frac{e^{-x/y}e^{-y}}{y}}{e^{-y}} = \frac{1}{y} e^{-x/y}, \quad x > 0
\]
This is the probability density function of an exponential random variable with mean parameters $\mu = y$. Therefore:
\[
E[X \mid Y = y] = \int_{0}^{\infty} x \left(\frac{1}{y} e^{-x/y}\right) \, dx = y
\]
}

\begin{exercise}
\label{ex:conditioning-expectation-13}
Let $X$ be exponential with mean $1/\lambda$; that is,
\[
f_X(x) = \lambda e^{-\lambda x}, \quad 0 < x < \infty
\]
Find $E[X \mid X > 1]$.
\end{exercise}

\solution{
By the memoryless property of the exponential distribution, given that $X > 1$, the excess value $X - 1$ follows the exact same exponential distribution as $X$ itself. That is:
\[
E[X - 1 \mid X > 1] = E[X] = \frac{1}{\lambda}
\]
Using the linearity of expectation:
\[
E[X \mid X > 1] - 1 = \frac{1}{\lambda} \implies E[X \mid X > 1] = 1 + \frac{1}{\lambda}
\]
}

\begin{exercise}
\label{ex:conditioning-expectation-14}
Let $X$ be uniform over $(0, 1)$. Find $E\left[X \mid X < \frac{1}{2}\right]$.
\end{exercise}

\solution{
Given that $X < \frac{1}{2}$, the random variable $X$ is restricted to the sub-interval $\left(0, \frac{1}{2}\right)$ and remains uniformly distributed over this new range.

The expected value of a uniform continuous distribution over an interval $(a, b)$ is simply its midpoint $\frac{a+b}{2}$. Setting $a = 0$ and $b = \frac{1}{2}$:
\[
E\left[X \mid X < \frac{1}{2}\right] = \frac{0 + 1/2}{2} = \frac{1}{4}
\]
}

\begin{exercise}
\label{ex:conditioning-expectation-15}
The joint density of $X$ and $Y$ is given by
\[
f(x,y) = \frac{e^{-y}}{y}, \quad 0 < x < y, \quad 0 < y < \infty
\]
Compute $E[X^2 \mid Y = y]$.
\end{exercise}

\solution{
First, compute the marginal density $f_Y(y)$ by integrating over the allowed range of $x$, which is $(0, y)$:
\[
f_Y(y) = \int_{0}^{y} \frac{e^{-y}}{y} \, dx = \frac{e^{-y}}{y} [x]_{0}^{y} = \frac{e^{-y}}{y} \cdot y = e^{-y}
\]
Now find the conditional density $f_{X|Y}(x|y)$:
\[
f_{X|Y}(x|y) = \frac{f(x,y)}{f_Y(y)} = \frac{\frac{e^{-y}}{y}}{e^{-y}} = \frac{1}{y}, \quad 0 < x < y
\]
This shows that given $Y = y$, $X$ is uniformly distributed over the interval $(0, y)$.

We can now compute the conditional expectation of $X^2$:
\[
E[X^2 \mid Y = y] = \int_{0}^{y} x^2 f_{X|Y}(x|y) \, dx = \int_{0}^{y} x^2 \frac{1}{y} \, dx = \frac{1}{y} \left[ \frac{x^3}{3} \right]_{0}^{y} = \frac{y^3}{3y} = \frac{y^2}{3}
\]
}

\setcounter{exercise}{18}
\begin{exercise}
\label{ex:conditioning-expectation-19}
Prove that if $X$ and $Y$ are jointly continuous, then
\[
E[X] = \int_{-\infty}^{\infty} E[X \mid Y = y] f_Y(y) \, dy
\]
\end{exercise}

\solution{
By definition, the conditional expectation of $X$ given $Y=y$ is:
\[
E[X \mid Y = y] = \int_{-\infty}^{\infty} x f_{X|Y}(x|y) \, dx
\]
Substituting this definition into the integral on the right-hand side gives:
\[
\int_{-\infty}^{\infty} E[X \mid Y = y] f_Y(y) \, dy = \int_{-\infty}^{\infty} \left( \int_{-\infty}^{\infty} x f_{X|Y}(x|y) \, dx \right) f_Y(y) \, dy
\]
Since $f_{X|Y}(x|y) f_Y(y) = f(x,y)$, where $f(x,y)$ is the joint probability density function of $X$ and $Y$, we can simplify the inner expression:
\[
= \int_{-\infty}^{\infty} \int_{-\infty}^{\infty} x f(x,y) \, dx \, dy
\]
Assuming $E[|X|]<\infty$, Fubini's theorem allows us to interchange the order of integration:
\[
= \int_{-\infty}^{\infty} x \left( \int_{-\infty}^{\infty} f(x,y) \, dy \right) dx
\]
Since $\int_{-\infty}^{\infty} f(x,y) \, dy = f_X(x)$ represents the marginal density function of $X$, the equation becomes:
\[
= \int_{-\infty}^{\infty} x f_X(x) \, dx = E[X]
\]
Thus, the identity is verified.
}

\begin{exercise}
\label{ex:conditioning-expectation-20}
An individual whose level of exposure to a certain pathogen is $x$ will contract the disease caused by this pathogen with probability $P(x)$. If the exposure level of a randomly chosen member of the population has probability density function $f$, determine the conditional probability density of the exposure level of that member given that he or she
\begin{enumerate}
    \item[(a)] has the disease.
    \item[(b)] does not have the disease.
    \item[(c)] Show that when $P(x)$ increases in $x$, then the ratio of the density of part (a) to that of part (b) also increases in $x$.
\end{enumerate}
\end{exercise}

\solution{
Let $D$ denote the event that the individual contracts the disease. We are given the conditional probability $P(D \mid X = x) = P(x)$ and the marginal density of exposure $f_X(x) = f(x)$.

\textbf{(a)} By Bayes' theorem for continuous random variables, the conditional density of $X$ given that the individual has the disease is:
\[
f_{X|D}(x) = \frac{P(D \mid X = x) f(x)}{P(D)} = \frac{P(x) f(x)}{\int_{-\infty}^{\infty} P(t) f(t) \, dt}
\]

\textbf{(b)} The probability of not contracting the disease given exposure level $x$ is $P(D^c \mid X = x) = 1 - P(x)$. Thus, the conditional density given that the individual is healthy is:
\[
f_{X|D^c}(x) = \frac{P(D^c \mid X = x) f(x)}{P(D^c)} = \frac{(1 - P(x)) f(x)}{\int_{-\infty}^{\infty} (1 - P(t)) f(t) \, dt}
\]

\textbf{(c)} Let us find the ratio of the conditional density from part (a) to that of part (b):
\[
\frac{f_{X|D}(x)}{f_{X|D^c}(x)} = \frac{\frac{P(x) f(x)}{P(D)}}{\frac{(1 - P(x)) f(x)}{P(D^c)}} = \frac{P(D^c)}{P(D)} \cdot \frac{P(x)}{1 - P(x)}
\]
The fraction $\frac{P(D^c)}{P(D)}$ is a constant value with respect to $x$. Therefore, the variation of the ratio depends entirely on the function:
\[
g(x) = \frac{P(x)}{1 - P(x)}
\]
If $P(x)$ is an increasing function of $x$, then as $x$ increases, the numerator $P(x)$ increases and the denominator $1 - P(x)$ decreases. Consequently, $g(x)$ is nondecreasing when $P(x)$ is nondecreasing (and strictly increasing when $P(x)$ is strictly increasing), on points where the density ratio is defined. Assume $0<P(D)<1$ and restrict the ratio to the common positive support. This proves that the ratio of the densities also increases in $x$.
}

\begin{exercise}
\label{ex:conditioning-expectation-21}
A miner is trapped in a mine with three doors. Door 1 leads to safety in 2 hours; doors 2 and 3 return the miner to the starting point after 3 and 5 hours, respectively. At each return, the miner independently chooses each door with probability $1/3$. Let $N$ denote the total number of doors selected before the miner reaches safety. Also, let $T_i$ denote the travel time corresponding to the $i$th choice, $i \ge 1$. Let $X$ denote the time when the miner reaches safety.
\begin{enumerate}
    \item[(a)] Give an identity that relates $X$ to $N$ and the $T_i$.
    \item[(b)] What is $E[N]$?
    \item[(c)] What is $E[T_N]$?
    \item[(d)] What is $E[\sum_{i=1}^N T_i \mid N = n]$?
    \item[(e)] Using the preceding, what is $E[X]$?
\end{enumerate}
\end{exercise}

\solution{
\textit{Setup:} A miner is trapped with 3 doors. Door 1 leads to safety in 2 hours; Door 2 returns him to the mine after 3 hours; Door 3 returns him to the mine after 5 hours. Each door is chosen at random with equal probability $1/3$.

\textbf{(a)} The total time $X$ required to escape is the sum of the travel times for all $N$ doors chosen sequentially:
\[
X = \sum_{i=1}^{N} T_i
\]

\textbf{(b)} On each independent attempt, the probability of selecting the safety door is $p = 1/3$. The total number of choices $N$ follows a geometric distribution with success parameter $p$. Thus:
\[
E[N] = \frac{1}{p} = 3
\]

\textbf{(c)} The final choice ($N$-th choice) must be the door that leads directly to safety. Based on the problem description, the travel time for the safety door is 2 hours. Thus:
\[
E[T_N] = 2
\]

\textbf{(d)} Given that $N = n$, we know that the first $n-1$ choices were incorrect doors (Door 2 or Door 3, each chosen with equal probability $1/2$), and the final $n$-th choice was the safety door (Door 1).
The expected time for any of the first $n-1$ incorrect choices is $\frac{3 + 5}{2} = 4$ hours.
The time for the last choice is 2 hours. By linearity of expectation:
\[
E\left[\sum_{i=1}^N T_i \;\middle|\; N = n\right] = \sum_{i=1}^{n-1} E[T_i \mid N=n] + E[T_n \mid N=n] = 4(n - 1) + 2 = 4n - 2
\]

\textbf{(e)} We use the law of total expectation, substituting $N$ into the conditional identity found in part (d):
\[
E[X] = E\left[E\left[\sum_{i=1}^N T_i \;\middle|\; N\right]\right] = E[4N - 2]
\]
Applying linearity of expectation:
\[
E[X] = 4E[N] - 2 = 4(3) - 2 = 12 - 2 = 10 \text{ hours}
\]
}

\setcounter{exercise}{22}
\begin{exercise}
\label{ex:conditioning-expectation-23}
A coin having probability $p$ of coming up heads is successively flipped until two of the most recent three flips are heads. Let $N$ denote the number of flips. (Note that if the first two flips are heads, then $N = 2$.) Find $E[N]$.
\end{exercise}

\solution{
Let $q = 1-p$. We can find $E[N]$ by setting up a system of equations based on the states of the recent flips:
\begin{itemize}
    \item $E_S$: Expected additional flips from the start.
    \item $E_H$: Expected additional flips when the last flip was $H$ (and the preceding was not $H$).
    \item $E_T$: Expected additional flips when the last flip was $T$ (with no useful history).
    \item $E_{HT}$: Expected additional flips when the last two flips were $H, T$.
    \item $E_{TH}$: Expected additional flips when the last two flips were $T, H$.
\end{itemize}

From each state, the transitions give the following equations:
\[
E_S = 1 + p E_H + q E_T
\]
\[
E_H = 1 + q E_{HT} \quad (\text{since an } H \text{ leads to } HH \text{ and stops})
\]
\[
E_T = 1 + p E_{TH} + q E_T \implies p E_T = 1 + p E_{TH} \implies E_T = \frac{1}{p} + E_{TH}
\]
\[
E_{HT} = 1 + q E_T \quad (\text{since an } H \text{ leads to } HTH \text{ and stops})
\]
\[
E_{TH} = 1 + q E_{HT} \quad (\text{since an } H \text{ leads to } THH \text{ and stops})
\]

Notice that $E_H = E_{TH} = 1 + q E_{HT}$. Substituting this into the expression for $E_T$:
\[
E_T = \frac{1}{p} + E_H
\]
Now substitute $E_T$ into the expression for $E_{HT}$:
\[
E_{HT} = 1 + q\left(\frac{1}{p} + E_H\right) = 1 + \frac{q}{p} + q E_H
\]
Substitute $E_{HT}$ back into the equation for $E_H$:
\[
E_H = 1 + q\left(1 + \frac{q}{p} + q E_H\right) = 1 + q + \frac{q^2}{p} + q^2 E_H
\]
\[
E_H(1 - q^2) = 1 + q + \frac{q^2}{p} \implies E_H \cdot p(1+q) = \frac{p+pq+q^2}{p} = \frac{p+q(p+q)}{p} = \frac{p+q}{p} = \frac{1}{p}
\]
\[
E_H = \frac{1}{p^2(1+q)} = \frac{1}{p^2(2-p)}
\]

Finally, we compute $E_S$:
\[
E_S = 1 + p E_H + q \left(\frac{1}{p} + E_H\right) = 1 + \frac{q}{p} + (p+q)E_H = \frac{1}{p} + E_H
\]
\[
E[N] = E_S = \frac{1}{p} + \frac{1}{p^2(2-p)} = \frac{p(2-p) + 1}{p^2(2-p)} = \frac{1 + 2p - p^2}{p^2(2-p)}
\]
}

\begin{exercise}
\label{ex:conditioning-expectation-24}
A coin, having probability $p$ of landing heads, is continually flipped until at least one head and one tail have been flipped.
\begin{enumerate}
    \item[(a)] Find the expected number of flips needed.
    \item[(b)] Find the expected number of flips that land on heads.
    \item[(c)] Find the expected number of flips that land on tails.
    \item[(d)] Repeat part (a) in the case where flipping is continued until a total of at least two heads and one tail have been flipped.
\end{enumerate}
\end{exercise}

\solution{
Let $q = 1-p$.

\textbf{(a)} Condition on the outcome of the first flip:
\[
E[N] = 1 + p E[\text{flips until T}] + q E[\text{flips until H}] = 1 + p\left(\frac{1}{q}\right) + q\left(\frac{1}{p}\right) = 1 + \frac{p}{q} + \frac{q}{p} = \frac{1}{pq} - 1
\]

\textbf{(b)} Let $H_N$ be the number of heads. Conditioning on the first flip:
\begin{itemize}
    \item If the first flip is $H$, we flip until the first $T$. The number of additional flips is Geometric($q$), which contains an expected $\frac{1}{q}-1$ heads. Thus, total expected heads is $1 + \frac{1}{q} - 1 = \frac{1}{q}$.
    \item If the first flip is $T$, we flip until the first $H$, which yields exactly $1$ head.
\end{itemize}
\[
E[H_N] = p\left(\frac{1}{q}\right) + q(1) = \frac{p + q^2}{q} = \frac{p + (1-p)^2}{q} = \frac{1 - p + p^2}{q} = 1 + \frac{p^2}{q}
\]

\textbf{(c)} By symmetric reasoning, swapping $p$ and $q$:
\[
E[T_N] = q\left(\frac{1}{p}\right) + p(1) = 1 + \frac{q^2}{p}
\]

\textbf{(d)} Let $N_{(2H,1T)}$ be the number of flips required. Condition on the first flip:
\begin{itemize}
    \item If the first is $T$ (prob $q$), we need $2$ more heads, requiring an expected $\frac{2}{p}$ additional flips.
    \item If the first is $H$ (prob $p$), we condition on the second flip:
    \begin{itemize}
        \item Second is $H$ (prob $p$): we have $2H$, so we need $1T$, requiring $\frac{1}{q}$ additional flips.
        \item Second is $T$ (prob $q$): we have $1H$ and $1T$, so we need $1H$, requiring $\frac{1}{p}$ additional flips.
    \end{itemize}
\end{itemize}
Combining these cases:
\[
E[N_{(2H,1T)}] = q\left(1 + \frac{2}{p}\right) + p\left(1 + p\left(1 + \frac{1}{q}\right) + q\left(1 + \frac{1}{p}\right)\right) = 2 + \frac{2q}{p} + \frac{p^2}{q}
\]
}

\begin{exercise}
\label{ex:conditioning-expectation-25}
Independent trials, resulting in one of the outcomes 1, 2, 3 with respective probabilities $p_1, p_2, p_3$, $\sum_{i=1}^3 p_i = 1$, are performed.
\begin{enumerate}
    \item[(a)] Let $N$ denote the number of trials needed until the initial outcome has occurred exactly 3 times. Find $E[N]$.
    \item[(b)] Find the expected number of trials needed until both outcome 1 and outcome 2 have occurred.
\end{enumerate}
\end{exercise}

\solution{
\textbf{(a)} Condition on the first outcome observed:
\begin{itemize}
    \item If the first outcome is $i$, then outcome $i$ is designated as the initial outcome. We require $2$ more occurrences of outcome $i$. The number of additional trials needed follows a Negative Binomial distribution with expected value $\frac{2}{p_i}$.
\end{itemize}
By the law of total expectation:
\[
E[N] = \sum_{i=1}^3 p_i \left(1 + \frac{2}{p_i}\right) = \sum_{i=1}^3 p_i + \sum_{i=1}^3 2 = 1 + 6 = 7
\]

\textbf{(b)} Let $M$ be the number of trials needed. Condition on the first trial:
\begin{itemize}
    \item First outcome is 1 (prob $p_1$): we only need outcome 2, which takes an expected $\frac{1}{p_2}$ additional steps.
    \item First outcome is 2 (prob $p_2$): we only need outcome 1, which takes an expected $\frac{1}{p_1}$ additional steps.
    \item First outcome is 3 (prob $p_3$): no progress is made towards either 1 or 2, resulting in a reset.
\end{itemize}
\[
E[M] = p_1\left(1 + \frac{1}{p_2}\right) + p_2\left(1 + \frac{1}{p_1}\right) + p_3(1 + E[M])
\]
\[
E[M](1 - p_3) = (p_1 + p_2 + p_3) + \frac{p_1}{p_2} + \frac{p_2}{p_1} \implies E[M](p_1 + p_2) = 1 + \frac{p_1^2 + p_2^2}{p_1 p_2}
\]
\[
E[M] = \frac{1}{p_1 + p_2} + \frac{p_1^2 + p_2^2}{p_1 p_2 (p_1 + p_2)} = \frac{p_1 p_2 + p_1^2 + p_2^2}{p_1 p_2 (p_1 + p_2)}
\]
}

\begin{exercise}
\label{ex:conditioning-expectation-26}
You have two opponents with whom you alternate play. Whenever you play $A$, you win with probability $p_A$; whenever you play $B$, you win with probability $p_B$, where $p_B > p_A$. If your objective is to minimize the expected number of games you need to play to win two in a row, should you start with $A$ or with $B$?
\end{exercise}

\solution{
Assume independent games and $0<p_A<p_B\le1$; if $p_A=0$, two consecutive wins are impossible from either starting order. Let $q_A = 1 - p_A$ and $q_B = 1 - p_B$. Following the hint, we write the conditional equations for $E[N_A]$ and $E[N_B]$:
\[
E[N_A] = p_A [ p_B \cdot 2 + q_B(2 + E[N_A]) ] + q_A [ 1 + E[N_B] ] = 1 + p_A + p_A q_B E[N_A] + q_A E[N_B]
\]
By symmetry, swapping $A$ and $B$:
\[
E[N_B] = 1 + p_B + p_B q_A E[N_B] + q_B E[N_A]
\]
Subtracting the second equation from the first:
\[
E[N_A] - E[N_B] = p_A - p_B + p_A q_B E[N_A] - q_B E[N_A] + q_A E[N_B] - p_B q_A E[N_B]
\]
\[
E[N_A] - E[N_B] = p_A - p_B + E[N_A] q_B(p_A - 1) + E[N_B] q_A(1 - p_B)
\]
Since $p_A - 1 = -q_A$ and $1 - p_B = q_B$:
\[
E[N_A] - E[N_B] = p_A - p_B - q_A q_B (E[N_A] - E[N_B])
\]
\[
(E[N_A] - E[N_B])(1 + q_A q_B) = p_A - p_B \implies E[N_A] - E[N_B] = \frac{p_A - p_B}{1 + q_A q_B}
\]
Since we are given $p_B > p_A$, the numerator $p_A - p_B$ is negative, while the denominator is positive. Therefore, $E[N_A] - E[N_B] < 0$, meaning $E[N_A] < E[N_B]$. To minimize the expected number of games, you should \textbf{start with A}.
}

\begin{exercise}
\label{ex:conditioning-expectation-27}
A coin that comes up heads with probability $p$ is continually flipped until the pattern T, T, H appears. Let $X$ denote the number of flips made, and find $E[X]$.
\end{exercise}

\solution{
Let $q = 1-p$. Let $E_0$ be the expected remaining flips from the start, $E_T$ after one $T$, and $E_{TT}$ after $TT$.
\[
E_0 = 1 + p E_0 + q E_T \implies q E_0 = 1 + q E_T \implies E_0 = \frac{1}{q} + E_T
\]
\[
E_T = 1 + p E_0 + q E_{TT}
\]
\[
E_{TT} = p(1) + q(1 + E_{TT}) \implies p E_{TT} = 1 \implies E_{TT} = \frac{1}{p}
\]
Substituting $E_{TT}$ into $E_T$:
\[
E_T = 1 + p E_0 + \frac{q}{p}
\]
Substituting $E_T$ into the simplified expression for $E_0$:
\[
E_0 = \frac{1}{q} + 1 + p E_0 + \frac{q}{p} \implies E_0(1-p) = \frac{1}{q} + 1 + \frac{q}{p}
\]
Since $1-p = q$:
\[
q E_0 = \frac{1}{q} + 1 + \frac{q}{p} \implies E[X] = E_0 = \frac{1}{q^2} + \frac{1}{q} + \frac{1}{p} = \frac{1}{(1-p)^2} + \frac{1}{1-p} + \frac{1}{p}
\]
}

\begin{exercise}
\label{ex:conditioning-expectation-28}
Polya's urn model supposes that an urn initially contains $r$ red and $b$ blue balls. At each stage a ball is randomly selected from the urn and is then returned along with $m$ other balls of the same color. Let $X_k$ be the number of red balls drawn in the first $k$ selections.
Find (a) $E[X_1]$, (b) $E[X_2]$, (c) $E[X_3]$. (d) Conjecture the value of $E[X_k]$ and verify it by conditioning. (e) Give an intuitive proof.
\end{exercise}

\solution{
Let $I_j$ be the indicator variable that equals 1 if the $j$-th draw is red, and 0 otherwise.

\textbf{(a)} $E[X_1] = P(I_1 = 1) = \frac{r}{r+b}$.

\textbf{(b)} Conditioning on the first draw:
\[
P(I_2 = 1) = \frac{r+m}{r+b+m}\left(\frac{r}{r+b}\right) + \frac{r}{r+b+m}\left(\frac{b}{r+b}\right) = \frac{r(r+m+b)}{(r+b)(r+b+m)} = \frac{r}{r+b}
\]
Thus, $E[X_2] = E[I_1] + E[I_2] = \frac{2r}{r+b}$.

\textbf{(c)} By the same expansion, $P(I_3 = 1) = \frac{r}{r+b}$, hence $E[X_3] = \frac{3r}{r+b}$.

\textbf{(d) Conjecture:} $E[X_k] = \frac{kr}{r+b}$.
\textit{Proof by induction:} Assume $P(I_j = 1) = \frac{r}{r+b}$ holds for the first $k$ trials under any starting configuration. Conditioning on the first draw's outcome:
\[
P(I_{k+1} = 1) = P(I_{k+1} = 1 \mid I_1 = 1) P(I_1 = 1) + P(I_{k+1} = 1 \mid I_1 = 0) P(I_1 = 0)
\]
By the induction hypothesis, the remaining $k$ trials behave like a fresh Polya urn sequence starting with parameters $(r+m, b)$ or $(r, b+m)$:
\[
P(I_{k+1} = 1) = \left(\frac{r+m}{r+b+m}\right)\frac{r}{r+b} + \left(\frac{r}{r+b+m}\right)\frac{b}{r+b} = \frac{r}{r+b}
\]
By linearity of expectation, $E[X_k] = \sum_{j=1}^k E[I_j] = \frac{kr}{r+b}$.

\textbf{(e) Intuitive proof:} By indexing each initial ball uniquely, symmetry dictates that every initial ball has an equal probability of being the 'ancestor' of the lineage chosen at any arbitrary draw stage $j$. Since $\frac{r}{r+b}$ of the initial balls are red, the probability of selecting a red ball at any stage remains constant at $\frac{r}{r+b}$.
}

\begin{exercise}
\label{ex:conditioning-expectation-29}
Two players take turns shooting at a target, with each shot by player $i$ hitting the target with probability $p_i$, $i = 1, 2$. Shooting ends when two consecutive shots hit the target. Let $\mu_i$ denote the mean number of shots taken when player $i$ shoots first, $i = 1, 2$.
Find (a) $\mu_1$ and $\mu_2$, (b) $h_1$ and $h_2$ (the mean number of hits).
\end{exercise}

\solution{
Assume $p_1,p_2>0$ and independent shots.\par
\textbf{(a)} Conditioning on the first results:
\[
\mu_1 = (1-p_1)(1 + \mu_2) + p_1 [ p_2 \cdot 2 + (1-p_2)(2 + \mu_1) ] = 1 + p_1 + p_1(1-p_2)\mu_1 + (1-p_1)\mu_2
\]
By symmetry, swapping indices 1 and 2:
\[
\mu_2 = 1 + p_2 + p_2(1-p_1)\mu_2 + (1-p_2)\mu_1
\]
Solving this linear system yields:
\[
\mu_1 = \frac{2 - p_1 p_2(1-p_1)}{p_1 p_2 [1 + (1-p_1)(1-p_2)]}, \quad \mu_2 = \frac{2 - p_1 p_2(1-p_2)}{p_1 p_2 [1 + (1-p_1)(1-p_2)]}
\]

\textbf{(b)} Let $h_i$ be the expected number of hits. Conditioning similarly:
\[
h_1 = (1-p_1)h_2 + p_1 [ p_2 \cdot 2 + (1-p_2)(1 + h_1) ] = p_1(1+p_2) + p_1(1-p_2)h_1 + (1-p_1)h_2
\]
\[
h_2 = p_2(1+p_1) + p_2(1-p_1)h_2 + (1-p_2)h_1
\]
Solving this system yields:
\[
h_1 = \frac{p_1 + p_2 - p_1 p_2^2(1-p_1)}{p_1 p_2 [1 + (1-p_1)(1-p_2)]}, \quad h_2 = \frac{p_1 + p_2 - p_1^2 p_2(1-p_2)}{p_1 p_2 [1 + (1-p_1)(1-p_2)]}
\]
}

\begin{exercise}
\label{ex:conditioning-expectation-30}
Let $X_i, i \ge 0$ be independent and identically distributed random variables with probability mass function $p(j) = P\{X_i = j\}, j = 1, \dots, m$. Find $E[N]$, where $N = \min\{n > 0 : X_n = X_0\}$.
\end{exercise}

\solution{
Condition on $X_0=j$ only for values with $p(j)>0$. Then $N$ is
geometric with mean $1/p(j)$, so
\[
E[N]=\sum_{j:p(j)>0}p(j)\frac1{p(j)}
=\#\{j:p(j)>0\}.
\]
In particular, the answer is $m$ if all $m$ listed probabilities are
positive. Zero-probability outcomes must not be counted.
}

\begin{exercise}
\label{ex:conditioning-expectation-31}
Each element in a sequence of binary data is either 1 with probability $p$ or 0 with probability $1 - p$. A maximal subsequence of consecutive values having identical outcomes is called a run. Find (a) the expected length of the first run, (b) the expected length of the second run.
\end{exercise}

\solution{
Assume independent data bits and $0<p<1$.\par
\textbf{(a)} Let $L_1$ be the length of the first run. Condition on the value of the first element:
\begin{itemize}
    \item If it is 1 (prob $p$), the run continues until the first 0 appears, which follows a Geometric($1-p$) distribution.
    \item If it is 0 (prob $1-p$), the run continues until the first 1 appears, which follows a Geometric($p$) distribution.
\end{itemize}
\[
E[L_1] = p\left(\frac{1}{1-p}\right) + (1-p)\left(\frac{1}{p}\right) = \frac{p}{1-p} + \frac{1-p}{p} = \frac{p^2 + (1-p)^2}{p(1-p)}
\]

\textbf{(b)} Let $L_2$ be the length of the second run. The second run is guaranteed to be composed of the opposite binary element compared to the first run. Conditioning on the first element of the sequence:
\begin{itemize}
    \item If the first element is 1, the second run consists entirely of 0s, terminating at the next 1 $\implies$ Expected length is $\frac{1}{p}$.
    \item If the first element is 0, the second run consists entirely of 1s, terminating at the next 0 $\implies$ Expected length is $\frac{1}{1-p}$.
\end{itemize}
\[
E[L_2] = p\left(\frac{1}{p}\right) + (1-p)\left(\frac{1}{1-p}\right) = 1 + 1 = 2
\]
}

\begin{exercise}
\label{ex:conditioning-expectation-32}
Independent trials, each resulting in success with probability $p$, are performed.
\begin{enumerate}
    \item[(a)] Find the expected number of trials needed for there to have been both at least $n$ successes and at least $m$ failures.
    \item[(b)] Find the expected number of trials needed for there to have been either at least $n$ successes or at least $m$ failures.
\end{enumerate}
\end{exercise}

\solution{
\textbf{(a)} Let $X$ be the number of trials until there are both at least $n$ successes and at least $m$ failures. Following the hint, let us look at the first $n+m$ trials. Let $S_{n+m}$ be the number of successes in these $n+m$ trials.
\begin{itemize}
    \item If $S_{n+m} \ge n$ and $(n+m) - S_{n+m} \ge m$ (which means $S_{n+m} = n$), then at trial $n+m$, we have exactly achieved both conditions.
    \item If $S_{n+m} = k > n$, we reached the $n$-th success earlier, but we still need more failures to reach $m$ failures. By the memoryless property of independent trials, the additional number of trials needed to reach $m$ failures given we already have $(n+m)-k$ failures is the expected number of trials to get the remaining $m - (n+m-k) = k-n$ failures, which is $\frac{k-n}{1-p}$.
    \item If $S_{n+m} = k < n$, we reached $m$ failures earlier, but we need more successes. The remaining successes needed are $n-k$, which takes an expected $\frac{n-k}{p}$ trials.
\end{itemize}
By the law of total expectation:
\[
E[X] = (n+m) + \sum_{k=n+1}^{n+m} \binom{n+m}{k} p^k (1-p)^{n+m-k} \frac{k-n}{1-p} + \sum_{k=0}^{n-1} \binom{n+m}{k} p^k (1-p)^{n+m-k} \frac{n-k}{p}
\]

\textbf{(b)} Let $Y$ be the number of trials until there are either at least $n$ successes or at least $m$ failures. This is the stopping time of a standard negative binomial race (Banach's matchbox-type problem). Notice that the total number of trials to achieve \textit{both} conditions ($X$) plus the number of trials to achieve the \textit{either} condition ($Y$) is equal to the sum of the two completion times (which are dependent): the time to get $n$ successes ($X_S$) and the time to get $m$ failures ($X_F$). Therefore, $X + Y = X_S + X_F$, which implies:
\[
E[Y] = E[X_S] + E[X_F] - E[X] = \frac{n}{p} + \frac{m}{1-p} - E[X]
\]
where $E[X]$ is the expression derived in part (a).
}

\begin{exercise}
\label{ex:conditioning-expectation-33}
If $R_i$ denotes the random amount that is earned in period $i$, then $\sum_{i=1}^{\infty} \beta^{i-1}R_i$, where $0 < \beta < 1$ is a specified constant, is called the total discounted reward with discount factor $\beta$. Let $T$ be a geometric random variable with parameter $1 - \beta$ that is independent of the $R_i$. Show that the expected total discounted reward is equal to the expected total (undiscounted) reward earned by time $T$. That is, show that
\[
E\left[\sum_{i=1}^{\infty} \beta^{i-1}R_i\right] = E\left[\sum_{i=1}^{T} R_i\right]
\]
\end{exercise}

\solution{
Assume the rewards are nonnegative, or $\sum_i\beta^{i-1}E[|R_i|]<\infty$. These conditions justify the exchanges below by Tonelli or Fubini. We can express the sum on the right-hand side using indicator functions:
\[
\sum_{i=1}^{T} R_i = \sum_{i=1}^{\infty} R_i \mathbf{1}[T \ge i]
\]
Taking expectations on both sides and using the independence of $T$ and $R_i$:
\[
E\left[\sum_{i=1}^{T} R_i\right] = \sum_{i=1}^{\infty} E\left[R_i \mathbf{1}[T \ge i]\right] = \sum_{i=1}^{\infty} E[R_i] P\{T \ge i\}
\]
Since $T$ is a geometric random variable with parameter $1-\beta$, the probability that $T$ is at least $i$ is the probability of having $i-1$ initial failures:
\[
P\{T \ge i\} = \beta^{i-1}
\]
Substituting this back into the sum yields:
\[
E\left[\sum_{i=1}^{T} R_i\right] = \sum_{i=1}^{\infty} E[R_i] \beta^{i-1} = E\left[\sum_{i=1}^{\infty} \beta^{i-1}R_i\right]
\]
The identity is thus proved.
}

\begin{exercise}
\label{ex:conditioning-expectation-34}
A set of $n$ dice is thrown. All those that land on six are put aside, and the others are again thrown. This is repeated until all the dice have landed on six. Let $N$ denote the number of throws needed. Let $m_n = E[N]$.
\begin{enumerate}
    \item[(a)] Derive a recursive formula for $m_n$ and use it to calculate $m_i, i = 2, 3, 4$ and to show that $m_5 \approx 13.024$.
    \item[(b)] Let $X_i$ denote the number of dice rolled on the $i$th throw. Find $E\left[\sum_{i=1}^N X_i\right]$.
\end{enumerate}
\end{exercise}

\solution{
\textbf{(a)} On the first throw of $n$ dice, the number of dice that land on six, $K$, follows a Binomial distribution with parameters $n$ and $1/6$. Conditioning on $K$:
\[
m_n = 1 + \sum_{k=0}^{n} \binom{n}{k} \left(\frac{1}{6}\right)^k \left(\frac{5}{6}\right)^{n-k} m_{n-k}
\]
Isolating $m_n$ (the $k=0$ term contains $m_n$):
\[
m_n = 1 + \left(\frac{5}{6}\right)^n m_n + \sum_{k=1}^{n} \binom{n}{k} \left(\frac{1}{6}\right)^k \left(\frac{5}{6}\right)^{n-k} m_{n-k}
\]
\[
m_n \left(1 - \left(\frac{5}{6}\right)^n\right) = 1 + \sum_{k=1}^{n} \binom{n}{k} \left(\frac{1}{6}\right)^k \left(\frac{5}{6}\right)^{n-k} m_{n-k}
\]
Starting with $m_0 = 0$:
\begin{itemize}
    \item For $n=1$: $m_1 (1 - 5/6) = 1 \implies m_1 = 6$.
    \item For $n=2$: $m_2 (1 - 25/36) = 1 + \binom{2}{1}\frac{5}{36}m_1 + \binom{2}{2}\frac{1}{36}m_0 \implies m_2 \frac{11}{36} = 1 + \frac{60}{36} \implies m_2 = \frac{96}{11} \approx 8.727$.
    \item Following this recursion for $m_3, m_4, m_5$, we find $m_3=10566/1001\approx10.555$, $m_4=728256/61061\approx11.927$, and $m_5 \approx 13.024$.
\end{itemize}

\textbf{(b)} Let $Y_j$ be the total number of throws required for die $j$ to land on six. Each $Y_j$ is independent and follows a Geometric($1/6$) distribution, so $\sum_{i=1}^N \mathbf{1}[\text{die } j \text{ is rolled in throw } i] = Y_j$. The total number of die rolls across all throws is $\sum_{i=1}^N X_i = \sum_{j=1}^n Y_j$. By linearity of expectation:
\[
E\left[\sum_{i=1}^N X_i\right] = \sum_{j=1}^n E[Y_j] = n \cdot 6 = 6n
\]
}

\begin{exercise}
\label{ex:conditioning-expectation-35}
Consider $n$ multinomial trials, where each trial independently results in outcome $i$ with probability $p_i$, $\sum_{i=1}^k p_i = 1$. With $X_i$ equal to the number of trials that result in outcome $i$, find $E[X_1 \mid X_2 > 0]$.
\end{exercise}

\solution{
Assume $n\ge1$ and $0<p_2<1$. Given $X_2=j$, the remaining $n-j$
trials have conditional type-$1$ probability $p_1/(1-p_2)$, so
$E[X_1\mid X_2=j]=(n-j)p_1/(1-p_2)$.
Also, $X_2\sim\operatorname{Binomial}(n,p_2)$ gives
\[
E[X_2\mid X_2>0]=\frac{np_2}{1-(1-p_2)^n}.
\]
Consequently
\begin{align*}
E[X_1\mid X_2>0]
&=\frac{p_1}{1-p_2}\left(n-\frac{np_2}{1-(1-p_2)^n}\right)\\
&=np_1\frac{1-(1-p_2)^{n-1}}{1-(1-p_2)^n}.
\end{align*}
If $p_2=1$, then $p_1=0$ and the answer is zero. If $p_2=0$ or $n=0$,
the conditioning event has probability zero.
}

\begin{exercise}
\label{ex:conditioning-expectation-36}
Let $p_0 = P\{X = 0\}$ and suppose that $0 < p_0 < 1$. Let $\mu = E[X]$ and $\sigma^2 = \mathrm{Var}(X)$.
\begin{enumerate}
    \item[(a)] Find $E[X \mid X \neq 0]$.
    \item[(b)] Find $\mathrm{Var}(X \mid X \neq 0)$.
\end{enumerate}
\end{exercise}

\solution{
\textbf{(a)} Since $X$ is zero off the conditioning event,
\[
\mu=E[X\mid X\ne0](1-p_0),\qquad
E[X\mid X\ne0]=\frac{\mu}{1-p_0}.
\]
\textbf{(b)} Similarly,
$E[X^2\mid X\ne0]=(\sigma^2+\mu^2)/(1-p_0)$. Therefore
\begin{align*}
\operatorname{Var}(X\mid X\ne0)
&=\frac{\sigma^2+\mu^2}{1-p_0}-\frac{\mu^2}{(1-p_0)^2}\\
&=\frac{\sigma^2}{1-p_0}-\frac{\mu^2p_0}{(1-p_0)^2}.
\end{align*}
The second term has a minus sign. For a Bernoulli variable, conditioning
on a nonzero value makes it constant and the formula gives variance zero.
}

\begin{exercise}
\label{ex:conditioning-expectation-37}
A manuscript is sent to a typing firm consisting of typists A, B, and C. If it is typed by A, then the number of errors made is a Poisson random variable with mean 2.6; if typed by B, then the number of errors is a Poisson random variable with mean 3; and if typed by C, then it is a Poisson random variable with mean 3.4. Let $X$ denote the number of errors in the typed manuscript. Assume that each typist is equally likely to do the work. Find (a) $E[X]$, (b) $\mathrm{Var}(X)$.
\end{exercise}

\solution{
Let $Y$ denote the typist chosen, with $P(A) = P(B) = P(C) = 1/3$.

\textbf{(a)} By the law of total expectation:
\[
E[X] = E[E[X|Y]] = \frac{1}{3}(2.6 + 3.0 + 3.4) = \frac{9}{3} = 3
\]

\textbf{(b)} We use the conditional variance formula $\mathrm{Var}(X) = E[\mathrm{Var}(X|Y)] + \mathrm{Var}(E[X|Y])$.
Since $X|Y$ is Poisson, its conditional variance equals its conditional mean: $\mathrm{Var}(X|Y) = E[X|Y]$. Thus, $E[\mathrm{Var}(X|Y)] = E[E[X|Y]] = 3$.
Next, compute $\mathrm{Var}(E[X|Y])$ where $E[X|Y]$ takes values 2.6, 3, 3.4 with probability 1/3:
\[
\mathrm{Var}(E[X|Y]) = \frac{1}{3}\left[(2.6-3)^2 + (3-3)^2 + (3.4-3)^2\right] = \frac{0.16 + 0 + 0.16}{3} = \frac{0.32}{3} \approx 0.1067
\]
Combining the parts:
\[
\mathrm{Var}(X) = 3 + \frac{0.32}{3} = \frac{9.32}{3} \approx 3.1067
\]
}

\begin{exercise}
\label{ex:conditioning-expectation-38}
Let $U$ be a uniform $(0, 1)$ random variable. Suppose that $n$ trials are to be performed and that conditional on $U = u$ these trials will be independent with a common success probability $u$. Compute the mean and variance of the number of successes that occur in these trials.
\end{exercise}

\solution{
Let $X$ be the number of successes. Given $U=u$, $X \mid U=u \sim \text{Binomial}(n, u)$. Thus, $E[X \mid U] = nU$ and $\mathrm{Var}(X \mid U) = nU(1-U)$.

\textbf{Mean:}
\[
E[X] = E[E[X \mid U]] = E[nU] = n E[U] = n\left(\frac{1}{2}\right) = \frac{n}{2}
\]

\textbf{Variance:}
Using the conditional variance formula:
\[
\mathrm{Var}(X) = E[\mathrm{Var}(X \mid U)] + \mathrm{Var}(E[X \mid U]) = E[nU(1-U)] + \mathrm{Var}(nU)
\]
We know that for a Uniform(0,1) variable, $E[U] = 1/2$, $E[U^2] = 1/3$, and $\mathrm{Var}(U) = 1/12$.
\[
E[nU(1-U)] = n(E[U] - E[U^2]) = n\left(\frac{1}{2} - \frac{1}{3}\right) = \frac{n}{6}
\]
\[
\mathrm{Var}(nU) = n^2 \mathrm{Var}(U) = \frac{n^2}{12}
\]
Thus:
\[
\mathrm{Var}(X) = \frac{n}{6} + \frac{n^2}{12} = \frac{2n + n^2}{12} = \frac{n(n+2)}{12}
\]
}

\begin{exercise}
\label{ex:conditioning-expectation-39}
A deck of \textit{n} cards, numbered $1$ through \textit{n}, is randomly shuffled so that all $n!$ possible
permutations are equally likely. The cards are then turned over one at a time until card number $1$ appears. These upturned cards constitute the first cycle. We now determine (by looking at the upturned cards) the lowest numbered card that has not yet appeared, and we continue to turn the cards face up until that card appears.
This new set of cards represents the second cycle. We again determine the lowest numbered of the remaining cards and turn the cards until it appears, and so on until
all cards have been turned over. Let $m_n$ denote the mean number of cycles.
\begin{enumerate}
    \item[(a)] Derive a recursive formula for $m_n$ in terms of $m_k, k = 1, \ldots, n - 1$.
    \item[(b)] Starting with $m_0 = 0$, find $m_1, m_2, m_3$, and $m_4$.
    \item[(c)] Conjecture a general formula for $m_n$.
    \item[(d)] Prove your formula by induction on $n$.
    \item[(e)] Let $X_i = 1$ if a cycle ends with card $i$, 0 otherwise. Express cycles using $X_i$.
    \item[(f)] Use part (e) to determine $m_n$.
    \item[(g)] Are the random variables $X_1, \ldots, X_n$ independent? Explain.
    \item[(h)] Find the variance of the number of cycles.
\end{enumerate}
\end{exercise}

\solution{
\textbf{(a)} The position of card $1$ is uniform on $\{1,\ldots,n\}$.
After its cycle ends in position $k$, the relative order of the remaining
$n-k$ cards is uniform. Consequently
\[
m_n=1+\frac1n\sum_{j=0}^{n-1}m_j,\qquad m_0=0.
\]
\textbf{(b)} This gives $m_1=1$, $m_2=3/2$, $m_3=11/6$, and $m_4=25/12$.

\textbf{(c)--(d)} Subtract the recurrence multiplied by $n-1$ at size
$n-1$ from the recurrence multiplied by $n$ at size $n$. The result is
$n(m_n-m_{n-1})=1$. Induction gives $m_n=H_n=\sum_{i=1}^n1/i$.

\textbf{(e)--(f)} A cycle ends at card $i$ exactly when card $i$ occurs
\emph{after every card with a smaller label}. Thus the number of cycles
is $\sum_iX_i$, and symmetry in the relative order of cards $1,\ldots,i$
gives $P(X_i=1)=1/i$. Hence the mean is again $H_n$.

\textbf{(g)} Construct a uniform permutation by inserting card $i$
uniformly into one of the $i$ positions among cards $1,\ldots,i-1$,
independently for each $i$. The event $X_i=1$ occurs exactly when card
$i$ is inserted at the end. Later insertions do not change its order
relative to smaller cards. The indicators $X_i$ are therefore mutually
independent.

\textbf{(h)} It follows that
\[
\operatorname{Var}\left(\sum_iX_i\right)
=\sum_{i=1}^n\frac1i\left(1-\frac1i\right)
=H_n-\sum_{i=1}^n\frac1{i^2}.
\]
}

\begin{exercise}
\label{ex:conditioning-expectation-40}
A prisoner is trapped in a cell containing three doors. The first door leads to a
tunnel that returns him to his cell after two days of travel. The second leads to a
tunnel that returns him to his cell after three days of travel. The third door leads
immediately to freedom.
\begin{enumerate}
    \item[(a)] Assuming that the prisoner will always select doors $1$, $2$, and $3$ with prob-
abilities $0.5$, $0.3$, $0.2$, what is the expected number of days until he reaches freedom?
    \item[(b)] Assuming that the prisoner is always equally likely to choose among those
doors that he has not used, what is the expected number of days until he reaches freedom? (In this version, for instance, if the prisoner initially tries
door 1, then when he returns to the cell, he will now select only from doors $2$ and $3$.)
    \item[(c)] For parts (a) and (b) find the variance of the number of days until the prisoner reaches freedom.
\end{enumerate}
\end{exercise}

\solution{
Let $m=E[X]$ and $v=E[X^2]$.
\textbf{(a)} With repeated independent door choices,
\[
m=0.5(2+m)+0.3(3+m),\qquad m=9.5\text{ days}.
\]
\textbf{(b)} Without reusing a door, the possible total times and
probabilities are $0$ with probability $1/3$, $2$ and $3$ with
probability $1/6$ each, and $5$ with probability $1/3$.
Their mean is $2.5$ days.

\textbf{(c)} For the repeated-choice policy, expanding the squared
first-step recurrence gives
\[
v=0.5(4+4m+v)+0.3(9+6m+v)=4.7+3.8m+0.8v.
\]
Thus $v=204$ and $\operatorname{Var}(X)=204-9.5^2=113.75$ days squared.
For the no-reuse policy, the distribution in (b) gives
\[
E[X^2]=\frac{2^2+3^2}{6}+\frac{5^2}{3}=10.5,
\qquad \operatorname{Var}(X)=10.5-2.5^2=4.25.
\]
}

\begin{exercise}
\label{ex:conditioning-expectation-41}
A rat is trapped in a maze. Initially it has to choose one of two directions. If it goes to the right, then it will wander around in the maze for three minutes and will
then return to its initial position. If it goes to the left, then with probability $\frac{1}{3}$ it will depart the maze after two minutes of traveling, and with probability $\frac{2}{3}$ it will return to its initial position after five minutes of traveling. Assuming that the rat is at all times equally likely to go to the left or the right, what is the expected number of minutes that it will be trapped in the maze?
\end{exercise}

\solution{
Let $E[T]$ be the expected time to escape. Conditioning on the initial choice (Left or Right with probability 1/2):
\[
E[T] = \frac{1}{2} E[T \mid \text{Right}] + \frac{1}{2} E[T \mid \text{Left}]
\]
where:
\[
E[T \mid \text{Right}] = 3 + E[T]
\]
\[
E[T \mid \text{Left}] = \frac{1}{3}(2) + \frac{2}{3}(5 + E[T]) = 4 + \frac{2}{3}E[T]
\]
Substituting these into the main relation:
\[
E[T] = \frac{1}{2}(3 + E[T]) + \frac{1}{2}\left(4 + \frac{2}{3}E[T]\right) = \frac{7}{2} + \left(\frac{1}{2} + \frac{1}{3}\right)E[T] = \frac{7}{2} + \frac{5}{6}E[T]
\]
\[
\left(1 - \frac{5}{6}\right) E[T] = \frac{7}{2} \implies \frac{1}{6} E[T] = \frac{7}{2} \implies E[T] = 21 \text{ minutes}
\]
}

\begin{exercise}
\label{ex:conditioning-expectation-42}
If $X_i$, $i = 1, \ldots, n$, are independent normal random variables, with $X_i$ having mean $\mu_i$ and variance $1$, then the random variable
\[
\sum_{i=1}^{n} X_i^2
\]
is said to be a \emph{noncentral chi-squared} random variable.

\begin{enumerate}
    \item[(a)] If $X$ is a normal random variable having mean $\mu$ and variance $1$, show that, for $|t|<\frac12$, the moment generating function of $X^2$ is
    \[
    M_{X^2}(t)
    =
    (1-2t)^{-1/2}
    \exp\!\left(\frac{t\mu^2}{1-2t}\right).
    \]

    \item[(b)] Derive the moment generating function of the noncentral chi-squared random variable
    \[
    \sum_{i=1}^{n} X_i^2,
    \]
    and show that its distribution depends on the sequence of means $\mu_1,\ldots,\mu_n$ only through the sum of their squares. As a result, show that
    \[
    \sum_{i=1}^{n} X_i^2
    \]
    is a noncentral chi-squared random variable with parameters $
    n
    ~\text{and}~
    \theta=\sum_{i=1}^{n}\mu_i^2.$

    \item[(c)] If all $\mu_i=0$, then
    \[
    \sum_{i=1}^{n} X_i^2
    \]
    is called a chi-squared random variable with $n$ degrees of freedom. Determine, by differentiating its moment generating function, its expected value and variance.

    \item[(d)] Let $K$ be a Poisson random variable with mean $\theta/2$, and suppose that, conditional on $K=k$, the random variable $W$ has a chi-squared distribution with $n+2k$ degrees of freedom. Show, by computing its moment generating function, that $W$ is a noncentral chi-squared random variable with parameters $n$ and $\theta$.

    \item[(e)] Find the expected value and variance of a noncentral chi-squared random variable with parameters $n$ and $\theta$.
\end{enumerate}
\end{exercise}

\solution{
\textbf{(a)} $M_{X^2}(t) = E[e^{tX^2}] = \int_{-\infty}^{\infty} \frac{1}{\sqrt{2\pi}} e^{tx^2} e^{-(x-\mu)^2/2} \, dx$. Completing the square in the exponent leads directly to the given formula: $(1-2t)^{-1/2}\exp\left(\frac{t\mu^2}{1-2t}\right)$.

\textbf{(b)} By independence, the MGF of the sum is the product of the individual MGFs:
\[
M_{\sum X_i^2}(t) = \prod_{i=1}^n (1-2t)^{-1/2} e^{\frac{t\mu_i^2}{1-2t}} = (1-2t)^{-n/2} e^{\frac{t \sum \mu_i^2}{1-2t}} = (1-2t)^{-n/2} e^{\frac{t\theta}{1-2t}}
\]
This depends only on $n$ and $\theta = \sum \mu_i^2$.

\textbf{(c)} For $\mu_i = 0 \implies \theta = 0$, $M(t) = (1-2t)^{-n/2}$.
$M'(t) = n(1-2t)^{-n/2-1} \implies E[W] = M'(0) = n$.
$M''(t) = n(n+2)(1-2t)^{-n/2-2} \implies E[W^2] = n(n+2) \implies \mathrm{Var}(W) = n(n+2) - n^2 = 2n$.

\textbf{(d)} Given $K=k$, $W \sim \chi^2_{n+2k}$, so $E[e^{tW} \mid K=k] = (1-2t)^{-(n+2k)/2}$.
\[
\begin{aligned}
M_W(t)&=E[E[e^{tW}\mid K]]\\
&=(1-2t)^{-n/2}e^{-\theta/2}
\sum_{k=0}^{\infty}\frac{[\theta/(2(1-2t))]^k}{k!}.
\end{aligned}
\]
\[
= (1-2t)^{-n/2} e^{-\theta/2} e^{\frac{\theta}{2(1-2t)}} = (1-2t)^{-n/2} e^{\frac{\theta}{2}\left(\frac{1}{1-2t}-1\right)} = (1-2t)^{-n/2} e^{\frac{t\theta}{1-2t}}
\]
This matches the noncentral chi-squared MGF perfectly.

\textbf{(e)} Using conditional formulas with $K \sim \text{Poisson}(\theta/2)$:
\[
E[W] = E[E[W \mid K]] = E[n + 2K] = n + 2(\theta/2) = n + \theta
\]
\[
\begin{aligned}
\operatorname{Var}(W)&=E[\operatorname{Var}(W\mid K)]+\operatorname{Var}(E[W\mid K])\\
&=2n+4E[K]+4\operatorname{Var}(K)=2n+4\theta.
\end{aligned}
\]
}

\begin{exercise}
\label{ex:conditioning-expectation-43}
The density function of a chi-squared random variable having $n$ degrees of freedom can be shown to be
\[
f(x)
=
\frac{\frac12 e^{-x/2}(x/2)^{n/2-1}}
{\Gamma(n/2)},
\qquad x>0,
\]
where the gamma function is defined by
\[
\Gamma(t)
=
\int_0^\infty e^{-x}x^{t-1}\,dx,
\qquad t>0.
\]

Integration by parts can be employed to show that
\[
\Gamma(t)=(t-1)\Gamma(t-1),
\qquad t>1.
\]

If $Z$ and $\chi_n^2$ are independent random variables with $Z$ having a standard normal distribution and $\chi_n^2$ having a chi-squared distribution with $n$ degrees of freedom, then the random variable
\[
T
=
\frac{Z}{\sqrt{\chi_n^2/n}}
\]
is said to have a $t$-distribution with $n$ degrees of freedom.

Compute the mean and variance of $T$ when $n>2$.
\end{exercise}

\solution{
\textbf{Mean:}
Using conditional expectation, conditioning on $\chi^2_n$:
\[
E[T] = E\left[ E\left[ \frac{Z}{\sqrt{\chi^2_n / n}} \;\middle|\; \chi^2_n \right] \right] = E\left[ \frac{1}{\sqrt{\chi^2_n / n}} E[Z \mid \chi^2_n] \right]
\]
Since $Z$ is independent of $\chi^2_n$, $E[Z \mid \chi^2_n] = E[Z] = 0$. Therefore, $E[T] = 0$.

\textbf{Variance:}
Since $E[T] = 0$, $\mathrm{Var}(T) = E[T^2]$.
\[
E[T^2] = E\left[ \frac{Z^2}{\chi^2_n / n} \right] = n E[Z^2] E\left[ \frac{1}{\chi^2_n} \right]
\]
Since $Z \sim N(0,1)$, $E[Z^2] = 1$. Now compute $E\left[\frac{1}{\chi^2_n}\right]$ using the density given in the text:
\[
E\left[ \frac{1}{\chi^2_n} \right] = \int_{0}^{\infty} \frac{1}{x} \frac{1}{2} \frac{e^{-x/2} (x/2)^{n/2 - 1}}{\Gamma(n/2)} \, dx = \frac{1}{2} \int_{0}^{\infty} \frac{e^{-x/2} (x/2)^{n/2 - 2} \cdot \frac{1}{2}}{\Gamma(n/2)} \, dx
\]
Substitute $u = x/2 \implies dx = 2du$:
\[
= \frac{1}{2} \int_{0}^{\infty} \frac{e^{-u} u^{n/2 - 2}}{\Gamma(n/2)} \, du = \frac{1}{2} \frac{\Gamma(n/2 - 1)}{\Gamma(n/2)}
\]
Using the recursive property $\Gamma(t) = (t-1)\Gamma(t-1)$, we have $\Gamma(n/2) = (n/2 - 1)\Gamma(n/2 - 1)$:
\[
= \frac{1}{2} \frac{1}{n/2 - 1} = \frac{1}{2} \frac{2}{n - 2} = \frac{1}{n - 2}
\]
Substituting this back into the formula for $E[T^2]$:
\[
\mathrm{Var}(T) = n \cdot 1 \cdot \frac{1}{n - 2} = \frac{n}{n - 2}
\]
}

\setcounter{exercise}{45}

\begin{exercise}
\label{ex:conditioning-expectation-46}
    (a) Show that
    \[
    \text{Cov}(X, Y) = \text{Cov}(X, E[Y \mid X])
    \]
    (b) Suppose that, for constants $a$ and $b$,
    \[
    E[Y \mid X] = a + bX
    \]
    Show that
    \[
    b = \text{Cov}(X, Y) / \text{Var}(X)
    \]
\end{exercise}

\solution{
    (a) By the definition of covariance, we know that:
    \[
    \text{Cov}(X, Y) = E[XY] - E[X]E[Y]
    \]
    Using the Law of Total Expectation, we can write $E[Y]$ as:
    \[
    E[Y] = E[E[Y \mid X]]
    \]
    Similarly, applying the conditioning property of total expectation to the term $E[XY]$, and recalling that given $X$, the variable $X$ acts as a constant and can be pulled out of the conditional expectation, we get:
    \[
    E[XY] = E[E[XY \mid X]] = E[X E[Y \mid X]]
    \]
    Substituting these two relationships into the initial covariance formula yields:
    \[
    \text{Cov}(X, Y) = E[X E[Y \mid X]] - E[X]E[E[Y \mid X]]
    \]
    The expression on the right matches exactly the definition of the covariance between $X$ and the random variable $E[Y \mid X]$. Consequently:
    \[
    \text{Cov}(X, Y) = \text{Cov}(X, E[Y \mid X])
    \]

    (b) Leveraging the result proven in part (a), we substitute the assumption $E[Y \mid X] = a + bX$ into the covariance:
    \[
    \text{Cov}(X, Y) = \text{Cov}(X, a + bX)
    \]
    Using the linearity properties of covariance, we know that adding a constant $a$ does not affect the covariance, and multiplicative coefficients can be pulled out:
    \[
    \text{Cov}(X, a + bX) = \text{Cov}(X, bX) = b\,\text{Cov}(X, X)
    \]
    Since the covariance of a variable with itself is its variance ($\text{Cov}(X, X) = \text{Var}(X)$), we get:
    \[
    \text{Cov}(X, Y) = b\,\text{Var}(X)
    \]
    Isolating the parameter $b$ (assuming $\text{Var}(X) > 0$), we arrive at the thesis:
    \[
    b = \frac{\text{Cov}(X, Y)}{\text{Var}(X)}
    \]
}

\begin{exercise}
\label{ex:conditioning-expectation-47}
    If $E[Y \mid X] = 1$, show that
    \[
    \text{Var}(X Y) \geqslant \text{Var}(X)
    \]
\end{exercise}

\solution{
    We use the Law of Total Variance, which states that for any two random variables, the variance can be decomposed as:
    \[
    \text{Var}(XY) = E[\text{Var}(XY \mid X)] + \text{Var}(E[XY \mid X])
    \]
    Let us analyze the two terms on the right-hand side individually:
    \begin{enumerate}
        \item For the second term, given $X$, the variable $X$ acts as a constant within the conditional expectation, so:
        \[
        E[XY \mid X] = X E[Y \mid X]
        \]
        Since by hypothesis $E[Y \mid X] = 1$, the term reduces to:
        \[
        E[XY \mid X] = X \cdot 1 = X
        \]
        Consequently, the variance of this term is:
        \[
        \text{Var}(E[XY \mid X]) = \text{Var}(X)
        \]

        \item For the first term, since the variance of any random variable (even conditional) is always non-negative, the expectation of its value will also be non-negative:
        \[
        E[\text{Var}(XY \mid X)] \geqslant 0
        \]
    \end{enumerate}
    Substituting these elements back into the initial decomposition formula, we get:
    \[
    \text{Var}(XY) = E[\text{Var}(XY \mid X)] + \text{Var}(X)
    \]
    Since $E[\text{Var}(XY \mid X)] \geqslant 0$, we can conclude that:
    \[
    \text{Var}(XY) \geqslant \text{Var}(X)
    \]
}

\setcounter{exercise}{50}

\begin{exercise}
\label{ex:conditioning-expectation-51}
    If $X$ is geometric with parameter $p$, find the probability that $X$ is even.
\end{exercise}

\solution{
    For $p=1$, $X=1$ and the probability of an even value is $0$. Assume $0<p<1$ for the following computation. A geometric random variable $X$ with parameter $p$ has a probability mass function given by $P(X = n) = (1-p)^{n-1}p$ for $n = 1, 2, \dots$. We want to find the probability that $X$ takes an even value ($X = 2k$ for $k = 1, 2, \dots$):
    \[
    P(X \text{ is even}) = \sum_{k=1}^{\infty} P(X = 2k) = \sum_{k=1}^{\infty} (1-p)^{2k-1}p
    \]
    We can rewrite the summation term by factoring out common terms:
    \[
    P(X \text{ is even}) = \frac{p}{1-p} \sum_{k=1}^{\infty} \left((1-p)^2\right)^k
    \]
    Using the formula for the sum of an infinite geometric series $\sum_{k=1}^{\infty} q^k = \frac{q}{1-q}$ with $q = (1-p)^2$:
    \[
    P(X~ \text {is even}) = \frac{p}{1-p} \cdot \frac{(1-p)^2}{1 - (1-p)^2} = \frac{p(1-p)}{1 - (1 - 2p + p^2)} = \frac{p(1-p)}{2p - p^2} = \frac{1-p}{2-p}
    \]
}

\begin{exercise}
\label{ex:conditioning-expectation-52}
    Suppose that $X$ and $Y$ are independent random variables with probability density functions $f_X$ and $f_Y$. Determine a one-dimensional integral expression for $P\{X + Y < x\}$.
\end{exercise}

\solution{
    Since $X$ and $Y$ are continuous and independent random variables, their joint density is the product of their individual densities, namely $f(s, t) = f_X(s)f_Y(t)$. The required probability is given by the double integral over the region $s + t < x$:
    \[
    P\{X + Y < x\} = \int_{-\infty}^{\infty} \int_{-\infty}^{x-s} f_X(s)f_Y(t) \, dt \, ds
    \]
    We can express this as a single one-dimensional integral by integrating with respect to $s$, where the inner integrand corresponds to the cumulative distribution function of $Y$ evaluated at $x-s$:
    \[
    P\{X + Y < x\} = \int_{-\infty}^{\infty} f_X(s) F_Y(x-s) \, ds
    \]
    Alternatively, integrating first with respect to $s$:
    \[
    P\{X + Y < x\} = \int_{-\infty}^{\infty} f_Y(t) F_X(x-t) \, dt
    \]
}

\begin{exercise}
\label{ex:conditioning-expectation-53}
    *Suppose $X$ is a Poisson random variable with mean $\lambda$. The parameter $\lambda$ is itself a random variable whose distribution is exponential with mean 1. Show that $P\{X = n\} = \left(\frac{1}{2}\right)^{n+1}$.
\end{exercise}

\solution{
    We use the Law of Total Probability by conditioning on the parameter $\Lambda$:
    \[
    P\{X = n\} = \int_{0}^{\infty} P\{X = n \mid \Lambda = \lambda\} f_{\Lambda}(\lambda) \, d\lambda
    \]
    We know that $X \mid \Lambda = \lambda \sim \text{Poisson}(\lambda)$ and that $\Lambda \sim \text{Exponential}(1)$, so $f_{\Lambda}(\lambda) = e^{-\lambda}$ for $\lambda > 0$. Substituting the densities, we get:
    \[
    P\{X = n\} = \int_{0}^{\infty} \frac{e^{-\lambda} \lambda^n}{n!} e^{-\lambda} \, d\lambda = \frac{1}{n!} \int_{0}^{\infty} \lambda^n e^{-2\lambda} \, d\lambda
    \]
    We perform the change of variable $u = 2\lambda$, which gives $\lambda = u/2$ and $d\lambda = du/2$:
    \[
    P\{X = n\} = \frac{1}{n!} \int_{0}^{\infty} \left(\frac{u}{2}\right)^n e^{-u} \frac{du}{2} = \frac{1}{n! \, 2^{n+1}} \int_{0}^{\infty} u^n e^{-u} \, du
    \]
    The integral corresponds to the definition of the Gamma function $\Gamma(n+1) = n!$. Consequently:
    \[
    P\{X = n\} = \frac{1}{n! \, 2^{n+1}} \cdot n! = \frac{1}{2^{n+1}} = \left(\frac{1}{2}\right)^{n+1}
    \]
}

\begin{exercise}
\label{ex:conditioning-expectation-54}
    A coin is randomly selected from a group of ten coins, the $n$th coin having a probability $n/10$ of coming up heads. The coin is then repeatedly flipped until a head appears. Let $N$ denote the number of flips necessary. What is the probability distribution of $N$? Is $N$ a geometric random variable? When would $N$ be a geometric random variable; that is, what would have to be done differently?
\end{exercise}

\solution{
    Let $P_h$ be the probability of getting heads for the chosen coin. Conditioning on the choice of the $n$-th coin (chosen uniformly with probability $1/10$ for $n=1,\dots,10$):
    \[
    P\{N = k\} = \sum_{n=1}^{10} P\{N = k \mid \text{coin } n\} \cdot P(\text{coin } n) = \sum_{n=1}^{10} \left(1 - \frac{n}{10}\right)^{k-1} \left(\frac{n}{10}\right) \cdot \frac{1}{10}
    \]
    Solving numerically or leaving it in closed form, this is the probability mass function of $N$.

    \textbf{Answers to the questions:}
    \par\noindent\textbullet\enspace $N$ \textbf{is not} a geometric random variable because its final distribution is a mixture of geometric distributions, which does not result in a pure geometric distribution (it lacks the memoryless property globally).
    \par\noindent\textbullet\enspace $N$ would be a geometric random variable if a coin were randomly chosen \textbf{afresh before each individual flip}, instead of repeatedly flipping the same single coin drawn at the very beginning.
}

\begin{exercise}
\label{ex:conditioning-expectation-55}
    You are invited to a party. Suppose the times at which invitees arrive are independent uniform $(0,1)$ random variables. Suppose that, aside from yourself, the number of other people who are invited is a Poisson random variable with mean 10.
    (a) Find the expected number of people who arrive before you.
    (b) Find the probability that you are the $n$th person to arrive.
\end{exercise}

\solution{
    Let $Y$ be the number of other invited people ($Y \sim \text{Poisson}(10)$) and let $T_0$ be your arrival time, with $T_0 \sim U(0,1)$. Each of the other guests arrives at a time $T_i \sim U(0,1)$.

    (a) Conditioning on your arrival time $T_0 = t$, each other person has an independent probability equal to $t$ of arriving before you. The number of people $X$ arriving before you, given $Y=y$ and $T_0=t$, is binomial $B(y, t)$. By the properties of the Poisson distribution, removing the conditioning on $Y$, we find that $X \mid T_0 = t \sim \text{Poisson}(10t)$. The conditional expectation is $E[X \mid T_0] = 10T_0$. Using the Law of Multiple Expectations:
    \[
    E[X] = E[E[X \mid T_0]] = E[10T_0] = 10 \cdot E[T_0] = 10 \cdot \frac{1}{2} = 5
    \]

    (b) Being the $n$-th person to arrive means that exactly $n-1$ people arrive before you. We compute the total probability by integrating against the density of $T_0$:
    \[
    P\{X = n-1\} = \int_{0}^{1} P\{X = n-1 \mid T_0 = t\} \, dt = \int_{0}^{1} \frac{e^{-10t} (10t)^{n-1}}{(n-1)!} \, dt
    \]
    Setting $u = 10t$, which gives $dt = du/10$:
    \[
    P\{X = n-1\} = \frac{1}{10 \cdot (n-1)!} \int_{0}^{10} e^{-u} u^{n-1} \, du
    \]
    This integral represents the lower incomplete Gamma function $\gamma(n, 10)$.
}

\begin{exercise}
\label{ex:conditioning-expectation-56}
    Data indicate that the number of traffic accidents in Berkeley on a rainy day is a Poisson random variable with mean 9, whereas on a dry day it is a Poisson random variable with mean 3. Let $X$ denote the number of traffic accidents tomorrow. If it will rain tomorrow with probability 0.6, find
    (a) $E[X]$;
    (b) $P\{X = 0\}$;
    (c) $\text{Var}(X)$.
\end{exercise}

\solution{
Let $R$ be the rain indicator, with $P(R=1)=0.6$. The conditional
Poisson mean and variance are both $\Lambda=3+6R$.
\textbf{(a)} $E[X]=E[\Lambda]=3+6(0.6)=6.6$.
\textbf{(b)} Conditioning on $R$ gives
$P(X=0)=0.6e^{-9}+0.4e^{-3}$.
\textbf{(c)} The conditional variance formula gives
\[
\operatorname{Var}(X)=E[\Lambda]+\operatorname{Var}(\Lambda)
=6.6+36(0.6)(0.4)=15.24.
\]
}

\begin{exercise}
\label{ex:conditioning-expectation-57}
    The number of storms in the upcoming rainy season is Poisson distributed but with a parameter value that is uniformly distributed over $(0,5)$. That is, $\Lambda$ is uniformly distributed over $(0,5)$, and given that $\Lambda = \lambda$, the number of storms is Poisson with mean $\lambda$. Find the probability there are at least three storms this season.
\end{exercise}

\solution{
    We want to compute $P\{X \geqslant 3\} = 1 - P\{X = 0\} - P\{X = 1\} - P\{X = 2\}$, knowing that $X \mid \Lambda = \lambda \sim \text{Poisson}(\lambda)$ and $\Lambda \sim U(0,5)$ with density $f_{\Lambda}(\lambda) = 1/5$ for $\lambda \in (0,5)$.
    We compute the general term $P\{X = k\}$ by integration:
    \[
    P\{X = k\} = \int_{0}^{5} \frac{e^{-\lambda}\lambda^k}{k!} \cdot \frac{1}{5} \, d\lambda = \frac{1}{5 \cdot k!} \int_{0}^{5} \lambda^k e^{-\lambda} \, d\lambda
    \]
    Let us evaluate the individual cases for $k=0,1,2$:
    \par\noindent\textbullet\enspace For $k=0$: $P\{X=0\} = \frac{1}{5}\int_{0}^{5} e^{-\lambda}d\lambda = \frac{1}{5}(1 - e^{-5})$.
    \par\noindent\textbullet\enspace For $k=1$: $P\{X=1\} = \frac{1}{5}\int_{0}^{5} \lambda e^{-\lambda}d\lambda = \frac{1}{5}[-e^{-\lambda}(\lambda+1)]_0^5 = \frac{1}{5}(1 - 6e^{-5})$.
    \par\noindent\textbullet\enspace For $k=2$: $P\{X=2\} = \frac{1}{10}\int_{0}^{5} \lambda^2 e^{-\lambda}d\lambda = \frac{1}{10}[-e^{-\lambda}(\lambda^2+2\lambda+2)]_0^5 = \frac{1}{10}(2 - 37e^{-5}) = \frac{1}{5}(1 - 18.5e^{-5})$.

    Summing the three probabilities:
    \[
    P\{X \leqslant 2\} = \frac{1}{5}\left[(1 - e^{-5}) + (1 - 6e^{-5}) + (1 - 18.5e^{-5})\right] = \frac{1}{5}(3 - 25.5e^{-5}) = 0.6 - 5.1e^{-5}
    \]
    Therefore, the required probability is:
    \[
    P\{X \geqslant 3\} = 1 - (0.6 - 5.1e^{-5}) = 0.4 + 5.1e^{-5}
    \]
}

\begin{exercise}
\label{ex:conditioning-expectation-58}
    A collection of $n$ coins is flipped. The outcomes are independent, and the $i$th coin comes up heads with probability $\alpha_i, i = 1, \dots, n$. Suppose that for some value of $j, 1 \leqslant j \leqslant n, \alpha_j = \frac{1}{2}$. Find the probability that the total number of heads to appear on the $n$ coins is an even number.
\end{exercise}

\solution{
    We condition on the outcomes of all coins except coin $j$. Let $Y$ be the total number of heads obtained from the remaining $n-1$ coins. The total number of heads $X$ will be even if:
    \par\noindent\textbf{1.}\enspace $Y$ is even and coin $j$ shows tails (probability $1 - \alpha_j = 1/2$).
    \par\noindent\textbf{2.}\enspace $Y$ is odd and coin $j$ shows heads (probability $\alpha_j = 1/2$).

    Using the Law of Total Probability by conditioning on $Y$:
    \[
    P(X \text{ is even}) = P(X \text{ is even} \mid Y \text{ is even})P(Y \text{ is even}) + P(X \text{ is even} \mid Y \text{ is odd})P(Y \text{ is odd})
    \]
    \[
    P(X \text{ is even}) = \frac{1}{2} P(Y \text{ is even}) + \frac{1}{2} P(Y \text{ is odd})
    \]
    Since the sum of the probabilities that $Y$ is even or odd must equal 1 ($P(Y \text{ is even}) + P(Y \text{ is odd}) = 1$), we conclude that:
    \[
    P(X \text{ is even}) = \frac{1}{2} (1) = \frac{1}{2}
    \]
}

\begin{exercise}
\label{ex:conditioning-expectation-59}
    Suppose each new coupon collected is, independent of the past, a type $i$ coupon with probability $p_i$. A total of $n$ coupons is to be collected. Let $A_i$ be the event that there is at least one type $i$ in this set. For $i \neq j$, compute $P(A_i A_j)$ by
    (a) conditioning on $N_i$, the number of type $i$ coupons in the set of $n$ coupons;
    (b) conditioning on $F_i$, the first time a type $i$ coupon is collected;
    (c) using the identity $P(A_i \cup A_j) = P(A_i) + P(A_j) - P(A_i A_j)$.
\end{exercise}

\solution{
    We first note that $P(A_i) = 1 - (1-p_i)^n$ and $P(A_i \cup A_j) = 1 - (1-p_i-p_j)^n$ (the complement event is that no coupon of type $i$ or $j$ is collected).

    (a) Conditioning on $N_i \sim \text{Binomial}(n, p_i)$:
    \[
    P(A_i A_j) = \sum_{k=1}^{n} P(A_j \mid N_i = k) P(N_i = k)
    \]
    Given $N_i = k$ (where $k \geqslant 1$, meaning event $A_i$ has certainly occurred), there remain $n-k$ slots available, for which the relative probability of obtaining a type $j$ coupon versus non-$i$ types is $\frac{p_j}{1-p_i}$. Consequently, the conditional probability of getting no type $j$ coupons is $\left(1 - \frac{p_j}{1-p_i}\right)^{n-k}$. Thus:
    \[
    P(A_i A_j) = \sum_{k=1}^{n} \left[1 - \left(1 - \frac{p_j}{1-p_i}\right)^{n-k}\right] \binom{n}{k} p_i^k (1-p_i)^{n-k}
    \]

    (b) Conditioning on the time of the first success $F_i = m$ ($m = 1, \dots, n$):
    If $F_i = m$, it means that type $i$ did not appear in the first $m-1$ trials. The event $A_i A_j$ occurs if at least one type $j$ is collected overall across all $n$ trials.
    \[
    P(A_i A_j) = \sum_{m=1}^{n} P(A_j \mid F_i = m) P(F_i = m)
    \]
    where $P(F_i = m) = (1-p_i)^{m-1}p_i$, and $P(A_j \mid F_i = m) = 1 - \left(1-\frac{p_j}{1-p_i}\right)^{m-1}(1-p_j)^{n-m}$.

    (c) Using the set identity, we isolate the intersection term:
    \[
    P(A_i A_j) = P(A_i) + P(A_j) - P(A_i \cup A_j)
    \]
    Substituting the respective known complement probabilities:
    \[
    P(A_i A_j) = \left[1 - (1-p_i)^n\right] + \left[1 - (1-p_j)^n\right] - \left[1 - (1-p_i-p_j)^n\right]
    \]
    \[
    P(A_i A_j) = 1 - (1-p_i)^n - (1-p_j)^n + (1-p_i-p_j)^n
    \]
}

\begin{exercise}
\label{ex:conditioning-expectation-60}
    *Two players alternate flipping a coin that comes up heads with probability $p$. The first player to obtain a head is declared the winner. We are interested in the probability that the first player to flip is the winner. Before determining this probability, which we will call $f(p)$, answer the following questions.
    (a) Do you think that $f(p)$ is a monotone function of $p$? If so, is it increasing or decreasing?
    (b) What do you think is the value of $\lim_{p \to 1} f(p)$?
    (c) What do you think is the value of $\lim_{p \to 0} f(p)$?
    (d) Find $f(p)$.
\end{exercise}

\solution{
Assume $0<p\le1$ and independent flips.
\textbf{(a)} The first-player winning probability is increasing in $p$:
a larger success probability strengthens the advantage of the first turn.
\textbf{(b)} Its limit as $p\to1$ is $1$.
\textbf{(c)} Its limit as $p\downarrow0$ is $1/2$; at $p=0$ the game
never ends, so this is a limit rather than a winning probability there.
\textbf{(d)} Conditioning on the first flip gives
\[
f(p)=p+(1-p)[1-f(p)],\qquad f(p)=\frac1{2-p}.
\]
Thus $f\prime(p)=1/(2-p)^2>0$, confirming the monotonicity and limits.
}

\begin{exercise}
\label{ex:conditioning-expectation-61}
    Suppose in Exercise 29 that the shooting ends when the target has been hit twice. Let $m_i$ denote the mean number of shots needed for the first hit when player $i$ shoots first, $i = 1, 2$. Also, let $P_i, i = 1, 2$, denote the probability that the first hit is by player 1, when player $i$ shoots first.
    (a) Find $m_1$ and $m_2$.
    (b) Find $P_1$ and $P_2$.
    For the remainder of the problem, assume that player 1 shoots first.
    (c) Find the probability that the final hit was by 1.
    (d) Find the probability that both hits were by 1.
    (e) Find the probability that both hits were by 2.
    (f) Find the mean number of shots taken.
\end{exercise}

\solution{
Put $q_i=1-p_i$ and $d=1-q_1q_2$, and assume $d>0$.
\textbf{(a)} First-step analysis gives
$m_1=1+q_1m_2$ and $m_2=1+q_2m_1$, so
\[
m_1=\frac{1+q_1}{d},\qquad m_2=\frac{1+q_2}{d}.
\]
\textbf{(b)} Similarly $P_1=p_1+q_1P_2$ and $P_2=q_2P_1$ give
\[
P_1=\frac{p_1}{d},\qquad P_2=\frac{q_2p_1}{d}.
\]
After player $1$ makes the first hit, player $2$ shoots next; after
player $2$ makes it, player $1$ shoots next. Restart the first-hit
calculation from that new turn order.

\textbf{(c)} The probability the final hit is by player $1$ is
$P_1P_2+(1-P_1)P_1$.
\textbf{(d)} The probability both hits are by player $1$ is $P_1P_2$.
\textbf{(e)} The probability both hits are by player $2$ is
$(1-P_1)^2$.
\textbf{(f)} The mean total number of shots is
\[
m_1+P_1m_2+(1-P_1)m_1.
\]
These expressions also cover one player never hitting, provided $d>0$.
}

\begin{exercise}
\label{ex:conditioning-expectation-62}
    A, B, and C are evenly matched tennis players. Initially A and B play a set, and the winner then plays C. This continues, with the winner always playing the waiting player, until one of the players has won two sets in a row. That player is then declared the overall winner. Find the probability that A is the overall winner.
\end{exercise}

\solution{
After the first set, let $c,h,w$ be A's winning probabilities when A is,
respectively, the player who just won, the next challenger, or the waiting
player. The other two players are interchangeable. Conditioning on the
next set gives
\[
c=\tfrac12+\tfrac12w,\qquad h=\tfrac12c,\qquad w=\tfrac12h.
\]
A loss by a challenger ends the tournament, while a win makes that
challenger the player with one consecutive win. Solving yields
$c=4/7$, $h=2/7$, and $w=1/7$.
In the first set A wins with probability $1/2$ and then has probability
$c$; if B wins, A waits and has probability $w$. Thus
\[
P(\text{A wins})=\frac12c+\frac12w=\frac5{14}.
\]
}

\begin{exercise}
\label{ex:conditioning-expectation-63}
    Suppose there are $n$ types of coupons, and that the type of each new coupon obtained is independent of past selections and is equally likely to be any of the $n$ types. Suppose one continues collecting until a complete set of at least one of each type is obtained.
    (a) Find the probability that there is exactly one type $i$ coupon in the final collection.
    Hint: Condition on $T$, the number of types that are collected before the first type $i$ appears.
    (b) Find the expected number of types that appear exactly once in the final collection.
\end{exercise}

\solution{
\textbf{(a)} Let $T$ count the distinct types discovered before the
first type $i$. By symmetry $T$ is uniform on $\{0,\ldots,n-1\}$.
Conditional on $T=k$, there are $n-k-1$ missing types after that first
$i$. For $i$ to occur exactly once, all those missing types must appear
before the next $i$. Ignoring already collected types other than $i$,
the discovery order of these $n-k$ relevant types is uniform. Type $i$
must be last, which has probability $1/(n-k)$. Consequently
\[
P(\text{exactly one type }i)
=\frac1n\sum_{k=0}^{n-1}\frac1{n-k}=\frac{H_n}{n}.
\]
Type $i$ need not be the last new type in the original collection:
it can appear earlier and then never repeat before completion.

\textbf{(b)} Summing the singleton-type indicators gives expected number
$n(H_n/n)=H_n$.
}

\begin{exercise}
\label{ex:conditioning-expectation-64}
    $A$ and $B$ roll a pair of dice in turn, with $A$ rolling first. $A$'s objective is to obtain a sum of 6, and $B$'s is to obtain a sum of 7. The game ends when either player reaches his or her objective, and that player is declared the winner.
    (a) Find the probability that $A$ is the winner.
    (b) Find the expected number of rolls of the dice.
    (c) Find the variance of the number of rolls of the dice.
\end{exercise}

\solution{
Let $a=5/36$ and $b=1/6$ be the success probabilities for A and B,
and let $q=(1-a)(1-b)=155/216$ be the probability a full round fails.
\textbf{(a)} Renewal after a failed round gives
\[
P(\text{A wins})=\frac{a}{1-q}=\frac{30}{61}.
\]
\textbf{(b)} A round uses one roll with probability $a$, and otherwise
uses two. First-step analysis yields
\[
E[N]=a+2(1-a)b+q(2+E[N]),
\qquad E[N]=\frac{2-a}{1-q}=\frac{402}{61}.
\]
\textbf{(c)} Let $G$ be the number of failed full rounds before the
successful round, and $D$ the number of rolls in that last round.
Then $N=2G+D$, where $G$ and $D$ are independent,
\[
P(G=g)=q^g(1-q),\qquad P(D=1)=\frac{30}{61},\quad P(D=2)=\frac{31}{61}.
\]
The independence follows from
$P(G=g,D=1)=q^g a$ and $P(G=g,D=2)=q^g(1-a)b$.
Hence
\[
\operatorname{Var}(N)
=\frac{4q}{(1-q)^2}+\frac{30\cdot31}{61^2}
=\frac{134850}{3721}.
\]
}

\begin{exercise}
\label{ex:conditioning-expectation-65}
    The number of red balls in an urn that contains $n$ balls is a random variable that is equally likely to be any of the values $0, 1, \dots, n$. That is,
    \[
    P\{i \text{ red}, n - i \text{ non-red}\} = \frac{1}{n+1}, \quad i = 0, 1, \dots, n
    \]
    The $n$ balls are then randomly removed one at a time. Let $Y_k$ denote the number of red balls in the first $k$ selections, $k = 1, \dots, n$.
    (a) Find $P\{Y_n = j\}, j = 0, \dots, n$.
    (b) Find $P\{Y_{n-1} = j\}, j = 0, \dots, n-1$.
    (c) What do you think is the value of $P\{Y_k = j\}, j = 0, \dots, k$?
    (d) Verify your answer to part (c) by a backwards induction argument. That is, check that your answer is correct when $k = n$, and then show that whenever it is true for $k$ it is also true for $k - 1, k = 1, \dots, n$.
\end{exercise}

\solution{
    (a) $Y_n$ represents the total number of red balls originally placed in the urn since all $n$ balls are eventually removed. Therefore, its probability distribution is identical to the prior distribution given:
    \[
    P\{Y_n = j\} = \frac{1}{n+1}, \quad j = 0, 1, \dots, n
    \]

    (b) We can condition on the total number of red balls $Y_n$:
    \[
    P\{Y_{n-1} = j\} = \sum_{i=0}^{n} P\{Y_{n-1} = j \mid Y_n = i\} P\{Y_n = i\}
    \]
    The conditional probability $P\{Y_{n-1} = j \mid Y_n = i\}$ is only non-zero if $i = j$ (the last ball removed was non-red) or $i = j+1$ (the last ball removed was red):
    \[
    P\{Y_{n-1} = j\} = \left(\frac{n-j}{n}\right)\left(\frac{1}{n+1}\right) + \left(\frac{j+1}{n}\right)\left(\frac{1}{n+1}\right) = \frac{n - j + j + 1}{n(n+1)} = \frac{n+1}{n(n+1)} = \frac{1}{n}
    \]

    (c) Following the pattern from parts (a) and (b), the distribution remains uniform for any sample size $k$:
    \[
    P\{Y_k = j\} = \frac{1}{k+1}, \quad j = 0, 1, \dots, k
    \]

    (d) \textbf{Proof by Backward Induction:}
    \par\noindent\textbullet\enspace \textbf{Base Case ($k=n$):} As verified in part (a), $P\{Y_n = j\} = \frac{1}{n+1}$, which matches the formula.
    \par\noindent\textbullet\enspace \textbf{Inductive Step:} Assume the hypothesis holds for $k$, meaning $P\{Y_k = m\} = \frac{1}{k+1}$. We show it must hold for $k-1$:
    \[
    P\{Y_{k-1} = j\} = P\{Y_{k-1} = j \mid Y_k = j\}P\{Y_k = j\} + P\{Y_{k-1} = j \mid Y_k = j+1\}P\{Y_k = j+1\}
    \]
    Given $Y_k$ balls containing $m$ red ones, the probability that the $k$-th ball removed was non-red is $\frac{k-j}{k}$, and the probability it was red is $\frac{j+1}{k}$:
    \[
    P\{Y_{k-1} = j\} = \left(\frac{k-j}{k}\right)\left(\frac{1}{k+1}\right) + \left(\frac{j+1}{k}\right)\left(\frac{1}{k+1}\right) = \frac{k - j + j + 1}{k(k+1)} = \frac{k+1}{k(k+1)} = \frac{1}{k}
    \]
    This matches our formula for $k-1$, completing the induction proof.
}

\begin{exercise}
\label{ex:conditioning-expectation-66}
    The opponents of soccer team A are of two types: either they are a class 1 or a class 2 team. The number of goals team A scores against a class $i$ opponent is a Poisson random variable with mean $\lambda_i$, where $\lambda_1 = 2, \lambda_2 = 3$. This weekend the team has two games against teams they are not very familiar with. Assuming that the first team they play is a class 1 team with probability 0.6 and the second is, independently of the class of the first team, a class 1 team with probability 0.3, determine
    (a) the expected number of goals team A will score this weekend.
    (b) the probability that team A will score a total of five goals.
\end{exercise}

\solution{
    Let $X_1$ and $X_2$ be the number of goals scored in the first and second matches, respectively. The total goals scored is $X = X_1 + X_2$.

    (a) Using the Law of Total Expectation for each match:
    \[
    E[X_1] = \lambda_1(0.6) + \lambda_2(0.4) = 2(0.6) + 3(0.4) = 1.2 + 1.2 = 2.4
    \]
    \[
    E[X_2] = \lambda_1(0.3) + \lambda_2(0.7) = 2(0.3) + 3(0.7) = 0.6 + 2.1 = 2.7
    \]
    By linearity of expectation:
    \[
    E[X] = E[X_1] + E[X_2] = 2.4 + 2.7 = 5.1
    \]

    (b) We condition on the four independent combinations of opponent classes $(i, j)$ for the two games:
    \begin{itemize}
        \item Class (1,1) with prob $0.6 \times 0.3 = 0.18$. $X \sim \text{Poisson}(\lambda_1 + \lambda_1 = 4)$
        \item Class (1,2) with prob $0.6 \times 0.7 = 0.42$. $X \sim \text{Poisson}(\lambda_1 + \lambda_2 = 5)$
        \item Class (2,1) with prob $0.4 \times 0.3 = 0.12$. $X \sim \text{Poisson}(\lambda_2 + \lambda_1 = 5)$
        \item Class (2,2) with prob $0.4 \times 0.7 = 0.28$. $X \sim \text{Poisson}(\lambda_2 + \lambda_2 = 6)$
    \end{itemize}
    Using the Poisson PMF $P(X = 5) = \frac{e^{-\lambda}\lambda^5}{5!}$:
    \[
    P(X = 5) = 0.18\left(\frac{e^{-4}4^5}{5!}\right) + (0.42 + 0.12)\left(\frac{e^{-5}5^5}{5!}\right) + 0.28\left(\frac{e^{-6}6^5}{5!}\right)
    \]
    \[
    P(X = 5) = 0.18\left(\frac{1024 e^{-4}}{120}\right) + 0.54\left(\frac{3125 e^{-5}}{120}\right) + 0.28\left(\frac{7776 e^{-6}}{120}\right)
    \]
}

\begin{exercise}
\label{ex:conditioning-expectation-67}
    *A coin having probability $p$ of coming up heads is continually flipped. Let $P_j(n)$ denote the probability that a run of $j$ successive heads occurs within the first $n$ flips.
    (a) Argue that
    \[
    P_j(n) = P_j(n - 1) + p^j(1 - p)[1 - P_j(n - j - 1)]
    \]
    (b) By conditioning on the first non-head to appear, derive another equation relating $P_j(n)$ to the quantities $P_j(n - k), k = 1, \dots, j$.
\end{exercise}

\solution{
    (a) The event that a run of $j$ successive heads occurs within the first $n$ flips can happen in two mutually exclusive ways:
    \par\noindent\textbf{1.}\enspace It already occurred within the first $n-1$ flips, which has probability $P_j(n-1)$.
    \par\noindent\textbf{2.}\enspace It occurs for the very first time ending exactly at the $n$-th flip. This requires that flips from $(n-j+1)$ to $n$ are all heads (probability $p^j$), the $(n-j)$-th flip is a tail (probability $1-p$), and no run of $j$ heads occurred in the first $n-j-1$ flips (probability $1 - P_j(n-j-1)$).
    Adding these probabilities gives the recurrence for $n\ge j+1$. The boundary values are $P_j(n)=0$ for $n<j$ and $P_j(j)=p^j$.

    (b) Let us condition on the position $k$ of the first tail (non-head) to appear ($k = 1, 2, \dots, j, \dots$):
    \par\noindent\textbullet\enspace If the first tail occurs at step $k \leqslant j$ (with probability $p^{k-1}(1-p)$), the sequence resets and we need a run of $j$ heads in the remaining $n-k$ flips.
    \par\noindent\textbullet\enspace If no tail occurs in the first $j$ flips (with probability $p^j$), the run of $j$ successive heads has already been achieved immediately.
    Combining these scenarios by the Law of Total Probability yields:
    \[
    P_j(n) = \sum_{k=1}^{j} p^{k-1}(1-p)P_j(n-k) + p^j
    \]
}

\begin{exercise}
\label{ex:conditioning-expectation-68}
    In a knockout tennis tournament of $2^n$ contestants, the players are paired and play a match. The losers depart, the remaining $2^{n-1}$ players are paired, and they play a match. This continues for $n$ rounds, after which a single player remains unbeaten and is declared the winner. Suppose that the contestants are numbered 1 through $2^n$, and that whenever two players contest a match, the lower numbered one wins with probability $p$. Also suppose that the pairings of the remaining players are always done at random so that all possible pairings for that round are equally likely.
    (a) What is the probability that player 1 wins the tournament?
    (b) What is the probability that player 2 wins the tournament?
\end{exercise}

\solution{
Assume independent match outcomes conditional on the players.
\textbf{(a)} Player $1$ is favoured against every opponent and must win
$n$ matches, so the probability is $p^n$.

\textbf{(b)} Let $b_n$ be player $2$'s winning probability with $2^n$
players. Put $q=1-p$. For $n=1$, $b_1=q$. In round one, player $2$
meets player $1$ with probability $1/(2^n-1)$. If that match is won,
player $2$ is the best remaining player, with remaining winning
probability $p^{n-1}$. Otherwise player $2$ wins the first match with
probability $p$, and player $1$ independently survives with probability
$p$. Thus, for $n\ge2$,
\[
b_n=\frac{q p^{n-1}}{2^n-1}
+\frac{2^n-2}{2^n-1}\,p\left(qp^{n-1}+p b_{n-1}\right).
\]
Induction simplifies this to
\[
b_n=p^n-\frac{p^{n-1}[(2p)^n-1]}{2^n-1}.
\]
In particular, for a two-player tournament this gives $1-p$, and for
$p=1/2$ it gives $2^{-n}$, as required by symmetry.
}

\begin{exercise}
\label{ex:conditioning-expectation-69}
    In the match problem, say that $(i, j), i < j$, is a pair if $i$ chooses $j$'s hat and $j$ chooses $i$'s hat.
    (a) Find the expected number of pairs.
    (b) Let $Q_n$ denote the probability that there are no pairs, and derive a recursive formula for $Q_n$ in terms of $Q_j, j < n$.
    (c) Use the recursion of part (b) to find $Q_8$.
\end{exercise}

\solution{
\textbf{(a)} For $n\ge2$, a specified pair forms a two-cycle with
probability $(n-2)!/n!=1/[n(n-1)]$. Summing over the $\binom n2$
pairs gives expected number $1/2$. For $n=0,1$, the expectation is zero.

\textbf{(b)} The length $j$ of the cycle containing person $1$ is
uniform on $\{1,\ldots,n\}$: there are
$\binom{n-1}{j-1}(j-1)!(n-j)!=(n-1)!$ permutations for each length.
Length $2$ is forbidden; for any other length, the remaining $n-j$
people must have no two-cycles. Therefore
\[
Q_0=1,\qquad
Q_n=\frac1n\sum_{\substack{1\le j\le n\\j\ne2}}Q_{n-j}.
\]
\textbf{(c)} Recursion gives
\[
Q_1=1,\quad Q_2=Q_3=\frac12,\quad Q_4=Q_5=\frac58,
\quad Q_6=Q_7=\frac{29}{48},\quad Q_8=\frac{233}{384}\approx0.606771.
\]
}

\begin{exercise}
\label{ex:conditioning-expectation-70}
    Let $N$ denote the number of cycles that result in the match problem.
    (a) Let $M_n = E[N]$, and derive an equation for $M_n$ in terms of $M_1, \dots, M_{n-1}$.
    (b) Let $C_j$ denote the size of the cycle that contains person $j$. Argue that
    \[
    N = \sum_{j=1}^{n} 1/C_j
    \]
    and use the preceding to determine $E[N]$.
    (c) Find the probability that persons $1, 2, \dots, k$ are all in the same cycle.
    (d) Find the probability that $1, 2, \dots, k$ is a cycle.
\end{exercise}

\solution{
    (a) By conditioning on the size of the cycle containing person 1, which is equally likely to be of any length $k \in \{1, \dots, n\}$:
    \[
    M_n = 1 + \frac{1}{n} \sum_{i=1}^{n-1} M_i
    \]

    (b) Each cycle of length $L$ contributes exactly $L$ terms of value $1/L$ to the summation, which sums up to exactly $1$ for that cycle. Summing over all elements thus counts the total number of distinct cycles. Since $C_j$ is uniformly distributed over $\{1, \dots, n\}$, $E[1/C_j] = \frac{1}{n}\sum_{m=1}^{n}\frac{1}{m}$. Summing across all $n$ individuals yields the harmonic number:
    \[
    E[N] = \sum_{m=1}^{n} \frac{1}{m}
    \]

    (c) Delete all labels other than $1,\ldots,k$ from the cycle notation of a uniform permutation. The induced permutation on these $k$ labels is uniform, and the specified labels share a cycle exactly when that induced permutation is one cycle. There are $(k-1)!$ such permutations out of $k!$, giving:
    \[
    P(\text{All in same cycle}) = \frac{1}{k}
    \]

    (d) For the specific elements $1, 2, \dots, k$ to form their own isolated cycle of length $k$, they must permute among themselves into a single cycle (with $(k-1)!$ options), while the remaining $n-k$ elements can permute freely in $(n-k)!$ ways:
    \[
    P(\{1, \dots, k\} \text{ is a cycle}) = \frac{(k-1)!(n-k)!}{n!} = \frac{1}{k \binom{n}{k}}
    \]
}

\begin{exercise}
\label{ex:conditioning-expectation-71}
Let $M_0=0$ and suppose that, for $n\ge1$,
\[
M_n=1+\frac1n\sum_{i=0}^{n-1}M_i.
\]
By subtracting suitable equations for $n$ and $n-1$, show that
$M_n=\sum_{i=1}^n1/i$.
\end{exercise}

\solution{
    Multiply the stated recurrence by $n$:
    \[
    n M_n = n + \sum_{i=1}^{n-1} M_i
    \]
    Next, we write down the equivalent relation for $n-1$:
    \[
    (n-1) M_{n-1} = (n-1) + \sum_{i=1}^{n-2} M_i
    \]
    Subtracting the second equation from the first isolates the final term of the summation:
    \[
    n M_n - (n-1) M_{n-1} = 1 + M_{n-1}
    \]
    \[
    n M_n = n M_{n-1} + 1 \implies M_n = M_{n-1} + \frac{1}{n}
    \]
    Expanding this recurrence relation downwards directly yields the harmonic-number formula:
    \[
    M_n = \sum_{i=1}^{n} \frac{1}{i}
    \]
}

\begin{exercise}
\label{ex:conditioning-expectation-72}
Let $U_1,U_2,\ldots$ be independent uniform random variables on $(0,1)$,
and define
\[
N=\min\{n\ge2:U_n>U_{n-1}\},\qquad
M=\min\{n\ge1:U_1+\cdots+U_n>1\}.
\]
For $0<y<1$, show that the conditional distribution of $N$ given
$U_1=y$ equals that of $M$ given $U_1=1-y$, and show that both
conditional expectations equal $1+e^y$.
\end{exercise}

\solution{
Given $U_1=y$, for $n\ge1$ the event $N>n$ means
$y>U_2>\cdots>U_n$, up to zero-probability ties. Its conditional
probability is $y^{n-1}/(n-1)!$: all $n-1$ values lie below $y$,
and one of their $(n-1)!$ orders is decreasing.
Given $U_1=1-y$, the event $M>n$ means
$U_2+\cdots+U_n\le y$. For $0\le y\le1$, this simplex has volume
$y^{n-1}/(n-1)!$. Thus the two conditional survival functions agree,
which proves equality of the conditional distributions. The tail sum gives
\[
E[N\mid U_1=y]=E[M\mid U_1=1-y]
=1+\sum_{n=1}^{\infty}\frac{y^{n-1}}{(n-1)!}=1+e^y.
\]
}

\setcounter{exercise}{75}

\begin{exercise}
\label{ex:conditioning-expectation-76}
    Consider a gambler who on each bet either wins 1 with probability 18/38 or loses 1 with probability 20/38. (These are the probabilities if the bet is that a roulette wheel will land on a specified color.) The gambler will quit either when he or she is winning a total of 5 or after 100 plays. What is the probability he or she plays exactly 15 times?
\end{exercise}

\solution{
    For the gambler to stop playing exactly at the 15th play, one of the two stopping conditions must be reached at exactly $n = 15$. Since the 100 plays limit is far away, the game must end because the gambler reaches a total net gain of $+5$ for the very first time.

    Let $W$ be the number of wins and $L$ be the number of losses in the 15 plays. We have the system of equations:
    \[
    \begin{cases}
    W + L = 15 \\
    W - L = 5
    \end{cases}
    \]
    Adding the equations gives $2W = 20 \implies W = 10$, which implies $L = 5$.

    For the target of $+5$ to be reached for the \textit{very first time} on the 15th play, the 15th play must be a win, and the gambler must never have reached a net profit of $+5$ during the first 14 plays.
    \par\noindent\textbullet\enspace The first 14 plays must contain exactly 9 wins and 5 losses.
    \par\noindent\textbullet\enspace By the reflection principle (or ballot problem argument), the number of valid path sequences that never touch the barrier $+5$ before step 15 is given by:
    \[
    \binom{14}{9} - \binom{14}{10} = 2002 - 1001 = 1001
    \]
    Each specific sequence of 10 wins and 5 losses has an independent probability of $p^{10} q^5$, where $p = \frac{18}{38} = \frac{9}{19}$ and $q = \frac{20}{38} = \frac{10}{19}$.

    Therefore, the final probability is:
    \[
    P(\text{Plays exactly 15 times}) = 1001 \left(\frac{9}{19}\right)^{10} \left(\frac{10}{19}\right)^5
    \]
}

\begin{exercise}
\label{ex:conditioning-expectation-77}
    Show that
    (a) $E[XY \mid Y = y] = yE[X \mid Y = y]$
    (b) $E[g(X, Y) \mid Y = y] = E[g(X, y) \mid Y = y]$
    (c) $E[XY] = E[YE[X \mid Y]]$
\end{exercise}

\solution{
    (a) Given the condition $Y = y$, the random variable $Y$ behaves deterministically as a constant value $y$ inside the conditional expectation space. By the linearity property of expectations, constants can be factored outside the operator:
    \[
    E[XY \mid Y = y] = E[X \cdot y \mid Y = y] = yE[X \mid Y = y]
    \]

    (b) Since we are evaluating the expectation restricted entirely to the subset where $Y$ takes the specific fixed value $y$, replacing the random variable argument $Y$ with its conditional numeric value $y$ inside the function $g(X,Y)$ preserves pointwise equivalence over the support:
    \[
    E[g(X, Y) \mid Y = y] = E[g(X, y) \mid Y = y]
    \]

    (c) By applying the Law of Total Expectation (or Adam's Law), we can evaluate the global expectation of $XY$ by first conditioning on $Y$:
    \[
    E[XY] = E[\,E[XY \mid Y]\,]
    \]
    Using the pulling-out property established in part (a), where $E[XY \mid Y] = YE[X \mid Y]$ since $Y$ is known relative to itself, substituting this into the outer expectation yields:
    \[
    E[XY] = E[Y E[X \mid Y]]
    \]
}

\setcounter{exercise}{78}

\begin{exercise}
\label{ex:conditioning-expectation-79}
    An urn contains $n$ white and $m$ black balls that are removed one at a time. If $n > m$, show that the probability that there are always more white than black balls in the urn (until, of course, the urn is empty) equals $(n - m)/(n + m)$. Explain why this probability is equal to the probability that the set of withdrawn balls always contains more white than black balls. (This latter probability is $(n - m)/(n + m)$ by the ballot problem.)
\end{exercise}

\solution{
    Let us look at the sampling process in reverse order. Let the sequence of balls remaining in the urn be viewed from the final ball left to the full initial urn composition.

    Equivalently, consider the sequence of removals backwards from step $(n+m)$ down to 1. If we reverse the timeline, the condition that "the urn always contains more white than black balls throughout the removal process" translates precisely to stating that if we count the balls added back into an empty urn one by one, the number of white balls added must strictly exceed the number of black balls added at every single step.

    By reversing the sequence of extractions, the process maps directly onto Bertrand's ballot problem where:
    \begin{itemize}
        \item Candidate $A$ (White balls) receives a total of $n$ votes.
        \item Candidate $B$ (Black balls) receives a total of $m$ votes.
        \item Votes are drawn uniformly at random without replacement.
    \end{itemize}
    According to the classical Ballot Theorem, the probability that Candidate $A$ stays strictly ahead of Candidate $B$ throughout the entire counting sequence is:
    \[
    P = \frac{n - m}{n + m}
    \]
    Since every full removal path has a one-to-one symmetric mapping to a forward withdrawal sequence preserving the property structural inequality at each index, the probability that the set of withdrawn balls always contains more white than black balls is identically $(n-m)/(n+m)$.
}

\setcounter{exercise}{81}

\begin{exercise}
\label{ex:conditioning-expectation-82}
Let $X_1, X_2, \dots$ be independent continuous random variables with a common distribution function $F$ and density $f = F'$, and for $k \ge 1$ let
\[
N_k = \min\{n \ge k : X_n = k\text{th largest of } X_1, \dots, X_n\}
\]
(a) Show that $P\{N_k = n\} = \frac{k-1}{n(n-1)}$, $n \ge k$.\\
(b) Argue that
\[
f_{X_{N_k}}(x) = f(x)(\bar{F}(x))^{k-1} \sum_{i=0}^{\infty} \binom{i+k-2}{i} (F(x))^i
\]
(c) Prove the following identity:
\[
a^{1-k} = \sum_{i=0}^{\infty} \binom{i+k-2}{i} (1-a)^i, \quad 0 < a < 1, \, k \ge 2
\]
\textbf{Hint:} Use induction. First prove it when $k = 2$, and then assume it for $k$. To prove it for $k + 1$, use the fact that
\[
\sum_{i=1}^{\infty} \binom{i+k-1}{i}(1-a)^i = \sum_{i=1}^{\infty} \binom{i+k-2}{i}(1-a)^i + \sum_{i=1}^{\infty} \binom{i+k-2}{i-1}(1-a)^i
\]
where the preceding used the combinatorial identity $\binom{m}{i} = \binom{m-1}{i} + \binom{m-1}{i-1}$. Now, use the induction hypothesis to evaluate the first term on the right side of the preceding equation.\\
(d) Conclude that $X_{N_k}$ has distribution $F$.
\end{exercise}

\solution{
For $k=1$, $N_1=1$ and the selected value has distribution $F$
immediately. Assume $k\ge2$ below.

\textbf{(a)} In an i.i.d. continuous sequence, the successive insertion
ranks are independent, with the rank at time $j$ uniform on
$\{1,\ldots,j\}$. This follows by inserting successive labels into a
uniform random permutation. Therefore, for $n\ge k$,
\[
P(N_k=n)=\left[\prod_{j=k}^{n-1}\left(1-\frac1j\right)\right]\frac1n
=\frac{k-1}{n(n-1)}.
\]
\textbf{(b)} The stopping event depends only on ranks and is independent
of the ordered sample values. Conditional on $N_k=n$, $X_n$ is the
$k$th largest of $n$ observations, with density
\[
\frac{n!}{(k-1)!(n-k)!}\,f(x)\overline F(x)^{k-1}F(x)^{n-k}.
\]
Multiplying by the probability in (a), summing over $n$, and writing
$i=n-k$ gives
\[
f_{X_{N_k}}(x)=f(x)\overline F(x)^{k-1}
\sum_{i=0}^{\infty}\binom{i+k-2}{i}F(x)^i.
\]
\textbf{(c)} Put $z=1-a$ and
$S_k(z)=\sum_{i\ge0}\binom{i+k-2}{i}z^i$.
For $k=2$ this is the geometric series $1/(1-z)=a^{-1}$.
Pascal's identity and reindexing give
$S_{k+1}(z)=S_k(z)+zS_{k+1}(z)$, so
$S_{k+1}(z)=S_k(z)/(1-z)=a^{-k}$. The series converge absolutely for
$0<z<1$ (their coefficients grow polynomially).

\textbf{(d)} Substitute $a=\overline F(x)$ into (c). The factors cancel,
leaving $f_{X_{N_k}}(x)=f(x)$ almost everywhere, hence distribution $F$.
}

\begin{exercise}
\label{ex:conditioning-expectation-83}
An urn contains $n$ balls, with ball $i$ having weight $w_i, i = 1, \dots, n$. The balls are withdrawn from the urn one at a time according to the following scheme: When $S$ is the set of balls that remains, ball $i, i \in S$, is the next ball withdrawn with probability $w_i / \sum_{j \in S} w_j$. Find the expected number of balls that are withdrawn before ball $i, i = 1, \dots, n$.
\end{exercise}

\solution{
For each pair of balls $j \neq i$, let us define an indicator variable $I_j$ which takes the value 1 if ball $j$ is withdrawn before ball $i$, and 0 otherwise. The total number of balls selected before ball $i$ is represented by the random variable $X = \sum_{j \neq i} I_j$.
By the linearity of expectation, we have $E[X] = \sum_{j \neq i} E[I_j] = \sum_{j \neq i} P\{\text{ball } j \text{ is withdrawn before ball } i\}$.

Looking at the isolated subset $\{i, j\}$, the relative probability that $j$ is chosen before $i$ at any point during the sampling process depends solely on their respective weights scaled to their partial sum:
\[
P\{\text{ball } j \text{ is withdrawn before ball } i\} = \frac{w_j}{w_i + w_j}
\]
Consequently, the expected number of balls removed from the urn before encountering ball $i$ is equal to:
\[
E[\text{Number of balls before } i] = \sum_{j \neq i} \frac{w_j}{w_i + w_j}
\]
}

\setcounter{exercise}{85}

\begin{exercise}
\label{ex:conditioning-expectation-86}
On five labelled vertices, each vertex independently chooses one of the five vertices uniformly as its directed-arc destination; self-loops are allowed. A component means a weakly connected component of this random mapping. Compute the probability distribution of the number of components and verify your solution by using it to compute $E[C]$ and then comparing your solution with
\[
E[C] = \sum_{k=1}^{5} \binom{5}{k} \frac{(k-1)!}{5^k}
\]
\end{exercise}

\solution{
The stated model is a uniform random mapping on five
labelled vertices, allowing self-loops: each vertex independently
chooses one of the five vertices as its destination. Components are weak
components, each containing one directed cycle.

Let $c_j$ be the number of connected mappings on $j$ labelled vertices.
Choose the length $\ell$ of the unique cycle. The cycle can be chosen
and ordered in $\binom j\ell(\ell-1)!$ ways; the remaining vertices
form a forest rooted on that cycle, with $\ell j^{j-\ell-1}$ choices
when $\ell<j$, and one choice when $\ell=j$. This gives
\[
(c_1,c_2,c_3,c_4,c_5)=(1,3,17,142,1569).
\]
For $r$ components, sum $\prod_B c_{|B|}$ over set partitions into $r$
blocks. The resulting counts for $r=1,\ldots,5$ are
\[
1569,\quad
5(142)+10(3)(17)=1220,\quad
10(17)+15(3^2)=305,\quad
10(3)=30,\quad1.
\]
They sum to $5^5=3125$. Therefore
\[
P(C=r)=\frac1{3125}
\begin{cases}1569,&r=1,\\1220,&r=2,\\305,&r=3,\\30,&r=4,\\1,&r=5.\end{cases}
\]
The expectation is
\[
E[C]=\frac{1569+2(1220)+3(305)+4(30)+5}{3125}
=\frac{5049}{3125}=1.61568.
\]
The proposed check gives exactly
\[
\sum_{j=1}^5\binom5j\frac{(j-1)!}{5^j}
=1+\frac25+\frac4{25}+\frac6{125}+\frac{24}{3125}
=\frac{5049}{3125}.
\]
}

\setcounter{exercise}{90}

\begin{exercise}
\label{ex:conditioning-expectation-91}
Find the expected number of flips of a coin, which comes up heads with probability $p$, that are necessary to obtain the pattern $h, t, h, h, t, h, t, h$.
\end{exercise}

\solution{
Let $A$ be the target sequence $h, t, h, h, t, h, t, h$, with length $k = 8$. Let $q = 1-p$ denote the probability of getting tails. To find the expected number of trials $E[M]$ to observe this pattern, we apply the renewal theoretic method of overlapping patterns (or Markov chain embedding via martingale optional stopping).

We check all possible overlaps of the prefix of $A$ with its suffix:
\begin{itemize}
    \item Overlap of length 1: $h = h$ $\implies$ Valid overlap.
    \item Overlap of length 3: $h, t, h = h, t, h$ $\implies$ Valid overlap.
    \item Overlap of length 8: The entire string matches itself $\implies$ Always valid.
\end{itemize}
No other shifts yield a valid overlap match. Summing the reciprocal probabilities of the valid overlapping prefixes, we obtain:
\[
E[M] = \frac{1}{p^5 q^3} + \frac{1}{p^2 q} + \frac{1}{p}
\]
where $p^5 q^3$ corresponds to the full sequence (5 heads, 3 tails), $p^2 q$ corresponds to the suffix of length 3 ($h,t,h$), and $p$ corresponds to the suffix of length 1 ($h$).
}

\setcounter{exercise}{92}

\begin{exercise}
\label{ex:conditioning-expectation-93}
Consider a sequence of independent trials, each of which is equally likely to result in any of the outcomes $0, 1, \dots, m$. Say that a round begins with the first trial, and that a new round begins each time outcome $0$ occurs. Let $N$ denote the number of trials that it takes until all of the outcomes $1, \dots, m-1$ have occurred in the same round. Assume $m\ge2$. Outcome $m$ is neutral. Let $T_j$ denote the number of trials until $j$ distinct outcomes from $\{0,\ldots,m-1\}$ have occurred, and let $I_j$ be the $j$th distinct outcome from that relevant set. Neutral trials count toward elapsed time, but not toward discoveries.
(a) Argue that the random vectors $(I_1, \dots, I_m)$ and $(T_1, \dots, T_m)$ are independent.\\
(b) Define $X$ by letting $X = j$ if outcome $0$ is the $j$th distinct outcome to occur. Derive an equation for $E[N]$ in terms of $E[T_j], j = 1, \dots, m-1$ by conditioning on $X$.\\
(c) Determine $E[T_j], j = 1, \dots, m-1$.\\
(d) Find $E[N]$.
\end{exercise}

\solution{
Outcome $m$ is neutral: it neither supplies a required type nor starts
a new round. In this solution, use $T_j$ and $I_j$ for the times and
identities of the $j$ distinct \emph{relevant} outcomes from
$\{0,\ldots,m-1\}$; neutral trials still count toward elapsed time.
This makes explicit the interpretation needed for the requested
conditioning argument. Assume $m\ge2$.

\textbf{(a)} The order of relevant discoveries is a uniform permutation.
After $i-1$ discoveries, the probability of a new relevant outcome on
the next trial is $(m-i+1)/(m+1)$, independent of the identities already
seen. Conditional on a new relevant discovery, its identity is uniform
among the unseen relevant outcomes, independently of its waiting time.
Thus the discovery order and waiting-time vector are independent.

\textbf{(b)} The rank $X$ of outcome $0$ among the $m$ relevant types is
uniform on $\{1,\ldots,m\}$. If $X=j<m$, a reset occurs at time $T_j$.
If $X=m$, success occurs at $T_{m-1}$, before $0$ is observed. Hence
\[
E[N]=\frac1mE[T_{m-1}]
+\frac1m\sum_{j=1}^{m-1}(E[T_j]+E[N]),
\]
so $E[N]=E[T_{m-1}]+\sum_{j=1}^{m-1}E[T_j]$.

\textbf{(c)} Summing geometric waiting times gives
\[
E[T_j]=(m+1)\sum_{r=m-j+1}^m\frac1r,\qquad1\le j\le m-1.
\]
\textbf{(d)} Substitution and reversing the finite sum yield
\[
E[N]=(m+1)\sum_{r=2}^m\frac{(r-1)+1}{r}
=(m+1)(m-1).
\]
For $m=1$, no types are required and the natural convention is $N=0$.
}

\begin{exercise}
\label{ex:conditioning-expectation-94}
Let $N$ be a hypergeometric random variable having the distribution of the number of white balls in a random sample of size $r$ from a set of $w$ white and $b$ blue balls.
That is,

\[
P\{ N =n\} = \frac{\binom{w}{n} \binom{b}{r-n}}{\binom{w+b}{r}}
\]
where we use the convention $\binom{m}{j}=0 ~\text{if either}~ j<0 ~~or~~j>m. $

Assume $w\ge1$, $b\ge0$, and $1\le r\le w+b$. Consider $S_N=\sum_{i=1}^N X_i$, where the $X_i$ are i.i.d. positive integer-valued random variables, independent of $N$, with $\alpha_j = P\{X_i = j\}$.\\
(a) Define $M$ by the size-biased law $P(M=n)=nP(N=n)/E[N]$. Find the distribution of $M-1$.\\
(b) Suppressing its dependence on $b$, let $P_{w,r}(k) = P\{S_N = k\}$, and derive a recursion equation for $P_{w,r}(k)$.\\
(c) Use the recursion of (b) to find $P_{w,r}(2)$.
\end{exercise}

\solution{
Assume the $X_i$ are i.i.d. and independent of $N$. Write
$p_n=P(N=n)$ and $h=E[N]=rw/(w+b)$.
\textbf{(a)} Under the size-biased definition
$P(M=n)=np_n/h$, we have
\[
P(M-1=j)
=\frac{\binom{w-1}{j}\binom b{r-1-j}}
{\binom{w+b-1}{r-1}}.
\]
Thus $M-1$ is hypergeometric with sample size $r-1$ from $w-1$ white
and $b$ blue balls. This follows by using
$n\binom wn=w\binom{w-1}{n-1}$ in the size-biased mass function.

\textbf{(b)} The dependence on $b$ must be restored for a first-draw
recursion. Put $P_{w,b,r}(k)=P(S_N=k)$. Then, for $r\ge1$,
\[
P_{w,b,r}(k)=\frac{w}{w+b}\sum_{j=1}^{k}\alpha_jP_{w-1,b,r-1}(k-j)
+\frac{b}{w+b}P_{w,b-1,r-1}(k).
\]
The boundary value is $P_{w,b,0}(k)=\mathbf1_{\{k=0\}}$; terms whose
coefficient is zero are omitted. In particular,
$P_{w,b,r}(0)=\binom b r/\binom{w+b}r$.

\textbf{(c)} Since summands are positive integers, $S_N=2$ arises only
from $N=1,X_1=2$ or $N=2,X_1=X_2=1$. Therefore
\[
P_{w,b,r}(2)
=\frac{w\binom b{r-1}\alpha_2+
\binom w2\binom b{r-2}\alpha_1^2}{\binom{w+b}r}.
\]
}

\begin{exercise}
\label{ex:conditioning-expectation-95}
Let $X_1,X_2,\ldots$ be i.i.d. integer-valued steps with $X_i\ge-1$ and $-\infty<E[X_i]<0$, and let $S_0=0$, $S_n=\sum_{i=1}^n X_i$. Define $\beta = P\{S_n \le 0 \text{ for all } n\}$ be the probability that the walk is never positive. Find $\beta$ when $E[X_i] < 0$.
\end{exercise}

\solution{
Assume the steps are i.i.d. integer-valued, $X_i\ge-1$, and
$-\infty<E[X_i]<0$. Put $v=-E[X_i]$ and $q=P(X_i=-1)>0$.
The answer is
\[
\beta=\frac vq=\frac{-E[X_i]}{P(X_i=-1)}.
\]
To prove it, let $A_n=\{S_1<0,\ldots,S_n<0\}$. The cyclic ballot
lemma for integer steps bounded below by $-1$ gives
\[
P(A_n\mid S_n=-k)=\frac kn,\qquad k=1,\ldots,n.
\]
Indeed, among the $n$ cyclic starting positions of such a sequence,
exactly $k$ have all partial sums strictly negative; i.i.d. sequences
are invariant under these cyclic shifts. Consequently
\[
P(A_n)=E\left[\left(-\frac{S_n}{n}\right)^+\right].
\]
By the strong law, $S_n/n\to-v$ almost surely. Also
$0\le(-S_n/n)^+\le1$, so bounded convergence gives
$P(S_j<0\text{ for all }j\ge1)=v$.
To be strictly negative forever, the first step must be $-1$.
Thereafter, from position $-1$, the independent remaining walk must
never exceed $0$ relative to its own starting level. This has probability
$\beta$, so $v=q\beta$.

A left skip-free walk can jump over state $1$ on an upward move;
``ever positive'' cannot in general be replaced by ``ever hits $1$''.
}

\begin{exercise}
\label{ex:conditioning-expectation-96}
Consider a large population of families, and suppose that the number of children in the different families are independent Poisson random variables with mean $\lambda$. Show that the number of siblings of a randomly chosen child is also Poisson distributed with mean $\lambda$.
\end{exercise}

\solution{
Let $X$ be the number of children in a randomly chosen family, so $P\{X = n\} = \frac{e^{-\lambda}\lambda^n}{n!}$ for $n \ge 0$.
Sampling a child at random biases the selection towards larger families. The probability that a chosen child comes from a family of size $n$ is proportional to $n \cdot P\{X = n\}$. Normalizing by the expected family size $E[X] = \lambda$, the size of the chosen child's family, $Y$, has the distribution:
\[
P\{Y = n\} = \frac{n P\{X = n\}}{E[X]} = \frac{n \frac{e^{-\lambda}\lambda^n}{n!}}{\lambda} = \frac{e^{-\lambda}\lambda^{n-1}}{(n-1)!}, \quad n \ge 1
\]

The number of siblings of this child is given by $S = Y - 1$. Let us find the distribution of $S$ by substituting $k = n - 1$:
\[
P\{S = k\} = P\{Y = k + 1\} = \frac{e^{-\lambda}\lambda^{(k+1)-1}}{((k+1)-1)!} = \frac{e^{-\lambda}\lambda^k}{k!}, \quad k \ge 0
\]
This expression is identically the probability mass function of a Poisson random variable with mean $\lambda$. Thus, the number of siblings of a randomly chosen child is also Poisson distributed with mean $\lambda$.
}

\section{Distribution Recursions}
\label{sec:distribution-recursions}

\setcounter{exercise}{88}
\begin{exercise}
\label{ex:distribution-recursions-89}
Let $I_1,\ldots,I_n$ be independent random variables, each of which is equally likely to be either $0$ or $1$. A well-known nonparametric statistical test, called the signed rank test, is concerned with determining $P_n(k)$ defined by
\[
P_n(k)=P\left\{\sum_{j=1}^n jI_j\le k\right\}.
\]
Justify the following formula:
\[
P_n(k)=\frac12 P_{n-1}(k)+\frac12 P_{n-1}(k-n).
\]
\end{exercise}

\solution{
Let
\[
S_n=\sum_{j=1}^n jI_j.
\]
Then
\[
S_n=S_{n-1}+nI_n,
\]
where
\[
S_{n-1}=\sum_{j=1}^{n-1}jI_j.
\]

Since $I_n$ is equally likely to be $0$ or $1$, we condition on the value of $I_n$:
\[
P_n(k)
=
P(S_n\le k).
\]
Thus,
\[
P_n(k)
=
P(S_n\le k\mid I_n=0)P(I_n=0)
+
P(S_n\le k\mid I_n=1)P(I_n=1).
\]

Since
\[
P(I_n=0)=P(I_n=1)=\frac12,
\]
we get
\[
P_n(k)
=
\frac12 P(S_{n-1}\le k)
+
\frac12 P(S_{n-1}+n\le k).
\]
The first term is
\[
P(S_{n-1}\le k)=P_{n-1}(k),
\]
and the second term is
\[
P(S_{n-1}+n\le k)=P(S_{n-1}\le k-n)=P_{n-1}(k-n).
\]
Therefore,
\[
P_n(k)=\frac12 P_{n-1}(k)+\frac12 P_{n-1}(k-n).
\]
}

\section{Monte Carlo and Simulation Methods}
\label{sec:simulation}

\setcounter{exercise}{4}
\begin{exercise}
\label{ex:simulation-5}
Suppose it is relatively easy to simulate from $F_i$ for each $i=1,\ldots,n$. How can we simulate from
\begin{enumerate}
\item $F(x)=\prod_{i=1}^n F_i(x)$?
\item $F(x)=1-\prod_{i=1}^n (1-F_i(x))$?
\item Give two methods for simulating from the distribution $F(x)=x^n$, $0<x<1$.
\end{enumerate}
\end{exercise}

\solution{
\textbf{(a)} Generate independent random variables
\[
X_1,\ldots,X_n,
\]
where
\[
X_i\sim F_i.
\]
Then set
\[
X=\max\{X_1,\ldots,X_n\}.
\]
Indeed,
\[
P(X\le x)
=
P(X_1\le x,\ldots,X_n\le x)
=
\prod_{i=1}^n F_i(x).
\]
Thus $X$ has distribution function
\[
F(x)=\prod_{i=1}^n F_i(x).
\]

\textbf{(b)} Generate independent random variables
\[
X_1,\ldots,X_n,
\]
where
\[
X_i\sim F_i.
\]
Then set
\[
X=\min\{X_1,\ldots,X_n\}.
\]
Then
\[
P(X>x)
=
P(X_1>x,\ldots,X_n>x)
=
\prod_{i=1}^n (1-F_i(x)).
\]
Therefore,
\[
P(X\le x)
=
1-P(X>x)
=
1-\prod_{i=1}^n (1-F_i(x)).
\]

\textbf{(c)} We want to simulate from
\[
F(x)=x^n,\qquad 0<x<1.
\]

\noindent
\textbf{Method 1: Maximum of uniforms.}

Generate independent random variables
\[
U_1,\ldots,U_n\sim \operatorname{Uniform}(0,1),
\]
and set
\[
X=\max\{U_1,\ldots,U_n\}.
\]
Then, for $0<x<1$,
\[
P(X\le x)
=
P(U_1\le x,\ldots,U_n\le x)
=
x^n.
\]

\noindent
\textbf{Method 2: Inverse transform method.}

Generate
\[
U\sim \operatorname{Uniform}(0,1),
\]
and set
\[
X=U^{1/n}.
\]
Then, for $0<x<1$,
\[
P(X\le x)
=
P(U^{1/n}\le x)
=
P(U\le x^n)
=
x^n.
\]
Hence $X$ has distribution function $F(x)=x^n$.
}

\setcounter{exercise}{5}
\begin{exercise}
\label{ex:simulation-6}
Consider simulating the absolute value of a standard normal random variable using rejection sampling with an exponential proposal. For a positive rate $\lambda$, the proposal density is
\[
g(x)=\lambda e^{-\lambda x}.
\]
Show that the mean number of iterations needed in the rejection scheme is minimized when $\lambda=1$.
\end{exercise}

\solution{
The density of the absolute value of a standard normal random variable is
\[
f(x)=\sqrt{\frac{2}{\pi}}e^{-x^2/2},
\qquad x>0.
\]
We use the exponential proposal density
\[
g(x)=\lambda e^{-\lambda x},
\qquad x>0.
\]

In the rejection method, the mean number of iterations is the rejection constant
\[
c_\lambda=\sup_{x>0}\frac{f(x)}{g(x)}.
\]
Now,
\[
\frac{f(x)}{g(x)}
=
\frac{\sqrt{\frac{2}{\pi}}e^{-x^2/2}}{\lambda e^{-\lambda x}}
=
\frac{\sqrt{\frac{2}{\pi}}}{\lambda}
e^{-x^2/2+\lambda x}.
\]
The exponent is
\[
-\frac{x^2}{2}+\lambda x.
\]
This is maximized at
\[
x=\lambda.
\]
Therefore,
\[
c_\lambda
=
\frac{\sqrt{\frac{2}{\pi}}}{\lambda}
e^{\lambda^2/2}.
\]

To minimize $c_\lambda$, minimize
\[
\log c_\lambda
=
\text{constant}-\log \lambda+\frac{\lambda^2}{2}.
\]
Differentiate:
\[
\frac{d}{d\lambda}\log c_\lambda
=
-\frac1\lambda+\lambda.
\]
Set this equal to $0$:
\[
-\frac1\lambda+\lambda=0.
\]
Hence
\[
\lambda^2=1.
\]
Since $\lambda>0$, we get
\[
\lambda=1.
\]
Also,
\[
\frac{d^2}{d\lambda^2}\log c_\lambda
=
\frac1{\lambda^2}+1>0,
\]
so this critical point is a minimum.

Therefore, the mean number of iterations in the rejection scheme is minimized when
\[
{\lambda=1.}
\]
}

\setcounter{exercise}{6}
\begin{exercise}
\label{ex:simulation-7}
Give an algorithm for simulating a random variable having density function
\[
f(x)=30(x^2-2x^3+x^4),\qquad 0<x<1.
\]
\end{exercise}

\solution{
First note that
\[
f(x)=30x^2(1-x)^2,
\qquad 0<x<1.
\]
This is the density of a Beta random variable with parameters $3$ and $3$.

A simple way to simulate such a random variable is to use order statistics.

Generate independent random variables
\[
U_1,U_2,U_3,U_4,U_5\sim \operatorname{Uniform}(0,1).
\]
Let
\[
U_{(1)}\le U_{(2)}\le U_{(3)}\le U_{(4)}\le U_{(5)}
\]
be their order statistics. Then set
\[
X=U_{(3)}.
\]

The density of the third order statistic from five independent uniforms is
\[
f_{U_{(3)}}(x)
=
\frac{5!}{2!2!}x^2(1-x)^2,
\qquad 0<x<1.
\]
Since
\[
\frac{5!}{2!2!}=30,
\]
we have
\[
f_{U_{(3)}}(x)=30x^2(1-x)^2.
\]
Thus the algorithm is:

\[
{
\text{Generate five independent uniforms and return their median.}
}
\]
}

\setcounter{exercise}{7}
\begin{exercise}
\label{ex:simulation-8}
Consider the technique of simulating a gamma $(n,\lambda)$ random variable by using the rejection method with $g$ being an exponential density with rate $\lambda/n$.
\begin{enumerate}
\item Show that the average number of iterations of the algorithm needed to generate a gamma is
\[
\frac{n^n e^{1-n}}{(n-1)!}.
\]
\item Use Stirling's approximation to show that for large $n$ the answer to part (a) is approximately equal to
\[
e\left[\frac{n-1}{2\pi}\right]^{1/2}.
\]
\item Show that the procedure is equivalent to the following:
\[
\text{Step 1: Generate }Y_1\text{ and }Y_2\text{ independent exponentials with rate }1.
\]
\[
\text{Step 2: If }Y_1<(n-1)[Y_2-\log(Y_2)-1],\text{ return to step 1.}
\]
\[
\text{Step 3: Set }X=\frac{nY_2}{\lambda}.
\]
\item Explain how to obtain an independent exponential along with a gamma from the preceding algorithm.
\end{enumerate}
\end{exercise}

\solution{
\textbf{(a)} The gamma density with parameters $n$ and $\lambda$ is
\[
f(x)=\frac{\lambda^n x^{n-1}e^{-\lambda x}}{(n-1)!},
\qquad x>0.
\]
The proposal density is exponential with rate $\lambda/n$, so
\[
g(x)=\frac{\lambda}{n}e^{-\lambda x/n},
\qquad x>0.
\]
In the rejection method, the average number of iterations is
\[
c=\sup_{x>0}\frac{f(x)}{g(x)}.
\]
Now
\[
\frac{f(x)}{g(x)}
=
\frac{\lambda^n x^{n-1}e^{-\lambda x}}{(n-1)!}
\cdot
\frac{n}{\lambda}e^{\lambda x/n}.
\]
Thus
\[
\frac{f(x)}{g(x)}
=
\frac{n\lambda^{n-1}x^{n-1}}{(n-1)!}
e^{-\lambda x(1-1/n)}.
\]
Let
\[
y=\frac{\lambda x}{n}.
\]
Then
\[
x=\frac{ny}{\lambda},
\]
and so
\[
\frac{f(x)}{g(x)}
=
\frac{n^n}{(n-1)!}y^{n-1}e^{-(n-1)y}.
\]
The function
\[
y^{n-1}e^{-(n-1)y}
\]
is maximized at
\[
y=1.
\]
Therefore,
\[
c
=
\frac{n^n}{(n-1)!}e^{-(n-1)}
=
\frac{n^n e^{1-n}}{(n-1)!}.
\]
Hence the average number of iterations is
\[
{
\frac{n^n e^{1-n}}{(n-1)!}
}.
\]

\textbf{(b)} By Stirling's approximation,
\[
(n-1)!\sim \sqrt{2\pi(n-1)}
\left(\frac{n-1}{e}\right)^{n-1}.
\]
Therefore,
\[
\frac{n^n e^{1-n}}{(n-1)!}
\sim
\frac{n^n e^{1-n}}
{\sqrt{2\pi(n-1)}
\left(\frac{n-1}{e}\right)^{n-1}}.
\]
Since
\[
\left(\frac{n-1}{e}\right)^{n-1}
=
(n-1)^{n-1}e^{-(n-1)},
\]
we get
\[
\frac{n^n e^{1-n}}{(n-1)!}
\sim
\frac{n^n}
{\sqrt{2\pi(n-1)}(n-1)^{n-1}}.
\]
Thus
\[
\frac{n^n e^{1-n}}{(n-1)!}
\sim
\frac{n}{\sqrt{2\pi(n-1)}}
\left(\frac{n}{n-1}\right)^{n-1}.
\]
For large $n$,
\[
\left(\frac{n}{n-1}\right)^{n-1}\to e,
\]
and
\[
\frac{n}{\sqrt{n-1}}\sim \sqrt{n-1}.
\]
Therefore,
\[
\frac{n^n e^{1-n}}{(n-1)!}
\approx
e\left[\frac{n-1}{2\pi}\right]^{1/2}.
\]

\textbf{(c)} In the rejection method, generate
\[
X\sim g,
\]
where $g$ is exponential with rate $\lambda/n$. Equivalently, generate
\[
Y_2\sim \operatorname{Exponential}(1)
\]
and set
\[
X=\frac{nY_2}{\lambda}.
\]

The acceptance probability is
\[
\frac{f(X)}{cg(X)}.
\]
From part (a),
\[
\frac{f(X)}{g(X)}
=
\frac{n^n}{(n-1)!}Y_2^{n-1}e^{-(n-1)Y_2},
\]
and
\[
c=\frac{n^n e^{1-n}}{(n-1)!}.
\]
Hence
\[
\frac{f(X)}{cg(X)}
=
Y_2^{n-1}e^{-(n-1)Y_2}e^{n-1}.
\]
Therefore,
\[
\frac{f(X)}{cg(X)}
=
\exp\{(n-1)(\log Y_2-Y_2+1)\}.
\]

In the usual rejection method, if $U\sim \operatorname{Uniform}(0,1)$, we accept when
\[
U
\le
\exp\{(n-1)(\log Y_2-Y_2+1)\}.
\]
Let
\[
Y_1=-\log U.
\]
Then
\[
Y_1\sim \operatorname{Exponential}(1).
\]
Taking negative logarithms, the acceptance condition becomes
\[
Y_1\ge (n-1)(Y_2-\log Y_2-1).
\]
Therefore, the rejection condition is
\[
Y_1<(n-1)(Y_2-\log Y_2-1).
\]
Thus the rejection method is equivalent to:

\[
\text{Generate }Y_1,Y_2\text{ independent exponentials with rate }1.
\]
Reject and start again if
\[
Y_1<(n-1)(Y_2-\log Y_2-1).
\]
Otherwise set
\[
X=\frac{nY_2}{\lambda}.
\]

\textbf{(d)} Suppose the algorithm accepts a pair $(Y_1,Y_2)$. Define
\[
h(Y_2)=(n-1)(Y_2-\log Y_2-1).
\]
The acceptance condition is
\[
Y_1\ge h(Y_2).
\]
By the memoryless property of the exponential distribution,
\[
Y_1-h(Y_2)\mid \{Y_1\ge h(Y_2)\}
\]
has an exponential distribution with rate $1$.

Thus, after acceptance, define
\[
Z=Y_1-h(Y_2).
\]
Then
\[
Z\sim \operatorname{Exponential}(1),
\]
and $Z$ is independent of the accepted value of $Y_2$, hence independent of
\[
X=\frac{nY_2}{\lambda}.
\]

Therefore, the algorithm produces both
\[
X\sim \operatorname{Gamma}(n,\lambda)
\]
and an independent exponential random variable
\[
Z\sim \operatorname{Exponential}(1).
\]
If an exponential with rate $\lambda$ is desired, use
\[
\frac{Z}{\lambda}.
\]
}

\setcounter{exercise}{12}
\begin{exercise}
\label{ex:simulation-13}
The Discrete Rejection Method. Suppose we want to simulate $X$ having probability mass function
\[
P\{X=i\}=P_i,\qquad i=1,\ldots,n,
\]
and suppose we can easily simulate from the probability mass function $Q_i$, where
\[
\sum_i Q_i=1,\qquad Q_i\ge 0.
\]
Let $C$ be such that
\[
P_i\le CQ_i,\qquad i=1,\ldots,n.
\]
Show that the following algorithm generates the desired random variable:

\[
\text{Step 1: Generate }Y\text{ having mass function }Q\text{ and }U\text{ an independent random number.}
\]
\[
\text{Step 2: If }U\le \frac{P_Y}{CQ_Y},\text{ set }X=Y.
\text{ Otherwise return to step 1.}
\]
\end{exercise}

\solution{
For proposal states with $Q_i>0$, assume that
\[
U\sim \operatorname{Uniform}(0,1)
\]
and that $U$ is independent of $Y$.

Let $A$ denote the event that the proposed value is accepted. Conditional on $Y=i$, the probability of acceptance is
\[
P(A\mid Y=i)=\frac{P_i}{CQ_i}.
\]
Therefore,
\[
P(A,Y=i)
=
P(A\mid Y=i)P(Y=i)
=
\frac{P_i}{CQ_i}Q_i
=
\frac{P_i}{C}.
\]
Thus the total probability of acceptance is
\[
P(A)
=
\sum_{i=1}^n P(A,Y=i)
=
\sum_{i=1}^n \frac{P_i}{C}
=
\frac{1}{C}\sum_{i=1}^n P_i
=
\frac1C.
\]

Now compute the conditional distribution of the accepted value:
\[
P(Y=i\mid A)
=
\frac{P(A,Y=i)}{P(A)}
=
\frac{\frac{P_i}{C}}{\frac1C}
=
P_i.
\]
States with $Q_i=0$ also have $P_i=0$ and are never proposed, so no acceptance ratio is evaluated there. Hence the distribution of the accepted value is exactly
\[
P\{X=i\}=P_i,\qquad i=1,\ldots,n.
\]
Therefore, the algorithm generates a random variable having the desired probability mass function.
}

\setcounter{exercise}{15}
\begin{exercise}
\label{ex:simulation-16}
Suppose $n$ balls having weights $w_1,w_2,\ldots,w_n$ are in an urn. These balls are sequentially removed in the following manner: At each selection, a given ball in the urn is chosen with probability equal to its weight divided by the sum of the weights of all balls still in the urn (including the chosen ball). Let $I_1,I_2,\ldots,I_n$ denote the order in which the balls are removed, so that $I_1,\ldots,I_n$ is a random permutation with weights.
\begin{enumerate}
\item Give a method for simulating $I_1,\ldots,I_n$.
\item Let $X_i$ be independent exponentials with rates $w_i$, $i=1,\ldots,n$. Explain how $X_i$ can be utilized to simulate $I_1,\ldots,I_n$.
\end{enumerate}
\end{exercise}

\solution{
\textbf{(a)} Assume all weights $w_i>0$. A direct simulation method is as follows.

At the first step, choose ball $i$ with probability
\[
\frac{w_i}{w_1+\cdots+w_n}.
\]
Remove that ball from the urn.

At a later step, if the remaining set of balls is $S$, choose ball $i\in S$ with probability
\[
\frac{w_i}{\sum_{j\in S}w_j}.
\]
Remove the chosen ball and continue until all balls have been removed.

The sequence of selected indices is
\[
I_1,I_2,\ldots,I_n.
\]

Equivalently, at each step one may generate a uniform random variable and use the cumulative probabilities
\[
\frac{w_i}{\sum_{j\in S}w_j},\qquad i\in S,
\]
to choose the next ball.

\textbf{(b)} Generate independent exponential random variables
\[
X_i\sim \operatorname{Exponential}(w_i),
\qquad i=1,\ldots,n.
\]
Then order them from smallest to largest:
\[
X_{I_1}<X_{I_2}<\cdots<X_{I_n}.
\]
The indices
\[
I_1,I_2,\ldots,I_n
\]
give the desired weighted random permutation.

To justify this, for any remaining set $S$ and any $i\in S$,
\[
P(X_i=\min_{j\in S}X_j)
=
\int_0^\infty
w_i e^{-w_it}
\prod_{\substack{j\in S\\ j\ne i}}
e^{-w_jt}\,dt.
\]
Hence
\[
P(X_i=\min_{j\in S}X_j)
=
\int_0^\infty
w_i e^{-\left(\sum_{j\in S}w_j\right)t}\,dt
=
\frac{w_i}{\sum_{j\in S}w_j}.
\]
Thus the first selected ball has the correct weighted probability.

After the smallest exponential time occurs, the memoryless property of exponentials implies that the remaining residual times are still independent exponentials with the same rates. Therefore, the same argument applies repeatedly. Hence ordering independent exponentials with rates $w_i$ simulates the desired sequence
\[
I_1,\ldots,I_n.
\]
}

\setcounter{exercise}{22}
\begin{exercise}
\label{ex:simulation-23}
For a nonhomogeneous Poisson process with intensity function $\lambda(t)$, $t\ge 0$, where
\[
\int_0^\infty \lambda(t)\,dt=\infty,
\]
let $X_1,X_2,\ldots$ denote the sequence of times at which events occur.
\begin{enumerate}
\item Show that
\[
\int_0^{X_1}\lambda(t)\,dt
\]
is exponential with rate $1$.
\item Show that
\[
\int_{X_{i-1}}^{X_i}\lambda(t)\,dt,\qquad i\ge 1,
\]
are independent exponentials with rate $1$, where $X_0=0$.
\end{enumerate}
In words, independent of the past, the additional amount of hazard that must be experienced until an event occurs is exponential with rate $1$.
\end{exercise}

\solution{
Assume $\lambda$ is nonnegative and locally integrable, and put
$\Lambda(t)=\int_0^t\lambda(s)\,ds$. It is continuous, nondecreasing,
and tends to infinity. Let $\tau(y)=\inf\{t:\Lambda(t)\ge y\}$.
For fixed hazard times $0\le a<b$, the count in the corresponding
physical-time interval is Poisson with mean
$\Lambda(\tau(b))-\Lambda(\tau(a))=b-a$. Disjoint hazard intervals
correspond to disjoint physical intervals, so their counts are independent.
Flat portions of $\Lambda$ have zero intensity and contain no events.
Thus the time-changed process is a rate-$1$ Poisson process.

\textbf{(a)} Its first arrival occurs at hazard time $\Lambda(X_1)$,
so $P(\Lambda(X_1)>y)=e^{-y}$ for $y\ge0$.

\textbf{(b)} Its successive interarrival times are
\[
\Lambda(X_i)-\Lambda(X_{i-1})
=\int_{X_{i-1}}^{X_i}\lambda(t)\,dt.
\]
As interarrival times of a rate-$1$ Poisson process, these are independent
exponentials with rate $1$. This argument uses deterministic transformed
intervals; independence on random intervals does not follow directly
from independent increments alone.
}

\setcounter{exercise}{24}
\begin{exercise}
\label{ex:simulation-25}
Let $(X,Y)$ be uniformly distributed in a circle of radius $r$ about the origin. That is, their joint density is given by
\[
f(x,y)=\frac1{\pi r^2},
\qquad 0\le x^2+y^2\le r^2.
\]
Let
\[
R=\sqrt{X^2+Y^2}
\qquad\text{and}\qquad
\theta=\operatorname{atan2}(Y,X)\pmod{2\pi}
\]
denote their polar coordinates. Show that $R$ and $\theta$ are independent with $\theta$ being uniform on $(0,2\pi)$ and
\[
P\{R<a\}=\frac{a^2}{r^2},
\qquad 0<a<r.
\]
\end{exercise}

\solution{
Use the polar coordinate transformation
\[
X=R\cos\theta,
\qquad
Y=R\sin\theta.
\]
The Jacobian of this transformation is
\[
\left|\frac{\partial(x,y)}{\partial(r,\theta)}\right|=R.
\]
Therefore, the joint density of $R$ and $\theta$ is
\[
f_{R,\theta}(\rho,\theta)
=
\frac1{\pi r^2}\rho,
\]
for
\[
0<\rho<r,\qquad 0<\theta<2\pi.
\]
Thus
\[
f_{R,\theta}(\rho,\theta)
=
\frac{\rho}{\pi r^2}.
\]
We can factor this as
\[
f_{R,\theta}(\rho,\theta)
=
\left(\frac{2\rho}{r^2}\right)
\left(\frac1{2\pi}\right),
\qquad 0<\rho<r,\quad 0<\theta<2\pi.
\]
Hence
\[
f_R(\rho)=\frac{2\rho}{r^2},
\qquad 0<\rho<r,
\]
and
\[
f_\theta(\theta)=\frac1{2\pi},
\qquad 0<\theta<2\pi.
\]
Since the joint density factors into a function of $\rho$ times a function of $\theta$, $R$ and $\theta$ are independent.

Also, $\theta$ is uniform on $(0,2\pi)$.

Finally, for $0<a<r$,
\[
P(R<a)
=
\int_0^a \frac{2\rho}{r^2}\,d\rho
=
\frac{2}{r^2}\cdot \frac{a^2}{2}
=
\frac{a^2}{r^2}.
\]
Therefore,
\[
{
P\{R<a\}=\frac{a^2}{r^2},\qquad 0<a<r.
}
\]
}

\setcounter{exercise}{30}
\begin{exercise}
\label{ex:simulation-31}
Consider a queueing system in which each service time, independent of the past, has mean $\mu$. Let $W_n$ and $D_n$ denote, respectively, the amounts of time customer $n$ spends in the system and in queue. Hence,
\[
D_n=W_n-S_n,
\]
where $S_n$ is the service time of customer $n$. Therefore,
\[
E[D_n]=E[W_n]-\mu.
\]
If we use simulation to estimate $E[D_n]$, should we
\begin{enumerate}
\item use the simulated data to determine $D_n$, which is then used as an estimate of $E[D_n]$, or
\item use the simulated data to determine $W_n$ and then use this quantity minus $\mu$ as an estimate of $E[D_n]$?
\end{enumerate}
Repeat for when we want to estimate $E[W_n]$.
\end{exercise}

\solution{
Since
\[
W_n=D_n+S_n,
\]
and the service time $S_n$ is independent of the past, $S_n$ is independent of $D_n$. Hence
\[
\operatorname{Cov}(D_n,S_n)=0.
\]
Therefore,
\[
\operatorname{Var}(W_n)
=
\operatorname{Var}(D_n+S_n)
=
\operatorname{Var}(D_n)+\operatorname{Var}(S_n).
\]
Thus,
\[
\operatorname{Var}(W_n)\ge \operatorname{Var}(D_n).
\]

\noindent
\textbf{Estimating $E[D_n]$}

There are two possible unbiased estimators:
\[
D_n
\]
or
\[
W_n-\mu.
\]
Since subtracting a constant does not change variance,
\[
\operatorname{Var}(W_n-\mu)=\operatorname{Var}(W_n).
\]
But
\[
\operatorname{Var}(W_n)\ge \operatorname{Var}(D_n).
\]
Therefore, it is better to estimate $E[D_n]$ directly using the simulated values of $D_n$.

\[
{\text{Use the simulated values of }D_n.}
\]

\noindent
\textbf{Estimating $E[W_n]$}

Now we can estimate $E[W_n]$ either by using
\[
W_n
\]
directly, or by using
\[
D_n+\mu.
\]
Both are unbiased because
\[
E[W_n]=E[D_n]+\mu.
\]
However,
\[
\operatorname{Var}(D_n+\mu)=\operatorname{Var}(D_n),
\]
whereas
\[
\operatorname{Var}(W_n)=\operatorname{Var}(D_n)+\operatorname{Var}(S_n).
\]
Thus
\[
\operatorname{Var}(D_n+\mu)\le \operatorname{Var}(W_n).
\]
Therefore, it is better to estimate $E[W_n]$ by using the simulated values of $D_n$ and then adding $\mu$.

\[
{\text{Use }D_n+\mu\text{ to estimate }E[W_n].}
\]
}

\begin{exercise}
\label{ex:simulation-32}
Show that if $X$ and $Y$ have the same distribution, then
\[
\operatorname{Var}\left(\frac{X+Y}{2}\right)\le \operatorname{Var}(X).
\]
Hence, conclude that the use of antithetic variables can never increase variance, though it need not be as efficient as generating an independent set of random numbers.
\end{exercise}

\solution{
Since $X$ and $Y$ have the same distribution,
\[
\operatorname{Var}(X)=\operatorname{Var}(Y).
\]
Now,
\[
\operatorname{Var}\left(\frac{X+Y}{2}\right)
=
\frac14\operatorname{Var}(X+Y).
\]
Expanding,
\[
\operatorname{Var}\left(\frac{X+Y}{2}\right)
=
\frac14\left[\operatorname{Var}(X)+\operatorname{Var}(Y)+2\operatorname{Cov}(X,Y)\right].
\]
Since $\operatorname{Var}(X)=\operatorname{Var}(Y)$,
\[
\operatorname{Var}\left(\frac{X+Y}{2}\right)
=
\frac12\operatorname{Var}(X)+\frac12\operatorname{Cov}(X,Y).
\]
By Cauchy's inequality,
\[
\operatorname{Cov}(X,Y)\le
\sqrt{\operatorname{Var}(X)\operatorname{Var}(Y)}
=
\operatorname{Var}(X).
\]
Therefore,
\[
\operatorname{Var}\left(\frac{X+Y}{2}\right)
\le
\frac12\operatorname{Var}(X)+\frac12\operatorname{Var}(X)
=
\operatorname{Var}(X).
\]

Hence,
\[
{
\operatorname{Var}\left(\frac{X+Y}{2}\right)\le \operatorname{Var}(X).
}
\]

Thus, averaging antithetic variables cannot increase the variance compared with using just one of them. However, it need not be as efficient as using two independent random variables. If $X_1$ and $X_2$ are independent with the same distribution as $X$, then
\[
\operatorname{Var}\left(\frac{X_1+X_2}{2}\right)
=
\frac12\operatorname{Var}(X).
\]
For antithetic variables,
\[
\operatorname{Var}\left(\frac{X+Y}{2}\right)
=
\frac12\operatorname{Var}(X)+\frac12\operatorname{Cov}(X,Y),
\]
which may be larger than $\frac12\operatorname{Var}(X)$ if $\operatorname{Cov}(X,Y)>0$.
}

\begin{exercise}
\label{ex:simulation-33}
If $0\le X\le a$, show that
\begin{enumerate}
\item $E[X^2]\le aE[X]$,
\item $\operatorname{Var}(X)\le E[X](a-E[X])$,
\item $\operatorname{Var}(X)\le a^2/4$.
\end{enumerate}
\end{exercise}

\solution{
Since
\[
0\le X\le a,
\]
we have
\[
X^2\le aX.
\]
Taking expectations gives
\[
E[X^2]\le aE[X].
\]
This proves part (a).

For part (b), recall that
\[
\operatorname{Var}(X)=E[X^2]-(E[X])^2.
\]
Using part (a),
\[
\operatorname{Var}(X)
\le
aE[X]-(E[X])^2.
\]
Therefore,
\[
\operatorname{Var}(X)
\le
E[X](a-E[X]).
\]

For part (c), let
\[
m=E[X].
\]
Since $0\le X\le a$, we have
\[
0\le m\le a.
\]
From part (b),
\[
\operatorname{Var}(X)\le m(a-m).
\]
The quadratic function
\[
m(a-m)
\]
is maximized on $0\le m\le a$ at
\[
m=\frac a2.
\]
Thus,
\[
m(a-m)\le \frac a2\left(a-\frac a2\right)=\frac{a^2}{4}.
\]
Hence,
\[
{\operatorname{Var}(X)\le \frac{a^2}{4}.}
\]
}

\setcounter{exercise}{34}
\begin{exercise}
\label{ex:simulation-35}
Suppose we are able to simulate independent random variables $X$ and $Y$. If we simulate $2k$ independent random variables
\[
X_1,\ldots,X_k
\]
and
\[
Y_1,\ldots,Y_k,
\]
where the $X_i$ have the same distribution as does $X$ and the $Y_j$ have the same distribution as does $Y$, how would you use them to estimate
\[
P(X<Y)?
\]
\end{exercise}

\solution{
Use all possible pairs $(X_i,Y_j)$ rather than only the $k$ pairs $(X_i,Y_i)$.

Define
\[
I_{ij}=
\begin{cases}
1, & X_i<Y_j,\\
0, & X_i\ge Y_j.
\end{cases}
\]
Then
\[
E[I_{ij}]=P(X_i<Y_j).
\]
Since $X_i$ has the same distribution as $X$, $Y_j$ has the same distribution as $Y$, and they are independent,
\[
P(X_i<Y_j)=P(X<Y).
\]
Therefore,
\[
E[I_{ij}]=P(X<Y).
\]

A natural estimator is
\[
\widehat{p}
=
\frac{1}{k^2}\sum_{i=1}^k\sum_{j=1}^k I_{ij}.
\]
That is,
\[
{
\widehat{p}
=
\frac{1}{k^2}
\sum_{i=1}^k\sum_{j=1}^k
\mathbf{1}\{X_i<Y_j\}.
}
\]
This estimator is unbiased because
\[
E[\widehat{p}]
=
\frac{1}{k^2}
\sum_{i=1}^k\sum_{j=1}^k
E[I_{ij}]
=
\frac{1}{k^2}
\sum_{i=1}^k\sum_{j=1}^k
P(X<Y)
=
P(X<Y).
\]
}

\section{Markov Chains}
\label{sec:markov-chains}

\textbf{Notation and Setup}
Throughout, $P = (P_{ij})$ denotes a transition probability matrix,
$\pi = (\pi_j)$ the stationary distribution, and $P^{(n)}_{ij}$ the
$n$-step transition probability.  The following models and state orderings are used in this chapter:
\begin{itemize}
  \item \textbf{Gary's mood:} States $0=C$ (cheerful),
    $1=S$ (so-so), $2=G$ (glum).
    \[
      P=\begin{pmatrix}0.5&0.4&0.1\\0.3&0.4&0.3\\0.2&0.3&0.5\end{pmatrix}.
    \]
  \item \textbf{Rain depends on the last two days:} States
    $0=(RR)$, $1=(RN)$, $2=(NR)$, $3=(NN)$ (R = rain, N = no rain,
    ordered today/yesterday).
    \[
      P=\begin{pmatrix}0.7&0&0.3&0\\0.5&0&0.5&0\\0&0.4&0&0.6\\0&0.2&0&0.8\end{pmatrix}.
    \]
  \item \textbf{Two-ball urn in Exercise 9:} State = \# red balls.
    \[
      P=\begin{pmatrix}0.8&0.2&0\\0.1&0.8&0.1\\0&0.2&0.8\end{pmatrix}.
    \]
  \item \textbf{Random walk:}
    $P_{i,i+1}=p=1-P_{i,i-1}$, $i\in\mathbb{Z}$.
\end{itemize}

\setcounter{exercise}{0}

\begin{exercise}
\label{ex:markov-chains-1}
Three white and three black balls are distributed in two urns so that
each contains three balls.  The system is in state $i$, $i=0,1,2,3$, if
the first urn contains $i$ white balls.  At each step, one ball is drawn
from each urn and the two balls are swapped.  Let $X_n$ be the state after
the $n$-th step.  Explain why $\{X_n\}$ is a Markov chain and calculate
its transition probability matrix.
\end{exercise}
\solution{
\textbf{Markov property.} The state $X_n = i$ fully determines the
composition of both urns (urn~1 has $i$ white and $3-i$ black; urn~2 has
$3-i$ white and $i$ black).  Because the draw at step $n+1$ depends only
on the current ball counts, the future evolution depends only on $X_n$,
not on the earlier history.

\textbf{Transition probabilities.} From state $i$:
\begin{itemize}
  \item \emph{State decreases by 1} ($i \to i-1$): we draw a white ball
    from urn~1 (prob $i/3$) and a black ball from urn~2 (prob $i/3$):
    \[P_{i,i-1}=\frac{i^2}{9},\qquad i=1,2,3.\]
  \item \emph{State increases by 1} ($i \to i+1$): we draw black from
    urn~1 (prob $(3-i)/3$) and white from urn~2 (prob $(3-i)/3$):
    \[P_{i,i+1}=\frac{(3-i)^2}{9},\qquad i=0,1,2.\]
  \item \emph{State unchanged}: same colour drawn from both urns:
    $P_{i,i}=1-P_{i,i-1}-P_{i,i+1}=\dfrac{2i(3-i)}{9}$.
\end{itemize}
The transition matrix (states $0,1,2,3$) is therefore
\[
P=\begin{pmatrix}
0 & 1 & 0 & 0 \\[2pt]
\tfrac{1}{9} & \tfrac{4}{9} & \tfrac{4}{9} & 0 \\[2pt]
0 & \tfrac{4}{9} & \tfrac{4}{9} & \tfrac{1}{9} \\[2pt]
0 & 0 & 1 & 0
\end{pmatrix}.
\]
(Boundary check: from state $0$ all balls in urn~1 are black so we
necessarily draw black from urn~1 and white from urn~2, giving $P_{0,1}=1$;
similarly $P_{3,2}=1$.)
}

\begin{exercise}
\label{ex:markov-chains-2}
Suppose that whether or not it rains today depends on previous weather
conditions through the \emph{last three days}.  Show how this system may
be analysed by a Markov chain.  How many states are needed?
\end{exercise}
\solution{
Define the state at time $n$ as the triple
$(W_{n-2},W_{n-1},W_n)\in\{0,1\}^3$, where $1$ = rain and $0$ = no rain.
Because today's weather depends only on the last three days, and the
triple encodes exactly that information, the process on these triples
satisfies the Markov property.

There are $2^3 = \mathbf{8}$ states needed.
}

\begin{exercise}
\label{ex:markov-chains-3}
In Exercise~2, suppose: if it has rained for the past three days, it
will rain today with probability $0.8$; if it did not rain for any of the
past three days, it will rain today with probability $0.2$; in any other
case, the weather today will with probability $0.6$ be the same as
yesterday.  Determine $P$ for this Markov chain.
\end{exercise}
\solution{
Label states as binary triples $(a,b,c)$ ($a$ = two days ago, $b$ =
yesterday, $c$ = today) ordered $000,001,010,011,100,101,110,111$.
From state $(a,b,c)$ the chain moves to $(b,c,X)$ where
\[
P(X=1)=\begin{cases}0.8 & \text{if }a=b=c=1,\\0.2&\text{if }a=b=c=0,\\
0.6\cdot c + 0.4\cdot(1-c)&\text{otherwise.}\end{cases}
\]
Denoting $p_{abc}:=P(X=1\mid \text{state }(a,b,c))$:
\begin{align*}
p_{000}&=0.2,\quad p_{001}=0.6,\quad p_{010}=0.4,\quad p_{011}=0.6,\\
p_{100}&=0.4,\quad p_{101}=0.6,\quad p_{110}=0.4,\quad p_{111}=0.8.
\end{align*}
Each row of $P$ has exactly two non-zero entries (to states $(b,c,0)$ and
$(b,c,1)$).  The full $8\times 8$ matrix is ($q_{abc}=1-p_{abc}$):
\[
P=\bordermatrix{
 &000&001&010&011&100&101&110&111\cr
000&0.8&0.2&0&0&0&0&0&0\cr
001&0&0&0.4&0.6&0&0&0&0\cr
010&0&0&0&0&0.6&0.4&0&0\cr
011&0&0&0&0&0&0&0.4&0.6\cr
100&0.6&0.4&0&0&0&0&0&0\cr
101&0&0&0.4&0.6&0&0&0&0\cr
110&0&0&0&0&0.6&0.4&0&0\cr
111&0&0&0&0&0&0&0.2&0.8\cr
}.
\]
}

\begin{exercise}
\label{ex:markov-chains-4}
Consider $\{X_n,n\ge0\}$ taking values $0,1,2$, with
\[
P\{X_{n+1}=j\mid X_n=i,\ldots\}=
\begin{cases}P^{\mathrm{I}}_{ij},&n\text{ even},\\
             P^{\mathrm{II}}_{ij},&n\text{ odd,}\end{cases}
\]
where $\sum_j P^{\mathrm{I}}_{ij}=\sum_j P^{\mathrm{II}}_{ij}=1$.
Is $\{X_n\}$ a Markov chain? If not, show how to transform it into one
by enlarging the state space.
\end{exercise}
\solution{
The process satisfies the Markov property: conditional on $X_n$, the
next state is independent of earlier states. Its transition law depends
on $n$, so it is a \emph{time-inhomogeneous} Markov chain (unless the two
matrices coincide).

To obtain a time-homogeneous chain, define $Y_n=(X_n,n\bmod 2)$ on
$\{0,1,2\}\times\{0,1\}$. Its transitions are
\[
P\{Y_{n+1}=(j,1)\mid Y_n=(i,0)\}=P^{\mathrm I}_{ij},
\]
\[
P\{Y_{n+1}=(j,0)\mid Y_n=(i,1)\}=P^{\mathrm {II}}_{ij},
\]
with all other transitions zero. These probabilities do not depend on $n$.
}

\begin{exercise}
\label{ex:markov-chains-5}
A Markov chain $\{X_n,n\ge0\}$ with states $0,1,2$ has transition matrix
\[
P=\begin{pmatrix}\tfrac{1}{2}&\tfrac{1}{3}&\tfrac{1}{6}\\[2pt]
                  0&\tfrac{1}{3}&\tfrac{2}{3}\\[2pt]
                  \tfrac{1}{2}&0&\tfrac{1}{2}\end{pmatrix}.
\]
If $P\{X_0=0\}=P\{X_0=1\}=\tfrac{1}{4}$, find $\E[X_3]$.
\end{exercise}
\solution{
Compute $P^2$ then $P^3$ by matrix multiplication.

\textbf{$P^2$:}
\[
P^2=\begin{pmatrix}\tfrac{1}{3}&\tfrac{5}{18}&\tfrac{7}{18}\\[2pt]
     \tfrac{1}{3}&\tfrac{1}{9}&\tfrac{5}{9}\\[2pt]
     \tfrac{1}{2}&\tfrac{1}{6}&\tfrac{1}{3}\end{pmatrix}.
\]

\textbf{$P^3=P^2 P$:}
\[
P^3=\begin{pmatrix}
     \tfrac{39}{108}&\tfrac{22}{108}&\tfrac{47}{108}\\[4pt]
     \tfrac{4}{9}&\tfrac{4}{27}&\tfrac{11}{27}\\[4pt]
     \tfrac{5}{12}&\tfrac{2}{9}&\tfrac{13}{36}
     \end{pmatrix}.
\]

Since $P\{X_0=2\}=\tfrac{1}{2}$,
\begin{align*}
\E[X_3]
&=\tfrac{1}{4}\E[X_3\mid X_0=0]+\tfrac{1}{4}\E[X_3\mid X_0=1]
  +\tfrac{1}{2}\E[X_3\mid X_0=2]\\
&=\tfrac{1}{4}\!\cdot\!\frac{0\cdot39+1\cdot22+2\cdot47}{108}
 +\tfrac{1}{4}\!\cdot\!\frac{0\cdot4/9+1\cdot4+2\cdot11}{27}
 +\tfrac{1}{2}\!\cdot\!\frac{0\cdot5+1\cdot8+2\cdot13}{36}\\
&=\tfrac{1}{4}\cdot\frac{116}{108}
 +\tfrac{1}{4}\cdot\frac{26}{27}
 +\tfrac{1}{2}\cdot\frac{34}{36}\\
&=\frac{116}{432}+\frac{26}{108}+\frac{17}{36}
 =\frac{116}{432}+\frac{104}{432}+\frac{204}{432}
 =\frac{424}{432}
 ={\dfrac{53}{54}}.
\end{align*}
}

\begin{exercise}
\label{ex:markov-chains-6}
Let the transition matrix of a two-state Markov chain be
\[
\mathbf{P}=\begin{Vmatrix}p & 1-p \\ 1-p & p\end{Vmatrix}.
\]
Show by induction that
\[
\mathbf{P}^{(n)}=
\begin{Vmatrix}
\tfrac{1}{2}+\tfrac{1}{2}(2p-1)^n & \tfrac{1}{2}-\tfrac{1}{2}(2p-1)^n \\[4pt]
\tfrac{1}{2}-\tfrac{1}{2}(2p-1)^n & \tfrac{1}{2}+\tfrac{1}{2}(2p-1)^n
\end{Vmatrix}.
\]
\end{exercise}
\solution{
Let $c:=2p-1$.

\textbf{Base case $n=1$.} $P^{(1)}_{00}=p=\tfrac{1}{2}+\tfrac{c}{2}$ and
$P^{(1)}_{01}=1-p=\tfrac{1}{2}-\tfrac{c}{2}$. 

\textbf{Inductive step.} Assume the formula holds for $n$.  Then
\begin{align*}
P^{(n+1)}_{00}
&=P^{(n)}_{00}\cdot p + P^{(n)}_{01}\cdot(1-p)\\
&=\bigl(\tfrac{1}{2}+\tfrac{c^n}{2}\bigr)p
  +\bigl(\tfrac{1}{2}-\tfrac{c^n}{2}\bigr)(1-p)\\
&=\tfrac{1}{2}+c^n\!\left(\tfrac{p}{2}-\tfrac{1-p}{2}\right)
 =\tfrac{1}{2}+c^n\cdot\tfrac{c}{2}=\tfrac{1}{2}+\tfrac{c^{n+1}}{2}.
\end{align*}
The other three entries follow by an identical argument (or by symmetry),
completing the induction. 
}

\begin{exercise}
\label{ex:markov-chains-7}
In the rain model given in the chapter setup, suppose it has rained neither yesterday nor the day before
yesterday.  What is the probability that it will rain tomorrow?
\end{exercise}
\solution{
Use the state ordering in the setup: the first component is the more
recent day's weather. At the end of yesterday, the state was $NN$,
because both yesterday and the preceding day were dry. Tomorrow is
\emph{two} transitions later. Thus
\[
P(\text{rain tomorrow}\mid\text{last two days dry})
=(P^2)_{3,0}+(P^2)_{3,1}
=0.2(0.5)+0.8(0.2)=0.26.
\]
Here, today's weather has been averaged out rather than assumed known.
}

\begin{exercise}
\label{ex:markov-chains-8}
Coin~1 has $P(H)=0.7$ and coin~2 has $P(H)=0.6$.  If today's flip is
heads we use coin~1 tomorrow; if tails we use coin~2 tomorrow.  The
initial coin is equally likely to be coin~1 or coin~2.
\begin{enumerate}
  \item[(a)] What is the probability that the coin flipped on the third
    day after the initial flip is coin~1?
  \item[(b)] Suppose the coin flipped on Monday comes up heads.  What is
    the probability that the coin flipped on Friday of the same week also
    comes up heads?
\end{enumerate}
\end{exercise}
\solution{
Model the coin in use as a Markov chain with states $C_1,C_2$ and
transition matrix
\[
P=\begin{pmatrix}0.7&0.3\\0.6&0.4\end{pmatrix}.
\]
The stationary distribution satisfies $0.3\pi_1=0.6\pi_2$, giving
$\pi_1=\tfrac{2}{3}$, $\pi_2=\tfrac{1}{3}$.  The non-unit eigenvalue is $1-0.3-0.6=0.1$.

Hence
\[
P^n_{11}=\pi_1+(1-\pi_1)\cdot0.1^n=\tfrac{2}{3}+\tfrac{1}{3}\cdot0.1^n,
\quad
P^n_{21}=\pi_1-\pi_1\cdot0.1^n=\tfrac{2}{3}(1-0.1^n).
\]

\textbf{(a)} Three days after the initial flip ($n=3$):
\begin{align*}
P(X_3=C_1)
&=\tfrac{1}{2}P^3_{11}+\tfrac{1}{2}P^3_{21}
 =\tfrac{1}{2}\!\left[\tfrac{2}{3}+\tfrac{1}{3}(0.1)^3\right]
 +\tfrac{1}{2}\!\left[\tfrac{2}{3}(1-(0.1)^3)\right]\\
&=\tfrac{2}{3}+\tfrac{1}{6}(0.1)^3-\tfrac{1}{3}(0.1)^3
 =\tfrac{2}{3}-\tfrac{1}{6000}
 \approx 0.6665.
\end{align*}

\textbf{(b)} Monday heads $\Rightarrow$ Tuesday coin is $C_1$
(deterministically, since heads always selects $C_1$ next).  From Tuesday
to Friday is $3$ flips (Wed, Thu, Fri).  Let $X_{\mathrm{Tue}}=C_1$.
\begin{align*}
P(\text{Fri heads}\mid\text{Mon heads})
&=P^3_{11}\cdot0.7+P^3_{12}\cdot0.6\\
&=\!\left[\tfrac{2}{3}+\tfrac{1}{3}(0.1)^3\right]\!0.7
 +\!\left[\tfrac{1}{3}-\tfrac{1}{3}(0.1)^3\right]\!0.6\\
&=\tfrac{2}{3}\cdot0.7+\tfrac{1}{3}\cdot0.6+(0.1)^3\tfrac{0.1}{3}
 =\tfrac{14}{30}+\tfrac{6}{30}+\varepsilon
 =\tfrac{2}{3}+\frac{(0.1)^3}{30}\approx0.6667.
\end{align*}
}

\begin{exercise}
\label{ex:markov-chains-9}
An urn initially contains one red and one blue ball. At each step, select
one of its two balls uniformly, record its colour, and replace it. With
probability $0.8$ the replacement has the same colour; with probability
$0.2$ replace it by a ball of the opposite colour. These random choices
are independent between steps. Let $X_n=1$ if the $n$th selected ball is
red, and $X_n=0$ otherwise. Is $\{X_n:n\ge1\}$ a Markov chain?
\end{exercise}
\solution{
No. Let $C_n$ be the number of red balls after step $n$, with $C_0=1$.
Although $C_n$ is Markov, the observed colours retain information from
more than the most recent observation.

If $X_1=X_2=1$, the first red ball must have been returned red;
otherwise the second red selection would be impossible. Immediately
before step $2$ there is therefore one red ball. After observing that
red selection, the expected red count after replacement is $1-0.2$.
Consequently
\[
P(X_3=1\mid X_2=1,X_1=1)=\frac{1-0.2}{2}=\frac25.
\]
If $X_1=0$, then $C_1=1$ with probability $0.8$ and $C_1=2$ with
probability $0.2$. Conditioning further on $X_2=1$ gives
\[
P(C_1=1\mid X_1=0,X_2=1)=\frac{0.8(1/2)}{0.8(1/2)+0.2}=\frac23,
\]
and the complementary probability is $1/3$. Given $C_1=i$ and a red
selection, the next-red probability is $(i-0.2)/2$. Thus
\[
P(X_3=1\mid X_2=1,X_1=0)
=\frac23\frac25+\frac13\frac9{10}=\frac{17}{30}.
\]
Both conditioning histories have positive probability. They end in the
same current observation but give different next-step probabilities,
which disproves the Markov property for $X_n$.
}

\begin{exercise}
\label{ex:markov-chains-10}
In the mood model given in the chapter setup, Gary is currently in a cheerful mood.  What is the
probability that he is \emph{not} in a glum mood on any of the following
three days?
\end{exercise}
\solution{
We restrict attention to the sub-matrix $Q$ of states $\{C,S\}$ (avoiding
state $G$):
\[
Q=\begin{pmatrix}0.5&0.4\\0.3&0.4\end{pmatrix}.
\]
The desired probability equals the sum of the entries in row $C$ of $Q^3$
(probability of remaining in $\{C,S\}$ for three steps).

\textbf{$Q^2$:}
\[
Q^2=\begin{pmatrix}0.37&0.36\\0.27&0.28\end{pmatrix}.
\]

\textbf{$Q^3=Q^2Q$:}
\[
Q^3[C,C]=0.37\times0.5+0.36\times0.3=0.293,\quad
Q^3[C,S]=0.37\times0.4+0.36\times0.4=0.292.
\]

Therefore
\[
P(\text{not glum on days }1,2,3\mid X_0=C)=0.293+0.292={0.585}.
\]
}

\begin{exercise}
\label{ex:markov-chains-11}
In the mood model given in the chapter setup, Gary was in a glum mood four days ago.  Given that he
hasn't felt cheerful in a week, what is the probability he is feeling
glum today?
\end{exercise}
\solution{
``Hasn't felt cheerful in a week'' means the chain has stayed in
$\{S,G\}$ for the past four steps.  Define the restricted sub-matrix
\[
R=\begin{pmatrix}P_{SS}&P_{SG}\\P_{GS}&P_{GG}\end{pmatrix}
 =\begin{pmatrix}0.4&0.3\\0.3&0.5\end{pmatrix}.
\]
We need
\[
P(X_0=G\mid X_{-4}=G,\,X_{-3},\ldots,X_0\in\{S,G\})
=\frac{R^4_{G,G}}{R^4_{G,S}+R^4_{G,G}}.
\]

Compute $R^2=(R_{ij}^2)$:
\[R^2=\begin{pmatrix}0.25&0.27\\0.27&0.34\end{pmatrix}.\]

Then $R^4=(R^2)^2$:
\begin{align*}
R^4_{G,S}&=0.27\times0.25+0.34\times0.27=0.0675+0.0918=0.1593,\\
R^4_{G,G}&=0.27\times0.27+0.34\times0.34=0.0729+0.1156=0.1885.
\end{align*}

Therefore
\[
P(\text{glum today})=\frac{0.1885}{0.1593+0.1885}=\frac{0.1885}{0.3478}
\approx {0.542}.
\]
}

\setcounter{exercise}{13}
\begin{exercise}
\label{ex:markov-chains-14}
Specify the classes of the following Markov chains, and determine whether
they are transient or recurrent:
\[
P_1=\begin{Vmatrix}0&\tfrac{1}{2}&\tfrac{1}{2}\\\tfrac{1}{2}&0&\tfrac{1}{2}\\\tfrac{1}{2}&\tfrac{1}{2}&0\end{Vmatrix},\quad
P_2=\begin{Vmatrix}0&0&0&1\\0&0&0&1\\\tfrac{1}{2}&\tfrac{1}{2}&0&0\\0&0&1&0\end{Vmatrix},
\]
\[
P_3=\begin{Vmatrix}\tfrac{1}{2}&0&\tfrac{1}{2}&0&0\\\tfrac{1}{4}&\tfrac{1}{2}&\tfrac{1}{4}&0&0\\\tfrac{1}{2}&0&\tfrac{1}{2}&0&0\\0&0&0&\tfrac{1}{2}&\tfrac{1}{2}\\0&0&0&\tfrac{1}{2}&\tfrac{1}{2}\end{Vmatrix},\quad
P_4=\begin{Vmatrix}\tfrac{1}{4}&\tfrac{3}{4}&0&0&0\\\tfrac{1}{2}&\tfrac{1}{2}&0&0&0\\0&0&1&0&0\\0&0&\tfrac{1}{3}&\tfrac{2}{3}&0\\1&0&0&0&0\end{Vmatrix}.
\]
\end{exercise}
\solution{
Use states numbered from $0$.
\begin{itemize}
\item $P_1$: a single closed communicating class $\{0,1,2\}$.
All states are recurrent. Return paths of lengths $2$ and $3$ give
period $1$.
\item $P_2$: a single closed communicating class $\{0,1,2,3\}$.
All states are recurrent. Transitions cycle through the three groups
$\{0,1\}\to\{3\}\to\{2\}\to\{0,1\}$. Every return time is a
multiple of $3$, and a return in $3$ steps is possible, so the period is $3$.
\item $P_3$: the closed classes are $\{0,2\}$ and $\{3,4\}$,
and all their states are recurrent. The singleton class $\{1\}$ is
transient: it leaves for $\{0,2\}$ with probability $1/2$ per step,
and neither $0$ nor $2$ can reach it. In particular, $0$ and $1$
do not communicate.
\item $P_4$: $\{0,1\}$ and $\{2\}$ are closed recurrent classes.
The singleton classes $\{3\}$ and $\{4\}$ are transient, since they
can enter a closed class that cannot reach them again.
\end{itemize}
}

\setcounter{exercise}{15}
\begin{exercise}
\label{ex:markov-chains-16}
Show that if state $i$ is recurrent and state $i$ does not communicate
with state $j$, then $P_{ij}=0$.
\end{exercise}
\solution{
Since $i$ and $j$ do not communicate, either $i$ cannot reach $j$ or $j$
cannot reach $i$.  Suppose for contradiction that $P_{ij}>0$ (so $i$ can
reach $j$ in one step).  Then $j$ can be visited starting from $i$.
Because $i$ and $j$ do not communicate, $j$ cannot reach $i$.  Therefore,
once the chain visits $j$ it can never return to $i$.

Formally, let $f_{ii}=P(\text{ever return to }i\mid X_0=i)$.  By first-step
analysis
\[
f_{ii} \le 1 - P_{ij}\cdot P(\text{never reach }i\mid X_1=j)
       = 1 - P_{ij}\cdot 1 = 1 - P_{ij} < 1,
\]
where the last equality uses the fact that $j$ cannot reach $i$.  This
contradicts the recurrence of $i$ ($f_{ii}=1$).  Hence $P_{ij}=0$. 

\textit{Corollary.} Once a process enters a recurrent class it cannot
leave: for any $i$ in the class and any $j$ outside, $P_{ij}=0$.
}

\setcounter{exercise}{17}
\begin{exercise}
\label{ex:markov-chains-18}
Coin~1 comes up heads with probability $0.6$ and coin~2 with probability
$0.5$.  A coin is continually flipped until it comes up tails, at which
time that coin is put aside and we start flipping the other one.
\begin{enumerate}
\item[(a)] What proportion of flips use coin~1?
\item[(b)] If we start with coin~1, what is the probability that coin~2
  is used on the fifth flip?
\end{enumerate}
\end{exercise}
\solution{
Model the coin in use on each individual flip as a Markov chain with
states $\{C_1,C_2\}$.  On each flip with coin $i$:
\begin{itemize}
  \item Heads (coin unchanged): $P(C_1\!\to\!C_1)=0.6$, $P(C_2\!\to\!C_2)=0.5$.
  \item Tails (switch coin): $P(C_1\!\to\!C_2)=0.4$, $P(C_2\!\to\!C_1)=0.5$.
\end{itemize}

\textbf{(a)} Stationary distribution: $0.4\pi_1=0.5\pi_2$.  With
$\pi_1+\pi_2=1$: $\pi_1=\dfrac{5}{9}$, $\pi_2=\dfrac{4}{9}$.

The proportion of flips using coin~1 is ${\tfrac{5}{9}\approx55.6\%}$.

\textbf{(b)} Starting from $C_1$, the second eigenvalue is
$\lambda_2=1-0.4-0.5=0.1$.  Hence
\[
P^n_{C_1,C_2}=\pi_2\bigl(1-0.1^n\bigr)=\tfrac{4}{9}\bigl(1-0.1^n\bigr).
\]
For the fifth flip ($n=4$ transitions from the first flip):
\[
P(\text{5th flip uses }C_2\mid\text{start }C_1)
=\tfrac{4}{9}\!\left(1-(0.1)^4\right)
=\tfrac{4}{9}\cdot\tfrac{9999}{10000}
=\frac{4\times9999}{90000}\approx{0.4444}.
\]
}

\setcounter{exercise}{20}
\begin{exercise}
\label{ex:markov-chains-21}
A DNA nucleotide has any of four values.  In the standard mutation model,
from period to period a nucleotide at a specific location does not change
with probability $1-3\alpha$, and if it changes it is equally likely to
become any of the other three values ($0<\alpha<\tfrac{1}{3}$).
\begin{enumerate}
\item[(a)] Show that $P^n_{1,1}=\tfrac{1}{4}+\tfrac{3}{4}(1-4\alpha)^n$.
\item[(b)] What is the long-run proportion of time the chain is in each state?
\end{enumerate}
\end{exercise}
\solution{
\textbf{(a)} Let $J$ be the $4\times4$ all-ones matrix and put
$\Pi=J/4$. The transition matrix is
\[
P=\Pi+(1-4\alpha)(I-\Pi).
\]
Since $\Pi^2=\Pi$, $(I-\Pi)^2=I-\Pi$, and the two projections have
zero product,
\[
P^n=\Pi+(1-4\alpha)^n(I-\Pi).
\]
Taking the $(1,1)$ entry gives
\[
P^n_{1,1}=\frac14+\frac34(1-4\alpha)^n.
\]
\textbf{(b)} Because $0<\alpha<1/3$, we have $|1-4\alpha|<1$.
Thus $P^n\to\Pi$, and each state has long-run proportion $1/4$.
}

\setcounter{exercise}{22}
\begin{exercise}
\label{ex:markov-chains-23}
In a good year, storms are Poisson with mean $1$; in a bad year, Poisson
with mean $3$.  A good year is equally likely to be followed by a good or
bad year; a bad year is twice as likely to be followed by a bad year as by
a good year.  Last year (year~$0$) was good.
\begin{enumerate}
\item[(a)] Find the expected total number of storms in years $1$ and $2$.
\item[(b)] Find the probability of no storms in year~$3$.
\item[(c)] Find the long-run average storms per year.
\end{enumerate}
\end{exercise}
\solution{
\textbf{Weather Markov chain.} States: $G$ (good), $B$ (bad).
\[
P=\begin{pmatrix}1/2&1/2\\1/3&2/3\end{pmatrix}.
\]
Let $\mu(G)=1$, $\mu(B)=3$ denote expected storms.

\textbf{(a)} Given $X_0=G$: $P(X_1=G)=1/2$, $P(X_1=B)=1/2$, so
\[\E[\text{storms yr 1}]=\tfrac{1}{2}\cdot1+\tfrac{1}{2}\cdot3=2.\]
\[P^2_{G,G}=\tfrac{1}{2}\cdot\tfrac{1}{2}+\tfrac{1}{2}\cdot\tfrac{1}{3}=\tfrac{5}{12},\quad
P^2_{G,B}=\tfrac{7}{12}.\]
\[\E[\text{storms yr 2}]=\tfrac{5}{12}\cdot1+\tfrac{7}{12}\cdot3=\tfrac{5+21}{12}=\tfrac{26}{12}=\tfrac{13}{6}.\]
\[\E[\text{total storms yrs 1--2}]=2+\tfrac{13}{6}={\tfrac{25}{6}\approx4.17}.\]

\textbf{(b)} $P(N=0\mid\text{good})=e^{-1}$, $P(N=0\mid\text{bad})=e^{-3}$.
\[P^3_{G,G}=P^2_{G,G}\cdot\tfrac{1}{2}+P^2_{G,B}\cdot\tfrac{1}{3}
=\tfrac{5}{24}+\tfrac{7}{36}=\tfrac{15+14}{72}=\tfrac{29}{72},\quad
P^3_{G,B}=\tfrac{43}{72}.\]
\[P(\text{no storms yr 3})=\tfrac{29}{72}e^{-1}+\tfrac{43}{72}e^{-3}.\]

\textbf{(c)} Stationary distribution: $\tfrac{1}{2}\pi_G=\tfrac{1}{3}\pi_B$,
so $\pi_B=\tfrac{3}{2}\pi_G$; normalising gives $\pi_G=\tfrac{2}{5}$,
$\pi_B=\tfrac{3}{5}$.  Long-run average:
\[{\pi_G\cdot1+\pi_B\cdot3=\tfrac{2}{5}+\tfrac{9}{5}=\tfrac{11}{5}=2.2 \text{ storms/year.}}\]
}

\begin{exercise}
\label{ex:markov-chains-24}
Consider three urns: red ($1$ red, $4$ blue), white ($3$ white, $2$ red,
$2$ blue), blue ($4$ white, $3$ red, $2$ blue).  At each stage a ball is
selected from the urn whose colour matches the previously selected ball.
In the long run, what proportion of the selected balls are red / white /
blue?
\end{exercise}
\solution{
The state is the colour of the urn currently in use
($R$, $W$, $B$).  Transition matrix:
\[
P=\bordermatrix{&R&W&B\cr
  R&\tfrac{1}{5}&0&\tfrac{4}{5}\cr
  W&\tfrac{2}{7}&\tfrac{3}{7}&\tfrac{2}{7}\cr
  B&\tfrac{3}{9}&\tfrac{4}{9}&\tfrac{2}{9}\cr}.
\]
Stationary equations $\pi P=\pi$:

From column $W$: $(3/7)\pi_W+(4/9)\pi_B=\pi_W \;\Rightarrow\;
\pi_W=\tfrac{7}{9}\pi_B$.

From column $R$: $(1/5)\pi_R+(2/7)\pi_W+(1/3)\pi_B=\pi_R \;\Rightarrow\;
(4/5)\pi_R=(2/7)\cdot(7/9)\pi_B+(1/3)\pi_B=(2/9+3/9)\pi_B=(5/9)\pi_B
\;\Rightarrow\;\pi_R=\tfrac{25}{36}\pi_B$.

Normalising ($\pi_R+\pi_W+\pi_B=1$):
\[\tfrac{25}{36}\pi_B+\tfrac{7}{9}\pi_B+\pi_B
=\tfrac{25+28+36}{36}\pi_B=\tfrac{89}{36}\pi_B=1
\;\Rightarrow\;\pi_B=\tfrac{36}{89}.\]

\[{\pi_R=\frac{25}{89},\quad\pi_W=\frac{28}{89},\quad\pi_B=\frac{36}{89}.}\]

The long-run proportion of balls that are red is $25/89\approx28.1\%$,
white is $28/89\approx31.5\%$, and blue is $36/89\approx40.4\%$.
}

\begin{exercise}
\label{ex:markov-chains-25}
An individual owns $k$ pairs of running shoes.  Each morning he leaves by
either door (front or back) with equal probability $1/2$, taking shoes if
available; on return he enters by either door with equal probability $1/2$,
leaving shoes there.  What proportion of the time does he run barefooted?
\end{exercise}
\solution{
State $i=$ number of pairs at the front door, $i=0,1,\ldots,k$.

Each morning the individual:
\begin{itemize}
  \item Leaves front (prob $1/2$): if $i>0$, takes a pair (later
    returns it to front w.p.\ $1/2$ or back w.p.\ $1/2$); if $i=0$,
    goes barefoot (no shoes move).
  \item Leaves back (prob $1/2$): if $k-i>0$, takes a pair; if $k-i=0$,
    goes barefoot.
\end{itemize}

For $0<i<k$:
$P_{i,i-1}=\tfrac{1}{4}$, $P_{i,i}=\tfrac{1}{2}$, $P_{i,i+1}=\tfrac{1}{4}$.
At boundaries: $P_{0,0}=\tfrac{3}{4}$, $P_{0,1}=\tfrac{1}{4}$;
$P_{k,k-1}=\tfrac{1}{4}$, $P_{k,k}=\tfrac{3}{4}$.

This is a birth--death chain with equal birth and death rates
$\tfrac{1}{4}$ for interior states.  By detailed balance the stationary
distribution is \emph{uniform}:
\[\pi_i=\frac{1}{k+1},\quad i=0,1,\ldots,k.\]

Barefoot occurs if and only if he leaves front with $i=0$ shoes there, or
leaves back with $k-i=0$ shoes there:
\[
P(\text{barefoot})
=\pi_0\cdot\tfrac{1}{2}+\pi_k\cdot\tfrac{1}{2}
=\frac{1}{k+1}\cdot\frac{1}{2}+\frac{1}{k+1}\cdot\frac{1}{2}
={\frac{1}{k+1}}.
\]
}

\begin{exercise}
\label{ex:markov-chains-26}
Consider the following approach to shuffling a deck of $n$ cards: choose
one of the positions $1,2,\ldots,n$ uniformly at random and move the card
in that position to the top of the deck (``move to front'').  Show that
in the limit the deck is perfectly shuffled, i.e.\ every ordering is
equally likely.
\end{exercise}
\solution{
\textbf{Irreducibility.}  Starting from any permutation, we can reach any
target permutation $\sigma=(\sigma_1,\ldots,\sigma_n)$ by performing at
most $n$ MTF moves: first move $\sigma_n$ to position~1, then move
$\sigma_{n-1}$ to position~1 (pushing $\sigma_n$ down), and so on.
Hence the chain is irreducible over all $n!$ orderings.

\textbf{Uniform stationary distribution.}  We verify that the uniform
distribution $\pi(\sigma)=1/n!$ is stationary.  For any target ordering
$\tau$, the number of orderings $\sigma$ that can produce $\tau$ in one
MTF step equals $n$ (move any of the $n$ cards to front); each such
$\sigma$ is reached with probability $1/n$ from a uniform starting
distribution.  Since there are exactly $n$ orderings mapping to $\tau$
in one step, and each contributes $\frac{1}{n!}\cdot\frac{1}{n}$:
\[\sum_\sigma\pi(\sigma)P_{\sigma,\tau}=n\cdot\frac{1}{n!}\cdot\frac{1}{n}=\frac{1}{n!}=\pi(\tau).\]

The self-loop obtained by selecting the top card makes the chain aperiodic. By irreducibility and finiteness, the unique stationary distribution is
the uniform distribution, so every ordering becomes equally likely in the
limit. 
}

\begin{exercise}
\label{ex:markov-chains-27}
Each of $N$ individuals is either active or inactive each period.
An active individual stays active w.p.\ $\alpha$; an inactive individual
stays inactive w.p.\ $\beta$.  Let $X_n=$ number of active individuals in
period $n$.
\begin{enumerate}
\item[(a)] Argue that $X_n$ is a Markov chain.
\item[(b)] Find $\E[X_n\mid X_0=i]$.
\item[(c)] Derive an expression for the transition probabilities.
\item[(d)] Find the long-run proportion of time that exactly $j$ individuals
  are active.
\end{enumerate}
\end{exercise}
\solution{
Assume individuals evolve independently and $0<\alpha,\beta<1$.\par
\textbf{(a)}  Given $X_n=i$, the next-period count $X_{n+1}$ is the sum
of $i$ independent Bernoulli($\alpha$) variables (one per active
individual) plus $N-i$ independent Bernoulli($1-\beta$) variables (one
per inactive individual).  The distribution of $X_{n+1}$ depends only on
$i=X_n$, not on the history, so $\{X_n\}$ is a Markov chain.

\textbf{(b)}  Let $\mu_n=\E[X_n\mid X_0=i]$.  Since
$\E[X_{n+1}\mid X_n]=\alpha X_n+(1-\beta)(N-X_n)$, taking expectations:
\[\mu_{n+1}=(\alpha+\beta-1)\mu_n+(1-\beta)N.\]
This linear recurrence (with $r=\alpha+\beta-1$) solves to
\[
\E[X_n\mid X_0=i]=\frac{(1-\beta)N}{2-\alpha-\beta}
+\left(i-\frac{(1-\beta)N}{2-\alpha-\beta}\right)(\alpha+\beta-1)^n.
\]

\textbf{(c)}  From state $i$, the next state $j$ is the sum of
$\mathrm{Bin}(i,\alpha)$ and $\mathrm{Bin}(N-i,1-\beta)$ random
variables, so
\[P_{ij}=\sum_{k=\max(0,j-N+i)}^{\min(i,j)}\binom{i}{k}\alpha^k(1-\alpha)^{i-k}
\binom{N-i}{j-k}(1-\beta)^{j-k}\beta^{N-i-(j-k)}.\]

\textbf{(d)}  Consider the $N=1$ case: a single individual has stationary
active probability $\pi_{\mathrm{act}}=(1-\beta)/(2-\alpha-\beta)$.  For
general $N$, because each individual evolves independently with the same
stationary distribution, in the steady state the number of active
individuals is $\mathrm{Binomial}(N,p^*)$ where
$p^*=(1-\beta)/(2-\alpha-\beta)$.  Hence
\[
\pi_j=\binom{N}{j}(p^*)^j(1-p^*)^{N-j},\quad j=0,1,\ldots,N.
\]
}

\begin{exercise}
\label{ex:markov-chains-28}
Every time the team wins a game it wins the next with probability $0.8$;
every time it loses it wins the next with probability $0.3$.  If it wins
it has dinner with probability $0.7$; if it loses, with probability $0.2$.
What proportion of games result in a team dinner?
\end{exercise}
\solution{
States are $W$ (win) and $L$ (loss). The stationary equation is
\[
0.2\pi_W=0.3\pi_L,\qquad \pi_W+\pi_L=1,
\]
so $\pi_W=3/5$ and $\pi_L=2/5$. Hence the long-run dinner proportion is
\[
\pi_W(0.7)+\pi_L(0.2)=\frac35(0.7)+\frac25(0.2)=0.50.
\]
}

\begin{exercise}
\label{ex:markov-chains-29}
An organisation has $N$ employees with three job classifications, changing
independently according to $P=\begin{pmatrix}0.7&0.2&0.1\\0.2&0.6&0.2\\0.1&0.4&0.5\end{pmatrix}$.
What percentage of employees are in each classification?
\end{exercise}
\solution{
Let $\pi_1,\pi_2,\pi_3$ be the stationary proportions. The first and
third balance equations give
\[
3\pi_1=2\pi_2+\pi_3,\qquad
5\pi_3=\pi_1+2\pi_2.
\]
Substituting the first into the second yields
$\pi_3=4\pi_2/7$, and then $\pi_1=6\pi_2/7$.
Normalizing gives
\[
(\pi_1,\pi_2,\pi_3)=\left(\frac6{17},\frac7{17},\frac4{17}\right).
\]
All entries of $P$ are positive, so these are also the limiting
proportions: approximately $35.3\%$, $41.2\%$, and $23.5\%$.
}

\begin{exercise}
\label{ex:markov-chains-30}
Three out of every four trucks are followed by a car; one out of every
five cars is followed by a truck.  What fraction of vehicles on the road
are trucks?
\end{exercise}
\solution{
States: $T$ (truck), $C$ (car).
$P_{T,C}=3/4$, $P_{T,T}=1/4$, $P_{C,T}=1/5$, $P_{C,C}=4/5$.

Stationary: $(3/4)\pi_T=(1/5)\pi_C \Rightarrow\pi_T/\pi_C=4/15$.
With $\pi_T+\pi_C=1$:
\[\pi_T=\frac{4}{4+15}={\frac{4}{19}\approx21.1\%}.\]
}

\begin{exercise}
\label{ex:markov-chains-31}
A town never has two sunny days in a row.  Days are sunny ($S$), cloudy
($C$), or rainy ($R$).  If sunny, next day is equally likely $C$ or $R$.
If rainy or cloudy, probability $1/2$ of staying the same; if it changes,
it is equally likely to be either of the other two.  What proportion of
days are sunny / cloudy?
\end{exercise}
\solution{
Transition matrix:
\[
P=\bordermatrix{&S&C&R\cr
  S&0&\tfrac{1}{2}&\tfrac{1}{2}\cr
  C&\tfrac{1}{4}&\tfrac{1}{2}&\tfrac{1}{4}\cr
  R&\tfrac{1}{4}&\tfrac{1}{4}&\tfrac{1}{2}\cr}.
\]
Stationary equations:
$\pi_S=(1/4)\pi_C+(1/4)\pi_R$;
$\pi_C=(1/2)\pi_S+(1/2)\pi_C+(1/4)\pi_R \Rightarrow (1/2)\pi_C=(1/2)\pi_S+(1/4)\pi_R$;
$\pi_R=(1/2)\pi_S+(1/4)\pi_C+(1/2)\pi_R$.

From eq.~2: $\pi_C=\pi_S+(1/2)\pi_R$.
From eq.~3: $(1/2)\pi_R=(1/2)\pi_S+(1/4)\pi_C=(1/2)\pi_S+(1/4)(\pi_S+(1/2)\pi_R)
\Rightarrow(3/8)\pi_R=(3/4)\pi_S \Rightarrow\pi_R=2\pi_S$.
Then $\pi_C=\pi_S+\pi_S=2\pi_S$.

Normalising: $\pi_S+2\pi_S+2\pi_S=5\pi_S=1\Rightarrow\pi_S=1/5$.
\[{\pi_S=\frac{1}{5},\quad\pi_C=\frac{2}{5},\quad\pi_R=\frac{2}{5}.}\]
}

\begin{exercise}
\label{ex:markov-chains-32}
Each of two switches is either on or off each day.  On day $n$, each
switch independently is on with probability
$[1+\text{(number on during day }n-1)]/4$.
What fraction of days are both switches on? Both off?
\end{exercise}
\solution{
State $k=$ number of switches on ($k=0,1,2$).  Conditional on $k$
switches being on, each switch on day $n$ is on independently with prob
$p=(1+k)/4$.

\[
P_{0\to j}=\binom{2}{j}\!\left(\tfrac{1}{4}\right)^j\!\left(\tfrac{3}{4}\right)^{2-j},\;
P_{1\to j}=\binom{2}{j}\!\left(\tfrac{1}{2}\right)^2,\;
P_{2\to j}=\binom{2}{j}\!\left(\tfrac{3}{4}\right)^j\!\left(\tfrac{1}{4}\right)^{2-j}.
\]
\[
P=\begin{pmatrix}\tfrac{9}{16}&\tfrac{6}{16}&\tfrac{1}{16}\\[3pt]
  \tfrac{1}{4}&\tfrac{1}{2}&\tfrac{1}{4}\\[3pt]
  \tfrac{1}{16}&\tfrac{6}{16}&\tfrac{9}{16}\end{pmatrix}.
\]
By the symmetry $P_{0\to j}=P_{2\to2-j}$, we have $\pi_0=\pi_2$.

From the balance equation for state $1$: $\pi_1=(6/16)\pi_0+(1/2)\pi_1+(6/16)\pi_2=(3/4)\pi_0+(1/2)\pi_1$
$\Rightarrow (1/2)\pi_1=(3/4)\pi_0\Rightarrow\pi_1=(3/2)\pi_0$.

Normalising: $\pi_0+(3/2)\pi_0+\pi_0=(7/2)\pi_0=1\Rightarrow\pi_0=2/7$.
\[{\text{Both on}: \pi_2=\frac{2}{7},\quad\text{Both off}: \pi_0=\frac{2}{7}.}\]
}

\begin{exercise}
\label{ex:markov-chains-33}
A professor gives three types of exams.  $p_i=$ probability of doing
well on type $i$ ($p_1=0.3,\,p_2=0.6,\,p_3=0.9$).  If the class does
well, the next exam is equally likely to be any type; if the class does
badly, the next exam is always type~1.  What proportion of exams are of
type $i$, $i=1,2,3$?
\end{exercise}
\solution{
Let the state be the type of the current exam. A successful exam leads
to each type with probability $1/3$, while failure leads to type $1$.
Thus $P_{ij}=p_i/3+(1-p_i)\mathbf{1}_{\{j=1\}}$, giving
\[
P=\begin{pmatrix}0.8&0.1&0.1\\0.6&0.2&0.2\\0.4&0.3&0.3\end{pmatrix}.
\]
The second and third columns are identical, so $\pi_2=\pi_3$.
The balance equation for state $2$ is then
\[
\pi_2=0.1\pi_1+0.2\pi_2+0.3\pi_2,
\]
which gives $\pi_1=5\pi_2$. Normalizing yields
\[
(\pi_1,\pi_2,\pi_3)=\left(\frac57,\frac17,\frac17\right).
\]
}

\begin{exercise}
\label{ex:markov-chains-34}
A flea moves on a triangle with vertices $1,2,3$.  At vertex $i$ it moves
clockwise (CW) w.p.\ $p_i$ and counterclockwise (CCW) w.p.\ $q_i=1-p_i$.
\begin{enumerate}
\item[(a)] Find the proportion of time at each vertex.
\item[(b)] How often does a CCW move immediately precede five consecutive
  CW moves?
\end{enumerate}
\end{exercise}
\solution{
Assume $0<p_i<1$ so the chain is irreducible. Label clockwise
neighbours $1\to2\to3\to1$. The balance equations are
\[
\pi_1=q_2\pi_2+p_3\pi_3,\quad
\pi_2=p_1\pi_1+q_3\pi_3,\quad
\pi_3=q_1\pi_1+p_2\pi_2.
\]
\textbf{(a)} Define
\[
w_1=q_2p_3+q_2q_3+p_2p_3,\qquad
w_2=p_1q_3+p_1p_3+q_1q_3,
\]
\[
w_3=q_1p_2+p_1p_2+q_1q_2.
\]
Direct substitution into the balance equations verifies that
\[
\pi_i=\frac{w_i}{w_1+w_2+w_3}.
\]
For example, $w_1$ is the sum of the weights of the three directed
spanning trees pointing toward vertex $1$.

\textbf{(b)} Count a six-move window consisting of one counterclockwise
move followed by five clockwise moves. Its long-run frequency is
\[
\pi_1q_1p_3p_1p_2p_3p_1
+\pi_2q_2p_1p_2p_3p_1p_2
+\pi_3q_3p_2p_3p_1p_2p_3.
\]
No condition is imposed on the move after these five clockwise moves.
}

\begin{exercise}
\label{ex:markov-chains-35}
Consider a Markov chain with states $\{0,1,2,3,4\}$, where $P_{0,4}=1$
and from state $i>0$ the next state is equally likely to be any of
$0,1,\ldots,i-1$.  Find the limiting probabilities.
\end{exercise}
\solution{
Stationary equations $\pi=\pi P$:
\begin{align*}
\pi_0 &= \pi_1+\tfrac{1}{2}\pi_2+\tfrac{1}{3}\pi_3+\tfrac{1}{4}\pi_4,\\
\pi_1 &= \tfrac{1}{2}\pi_2+\tfrac{1}{3}\pi_3+\tfrac{1}{4}\pi_4,\\
\pi_2 &= \tfrac{1}{3}\pi_3+\tfrac{1}{4}\pi_4,\\
\pi_3 &= \tfrac{1}{4}\pi_4,\\
\pi_4 &= \pi_0.
\end{align*}

Set $\pi_0=a$.  Then $\pi_4=a$, $\pi_3=a/4$, $\pi_2=a/12+a/4=a/3$,
$\pi_1=a/6+a/12+a/4=a/2$.  (Verify: $\pi_0=a/2+a/6+a/12+a/4=a$. )

Normalising:
$a+\tfrac{a}{2}+\tfrac{a}{3}+\tfrac{a}{4}+a
=a\bigl(2+\tfrac{1}{2}+\tfrac{1}{3}+\tfrac{1}{4}\bigr)
=a\cdot\tfrac{37}{12}=1\Rightarrow a=\tfrac{12}{37}$.

\[{\pi_0=\frac{12}{37},\;\pi_1=\frac{6}{37},\;\pi_2=\frac{4}{37},\;\pi_3=\frac{3}{37},\;\pi_4=\frac{12}{37}.}\]
}

\begin{exercise}
\label{ex:markov-chains-36}
A process changes state daily according to the two-state Markov chain
with $P_{0,0}=0.4$, $P_{0,1}=0.6$, $P_{1,0}=0.2$, $P_{1,1}=0.8$.  Each
day a message is sent: ``good'' w.p.\ $p_i$ or ``bad'' w.p.\ $q_i=1-p_i$
if the chain is in state $i$.
\begin{enumerate}
\item[(a)] State~$0$ on Monday: probability of a good message on Tuesday?
\item[(b)] State~$0$ on Monday: probability of a good message on Friday?
\item[(c)] In the long run, what proportion of messages are good?
\item[(d)] Let $Y_n=1$ if a good message is sent on day $n$, $Y_n=2$
  otherwise.  Is $\{Y_n\}$ a Markov chain?
\end{enumerate}
\end{exercise}
\solution{
\textbf{Stationary distribution.}  $0.6\pi_0=0.2\pi_1\Rightarrow
\pi_0=1/4$, $\pi_1=3/4$.  Second eigenvalue: $\lambda_2=1-0.6-0.2=0.2$.

\[
P^n_{0,0}=\tfrac{1}{4}+\tfrac{3}{4}(0.2)^n,\quad
P^n_{0,1}=\tfrac{3}{4}-\tfrac{3}{4}(0.2)^n=\tfrac{3}{4}[1-(0.2)^n].
\]

\textbf{(a)} (Tuesday, $n=1$):
\[P(\text{good Tue})=P_{0,0}\cdot p_0+P_{0,1}\cdot p_1=0.4p_0+0.6p_1.\]

\textbf{(b)} (Friday, $n=4$ days after Monday):
\begin{align*}
P(\text{good Fri})
&=P^4_{0,0}\cdot p_0+P^4_{0,1}\cdot p_1\\
&=\bigl(\tfrac{1}{4}+\tfrac{3}{4}(0.2)^4\bigr)p_0
 +\tfrac{3}{4}\bigl(1-(0.2)^4\bigr)p_1.
\end{align*}
Numerically $(0.2)^4=0.0016$:
$P(\text{good Fri})\approx0.2512\,p_0+0.7488\,p_1$.

\textbf{(c)} Long-run:
$\pi_0 p_0+\pi_1 p_1=\dfrac{p_0}{4}+\dfrac{3p_1}{4}$.

\textbf{(d)}  In general, $\{Y_n\}$ is \textbf{not} a Markov chain. Special choices, such as $p_0=p_1$ with independent message draws, give a Markov output process.  Knowing $Y_n=1$
(good message) does not determine the underlying state: both states can
produce a good message, with different probabilities.  Hence $P(Y_{n+1}=1
\mid Y_n=1)$ depends on the probability of being in each underlying state
given the history, which is not determined by $Y_n$ alone.
}

\setcounter{exercise}{37}

\begin{exercise}%[Exercise 38]
\label{ex:markov-chains-38}
Recall that state $i$ is said to be positive recurrent if $m_{i,i} < \infty$, where $m_{i,i}$ is the expected number of transitions until the Markov chain, starting in state $i$, makes a transition back into that state. Because $\pi_i$, the long-run proportion of time the Markov chain, starting in state $i$, spends in state $i$, satisfies
\[
\pi_i = \frac{1}{m_{i,i}}
\]
it follows that state $i$ is positive recurrent if and only if $\pi_i > 0$. Suppose that state $i$ is positive recurrent and that state $i$ communicates with state $j$. Show that state $j$ is also positive recurrent by arguing that there is an integer $n$ such that
\[
\pi_j \ge \pi_i P_{i,j}^n > 0
\]
\end{exercise}
\solution{
Since state $i$ and state $j$ communicate ($i \leftrightarrow j$), there exists an integer $n \ge 0$ such that the $n$-step transition probability from $i$ to $j$ is strictly positive, i.e., $P_{i,j}^n > 0$.

By the properties of stationary distributions and the law of total probability applied over $n$ steps, we have:
\[
\pi_j = \sum_{k} \pi_k P_{k,j}^n
\]
Since all terms in the summation are non-negative, we can lower-bound $\pi_j$ by keeping only the single term corresponding to $k = i$:
\[
\pi_j \ge \pi_i P_{i,j}^n
\]
We are given that state $i$ is positive recurrent, which implies $\pi_i > 0$. Since we also established that $P_{i,j}^n > 0$, their product must be strictly greater than zero:
\[
\pi_j \ge \pi_i P_{i,j}^n > 0
\]
Since $\pi_j > 0$, it follows directly from the given definition that state $j$ is also positive recurrent.
}

\begin{exercise}%[Exercise 39]
\label{ex:markov-chains-39}
Recall that a recurrent state that is not positive recurrent is called null recurrent. Use the result of Exercise 38 to prove that null recurrence is a class property. That is, if state $i$ is null recurrent and state $i$ communicates with state $j$, show that state $j$ is also null recurrent.
\end{exercise}
\solution{
We are given that state $i$ is null recurrent and communicates with state $j$ ($i \leftrightarrow j$). Since recurrence is a class property and state $i$ is recurrent, state $j$ must also be recurrent.

Now, we proceed by contradiction to show that state $j$ must be null recurrent. Suppose that state $j$ is \emph{not} null recurrent. Since it is recurrent, it would therefore have to be positive recurrent.

However, since $i \leftrightarrow j$, we can apply the result from Exercise 38: if state $j$ is positive recurrent, then state $i$ must also be positive recurrent. This directly contradicts our initial assumption that state $i$ is null recurrent.

Thus, state $j$ cannot be positive recurrent, meaning it must be null recurrent. Therefore, null recurrence is a class property.
}

\begin{exercise}%[Exercise 40]
\label{ex:markov-chains-40}
It follows from the argument made in Exercise 38 that state $i$ is null recurrent if it is recurrent and $\pi_i = 0$. Consider the symmetric random walk on $\mathbb Z$, with independent steps $+1$ and $-1$ of probability $1/2$ each.
\begin{enumerate}
\item[(a)] Argue that $\pi_i = \pi_0$ for all $i$.
\item[(b)] Argue that all states are null recurrent.
\end{enumerate}
\end{exercise}
\solution{
\textbf{(a)} All states communicate, and the symmetric random walk is
recurrent. If it were positive recurrent, its stationary probability
vector would satisfy
\[
2\pi_i=\pi_{i-1}+\pi_{i+1},\qquad i\in\mathbb Z.
\]
Thus consecutive differences are constant, so $\pi_i=a+bi$.
Nonnegativity for every integer $i$ forces $b=0$, hence $\pi_i=\pi_0$.

\textbf{(b)} No constant nonnegative sequence on $\mathbb Z$ sums to
$1$. Thus a stationary probability vector cannot exist, and the
irreducible recurrent chain is null recurrent. Its long-run occupation
proportions are all zero, so the equality in (a) also holds for these
proportions. These zero proportions are not a stationary probability
\emph{distribution}.
}

\begin{exercise}%[Exercise 41]
\label{ex:markov-chains-41}
Let $\pi_i$ denote the long-run proportion of time a given irreducible Markov chain is in state $i$.
\begin{enumerate}
\item[(a)] Explain why $\pi_i$ is also the proportion of transitions that are into state $i$ as well as being the proportion of transitions that are from state $i$.
\item[(b)] $\pi_i P_{i,j}$ represents the proportion of transitions that satisfy what property?
\item[(c)] $\sum_{i} \pi_i P_{i,j}$ represent the proportion of transitions that satisfy what property?
\item[(d)] Using the preceding explain why
\[
\pi_j = \sum_{i} \pi_i P_{i,j}
\]
\end{enumerate}
\end{exercise}
\solution{
\textbf{(a)} In the long run, the number of entries into state $i$ must essentially equal the number of departures from state $i$ (differing by at most 1 at any finite step). Since $\pi_i$ is the long-run proportion of time spent in state $i$, each unit of time corresponds to exactly one transition out of state $i$. Thus, $\pi_i$ represents the long-run proportion of transitions originating from state $i$. By balance, it must also equal the long-run proportion of transitions entering state $i$.

\textbf{(b)} $\pi_i P_{i,j}$ represents the long-run proportion of all transitions that start in state $i$ and move directly into state $j$ in the next step.

\textbf{(c)} $\sum_{i} \pi_i P_{i,j}$ sums the proportions of transitions moving into state $j$ from any possible starting state $i$. Therefore, it represents the total long-run proportion of transitions that end in state $j$.

\textbf{(d)} From part (a), $\pi_j$ is the long-run proportion of transitions that are into state $j$. From part (c), $\sum_{i} \pi_i P_{i,j}$ is also the total long-run proportion of transitions that enter state $j$. Equating these two expressions for the same physical quantity yields the stationary identity:
\[
\pi_j = \sum_{i} \pi_i P_{i,j}
\]
}

\begin{exercise}%[Exercise 42]
\label{ex:markov-chains-42}
Let $A$ be a set of states, and let $A^c$ be the remaining states.
\begin{enumerate}
\item[(a)] What is the interpretation of
\[
\sum_{i \in A} \sum_{j \in A^c} \pi_i P_{i,j}?
\]
\item[(b)] What is the interpretation of
\[
\sum_{i \in A^c} \sum_{j \in A} \pi_i P_{i,j}?
\]
\item[(c)] Explain the identity
\[
\sum_{i \in A} \sum_{j \in A^c} \pi_i P_{i,j} = \sum_{i \in A^c} \sum_{j \in A} \pi_i P_{i,j}
\]
\end{enumerate}
\end{exercise}
\solution{
\textbf{(a)} This expression represents the long-run proportion of all transitions that start from a state inside set $A$ and end in a state outside set $A$ (in $A^c$). In other words, it is the long-run rate of escaping or leaving set $A$.

\textbf{(b)} This expression represents the long-run proportion of all transitions that start from a state outside set $A$ (in $A^c$) and end in a state inside set $A$. In other words, it is the long-run rate of entering set $A$.

\textbf{(c)} For the Markov chain to remain in a stationary regime, the system cannot indefinitely pile up inside or outside of set $A$. Over the long run, every time the process exits set $A$, it must eventually return to set $A$ to exit again. Thus, the total number of transitions leaving $A$ must balance the total number of transitions entering $A$ (differing by at most 1 at any given time). Therefore, the long-run rate of leaving $A$ must equal the long-run rate of entering $A$:
\[
\sum_{i \in A} \sum_{j \in A^c} \pi_i P_{i,j} = \sum_{i \in A^c} \sum_{j \in A} \pi_i P_{i,j}
\]
}

\begin{exercise}%[Exercise 43]
\label{ex:markov-chains-43}
Each day, one of $n$ possible elements is requested, the $i$th one with probability $P_i, i \ge 1, \sum_1^n P_i = 1$. These elements are at all times arranged in an ordered list that is revised as follows: The element selected is moved to the front of the list with the relative positions of all the other elements remaining unchanged. Define the state at any time to be the list ordering at that time and note that there are $n!$ possible states.
\begin{enumerate}
\item[(a)] Argue that the preceding is a Markov chain.
\item[(b)] For any state $i_1, \dots, i_n$ (which is a permutation of $1, 2, \dots, n$), let $\pi(i_1, \dots, i_n)$ denote the limiting probability. In order for the state to be $i_1, \dots, i_n$, it is necessary for the last request to be for $i_1$, the last non-$i_1$ request for $i_2$, the last non-$i_1$ or $i_2$ request for $i_3$, and so on. Hence, it appears intuitive that
\[
\pi(i_1, \dots, i_n) = P_{i_1} \frac{P_{i_2}}{1 - P_{i_1}} \frac{P_{i_3}}{1 - P_{i_1} - P_{i_2}} \dots \frac{P_{i_{n-1}}}{1 - P_{i_1} - \dots - P_{i_{n-2}}}
\]
Verify when $n = 3$ that the preceding are indeed the limiting probabilities.
\end{enumerate}
\end{exercise}
\solution{
\textbf{(a)} The next ordering of the list depends exclusively on two things: the current arrangement of the list and which element is selected next. The probability of choosing any element $i$ is constant ($P_i$) and independent of previous choices or past positions. Since the future list configuration depends only on the current configuration and the independent current selection, the process satisfies the Markov property and is a Markov chain.

Assume $P_i>0$ for all items.\par
\textbf{(b)} Let $n = 3$. We wish to verify that the proposed formula satisfies the stationary equations for a state, say $\pi(1, 2, 3)$. According to the formula:
\[
\pi(1, 2, 3) = P_1 \frac{P_2}{1 - P_1}
\]
Note that since $P_1 + P_2 + P_3 = 1$, the final implied factor $\frac{P_3}{1-P_1-P_2} = \frac{P_3}{P_3} = 1$, so the expression is complete.

To verify stationarity, we look at all states that can transition into $(1, 2, 3)$ in one step. A transition into $(1, 2, 3)$ occurs if and only if element 1 is requested while the relative order of the remaining elements is $(2, 3)$. The states that have the relative order $(2, 3)$ for the remaining elements are $(1, 2, 3)$, $(2, 1, 3)$, and $(2, 3, 1)$. In each case, requesting element 1 (with probability $P_1$) moves 1 to the front, yielding $(1, 2, 3)$.

Thus, the balance equation for state $(1, 2, 3)$ is:
\[
\pi(1, 2, 3) = \Big( \pi(1, 2, 3) + \pi(2, 1, 3) + \pi(2, 3, 1) \Big) \cdot P_1
\]
Let us substitute the proposed formula for each of these three states on the right-hand side:
\begin{align*}
\pi(1, 2, 3) &= P_1 \frac{P_2}{1 - P_1} \\
\pi(2, 1, 3) &= P_2 \frac{P_1}{1 - P_2} \\
\pi(2, 3, 1) &= P_2 \frac{P_3}{1 - P_2}
\end{align*}
Summing the terms inside the parentheses:
\begin{align*}
\pi(1, 2, 3) + \pi(2, 1, 3) + \pi(2, 3, 1) &= P_1 \frac{P_2}{1 - P_1} + P_2 \frac{P_1}{1 - P_2} + P_2 \frac{P_3}{1 - P_2} \\
&= P_1 \frac{P_2}{1 - P_1} + \frac{P_2}{1 - P_2} (P_1 + P_3)
\end{align*}
Since $P_1 + P_2 + P_3 = 1$, we have $P_1 + P_3 = 1 - P_2$. Substituting this back:
\[
= P_1 \frac{P_2}{1 - P_1} + \frac{P_2}{1 - P_2} (1 - P_2) = P_1 \frac{P_2}{1 - P_1} + P_2
\]
Now, find a common denominator for these two remaining terms:
\[
= P_2 \left( \frac{P_1}{1 - P_1} + 1 \right) = P_2 \left( \frac{P_1 + 1 - P_1}{1 - P_1} \right) = \frac{P_2}{1 - P_1}
\]
Finally, multiply this sum by the transition probability $P_1$ as required by the balance equation:
\[
\Big( \pi(1, 2, 3) + \pi(2, 1, 3) + \pi(2, 3, 1) \Big) \cdot P_1 = \frac{P_2}{1 - P_1} \cdot P_1 = P_1 \frac{P_2}{1 - P_1} = \pi(1,2,3)
\]
The equation holds. By symmetry, this verification applies to all other permutations, confirming that the proposed formula gives the true limiting probabilities.
}

\begin{exercise}%[Exercise 44]
\label{ex:markov-chains-44}
Suppose that a population consists of a fixed number, say, $m$, of genes in any generation. Each gene is one of two possible genetic types. If exactly $i$ (of the $m$) genes of any generation are of type 1, then the next generation will have $j$ type 1 genes (and $m - j$ type 2 genes) with probability
\[
\binom{m}{j} \left(\frac{i}{m}\right)^j \left(\frac{m-i}{m}\right)^{m-j}, \quad j = 0, 1, \dots, m
\]
Let $X_n$ denote the number of type 1 genes in the $n$th generation, and assume that $X_0 = i$.
\begin{enumerate}
\item[(a)] Find $E[X_n]$.
\item[(b)] What is the probability that eventually all the genes will be type 1?
\end{enumerate}
\end{exercise}
\solution{
\textbf{(a)} Given $X_n = i$, the next state $X_{n+1}$ follows a binomial distribution with parameters $m$ (number of trials) and $p = \frac{i}{m} = \frac{X_n}{m}$ (success probability).

The conditional expectation is:
\[
E[X_{n+1} \mid X_n] = m \cdot \left(\frac{X_n}{m}\right) = X_n
\]
Taking the total expectation on both sides yields:
\[
E[X_{n+1}] = E[E[X_{n+1} \mid X_n]] = E[X_n]
\]
By induction, the expected value remains constant across all generations:
\[
E[X_n] = E[X_0] = i
\]

\textbf{(b)} The states $j = 0$ and $j = m$ are absorbing states because if the population consists entirely of one type of gene, no genes of the other type can be generated. Since the chain is finite and every other state can reach these absorbing states, the process will eventually be absorbed into either state 0 (extinction of type 1) or state $m$ (fixation of type 1) with probability 1.

Let $P_{\text{all}}$ be the probability that the population eventually consists entirely of type 1 genes (absorption at $m$). As $n \to \infty$, the random variable $X_n$ converges to a limiting random variable $X_\infty$ which takes value $m$ with probability $P_{\text{all}}$ and value $0$ with probability $1 - P_{\text{all}}$.

Using the property from part (a) that $E[X_n] = i$ for all $n$, we can take the limit by bounded convergence, since $0\le X_n\le m$:
\[
E[X_\infty] = m \cdot P_{\text{all}} + 0 \cdot (1 - P_{\text{all}}) = i
\]
Solving for $P_{\text{all}}$ gives:
\[
P_{\text{all}} = \frac{i}{m}
\]
}

\setcounter{exercise}{45}

\begin{exercise}%[Exercise 46]
\label{ex:markov-chains-46}
An individual possesses $r$ umbrellas that he employs in going from his home to office, and vice versa. If he is at home (the office) at the beginning (end) of a day and it is raining, then he will take an umbrella with him to the office (home), provided there is one to be taken. If it is not raining, then he never takes an umbrella. Assume that, independent of the past, it rains at the beginning (end) of a day with probability $p$, independent of the past.
\begin{enumerate}
\item[(a)] Define a Markov chain with $r + 1$ states, which will help us to determine the proportion of time that our man gets wet. (Note: He gets wet if it is raining, and all umbrellas are at his other location.)
\item[(b)] Show that the limiting probabilities are given by
\[
\pi_i = \begin{cases}
\dfrac{q}{r + q}, & \text{if } i = 0 \\
\dfrac{1}{r + q}, & \text{if } i = 1, \dots, r
\end{cases} \quad \text{where } q = 1 - p
\]
\item[(c)] What fraction of time does our man get wet?
\item[(d)] When $r = 3$, what value of $p$ maximizes the fraction of time he gets wet?
\end{enumerate}
\end{exercise}
\solution{
Assume $r\ge1$ and $0<p<1$ for the limiting-distribution argument.\par
\textbf{(a)} Let the state $X_n$ be the number of umbrellas available at the individual's current location (home or office) before his trip. There are $r+1$ possible states, $i \in \{0, 1, \dots, r\}$. Let $q = 1 - p$.
If the current state is $i$ ($0 < i \le r$):
\begin{itemize}
\item With probability $p$, it rains, he takes one umbrella, and leaves $i-1$ umbrellas behind. The next location will have the remaining $r - (i - 1) = r - i + 1$ umbrellas. Thus, $P_{i, r-i+1} = p$.
\item With probability $q$, it does not rain, he takes no umbrella, leaving $i$ umbrellas behind. The next location will have $r - i$ umbrellas. Thus, $P_{i, r-i} = q$.
\end{itemize}
If the current state is $0$, he cannot take any umbrella regardless of the weather. Thus, he leaves $0$ umbrellas behind, and the next location will have all $r$ umbrellas. Hence, $P_{0, r} = 1$.

\textbf{(b)} The stationary equations are given by $\pi_j = \sum_i \pi_i P_{i,j}$. We can guess the form $\pi_i = c$ for $i=1,\dots,r$ and $\pi_0 = c q$. Let us verify this balance by verifying $\sum_{i=0}^r \pi_i = 1$:
\[
\pi_0 + \sum_{i=1}^r \pi_i = cq + rc = c(r + q) = 1 \implies c = \frac{1}{r+q}
\]
Thus, $\pi_0 = \frac{q}{r+q}$ and $\pi_i = \frac{1}{r+q}$ for $i=1,\dots,r$. We check equation for state $r$:
\[
\pi_r = \pi_0 P_{0,r} + \pi_1 P_{1,r} = \pi_0(1) + \pi_1(p) = \frac{q}{r+q} + \frac{p}{r+q} = \frac{1}{r+q}
\]
This matches perfectly. Similar checks verify the other states, confirming the unique stationary distribution.

\textbf{(c)} The individual gets wet if he is at a location with 0 umbrellas and it is raining. Since the weather is independent of the location's state, this long-run fraction of time is:
\[
\text{Fraction} = \pi_0 \cdot p = \frac{p q}{r + q} = \frac{p(1-p)}{r + 1 - p}
\]

\textbf{(d)} For $r = 3$, the fraction is $f(p) = \frac{p - p^2}{4 - p}$. To find the maximum, we differentiate with respect to $p$:
\[
f'(p) = \frac{(1 - 2p)(4 - p) - (p - p^2)(-1)}{(4 - p)^2} = \frac{4 - p - 8p + 2p^2 + p - p^2}{(4 - p)^2} = \frac{p^2 - 8p + 4}{(4 - p)^2}
\]
Setting $f'(p) = 0$ yields $p^2 - 8p + 4 = 0$. Solving the quadratic equation for $p \in [0, 1]$:
\[
p = \frac{8 - \sqrt{64 - 16}}{2} = 4 - \sqrt{12} = 4 - 2\sqrt{3} \approx 0.5359
\]
}

\begin{exercise}%[Exercise 47]
\label{ex:markov-chains-47}
Let $\{X_n, n \ge 0\}$ denote an ergodic Markov chain with limiting probabilities $\pi_i$. Define the process $\{Y_n, n \ge 1\}$ by $Y_n = (X_{n-1}, X_n)$. That is, $Y_n$ keeps track of the last two states of the original chain. Is $\{Y_n, n \ge 1\}$ a Markov chain? If so, determine its transition probabilities and find $\lim_{n \to \infty} P\{Y_n = (i,j)\}$.
\end{exercise}
\solution{
Yes, $\{Y_n\}$ is a Markov chain. The value of $Y_{n+1} = (X_n, X_{n+1})$ depends on the past history only through its immediate predecessor $Y_n = (X_{n-1}, X_n)$, because $X_{n+1}$ depends only on $X_n$ via the Markov property of $\{X_n\}$.

The transition probabilities are given by:
\[
P_{Y_n, Y_{n+1}} = P\{(X_n, X_{n+1}) = (k, l) \mid (X_{n-1}, X_n) = (i, j)\}
\]
This probability is $0$ unless the current second component matches the next first component (i.e., $j = k$). If $j = k$, then:
\[
P\{(j, l) \mid (i, j)\} = P\{X_{n+1} = l \mid X_n = j, X_{n-1} = i\} = P_{j,l}
\]
The limiting probabilities are found directly from the joint distribution under stationarity:
\[
\lim_{n \to \infty} P\{Y_n = (i,j)\} = \lim_{n \to \infty} P\{X_{n-1} = i, X_n = j\} = \lim_{n \to \infty} P\{X_{n-1} = i\} P_{i,j} = \pi_i P_{i,j}
\]
}

\begin{exercise}%[Exercise 48]
\label{ex:markov-chains-48}
Consider a Markov chain in steady state. Say that a $k$ length run of zeroes ends at time $m$ if
\[
X_{m-k-1} \neq 0, \quad X_{m-k} = X_{m-k+1} = \dots = X_{m-1} = 0, \quad X_m \neq 0
\]
Show that the probability of this event is $\pi_0 (P_{0,0})^{k-1} (1 - P_{0,0})^2$, where $\pi_0$ is the limiting probability of state 0.
\end{exercise}
\solution{
By stationarity, the probability of entering state $0$ from outside it
in any given step is
\[
\sum_{i\ne0}\pi_iP_{i0}
=\pi_0-\pi_0P_{00}=\pi_0(1-P_{00}).
\]
After entering, a run of exactly $k\ge1$ zeros requires $k-1$ transitions
from $0$ to itself, followed by a transition out of $0$. By the Markov
property, its probability at the specified time is therefore
\[
\pi_0(1-P_{00})P_{00}^{k-1}(1-P_{00})
=\pi_0P_{00}^{k-1}(1-P_{00})^2.
\]
}

\begin{exercise}%[Exercise 49]
\label{ex:markov-chains-49}
Let $P^{(1)}$ and $P^{(2)}$ denote transition probability matrices for ergodic Markov chains having the same state space. Let $\pi^1$ and $\pi^2$ denote the stationary (limiting) probability vectors for the two chains. Consider a process defined as follows:
\begin{enumerate}
\item[(a)] $X_0 = 1$. A coin is then flipped and if it comes up heads, then the remaining states $X_1, \dots$ are obtained from the transition probability matrix $P^{(1)}$ and if tails from the matrix $P^{(2)}$. Is $\{X_n, n \ge 0\}$ a Markov chain? If $p = P\{\text{coin comes up heads}\}$, what is $\lim_{n \to \infty} P(X_n = i)$?
\item[(b)] $X_0 = 1$. At each stage the coin is flipped and if it comes up heads, then the next state is chosen according to $P^{(1)}$ and if tails comes up, then it is chosen according to $P^{(2)}$. In this case do the successive states constitute a Markov chain? If so, determine the transition probabilities. Show by a counterexample that the limiting probabilities are not the same as in part (a).
\end{enumerate}
\end{exercise}
\solution{
\textbf{(a)} In general, the process is not Markov: the past path
contains information about the transition matrix chosen initially.
There are exceptions, such as $P^{(1)}=P^{(2)}$. Conditioning on the
initial choice gives
\[
\lim_{n\to\infty}P(X_n=i)=p\pi_i^1+(1-p)\pi_i^2.
\]
\textbf{(b)} Fresh independent choices at each step give a Markov chain
with matrix $P=pP^{(1)}+(1-p)P^{(2)}$.
For a counterexample satisfying the ergodicity assumption, take $p=1/2$
and
\[
P^{(1)}=\begin{pmatrix}9/10&1/10\\1/5&4/5\end{pmatrix},
\qquad
P^{(2)}=\begin{pmatrix}1/5&4/5\\2/5&3/5\end{pmatrix}.
\]
Both matrices have strictly positive entries. Their stationary vectors
are $(2/3,1/3)$ and $(1/3,2/3)$, whose equally weighted mixture is
$(1/2,1/2)$. In contrast,
\[
P=\begin{pmatrix}11/20&9/20\\3/10&7/10\end{pmatrix}
\]
has stationary vector $(2/5,3/5)$. Thus the two limiting distributions
differ even when both original chains are ergodic.
}

\begin{exercise}%[Exercise 50]
\label{ex:markov-chains-50}
In Exercise 8, if today's flip lands heads, what is the expected number of additional flips needed until the pattern $t, t, h, t, h, t, t$ occurs?
\end{exercise}
\solution{
The flips follow Exercise~8: the probability of heads next is $0.7$
after heads and $0.6$ after tails. Consequently the formula for a
pattern in independent fair flips does not apply.

Track the longest suffix of the observed flips that is a proper prefix
of $TTHTHTT$. The transient states are
\[
\varnothing,\ T,\ TT,\ TTH,\ TTHT,\ TTHTH,\ TTHTHT.
\]
The empty-prefix state ends in heads, so today's heads starts in that
state. Let $e_0,\ldots,e_6$ be the expected additional numbers of flips
from these states; reaching the full pattern has remaining time zero.
First-step analysis gives
\begin{align*}
e_0&=1+0.7e_0+0.3e_1,\\
e_1&=1+0.6e_0+0.4e_2,\\
e_2&=1+0.4e_2+0.6e_3,\\
e_3&=1+0.7e_0+0.3e_4,\\
e_4&=1+0.4e_2+0.6e_5,\\
e_5&=1+0.7e_0+0.3e_6,\\
e_6&=1+0.6e_0.
\end{align*}
For example, a tail after $TTHT$ leaves the suffix $TT$, and a tail
after $TTHTHT$ completes the target. Solving gives
\[
e_0=\frac{31835}{54}\approx589.54.
\]
This is the expected number of additional flips after today's heads.
}

\begin{exercise}%[Exercise 51]
\label{ex:markov-chains-51}
In the mood model given in the chapter setup, Gary is in a cheerful mood today. Find the expected number of days until he has been glum for three consecutive days.
\end{exercise}
\solution{
The states must distinguish cheerful and so-so moods, since they have
different probabilities of becoming glum. Use states $C,S,G_1,G_2$,
where $G_k$ records a current run of $k$ glum days; three glum days is
the absorbing target. Let $c,s,g_1,g_2$ be the expected additional days.
Using the transition matrix given in the setup,
\begin{align*}
c&=1+0.5c+0.4s+0.1g_1,\\
s&=1+0.3c+0.4s+0.3g_1,\\
g_1&=1+0.2c+0.3s+0.5g_2,\\
g_2&=1+0.2c+0.3s.
\end{align*}
The last equation has no continuation term when the next day is glum,
because that completes the three-day run. Solving the system yields
\[
(c,s,g_1,g_2)=\left(\frac{236}{9},\frac{226}{9},\frac{62}{3},\frac{124}{9}\right).
\]
Starting cheerful, the expected waiting time is $236/9\approx26.22$ days.
}

\begin{exercise}%[Exercise 52]
\label{ex:markov-chains-52}
A taxi driver provides service in two zones of a city. Fares picked up in zone $A$ will have destinations in zone $A$ with probability $0.6$ or in zone $B$ with probability $0.4$. Fares picked up in zone $B$ will have destinations in zone $A$ with probability $0.3$ or in zone $B$ with probability $0.7$. The driver's expected profit for a trip entirely in zone $A$ is $6$; for a trip entirely in zone $B$ is $8$; and for a trip that involves both zones is $12$. Find the taxi driver's average profit per trip.
\end{exercise}
\solution{
First, we find the stationary probabilities of being in each zone by modeling the location of transitions as a two-state Markov chain with matrix:
\[
P = \begin{pmatrix} 0.6 & 0.4 \\ 0.3 & 0.7 \end{pmatrix}
\]
The stationary equations $\pi_A P_{A,B} = \pi_B P_{B,A}$ give:
\[
0.4 \pi_A = 0.3 \pi_B \implies \pi_A = \frac{3}{7}, \quad \pi_B = \frac{4}{7}
\]
Next, we calculate the expected profit generated from a trip originating in each zone:
\begin{align*}
R_A &= P_{A,A}(6) + P_{A,B}(12) = 0.6(6) + 0.4(12) = 3.6 + 4.8 = 8.4 \\
R_B &= P_{B,A}(12) + P_{B,B}(8) = 0.3(12) + 0.7(8) = 3.6 + 5.6 = 9.2
\end{align*}
The long-run average profit per trip is the expected value across the stationary distribution:
\[
\text{Average Profit} = \pi_A R_A + \pi_B R_B = \frac{3}{7}(8.4) + \frac{4}{7}(9.2) = \frac{25.2 + 36.8}{7} = \frac{62}{7} \approx 8.86
\]
}

\begin{exercise}%[Exercise 53]
\label{ex:markov-chains-53}
An insurer assigns each policyholder an annual premium level $0$, $1$,
or $2$, charging $500$, $750$, or $1000$ monetary units, respectively.
The level at the start of a year determines that year's premium. For a
policyholder of type $\lambda$, annual claim counts are independent
Poisson random variables with mean $\lambda$. A claim-free year lowers
the next year's level by one, stopping at level $0$; a year with one or
more claims sets the next year's level to $2$. A policyholder's type
remains fixed. One-third of policyholders have $\lambda=1/4$ and
two-thirds have $\lambda=1/2$. Find the long-run average annual premium
per policyholder.
\end{exercise}
\solution{
For a fixed type, put $q=e^{-\lambda}$, the probability of a claim-free
year. With rows and columns ordered $0,1,2$, the transition matrix is
\[
P(\lambda)=\begin{pmatrix}q&0&1-q\\q&0&1-q\\0&q&1-q\end{pmatrix}.
\]
Since $0<q<1$, the chain is irreducible and aperiodic. Every row sends
probability $1-q$ to state $2$, so its stationary probability is
$\pi_2=1-q$. The balance equation for state $1$ gives
$\pi_1=q\pi_2=q(1-q)$, and normalization gives $\pi_0=q^2$.
The mean premium for this type is therefore
\begin{align*}
R(\lambda)&=500q^2+750q(1-q)+1000(1-q)\\
&=1000-250e^{-\lambda}-250e^{-2\lambda}.
\end{align*}
Average over the fixed policyholder types, using their population
fractions rather than their claim frequencies:
\begin{align*}
\overline R&=\frac13R(1/4)+\frac23R(1/2)\\
&=1000-\frac{250}{3}\left(e^{-1/4}+3e^{-1/2}+2e^{-1}\right)\\
&\approx722.15\text{ monetary units per year}.
\end{align*}
}

\begin{exercise}%[Exercise 54]
\label{ex:markov-chains-54}
Consider the Ehrenfest urn model in which $M$ molecules are distributed between two urns, and at each time point one of the molecules is chosen at random and is then removed from its urn and placed in the other one. Let $X_n$ denote the number of molecules in urn 1 after the $n$th switch and let $\mu_n = E[X_n]$. Show that
\begin{enumerate}
\item[(a)] $\mu_{n+1} = 1 + (1 - 2/M)\mu_n$.
\item[(b)] Use (a) to prove that
\[
\mu_n = \frac{M}{2} + \left(\frac{M - 2}{M}\right)^n \left(E[X_0] - \frac{M}{2}\right)
\]
\end{enumerate}
\end{exercise}
\solution{
\textbf{(a)} Given $X_n$, the number of molecules in urn 1 at step $n+1$ increases by 1 if a molecule from urn 2 is chosen, and decreases by 1 if a molecule from urn 1 is chosen. Thus:
\[
X_{n+1} = \begin{cases}
X_n + 1, & \text{with prob } \dfrac{M - X_n}{M} \\
X_n - 1, & \text{with prob } \dfrac{X_n}{M}
\end{cases}
\]
Taking the conditional expectation:
\[
E[X_{n+1} \mid X_n] = (X_n + 1)\frac{M - X_n}{M} + (X_n - 1)\frac{X_n}{M} = X_n + 1 - \frac{2X_n}{M} = 1 + \left(1 - \frac{2}{M}\right)X_n
\]
Taking total expectation yields the desired relation:
\[
\mu_{n+1} = 1 + \left(1 - \frac{2}{M}\right)\mu_n
\]

\textbf{(b)} We rewrite the relation as $\mu_{n+1} - \frac{M}{2} = \left(1 - \frac{2}{M}\right)\left(\mu_n - \frac{M}{2}\right)$. Let $Y_n = \mu_n - \frac{M}{2}$. Then:
\[
Y_{n+1} = \left(\frac{M - 2}{M}\right) Y_n \implies Y_n = \left(\frac{M - 2}{M}\right)^n Y_0
\]
Substituting back $Y_n$:
\[
\mu_n - \frac{M}{2} = \left(\frac{M - 2}{M}\right)^n \left(E[X_0] - \frac{M}{2}\right) \implies \mu_n = \frac{M}{2} + \left(\frac{M - 2}{M}\right)^n \left(E[X_0] - \frac{M}{2}\right)
\]
}

\begin{exercise}%[Exercise 55]
\label{ex:markov-chains-55}
Consider a population of individuals each of whom possesses two genes that can be either type $A$ or type $a$, with $A$ dominant. For randomly and independently chosen parents, let $S_{10}$ be the probability of a recessive child given one dominant and one recessive parent, and $S_{11}$ the probability given two dominant parents. Compute $S_{11}$ and $S_{10}$ to show that $S_{11} = S_{10}^2$. (The quantities $S_{10}$ and $S_{11}$ are known in the genetics literature as Snyder's ratios.)
\end{exercise}
\solution{
Let $a,b,c$ be the proportions of genotypes $AA,aa,Aa$ in the parent
population, and assume $a+c>0$. Given that a parent is dominant, it is
heterozygous with probability $c/(a+c)$. Therefore its probability of
transmitting $a$ is $u=c/[2(a+c)]$. A recessive parent always transmits
$a$, and independently sampled dominant parents transmit independently.
Consequently
\[
S_{10}=u=\frac{c}{2(a+c)},\qquad
S_{11}=u^2=S_{10}^2.
\]
Hardy--Weinberg equilibrium is not needed for this relation. If it is
additionally assumed and the recessive allele frequency is $z$, then
$u=z/(1+z)$.
}

\begin{exercise}%[Exercise 56]
\label{ex:markov-chains-56}
Suppose that on each play of the game a gambler either wins $1$ with probability $p$ or loses $1$ with probability $1 - p$. The gambler continues betting until she or he is either up $n$ or down $m$. What is the probability that the gambler quits a winner?
\end{exercise}
\solution{
This is the classical Gambler's Ruin problem. Let $P_i$ be the probability that the gambler reaches total profit $n$ starting from a current relative position $i$ (where the boundaries are $-m$ and $n$). Let $q = 1-p$.
The difference equation is:
\[
P_i = p P_{i+1} + q P_{i-1}
\]
with boundary conditions $P_{-m} = 0$ and $P_n = 1$.
The solution to this standard system depends on whether $p = 0.5$:
\[
P_i = \begin{cases}
\dfrac{1 - (q/p)^{i+m}}{1 - (q/p)^{n+m}}, & \text{if } p \neq q \\
\dfrac{i + m}{n + m}, & \text{if } p = q = 0.5
\end{cases}
\]
Starting at position $i = 0$ (the initial state), the probability that the gambler quits a winner is:
\[
P_0 = \begin{cases}
\dfrac{1 - (q/p)^m}{1 - (q/p)^{n+m}}, & \text{if } p \neq 0.5 \\
\dfrac{m}{n + m}, & \text{if } p = 0.5
\end{cases}
\]
}

\begin{exercise}%[Exercise 57]
\label{ex:markov-chains-57}
A particle moves among $n+1$ vertices on a circle, moving clockwise with probability $p$ and counterclockwise with probability $1-p$ at each independent step. Starting at a specified state, call it state $0$, let $T$ be the time of the first return to state $0$. Find the probability that all states have been visited by time $T$.
\hfill \textbf{Hint:} Condition on the initial transition and then use results from the gambler's ruin problem.
\end{exercise}
\solution{
Assume $n\ge1$ and $0<p<1$. By symmetry, on the first step the particle moves either to state 1 (clockwise) with probability $p$, or to state $n$ (counterclockwise) with probability $q = 1-p$.

Suppose it moves to state 1. For all states to be visited before returning to 0, the particle must reach state $n$ before hitting state 0. This is exactly equivalent to a Gambler's Ruin problem on the line $\{0, 1, \dots, n\}$ starting at position 1, where 0 and $n$ are the boundaries. It wins if it hits $n$ first.
From the Gambler's ruin formulas:
\[
P(\text{reach } n \text{ before } 0 \mid \text{start at } 1) = \begin{cases}
\dfrac{1 - (q/p)^1}{1 - (q/p)^n}, & \text{if } p \neq q \\
\dfrac{1}{n}, & \text{if } p = q = 0.5
\end{cases}
\]

Conversely, if the first step was to state $n$ (probability $q$), it must reach state 1 before hitting 0 to visit all states. This is symmetric but with the roles of $p$ and $q$ reversed:
\[
P(\text{reach } 1 \text{ before } 0 \mid \text{start at } n) = \begin{cases}
\dfrac{1 - (p/q)^1}{1 - (p/q)^n}, & \text{if } p \neq q \\
\dfrac{1}{n}, & \text{if } p = q = 0.5
\end{cases}
\]

Using the law of total probability, the overall probability is:
\[
\text{Total Prob} = \begin{cases}
p \cdot \dfrac{1 - q/p}{1 - (q/p)^n} + q \cdot \dfrac{1 - p/q}{1 - (p/q)^n}, & \text{if } p \neq 0.5 \\
\dfrac{1}{n}, & \text{if } p = 0.5
\end{cases}
\]
}

\begin{exercise}%[Exercise 58]
\label{ex:markov-chains-58}
In a gambler's ruin game with fortunes $0,1,\ldots,N$, each independent bet increases the fortune by $1$ with probability $p\in(0,1)$ or decreases it by $1$ with probability $q=1-p$; play stops at $0$ or $N$. Suppose the gambler's fortune is presently $i$, and suppose that we know that the gambler's fortune will eventually reach $N$ (before it goes to 0). Given this information, show that the probability he wins the next gamble is
\[
\begin{cases}
\dfrac{p[1 - (q/p)^{i+1}]}{1 - (q/p)^i}, & \text{if } p \neq \frac{1}{2} \\
\dfrac{i + 1}{2i}, & \text{if } p = \frac{1}{2}
\end{cases}
\]
\end{exercise}
\solution{
Let $P_k$ denote the probability that the gambler reaches $N$ before 0 starting from fortune $k$. From the classical Gambler's Ruin formulas, we have:
\[
P_k = \begin{cases}
\dfrac{1 - (q/p)^k}{1 - (q/p)^N}, & \text{if } p \neq \frac{1}{2} \\
\dfrac{k}{N}, & \text{if } p = \frac{1}{2}
\end{cases}
\]
Following the hint, we use the definition of conditional probability:
\[
P\{X_{n+1} = i+1 \mid X_n = i, \lim_{m \to \infty} X_m = N\} = \frac{P\{X_{n+1} = i+1, \lim_{m \to \infty} X_m = N \mid X_n = i\}}{P\{\lim_{m \to \infty} X_m = N \mid X_n = i\}}
\]
Applying the Markov property, the numerator is the probability of moving from $i$ to $i+1$ in one step, times the probability of eventually reaching $N$ starting from $i+1$:
\[
\text{Numerator} = P\{X_{n+1} = i+1 \mid X_n = i\} \cdot P\{\lim_{m \to \infty} X_m = N \mid X_{n+1} = i+1\} = p \cdot P_{i+1}
\]
The denominator is simply $P_i$. Therefore, the conditional probability is $\dfrac{p P_{i+1}}{P_i}$.

\textbf{Case 1:} $p \neq \frac{1}{2}$
\[
\frac{p P_{i+1}}{P_i} = p \cdot \frac{\frac{1 - (q/p)^{i+1}}{1 - (q/p)^N}}{\frac{1 - (q/p)^i}{1 - (q/p)^N}} = \frac{p[1 - (q/p)^{i+1}]}{1 - (q/p)^i}
\]

\textbf{Case 2:} $p = \frac{1}{2}$
\[
\frac{p P_{i+1}}{P_i} = \frac{1}{2} \cdot \frac{\frac{i+1}{N}}{\frac{i}{N}} = \frac{i+1}{2i}
\]
The claims are shown.
}

\begin{exercise}%[Exercise 59]
\label{ex:markov-chains-59}
For independent bets that change the fortune by $+1$ with probability $p\in(0,1)$ and $-1$ with probability $q=1-p$, stopping at $0$ or $N$, let $M_i$ denote the mean number of games that must be played until the gambler either goes broke or reaches a fortune of $N$, given that he starts with $i, i = 0, 1, \dots, N$. Show that $M_i$ satisfies
\[
M_0 = M_N = 0; \quad M_i = 1 + p M_{i+1} + q M_{i-1}, \quad i = 1, \dots, N - 1
\]
\end{exercise}
\solution{
If the gambler starts at $i = 0$ or $i = N$, the process terminates instantly because the game boundaries are already reached. Thus, $M_0 = 0$ and $M_N = 0$.

For any intermediate fortune $i \in \{1, \dots, N-1\}$, the gambler plays exactly 1 game.
\begin{itemize}
\item With probability $p$, he wins, increasing his fortune to $i+1$, from which the expected remaining games is $M_{i+1}$.
\item With probability $q$, he loses, decreasing his fortune to $i-1$, from which the expected remaining games is $M_{i-1}$.
\end{itemize}
By conditioning on the outcome of this first step (using the law of total expectation), we obtain:
\[
M_i = (1 + M_{i+1})p + (1 + M_{i-1})q = 1 + p M_{i+1} + q M_{i-1}
\]
since $p + q = 1$. This yields the given system of difference equations.
}

\begin{exercise}%[Exercise 60]
\label{ex:markov-chains-60}
Solve the equations given in Exercise 59 to obtain
\[
M_i = \begin{cases}
i(N - i), & \text{if } p = \frac{1}{2} \\
\dfrac{i}{q - p} - \dfrac{N}{q - p} \left( \dfrac{1 - (q/p)^i}{1 - (q/p)^N} \right), & \text{if } p \neq \frac{1}{2}
\end{cases}
\]
\end{exercise}
\solution{
We can rewrite the second-order difference equation $M_i = 1 + p M_{i+1} + q M_{i-1}$ as:
\[
p(M_{i+1} - M_i) - q(M_i - M_{i-1}) = -1
\]

\textbf{Case 1:} $p = q = \frac{1}{2}$
The equation simplifies to $(M_{i+1} - M_i) - (M_i - M_{i-1}) = -2$. Let $d_i = M_i - M_{i-1}$. Then $d_{i+1} - d_i = -2$, which implies $d_i = d_1 - 2(i-1)$. Summing these differences from $1$ to $i$ gives:
\[
M_i = M_0 + \sum_{k=1}^i d_k = 0 + i d_1 - 2\frac{i(i-1)}{2} = i d_1 - i^2 + i
\]
Using the boundary condition $M_N = 0$:
\[
N d_1 - N^2 + N = 0 \implies d_1 = N - 1
\]
Substituting $d_1$ back into the formula for $M_i$:
\[
M_i = i(N - 1) - i^2 + i = iN - i - i^2 + i = i(N - i)
\]

\textbf{Case 2:} $p \neq q$
The general solution to the non-homogeneous equation consists of a homogeneous part and a particular solution. The homogeneous recurrence equation is $p M_{i+1} - M_i + q M_{i-1} = 0$, which has the standard solution $A + B(q/p)^i$. For the non-homogeneous part, a linear particular solution of the form $C \cdot i$ works. Substituting $M_i = C \cdot i$ into the original recurrence gives:
\[
C \cdot i = 1 + p C(i+1) + q C(i-1) \implies 0 = 1 + C(p - q) \implies C = \frac{1}{q - p}
\]
Thus, the general solution is:
\[
M_i = A + B\left(\frac{q}{p}\right)^i + \frac{i}{q - p}
\]
Using the boundary conditions $M_0 = 0$ and $M_N = 0$:
\[
M_0 = A + B = 0 \implies B = -A
\]
\[
M_N = A - A\left(\frac{q}{p}\right)^N + \frac{N}{q - p} = 0 \implies A \left[1 - \left(\frac{q}{p}\right)^N \right] = -\frac{N}{q - p} \implies A = -\frac{N}{q - p} \frac{1}{1 - (q/p)^N}
\]
Substituting $A$ and $B = -A$ back into our expression for $M_i$:
\[
M_i = \frac{i}{q - p} - \frac{N}{q - p} \frac{1}{1 - (q/p)^N} + \frac{N}{q - p} \frac{(q/p)^i}{1 - (q/p)^N} = \frac{i}{q - p} - \frac{N}{q - p} \left( \frac{1 - (q/p)^i}{1 - (q/p)^N} \right)
\]
This completes the verification.
}

\begin{exercise}%[Exercise 61]
\label{ex:markov-chains-61}
Suppose in the gambler's ruin problem that the probability of winning a bet depends on the gambler's present fortune. Specifically, suppose that $\alpha_i$ is the probability that the gambler wins a bet when his or her fortune is $i$. Given that the gambler's initial fortune is $i$, let $P(i)$ denote the probability that the gambler's fortune reaches $N$ before 0.
\begin{enumerate}
\item[(a)] Derive a formula that relates $P(i)$ to $P(i - 1)$ and $P(i + 1)$.
\item[(b)] Using the same approach as in the gambler's ruin problem, solve the equation of part (a) for $P(i)$.
\item[(c)] Suppose that $i$ balls are initially in urn 1 and $N - i$ are in urn 2, and suppose that at each stage one of the $N$ balls is randomly chosen, taken from whichever urn it is in, and placed in the other urn. Find the probability that the first urn becomes empty before the second.
\end{enumerate}
\end{exercise}
\solution{
\textbf{(a)} By conditioning on the very first bet from state $i$:
\[
P(i) = \alpha_i P(i+1) + (1 - \alpha_i) P(i-1), \quad i = 1, \dots, N-1
\]
with boundary conditions $P(0) = 0$ and $P(N) = 1$.

\textbf{(b)} Rearranging the relation gives $\alpha_i(P(i+1) - P(i)) = (1 - \alpha_i)(P(i) - P(i-1))$, which means:
\[
P(i+1) - P(i) = \frac{1 - \alpha_i}{\alpha_i} (P(i) - P(i-1))
\]
Let $\beta_i = \frac{1 - \alpha_i}{\alpha_i}$. By iterating this relationship downwards to $P(1)-P(0)$, we find:
\[
P(i+1) - P(i) = \left( \prod_{k=1}^i \beta_k \right) P(1)
\]
Summing these intervals from $k=0$ to $i-1$, and defining the empty product as 1:
\[
P(i) = P(1) \sum_{j=0}^{i-1} \left( \prod_{k=1}^j \beta_k \right)
\]
Using the upper boundary condition $P(N) = 1$, we can find $P(1)$:
\[
1 = P(1) \sum_{j=0}^{N-1} \left( \prod_{k=1}^j \beta_k \right) \implies P(1) = \left[ \sum_{j=0}^{N-1} \left( \prod_{k=1}^j \beta_k \right) \right]^{-1}
\]
Substituting this back yields the final explicit expression for $P(i)$:
\[
P(i) = \frac{\sum_{j=0}^{i-1} \left( \prod_{k=1}^j \beta_k \right)}{\sum_{j=0}^{N-1} \left( \prod_{k=1}^j \beta_k \right)}, \quad \text{where } \beta_k = \frac{1 - \alpha_k}{\alpha_k}
\]

\textbf{(c)} This system is the Ehrenfest urn model. Let $X_n$ be the number of balls in urn 1. At state $i$, a ball is selected from urn 1 with probability $i/N$, moving the chain to $i-1$. It is selected from urn 2 with probability $(N-i)/N$, moving it to $i+1$.
Thus, the probability of increasing the count in urn 1 is $\alpha_i = \frac{N-i}{N}$. The ratio $\beta_k$ becomes:
\[
\beta_k = \frac{1 - \alpha_k}{\alpha_k} = \frac{1 - \frac{N-k}{N}}{\frac{N-k}{N}} = \frac{k}{N-k}
\]
The product term inside our equation is:
\[
\prod_{k=1}^j \beta_k = \prod_{k=1}^j \frac{k}{N-k} = \frac{j!}{(N-1)(N-2)\cdots(N-j)} = \frac{j!(N-1-j)!}{(N-1)!} = \frac{1}{\binom{N-1}{j}}
\]
We want to find the probability that urn 1 becomes empty first, which means reaching state $0$ before state $N$. This is exactly $1 - P(i)$, where $P(i)$ is the probability of reaching $N$ first. Utilizing our formula from part (b):
\[
1 - P(i) = 1 - \frac{\sum_{j=0}^{i-1} \frac{1}{\binom{N-1}{j}}}{\sum_{j=0}^{N-1} \frac{1}{\binom{N-1}{j}}} = \frac{\sum_{j=i}^{N-1} \frac{1}{\binom{N-1}{j}}}{\sum_{j=0}^{N-1} \frac{1}{\binom{N-1}{j}}}
\]
}

\setcounter{exercise}{63}
\begin{exercise}%[Exercise 64]
\label{ex:markov-chains-64}
Consider a branching process having $\mu < 1$. Show that if $X_0 = 1$, then the expected number of individuals that ever exist in this population is given by $1/(1 - \mu)$. What if $X_0 = n$?
\end{exercise}
\solution{
Let $X_k$ be the population size of the $k$-th generation. The total number of individuals that will ever exist in the population is given by the infinite sum $T = \sum_{k=0}^{\infty} X_k$.

By linearity of expectation, we can express the expected total population as:
\[
E[T] = E\left[ \sum_{k=0}^{\infty} X_k \right] = \sum_{k=0}^{\infty} E[X_k]
\]
For a standard branching process starting with a single individual ($X_0 = 1$), the expected size of the $k$-th generation is $E[X_k] = \mu^k$. Since $\mu < 1$, this forms a convergent geometric series:
\[
E[T] = \sum_{k=0}^{\infty} \mu^k = \frac{1}{1 - \mu}
\]
If the process starts instead with $X_0 = n$ individuals, each of these $n$ initial individuals initiates an independent and identically distributed branching tree process. By linearity, the total expected population is simply $n$ times the expectation of a single lineage:
\[
E[T] = \frac{n}{1 - \mu}
\]
}

\setcounter{exercise}{65}

\begin{exercise}%[Exercise 66]
\label{ex:markov-chains-66}
For a branching process, calculate $\pi_0$ when
\begin{enumerate}
\item[(a)] $P_0 = \frac{1}{4}, P_2 = \frac{3}{4}$
\item[(b)] $P_0 = \frac{1}{4}, P_1 = \frac{1}{2}, P_2 = \frac{1}{4}$
\item[(c)] $P_0 = \frac{1}{6}, P_1 = \frac{1}{2}, P_3 = \frac{1}{3}$
\end{enumerate}
\end{exercise}
\solution{
The extinction probability $\pi_0$ is the smallest non-negative root of the generating function relation $\pi_0 = \sum_{j=0}^{\infty} P_j (\pi_0)^j$.

\textbf{(a)} The equation is $\pi_0 = P_0 + P_2 \pi_0^2 = \frac{1}{4} + \frac{3}{4}\pi_0^2$. Rearranging terms:
\[
3\pi_0^2 - 4\pi_0 + 1 = 0 \implies (3\pi_0 - 1)(\pi_0 - 1) = 0
\]
The roots are $\pi_0 = \frac{1}{3}$ and $\pi_0 = 1$. The smallest non-negative root is $\pi_0 = \frac{1}{3}$.

\textbf{(b)} The equation is $\pi_0 = P_0 + P_1 \pi_0 + P_2 \pi_0^2 = \frac{1}{4} + \frac{1}{2}\pi_0 + \frac{1}{4}\pi_0^2$. Rearranging terms:
\[
\frac{1}{4}\pi_0^2 - \frac{1}{2}\pi_0 + \frac{1}{4} = 0 \implies \frac{1}{4}(\pi_0 - 1)^2 = 0
\]
This quadratic equation has a single repeated root, which is $\pi_0 = 1$.

\textbf{(c)} The equation is $\pi_0 = P_0 + P_1 \pi_0 + P_3 \pi_0^3 = \frac{1}{6} + \frac{1}{2}\pi_0 + \frac{1}{3}\pi_0^3$. Rearranging terms:
\[
2\pi_0^3 - 3\pi_0 + 1 = 0
\]
We know that $\pi_0 = 1$ is always a root of this equation. Factoring out $(\pi_0 - 1)$ via polynomial division:
\[
(\pi_0 - 1)(2\pi_0^2 + 2\pi_0 - 1) = 0
\]
Solving the remaining quadratic term $2\pi_0^2 + 2\pi_0 - 1 = 0$ using the quadratic formula:
\[
\pi_0 = \frac{-2 \pm \sqrt{4 - 4(2)(-1)}}{4} = \frac{-2 \pm \sqrt{12}}{4} = \frac{-1 \pm \sqrt{3}}{2}
\]
Since the probability must be non-negative, we discard the negative root. The two valid roots are $\pi_0 = 1$ and $\pi_0 = \frac{\sqrt{3} - 1}{2} \approx 0.366$. The smallest root is $\pi_0 = \frac{\sqrt{3} - 1}{2}$.
}

\setcounter{exercise}{67}

\begin{exercise}%[Exercise 68]
\label{ex:markov-chains-68}
\begin{enumerate}
\item[(a)] Show that the limiting probabilities of the reversed Markov chain are the same as for the forward chain by showing that they satisfy the equations
\[
\pi_j = \sum_{i} \pi_i Q_{i,j}
\]
\item[(b)] Give an intuitive explanation for the result of part (a).
\end{enumerate}
\end{exercise}
\solution{
\textbf{(a)} By definition, the transition probabilities $Q_{i,j}$ of the reversed Markov chain are related to the transition probabilities $P_{j,i}$ of the forward chain via the relation:
\[
Q_{i,j} = \frac{\pi_j P_{j,i}}{\pi_i}
\]
To verify that $\pi$ is also the stationary distribution for the reversed chain, we substitute this definition into the right-hand side of the stationary equation for the reversed chain:
\[
\sum_{i} \pi_i Q_{i,j} = \sum_{i} \pi_i \left( \frac{\pi_j P_{j,i}}{\pi_i} \right) = \sum_{i} \pi_j P_{j,i}
\]
Since $\pi_j$ does not depend on the summation index $i$, we can factor it out of the sum:
\[
= \pi_j \sum_{i} P_{j,i}
\]
Because the row sums of any transition probability matrix must equal 1 ($\sum_{i} P_{j,i} = 1$), the expression simplifies to:
\[
= \pi_j \cdot 1 = \pi_j
\]
Thus, the equation $\pi_j = \sum_{i} \pi_i Q_{i,j}$ holds, meaning that the reversed chain shares the exact same stationary probabilities $\pi$ as the forward chain.

\textbf{(b)} The limiting probability $\pi_j$ represents the long-run proportion of time that the system spends in state $j$. When we observe a Markov chain in steady state and view the sequence of states in reverse chronological order, the proportion of time the process spends in any given state $j$ does not change. Time reversal alters the sequence of transitions between states, but it cannot alter the overall long-run occupancy of the states themselves.
}

\setcounter{exercise}{69}

\begin{exercise}%[Exercise 70]
\label{ex:markov-chains-70}
A total of $m$ white and $m$ black balls are distributed among two urns, with each urn containing $m$ balls. At each stage, a ball is randomly selected from each urn and the two selected balls are interchanged. Let $X_n$ denote the number of black balls in urn 1 after the $n$th interchange.
\begin{enumerate}
\item[(a)] Give the transition probabilities of the Markov chain $X_n, n \ge 0$.
\item[(b)] Without any computations, what do you think are the limiting probabilities of this chain?
\item[(c)] Find the limiting probabilities and show that the stationary chain is time reversible.
\end{enumerate}
\end{exercise}
\solution{
\textbf{(a)} Let the state $X_n = i$ be the number of black balls in urn 1. Since each urn must contain exactly $m$ balls in total, if urn 1 has $i$ black balls, it must also contain $m - i$ white balls. Consequently, urn 2 must contain the remaining $m - i$ black balls and $i$ white balls.

At each transition, one ball is drawn at random from each urn. The state can change by at most 1:
\begin{itemize}
\item \textbf{Transition $i \to i+1$:} We must select a white ball from urn 1 and a black ball from urn 2.
\[
P_{i, i+1} = \left(\frac{m - i}{m}\right) \left(\frac{m - i}{m}\right) = \frac{(m - i)^2}{m^2}
\]
\end{itemize}
\begin{itemize}
\item \textbf{Transition $i \to i-1$:} We must select a black ball from urn 1 and a white ball from urn 2.
\[
P_{i, i-1} = \left(\frac{i}{m}\right) \left(\frac{i}{m}\right) = \frac{i^2}{m^2}
\]
\end{itemize}
\begin{itemize}
\item \textbf{Transition $i \to i$:} We either select black from both urns, or white from both urns.
\[
P_{i,i} = \left(\frac{i}{m}\right)\left(\frac{m-i}{m}\right) + \left(\frac{m-i}{m}\right)\left(\frac{i}{m}\right) = \frac{2i(m-i)}{m^2}
\]
\end{itemize}

\textbf{(b)} By symmetry and randomness, any collection of $m$ balls chosen out of the total pool of $2m$ balls (containing $m$ black and $m$ white balls) is equally likely to end up in urn 1 in the long run. Therefore, the limiting distribution should be hyper-geometric, corresponding to choosing $m$ balls from a pool of $2m$ containing $m$ black balls.

\textbf{(c)} The proposed hyper-geometric stationary distribution is:
\[
\pi_i = \frac{\binom{m}{i}\binom{m}{m-i}}{\binom{2m}{m}} = \frac{\binom{m}{i}^2}{\binom{2m}{m}}, \quad i = 0, 1, \dots, m
\]
To show that the chain is time reversible, we verify that it satisfies the detailed balance equations $\pi_i P_{i, i+1} = \pi_{i+1} P_{i+1, i}$ for all $i$:
\[
\pi_i P_{i, i+1} = \frac{\binom{m}{i}^2}{\binom{2m}{m}} \cdot \frac{(m - i)^2}{m^2}
\]
Using the definition of the binomial coefficient, we note that $\binom{m}{i}(m-i) = \frac{m!}{i!(m-i-1)!} = m \binom{m-1}{i}$. Thus:
\[
\pi_i P_{i, i+1} = \frac{1}{\binom{2m}{m}} \left[ \frac{m \binom{m-1}{i}}{m} \right]^2 = \frac{\binom{m-1}{i}^2}{\binom{2m}{m}}
\]
Now we look at the reverse direction:
\[
\pi_{i+1} P_{i+1, i} = \frac{\binom{m}{i+1}^2}{\binom{2m}{m}} \cdot \frac{(i + 1)^2}{m^2} = \frac{1}{\binom{2m}{m}} \left[ \frac{\binom{m}{i+1}(i+1)}{m} \right]^2
\]
Since $\binom{m}{i+1}(i+1) = \frac{m!}{i!(m-i-1)!} = m \binom{m-1}{i}$, we obtain:
\[
\pi_{i+1} P_{i+1, i} = \frac{1}{\binom{2m}{m}} \left[ \frac{m \binom{m-1}{i}}{m} \right]^2 = \frac{\binom{m-1}{i}^2}{\binom{2m}{m}}
\]
Because $\pi_i P_{i, i+1} = \pi_{i+1} P_{i+1, i}$ holds for all $i$, the detailed balance condition is satisfied. Therefore, the chain is time reversible and the guessed distribution is its unique stationary distribution. For $m\ge2$ there is a self-loop in an interior state, so the chain is aperiodic and this is also the limiting distribution. For $m=1$, the chain alternates between $0$ and $1$; it has stationary vector $(1/2,1/2)$ but no limiting distribution from a fixed initial state.
}

\begin{exercise}%[Exercise 71]
\label{ex:markov-chains-71}
It follows from Theorem 4.2 that for a time reversible Markov chain
\[
P_{i,j} P_{j,k} P_{k,i} = P_{i,k} P_{k,j} P_{j,i}, \quad \text{for all } i,j,k
\]
It turns out that if the state space is finite and $P_{i,j} > 0$ for all $i,j$, then the preceding is also a sufficient condition for time reversibility. (That is, in this case, we need only check the equation for paths from $i$ to $i$ that have only two intermediate states.) Prove this.
\\ \textbf{Hint:} Fix $i$ and show that the equations
\[
\pi_j P_{j,k} = \pi_k P_{k,j}
\]
are satisfied by $\pi_j = c P_{i,j}/P_{j,i}$, where $c$ is chosen so that $\sum_j \pi_j = 1$.
\end{exercise}
\solution{
Following the hint, let us define a vector $\pi$ by setting $\pi_j = c \frac{P_{i,j}}{P_{j,i}}$ for a fixed state $i$, where $c$ is a normalization constant. We must show that this choice satisfies the detailed balance equations $\pi_j P_{j,k} = \pi_k P_{k,j}$ for any pair of states $j$ and $k$.

We substitute our definition of $\pi_j$ into the left-hand side of the detailed balance equation:
\[
\pi_j P_{j,k} = \left( c \frac{P_{i,j}}{P_{j,i}} \right) P_{j,k} = c \frac{P_{i,j} P_{j,k}}{P_{j,i}}
\]
From the given sufficient condition, we are allowed to use the cycle property $P_{i,j} P_{j,k} P_{k,i} = P_{i,k} P_{k,j} P_{j,i}$. Since $P_{k,i} > 0$, we can rearrange this identity to express the numerator product as:
\[
P_{i,j} P_{j,k} = \frac{P_{i,k} P_{k,j} P_{j,i}}{P_{k,i}}
\]
Substituting this substitution back into our simplified expression for $\pi_j P_{j,k}$:
\[
\pi_j P_{j,k} = c \frac{\frac{P_{i,k} P_{k,j} P_{j,i}}{P_{k,i}}}{P_{j,i}} = c \frac{P_{i,k} P_{k,j}}{P_{k,i}} = \left( c \frac{P_{i,k}}{P_{k,i}} \right) P_{k,j}
\]
By our original definition of the vector components, the term inside the parentheses is exactly $\pi_k$. Thus, we have:
\[
\pi_j P_{j,k} = \pi_k P_{k,j}
\]
Since this detailed balance relation holds for all choices of states $j$ and $k$, the Markov chain is by definition time reversible, and the normalized vector $\pi$ is its unique stationary probability distribution.
}

\setcounter{exercise}{73}
\begin{exercise}%[Exercise 74]
\label{ex:markov-chains-74}
A group of $n$ processors is arranged in an ordered list. When a job arrives, the first processor in line attempts it; if it is unsuccessful, then the next in line tries it; if it too is unsuccessful, then the next in line tries it, and so on. When the job is successfully processed or after all processors have been unsuccessful, the job leaves the system. At this point we are allowed to reorder the processors, and a new job appears. Suppose that we use the one-closer reordering rule, which moves the processor that was successful one closer to the front of the line by interchanging its position with the one in front of it. If all processors were unsuccessful (or if the processor in the first position was successful), then the ordering remains the same. Suppose that each time processor $i$ attempts a job then, independently of anything else, it is successful with probability $p_i$.
\begin{enumerate}
\item[(a)] Define an appropriate Markov chain to analyze this model.
\item[(b)] Show that this Markov chain is time reversible.
\item[(c)] Find the long-run probabilities.
\end{enumerate}
\end{exercise}
\solution{
\textbf{(a)} The appropriate state for this system at any step is the specific ordered list of the processors. Let the state be a permutation $\sigma = (i_1, i_2, \dots, i_n)$ of the set $\{1, 2, \dots, n\}$, where $i_k$ denotes the processor currently occupying position $k$. There are $n!$ possible states.

A transition occurs when a new job is processed. Processor $i_k$ in position $k$ is called upon if and only if all processors in front of it failed, which occurs with probability $\prod_{j=1}^{k-1}(1-p_{i_j})$. Then, it succeeds with probability $p_{i_k}$, causing it to swap places with processor $i_{k-1}$.

\textbf{(b) \& (c)} Assume $0<p_i<1$ for every processor, so all adjacent swaps have positive probability and the chain is irreducible. Let $\alpha_i = \frac{p_i}{1 - p_i}$. We propose that the long-run stationary probability $\pi(\sigma)$ of being in state $\sigma = (i_1, i_2, \dots, i_n)$ is proportional to the product of these values weighted by position:
\[
\pi(i_1, i_2, \dots, i_n) = C \prod_{k=1}^n (\alpha_{i_k})^{n-k}
\]
where $C$ is a normalization constant. To establish time reversibility and prove this is the correct distribution, we check detailed balance between any two states that can transition into one another in a single step. Let $\sigma$ and $\sigma'$ be two permutations that differ only by swapping two adjacent elements at positions $k$ and $k+1$:
\[
\sigma = (\dots, i_k, i_{k+1}, \dots), \quad \sigma' = (\dots, i_{k+1}, i_k, \dots)
\]
A transition from $\sigma \to \sigma'$ happens if processor $i_{k+1}$ (currently in position $k+1$) is successful. The probability of this single event is:
\[
P_{\sigma, \sigma'} = \left[ \prod_{j=1}^{k-1} (1 - p_{i_j}) \right] (1 - p_{i_k}) p_{i_{k+1}}
\]
Conversely, a transition from $\sigma' \to \sigma$ happens if processor $i_k$ (now in position $k+1$ in state $\sigma'$) succeeds:
\[
P_{\sigma', \sigma} = \left[ \prod_{j=1}^{k-1} (1 - p_{i_j}) \right] (1 - p_{i_{k+1}}) p_{i_k}
\]
Now we evaluate the detailed balance ratio using our proposed formula for $\pi$:
\[
\frac{\pi(\sigma)}{\pi(\sigma')} = \frac{(\alpha_{i_k})^{n-k} (\alpha_{i_{k+1}})^{n-(k+1)}}{(\alpha_{i_{k+1}})^{n-k} (\alpha_{i_k})^{n-(k+1)}} = \frac{\alpha_{i_k}}{\alpha_{i_{k+1}}} = \frac{p_{i_k}/(1-p_{i_k})}{p_{i_{k+1}}/(1-p_{i_{k+1}})} = \frac{p_{i_k}(1-p_{i_{k+1}})}{p_{i_{k+1}}(1-p_{i_k})}
\]
Next, we evaluate the ratio of the transition probabilities:
\[
\frac{P_{\sigma', \sigma}}{P_{\sigma, \sigma'}} = \frac{\left[ \prod_{j=1}^{k-1} (1 - p_{i_j}) \right] (1 - p_{i_{k+1}}) p_{i_k}}{\left[ \prod_{j=1}^{k-1} (1 - p_{i_j}) \right] (1 - p_{i_k}) p_{i_{k+1}}} = \frac{p_{i_k}(1-p_{i_{k+1}})}{p_{i_{k+1}}(1-p_{i_k})}
\]
Since $\frac{\pi(\sigma)}{\pi(\sigma')} = \frac{P_{\sigma', \sigma}}{P_{\sigma, \sigma'}}$, it follows that $\pi(\sigma) P_{\sigma, \sigma'} = \pi(\sigma') P_{\sigma', \sigma}$.

Detailed balance is satisfied for all adjacent transpositions, meaning the chain is time reversible and the long-run probabilities are given by this product form.
}

\section{Counting Processes}
\label{sec:counting-processes}

\setcounter{exercise}{0}
\begin{exercise}
\label{ex:counting-processes-1}
The time $T$ required to repair a machine is an exponentially distributed random variable with mean $\frac12$ hours.
\begin{enumerate}
\item What is the probability that a repair time exceeds $\frac12$ hour?
\item What is the probability that a repair takes at least $12\frac12$ hours given that its duration exceeds $12$ hours?
\end{enumerate}
\end{exercise}

\solution{
Since $T$ is exponential with mean $\frac12$, its rate is
\[
\lambda=2.
\]
Thus
\[
P(T>t)=e^{-2t}, \qquad t\ge 0.
\]

\noindent
\textbf{(a)}
\[
P\left(T>\frac12\right)=e^{-2(1/2)}=e^{-1}.
\]

\noindent
\textbf{(b)}
Using the memoryless property of the exponential distribution,
\[
P\left(T\ge 12\frac12 \mid T>12\right)
=
P\left(T\ge \frac12\right)
=
e^{-1}.
\]
}

\begin{exercise}
\label{ex:counting-processes-2}
Suppose that you arrive at a single-teller bank to find five other customers in the bank, one being served and the other four waiting in line. You join the end of the line. If the service times are all exponential with rate $\mu$, what is the expected amount of time you will spend in the bank?
\end{exercise}

\solution{
There are five customers ahead of you: one in service and four waiting. By the memoryless property, the remaining service time of the customer currently being served is exponential with rate $\mu$.

Thus, before you leave the bank, there must be six exponential service completions: the remaining service of the customer in service, the service times of the four customers ahead of you, and your own service time.

Each has mean
\[
\frac1\mu.
\]
Therefore, the expected time you spend in the bank is
\[
6\cdot \frac1\mu=\frac6\mu.
\]
}

\begin{exercise}
\label{ex:counting-processes-3}
Let $X$ be an exponential random variable. Without any computations, tell which one of the following is correct. Explain your answer.
\begin{enumerate}
\item $E[X^2\mid X>1]=E[(X+1)^2]$,
\item $E[X^2\mid X>1]=E[X^2]+1$,
\item $E[X^2\mid X>1]=(1+E[X])^2$.
\end{enumerate}
\end{exercise}

\solution{
The correct statement is
\[
{E[X^2\mid X>1]=E[(X+1)^2].}
\]

This follows from the memoryless property. Given that $X>1$, the excess lifetime
\[
X-1
\]
has the same distribution as $X$. Hence, conditionally on $X>1$, the random variable $X$ has the same distribution as
\[
1+X.
\]
Therefore,
\[
E[X^2\mid X>1]=E[(1+X)^2].
\]
}

\begin{exercise}
\label{ex:counting-processes-4}
Consider a post office with two clerks. Three people, A, B, and C, enter simultaneously. A and B go directly to the clerks, and C waits until either A or B leaves before he begins service. What is the probability that A is still in the post office after the other two have left when
\begin{enumerate}
\item the service time for each clerk is exactly ten minutes?
\item the service times are $i$ with probability $\frac13$, $i=1,2,3$?
\item the service times are exponential with mean $1/\mu$?
\end{enumerate}
\end{exercise}

\solution{
Let $S_A,S_B,S_C$ denote the service times of A, B, and C.

\noindent
\textbf{(a)}
If every service time is exactly ten minutes, then A and B both leave at time $10$, and C then begins service and leaves at time $20$. Hence A cannot be the last to leave. Therefore,
\[
P(\text{A is last})=0.
\]

\noindent
\textbf{(b)}
Here
\[
S_A,S_B,S_C\in\{1,2,3\}
\]
independently, each value having probability $\frac13$.

For A to be last, B must finish before A, and then C must finish before A. Thus we need
\[
S_B+S_C<S_A.
\]
Since $S_A\le 3$, the only possibility is
\[
S_A=3,\qquad S_B=1,\qquad S_C=1.
\]
Therefore,
\[
P(\text{A is last})
=
\left(\frac13\right)^3
=
\frac1{27}.
\]

\noindent
\textbf{(c)}
Suppose the service times are exponential with rate $\mu$. A can be last only if B finishes before A. This happens with probability
\[
\frac12.
\]
Once B finishes, C begins service. By the memoryless property, A's remaining service time is exponential with rate $\mu$, independent of C's service time. Thus, after C begins service, the probability that C finishes before A is again
\[
\frac12.
\]
Therefore,
\[
P(\text{A is last})
=
\frac12\cdot \frac12
=
\frac14.
\]
}

\begin{exercise}
\label{ex:counting-processes-5}
The lifetime of a radio is exponentially distributed with a mean of ten years. If Jones buys a ten-year-old radio, what is the probability that it will be working after an additional ten years?
\end{exercise}

\solution{
Let $T$ be the lifetime of the radio. Since $T$ is exponential with mean $10$, its rate is
\[
\lambda=\frac1{10}.
\]
We want
\[
P(T>20\mid T>10).
\]
By the memoryless property,
\[
P(T>20\mid T>10)=P(T>10).
\]
Hence
\[
P(T>10)=e^{-(1/10)10}=e^{-1}.
\]
Therefore,
\[
{P(T>20\mid T>10)=e^{-1}.}
\]
}

\begin{exercise}
\label{ex:counting-processes-6}
Smith arrives to find two servers busy and no one else waiting. The servers work independently, with exponential service times of rate $\lambda_i$ at server $i$. Smith enters service with the first server to become free. Show that, among Smith and the two customers already in service,
\[
P\{\text{Smith is not last}\}
=
\left(\frac{\lambda_1}{\lambda_1+\lambda_2}\right)^2
+
\left(\frac{\lambda_2}{\lambda_1+\lambda_2}\right)^2.
\]
\end{exercise}

\solution{
Smith begins service when one of the two servers first becomes free.

The probability that server $1$ becomes free before server $2$ is
\[
\frac{\lambda_1}{\lambda_1+\lambda_2}.
\]
If Smith is then served by server $1$, Smith is not last precisely if his service is completed before the remaining service at server $2$. By the memoryless property, this probability is again
\[
\frac{\lambda_1}{\lambda_1+\lambda_2}.
\]
Thus the contribution from server $1$ is
\[
\left(\frac{\lambda_1}{\lambda_1+\lambda_2}\right)^2.
\]

Similarly, the contribution from server $2$ is
\[
\left(\frac{\lambda_2}{\lambda_1+\lambda_2}\right)^2.
\]

Therefore,
\[
P\{\text{Smith is not last}\}
=
\left(\frac{\lambda_1}{\lambda_1+\lambda_2}\right)^2
+
\left(\frac{\lambda_2}{\lambda_1+\lambda_2}\right)^2.
\]
}

\begin{exercise}
\label{ex:counting-processes-7}
If $X_1$ and $X_2$ are independent nonnegative continuous random variables, show that
\[
P\{X_1<X_2\mid \min(X_1,X_2)=t\}
=
\frac{r_1(t)}{r_1(t)+r_2(t)},
\]
where $r_i(t)$ is the failure rate function of $X_i$.
\end{exercise}

\solution{
Let
\[
\overline F_i(t)=P(X_i>t)
\]
and
\[
r_i(t)=\frac{f_i(t)}{\overline F_i(t)}.
\]

For small $h>0$,
\[
P(t<X_1<t+h,\ X_2>t)
\approx f_1(t)\overline F_2(t)h,
\]
and
\[
P(t<X_2<t+h,\ X_1>t)
\approx f_2(t)\overline F_1(t)h.
\]
Therefore,
\[
P\{X_1<X_2\mid \min(X_1,X_2)=t\}
=
\frac{f_1(t)\overline F_2(t)}
{f_1(t)\overline F_2(t)+f_2(t)\overline F_1(t)}.
\]
Using
\[
f_i(t)=r_i(t)\overline F_i(t),
\]
we get
\[
P\{X_1<X_2\mid \min(X_1,X_2)=t\}
=
\frac{r_1(t)\overline F_1(t)\overline F_2(t)}
{r_1(t)\overline F_1(t)\overline F_2(t)+r_2(t)\overline F_1(t)\overline F_2(t)}.
\]
Canceling the common factor gives
\[
P\{X_1<X_2\mid \min(X_1,X_2)=t\}
=
\frac{r_1(t)}{r_1(t)+r_2(t)}.
\]
}

\begin{exercise}
\label{ex:counting-processes-8}
Let $X_i$, $i=1,\ldots,n$, be independent continuous random variables, with $X_i$ having failure rate function $r_i(t)$. Let $T$ be independent of this sequence, and suppose that
\[
\sum_{i=1}^n P\{T=i\}=1.
\]
Show that the failure rate function $r(t)$ of $X_T$ is given by
\[
r(t)=\sum_{i=1}^n r_i(t)P\{T=i\mid X_T>t\}.
\]
\end{exercise}

\solution{
Let
\[
p_i=P(T=i).
\]
The survival function of $X_T$ is
\[
P(X_T>t)
=
\sum_{i=1}^n P(T=i)P(X_i>t)
=
\sum_{i=1}^n p_i \overline F_i(t).
\]
The density of $X_T$ is
\[
f_{X_T}(t)
=
\sum_{i=1}^n p_i f_i(t).
\]
Hence its failure rate is
\[
r(t)
=
\frac{f_{X_T}(t)}{P(X_T>t)}
=
\frac{\sum_{i=1}^n p_i f_i(t)}
{\sum_{i=1}^n p_i \overline F_i(t)}.
\]
Since
\[
f_i(t)=r_i(t)\overline F_i(t),
\]
we have
\[
r(t)
=
\frac{\sum_{i=1}^n p_i r_i(t)\overline F_i(t)}
{\sum_{i=1}^n p_i \overline F_i(t)}.
\]
But
\[
P(T=i\mid X_T>t)
=
\frac{p_i\overline F_i(t)}
{\sum_{j=1}^n p_j\overline F_j(t)}.
\]
Therefore,
\[
r(t)
=
\sum_{i=1}^n r_i(t)P(T=i\mid X_T>t).
\]
}

\begin{exercise}
\label{ex:counting-processes-9}
Machine $1$ is currently working. Machine $2$ will be put in use at a time $t$ from now. If the lifetime of machine $i$ is exponential with rate $\lambda_i$, $i=1,2$, what is the probability that machine $1$ is the first machine to fail?
\end{exercise}

\solution{
Let
\[
X_1\sim \operatorname{Exponential}(\lambda_1)
\]
be the remaining lifetime of machine $1$, and let
\[
X_2\sim \operatorname{Exponential}(\lambda_2)
\]
be the lifetime of machine $2$ once it is put into use.

Machine $2$ fails at time
\[
t+X_2.
\]
Thus machine $1$ fails first if
\[
X_1<t+X_2.
\]
Therefore,
\[
P(X_1<t+X_2)
=
P(X_1<t)+P(t\le X_1<t+X_2).
\]
Now,
\[
P(X_1<t)=1-e^{-\lambda_1 t}.
\]
Also,
\[
P(t\le X_1<t+X_2)
=
\int_t^\infty \lambda_1 e^{-\lambda_1 x}P(X_2>x-t)\,dx.
\]
Since
\[
P(X_2>x-t)=e^{-\lambda_2(x-t)},
\]
we get
\[
P(t\le X_1<t+X_2)
=
\int_t^\infty \lambda_1 e^{-\lambda_1 x}e^{-\lambda_2(x-t)}\,dx.
\]
Thus
\[
P(t\le X_1<t+X_2)
=
\frac{\lambda_1}{\lambda_1+\lambda_2}e^{-\lambda_1 t}.
\]
Hence
\[
P(\text{machine 1 fails first})
=
1-e^{-\lambda_1 t}
+
\frac{\lambda_1}{\lambda_1+\lambda_2}e^{-\lambda_1 t}.
\]
Simplifying,
\[
{
P(\text{machine 1 fails first})
=
1-\frac{\lambda_2}{\lambda_1+\lambda_2}e^{-\lambda_1 t}.
}
\]
}

\begin{exercise}
\label{ex:counting-processes-10}
Let $X$ and $Y$ be independent exponential random variables with respective rates $\lambda$ and $\mu$. Let $M=\min(X,Y)$. Find
\begin{enumerate}
\item $E[MX\mid M=X]$,
\item $E[MX\mid M=Y]$,
\item $\operatorname{Cov}(X,M)$.
\end{enumerate}
\end{exercise}

\solution{
Let
\[
r=\lambda+\mu.
\]
Then
\[
M\sim \operatorname{Exponential}(r),
\]
and
\[
P(M=X)=\frac{\lambda}{r},\qquad P(M=Y)=\frac{\mu}{r}.
\]

\noindent
\textbf{(a)}
On the event $\{M=X\}$, we have $X=M$. Therefore,
\[
MX=M^2.
\]
Since $M$ is exponential with rate $r$,
\[
E[M^2]=\frac{2}{r^2}.
\]
Hence
\[
{
E[MX\mid M=X]=\frac{2}{(\lambda+\mu)^2}.
}
\]

\noindent
\textbf{(b)}
On the event $\{M=Y\}$, we have $M=Y<X$. By the memoryless property,
\[
X=M+Z,
\]
where
\[
Z\sim \operatorname{Exponential}(\lambda)
\]
is independent of $M$. Hence
\[
MX=M(M+Z)=M^2+MZ.
\]
Therefore,
\[
E[MX\mid M=Y]
=
E[M^2]+E[M]E[Z].
\]
Thus
\[
E[MX\mid M=Y]
=
\frac{2}{r^2}
+
\frac1r\cdot \frac1\lambda.
\]
So
\[
{
E[MX\mid M=Y]
=
\frac{2}{(\lambda+\mu)^2}
+
\frac{1}{\lambda(\lambda+\mu)}.
}
\]

\noindent
\textbf{(c)}
Using the law of total expectation,
\[
E[XM]
=
E[XM\mid M=X]P(M=X)
+
E[XM\mid M=Y]P(M=Y).
\]
Therefore,
\[
E[XM]
=
\frac{2}{r^2}\cdot \frac{\lambda}{r}
+
\left(\frac{2}{r^2}+\frac{1}{\lambda r}\right)\frac{\mu}{r}.
\]
This simplifies to
\[
E[XM]
=
\frac{2}{r^2}+\frac{\mu}{\lambda r^2}.
\]
Also,
\[
E[X]=\frac1\lambda,
\qquad
E[M]=\frac1r.
\]
Hence
\[
\operatorname{Cov}(X,M)
=
E[XM]-E[X]E[M].
\]
Thus
\[
\operatorname{Cov}(X,M)
=
\frac{2}{r^2}+\frac{\mu}{\lambda r^2}
-
\frac1{\lambda r}.
\]
Since $r=\lambda+\mu$, this becomes
\[
{
\operatorname{Cov}(X,M)=\frac1{(\lambda+\mu)^2}.
}
\]
}

\begin{exercise}
\label{ex:counting-processes-11}
Let $X,Y_1,\ldots,Y_n$ be independent exponential random variables, $X$ having rate $\lambda$ and $Y_i$ having rate $\mu$. Let $A_i$ be the event that the $i$th smallest of these $n+1$ random variables is one of the $Y_i$. Find
\[
p=P\{X>\max_i Y_i\}
\]
by using the identity
\[
p=P(A_1\cdots A_n)=P(A_1)P(A_2\mid A_1)\cdots P(A_n\mid A_1\cdots A_{n-1}).
\]
Verify your answer when $n=2$ by conditioning on $X$ to obtain $p$.
\end{exercise}

\solution{
The event
\[
\{X>\max_i Y_i\}
\]
means that the first $n$ smallest values are all among
\[
Y_1,\ldots,Y_n.
\]
Thus
\[
p=P(A_1A_2\cdots A_n).
\]

Initially, the total rate of the $Y_i$ variables is $n\mu$, while the rate of $X$ is $\lambda$. Hence
\[
P(A_1)=\frac{n\mu}{\lambda+n\mu}.
\]
Given that the first smallest value was one of the $Y_i$, there remain $n-1$ variables of type $Y$ and the variable $X$. Therefore,
\[
P(A_2\mid A_1)=\frac{(n-1)\mu}{\lambda+(n-1)\mu}.
\]
Continuing in this way,
\[
P(A_j\mid A_1\cdots A_{j-1})
=
\frac{(n-j+1)\mu}{\lambda+(n-j+1)\mu}.
\]
Therefore,
\[
p
=
\prod_{j=1}^n
\frac{j\mu}{\lambda+j\mu}.
\]
Hence
\[
{
P\{X>\max_i Y_i\}
=
\prod_{j=1}^n
\frac{j\mu}{\lambda+j\mu}.
}
\]

For $n=2$, this gives
\[
p=
\frac{\mu}{\lambda+\mu}
\cdot
\frac{2\mu}{\lambda+2\mu}
=
\frac{2\mu^2}{(\lambda+\mu)(\lambda+2\mu)}.
\]

Now verify by conditioning on $X$. Since
\[
P(Y_1<X,Y_2<X\mid X=x)
=
(1-e^{-\mu x})^2,
\]
we have
\[
p
=
\int_0^\infty
(1-e^{-\mu x})^2\lambda e^{-\lambda x}\,dx.
\]
Expanding,
\[
p
=
\int_0^\infty
\lambda e^{-\lambda x}
\left(1-2e^{-\mu x}+e^{-2\mu x}\right)\,dx.
\]
Thus
\[
p
=
1-\frac{2\lambda}{\lambda+\mu}
+
\frac{\lambda}{\lambda+2\mu}.
\]
Simplifying,
\[
p=
\frac{2\mu^2}{(\lambda+\mu)(\lambda+2\mu)},
\]
which agrees with the formula above.
}

\begin{exercise}
\label{ex:counting-processes-12}
If $X_i$, $i=1,2,3$, are independent exponential random variables with rates $\lambda_i$, $i=1,2,3$, find
\begin{enumerate}
\item $P\{X_1<X_2<X_3\}$,
\item $P\{X_1<X_2\mid \max(X_1,X_2,X_3)=X_3\}$,
\item $E[\max X_i\mid X_1<X_2<X_3]$,
\item $E[\max X_i]$.
\end{enumerate}
\end{exercise}

\solution{
Let
\[
\Lambda=\lambda_1+\lambda_2+\lambda_3.
\]

\noindent
\textbf{(a)}
The probability that $X_1$ is the first to occur is
\[
\frac{\lambda_1}{\Lambda}.
\]
Given that $X_1$ is first, the probability that $X_2$ occurs before $X_3$ is
\[
\frac{\lambda_2}{\lambda_2+\lambda_3}.
\]
Therefore,
\[
{
P(X_1<X_2<X_3)
=
\frac{\lambda_1}{\Lambda}
\cdot
\frac{\lambda_2}{\lambda_2+\lambda_3}.
}
\]

\noindent
\textbf{(b)}
The event $\{\max(X_1,X_2,X_3)=X_3\}$ means that $X_3$ is last. The possible orders are
\[
X_1<X_2<X_3
\qquad\text{or}\qquad
X_2<X_1<X_3.
\]
Thus
\[
P(X_1<X_2\mid X_3\text{ is last})
=
\frac{P(X_1<X_2<X_3)}
{P(X_1<X_2<X_3)+P(X_2<X_1<X_3)}.
\]
Now,
\[
P(X_1<X_2<X_3)
=
\frac{\lambda_1}{\Lambda}
\frac{\lambda_2}{\lambda_2+\lambda_3},
\]
and
\[
P(X_2<X_1<X_3)
=
\frac{\lambda_2}{\Lambda}
\frac{\lambda_1}{\lambda_1+\lambda_3}.
\]
Therefore,
\[
P(X_1<X_2\mid X_3\text{ is last})
=
\frac{\frac{1}{\lambda_2+\lambda_3}}
{\frac{1}{\lambda_2+\lambda_3}+\frac{1}{\lambda_1+\lambda_3}}.
\]
Hence
\[
{
P(X_1<X_2\mid \max(X_1,X_2,X_3)=X_3)
=
\frac{\lambda_1+\lambda_3}{\lambda_1+\lambda_2+2\lambda_3}.
}
\]

\noindent
\textbf{(c)}
Given the order
\[
X_1<X_2<X_3,
\]
the spacings are independent exponentials with rates
\[
\Lambda,\qquad \lambda_2+\lambda_3,\qquad \lambda_3.
\]
Thus
\[
X_3
=
X_1+(X_2-X_1)+(X_3-X_2).
\]
Therefore,
\[
E[\max X_i\mid X_1<X_2<X_3]
=
E[X_3\mid X_1<X_2<X_3].
\]
Hence
\[
{
E[\max X_i\mid X_1<X_2<X_3]
=
\frac1{\Lambda}
+
\frac1{\lambda_2+\lambda_3}
+
\frac1{\lambda_3}.
}
\]

\noindent
\textbf{(d)}
For $t\ge 0$,
\[
P(\max X_i\le t)
=
\prod_{i=1}^3(1-e^{-\lambda_i t}).
\]
Thus
\[
E[\max X_i]
=
\int_0^\infty P(\max X_i>t)\,dt.
\]
Now,
\[
P(\max X_i>t)
=
1-\prod_{i=1}^3(1-e^{-\lambda_i t}).
\]
Expanding and integrating,
\[
{
E[\max X_i]
=
\frac1{\lambda_1}
+\frac1{\lambda_2}
+\frac1{\lambda_3}
-\frac1{\lambda_1+\lambda_2}
-\frac1{\lambda_1+\lambda_3}
-\frac1{\lambda_2+\lambda_3}
+\frac1{\lambda_1+\lambda_2+\lambda_3}.
}
\]
}

\begin{exercise}
\label{ex:counting-processes-13}
A single server is busy and there are $n\ge1$ customers waiting in
first-come, first-served order. Service times are independent exponentials
with rate $\mu$. Independently, each waiting customer abandons the queue
after an exponential patience time with rate $\theta>0$. Find the
expected time until the customer initially in position $n$ in the queue
leaves the queue, either by starting service or by abandoning it.
\end{exercise}

\solution{
Under the stated assumptions, service times are exponential with rate $\mu$, and each person waiting in line independently departs without service at rate $\theta$.

Let $E_k$ be the expected time until the tagged customer leaves the line when he is in position $k$ in line. Thus the desired quantity is $E_n$.

When the tagged customer is in position $k$, the next relevant event occurs at rate
\[
\mu+k\theta.
\]
This event is either a service completion, an abandonment by one of the $k-1$ people ahead of him, or his own abandonment.

For $k=1$,
\[
E_1=\frac1{\mu+\theta}.
\]
For $k\ge 2$,
\[
E_k
=
\frac1{\mu+k\theta}
+
\frac{\mu+(k-1)\theta}{\mu+k\theta}E_{k-1}.
\]
Multiplying by $\mu+k\theta$ gives
\[
(\mu+k\theta)E_k
=
1+(\mu+(k-1)\theta)E_{k-1}.
\]
Since
\[
(\mu+\theta)E_1=1,
\]
it follows by induction that
\[
(\mu+k\theta)E_k=k.
\]
Therefore,
\[
{
E_n=\frac{n}{\mu+n\theta}.
}
\]
}

\begin{exercise}
\label{ex:counting-processes-14}
Let $X$ be an exponential random variable with rate $\lambda$.
\begin{enumerate}
\item Use the definition of conditional expectation to determine $E[X\mid X<c]$.
\item Now determine $E[X\mid X<c]$ by using the following identity:
\[
E[X]=E[X\mid X<c]P(X<c)+E[X\mid X>c]P(X>c).
\]
\end{enumerate}
\end{exercise}

\solution{
\textbf{(a)}
By definition,
\[
E[X\mid X<c]
=
\frac{\int_0^c x\lambda e^{-\lambda x}\,dx}{P(X<c)}.
\]
Since
\[
P(X<c)=1-e^{-\lambda c},
\]
and
\[
\int_0^c x\lambda e^{-\lambda x}\,dx
=
\frac1\lambda-e^{-\lambda c}\left(c+\frac1\lambda\right),
\]
we get
\[
E[X\mid X<c]
=
\frac{\frac1\lambda-e^{-\lambda c}\left(c+\frac1\lambda\right)}
{1-e^{-\lambda c}}.
\]
Equivalently,
\[
{
E[X\mid X<c]
=
\frac1\lambda-\frac{c e^{-\lambda c}}{1-e^{-\lambda c}}
=
\frac1\lambda-\frac{c}{e^{\lambda c}-1}.
}
\]

\textbf{(b)}
We use
\[
E[X]=E[X\mid X<c]P(X<c)+E[X\mid X>c]P(X>c).
\]
Now,
\[
E[X]=\frac1\lambda,
\qquad
P(X<c)=1-e^{-\lambda c},
\qquad
P(X>c)=e^{-\lambda c}.
\]
By the memoryless property,
\[
E[X\mid X>c]=c+\frac1\lambda.
\]
Therefore,
\[
\frac1\lambda
=
E[X\mid X<c](1-e^{-\lambda c})
+
\left(c+\frac1\lambda\right)e^{-\lambda c}.
\]
Solving for $E[X\mid X<c]$ gives
\[
{
E[X\mid X<c]
=
\frac1\lambda-\frac{c e^{-\lambda c}}{1-e^{-\lambda c}}.
}
\]
}

\begin{exercise}
\label{ex:counting-processes-15}
One hundred items are simultaneously put on a life test. Suppose the lifetimes of the individual items are independent exponential random variables with mean $200$ hours. The test will end when there have been a total of $5$ failures. If $T$ is the time at which the test ends, find $E[T]$ and $\operatorname{Var}(T)$.
\end{exercise}

\solution{
Each lifetime is exponential with mean $200$, so the rate is
\[
\lambda=\frac1{200}.
\]
The time until the first failure is exponential with rate
\[
100\lambda.
\]
After the first failure, $99$ items remain, so the time until the next failure is exponential with rate
\[
99\lambda.
\]
Continuing, the time until the fifth failure can be written as a sum of independent exponential random variables with rates
\[
100\lambda,\ 99\lambda,\ 98\lambda,\ 97\lambda,\ 96\lambda.
\]
Therefore,
\[
E[T]
=
\sum_{j=96}^{100}\frac1{j\lambda}.
\]
Since $\lambda=\frac1{200}$,
\[
{
E[T]
=
200\left(
\frac1{100}+\frac1{99}+\frac1{98}+\frac1{97}+\frac1{96}
\right).
}
\]

Also,
\[
\operatorname{Var}(T)
=
\sum_{j=96}^{100}\frac1{(j\lambda)^2}.
\]
Hence
\[
{
\operatorname{Var}(T)
=
200^2\left(
\frac1{100^2}+\frac1{99^2}+\frac1{98^2}+\frac1{97^2}+\frac1{96^2}
\right).
}
\]
}

\begin{exercise}
\label{ex:counting-processes-16}
There are three jobs that need to be processed, with the processing time of job $i$ being exponential with rate $\mu_i$. There are two processors available, so processing on two of the jobs can immediately start, with processing on the final job to start when one of the initial ones is finished.
\begin{enumerate}
\item Let $T_i$ denote the time at which the processing of job $i$ is completed. If the objective is to minimize $E[T_1+T_2+T_3]$, which jobs should be initially processed if $\mu_1<\mu_2<\mu_3$?
\item Let $M$, called the makespan, be the time until all three jobs have been processed. With $S$ equal to the time that there is only a single processor working, show that
\[
2E[M]=E[S]+\sum_{i=1}^3\frac1{\mu_i}.
\]
\end{enumerate}
For the rest of this problem, suppose that $\mu_1=\mu_2=\mu$, $\mu_3=\lambda$. Also, let $P(\mu)$ be the probability that the last job to finish is either job $1$ or job $2$, and let $P(\lambda)=1-P(\mu)$ be the probability that the last job to finish is job $3$.
\begin{enumerate}
\setcounter{enumi}{2}
\item Express $E[S]$ in terms of $P(\mu)$ and $P(\lambda)$.
\item Let $P_{i,j}(\mu)$ be the value of $P(\mu)$ when $i$ and $j$ are the jobs that are initially started. Show that $P_{1,2}(\mu)\le P_{1,3}(\mu)$.
\item If $\mu>\lambda$, show that $E[M]$ is minimized when job $3$ is one of the jobs that is initially started.
\item If $\mu<\lambda$, show that $E[M]$ is minimized when processing is initially started on jobs $1$ and $2$.
\end{enumerate}
\end{exercise}

\solution{
\textbf{(a)}
Suppose jobs $i$ and $j$ are initially processed, and job $k$ waits. Then
\[
E[T_i]=\frac1{\mu_i},
\qquad
E[T_j]=\frac1{\mu_j}.
\]
The waiting job $k$ starts when the first of jobs $i$ and $j$ finishes. The expected time until that first completion is
\[
\frac1{\mu_i+\mu_j}.
\]
Therefore,
\[
E[T_k]=\frac1{\mu_i+\mu_j}+\frac1{\mu_k}.
\]
Hence
\[
E[T_1+T_2+T_3]
=
\sum_{\ell=1}^3\frac1{\mu_\ell}
+
\frac1{\mu_i+\mu_j}.
\]
The first term is fixed, so we minimize the expectation by maximizing
\[
\mu_i+\mu_j.
\]
Since
\[
\mu_1<\mu_2<\mu_3,
\]
we should initially process jobs $2$ and $3$.

\[
{\text{Initially process jobs }2\text{ and }3.}
\]

\textbf{(b)}
Let the total amount of required processing work be
\[
W=\sum_{i=1}^3 V_i,
\]
where $V_i$ is the processing time of job $i$. During the interval $[0,M]$, two processors are working except during the time interval of length $S$, when only one processor is working.

Thus the total processor work done is
\[
2(M-S)+S=2M-S.
\]
This must equal the total required work:
\[
2M-S=\sum_{i=1}^3 V_i.
\]
Taking expectations,
\[
2E[M]-E[S]=\sum_{i=1}^3 E[V_i].
\]
Since
\[
E[V_i]=\frac1{\mu_i},
\]
we get
\[
2E[M]-E[S]=\sum_{i=1}^3\frac1{\mu_i}.
\]
Therefore,
\[
{
2E[M]=E[S]+\sum_{i=1}^3\frac1{\mu_i}.
}
\]

\textbf{(c)}
The time $S$ is the final period during which only the last remaining job is being processed.

If the last job is either job $1$ or job $2$, its rate is $\mu$, so the expected value of $S$ is $\frac1\mu$. If the last job is job $3$, its rate is $\lambda$, so the expected value of $S$ is $\frac1\lambda$.

Therefore,
\[
{
E[S]=\frac{P(\mu)}{\mu}+\frac{P(\lambda)}{\lambda}.
}
\]

\textbf{(d)}
First suppose jobs $1$ and $2$ are initially started. Then job $3$ waits. After the first of jobs $1$ and $2$ finishes, job $3$ starts. At that point, one job of rate $\mu$ and one job of rate $\lambda$ remain. The last job has rate $\mu$ precisely if the job of rate $\lambda$ finishes first. Hence
\[
P_{1,2}(\mu)=\frac{\lambda}{\mu+\lambda}.
\]

Now suppose jobs $1$ and $3$ are initially started. Job $2$ waits. If job $1$ finishes first, which occurs with probability
\[
\frac{\mu}{\mu+\lambda},
\]
then job $2$ starts, and the remaining jobs have rates $\mu$ and $\lambda$. The probability that the last job has rate $\mu$ is then
\[
\frac{\lambda}{\mu+\lambda}.
\]
If job $3$ finishes first, which occurs with probability
\[
\frac{\lambda}{\mu+\lambda},
\]
then jobs $1$ and $2$ remain, both having rate $\mu$, so the last job certainly has rate $\mu$.

Therefore,
\[
P_{1,3}(\mu)
=
\frac{\mu}{\mu+\lambda}\cdot \frac{\lambda}{\mu+\lambda}
+
\frac{\lambda}{\mu+\lambda}.
\]
Thus
\[
P_{1,3}(\mu)
=
\frac{\mu\lambda}{(\mu+\lambda)^2}
+
\frac{\lambda}{\mu+\lambda}.
\]
Since
\[
P_{1,2}(\mu)=\frac{\lambda}{\mu+\lambda},
\]
we have
\[
P_{1,3}(\mu)-P_{1,2}(\mu)
=
\frac{\mu\lambda}{(\mu+\lambda)^2}\ge 0.
\]
Hence
\[
{
P_{1,2}(\mu)\le P_{1,3}(\mu).
}
\]

\textbf{(e)}
From part (b), minimizing $E[M]$ is equivalent to minimizing $E[S]$, because
\[
\sum_{i=1}^3\frac1{\mu_i}
\]
is fixed.

From part (c),
\[
E[S]=\frac{P(\mu)}{\mu}+\frac{1-P(\mu)}{\lambda}.
\]
Thus
\[
E[S]
=
\frac1\lambda
+
P(\mu)\left(\frac1\mu-\frac1\lambda\right).
\]
If
\[
\mu>\lambda,
\]
then
\[
\frac1\mu-\frac1\lambda<0.
\]
Therefore, $E[S]$ is minimized by making $P(\mu)$ as large as possible. By part (d), $P(\mu)$ is larger when job $3$ is one of the jobs initially started. Hence
\[
{
\text{if }\mu>\lambda,\text{ then }E[M]\text{ is minimized by initially starting job }3
}
\]
together with either job $1$ or job $2$.

\textbf{(f)}
Again,
\[
E[S]
=
\frac1\lambda
+
P(\mu)\left(\frac1\mu-\frac1\lambda\right).
\]
If
\[
\mu<\lambda,
\]
then
\[
\frac1\mu-\frac1\lambda>0.
\]
Therefore, $E[S]$ is minimized by making $P(\mu)$ as small as possible. By part (d), this occurs when jobs $1$ and $2$ are initially started.

Hence
\[
{
\text{if }\mu<\lambda,\text{ then }E[M]\text{ is minimized by initially starting jobs }1\text{ and }2.
}
\]
}

\setcounter{exercise}{19}
\begin{exercise}
\label{ex:counting-processes-20}
Consider a two-server system in which a customer is served first by server $1$, then by server $2$, and then departs. The service times at server $i$ are exponential random variables with rates $\mu_i$, $i=1,2$. When you arrive, you find server $1$ free and two customers at server $2$---customer A in service and customer B waiting in line.
\begin{enumerate}
\item Find $P_A$, the probability that A is still in service when you move over to server $2$.
\item Find $P_B$, the probability that B is still in the system when you move over to server $2$.
\item Find $E[T]$, where $T$ is the time that you spend in the system.
\end{enumerate}
\end{exercise}

\solution{
Let $S_1$ be your service time at server $1$, and let $R_A$ be the remaining service time of customer A at server $2$. By the memoryless property,
\[
S_1\sim \operatorname{Exp}(\mu_1),
\qquad
R_A\sim \operatorname{Exp}(\mu_2).
\]

\textbf{(a)} A is still in service when you move to server $2$ if
\[
S_1<R_A.
\]
Thus
\[
P_A=P(S_1<R_A)=\frac{\mu_1}{\mu_1+\mu_2}.
\]

\textbf{(b)} Let $S_B$ be B's service time at server $2$. Then B is still in the system when you move to server $2$ if
\[
S_1<R_A+S_B.
\]
Since $R_A$ and $S_B$ are independent exponentials with rate $\mu_2$,
\[
P_B=P(S_1<R_A+S_B).
\]
Using conditioning,
\[
P_B=1-E\left[e^{-\mu_1(R_A+S_B)}\right].
\]
Now
\[
E[e^{-\mu_1R_A}]
=
\frac{\mu_2}{\mu_1+\mu_2}.
\]
Therefore,
\[
P_B
=
1-\left(\frac{\mu_2}{\mu_1+\mu_2}\right)^2.
\]

\textbf{(c)} Write
\[
T=S_1+S_2+W_A+W_B,
\]
where $S_2$ is your own service time at server $2$, $W_A$ is the time you wait while A is being served, and $W_B$ is the time you wait while B is being served.

Clearly,
\[
E[S_1]=\frac1{\mu_1},
\qquad
E[S_2]=\frac1{\mu_2}.
\]
Also,
\[
W_A=(R_A-S_1)^+.
\]
By the memoryless property,
\[
E[W_A]
=
P(S_1<R_A)\frac1{\mu_2}
=
\frac{\mu_1}{\mu_1+\mu_2}\cdot \frac1{\mu_2}.
\]
Similarly, if B is still in the system when you arrive at server $2$, the expected amount of service of B that remains to be waited through is $\frac1{\mu_2}$. Hence
\[
E[W_B]
=
P_B\frac1{\mu_2}
=
\frac1{\mu_2}\left[
1-\left(\frac{\mu_2}{\mu_1+\mu_2}\right)^2
\right].
\]
Therefore,
\[
{
E[T]
=
\frac1{\mu_1}
+\frac1{\mu_2}
+
\frac{\mu_1}{\mu_2(\mu_1+\mu_2)}
+
\frac1{\mu_2}\left[
1-\left(\frac{\mu_2}{\mu_1+\mu_2}\right)^2
\right].
}
\]
}

\begin{exercise}
\label{ex:counting-processes-21}
In a certain system, a customer must first be served by server $1$ and then by server $2$. The service times at server $i$ are exponential with rate $\mu_i$, $i=1,2$. An arrival finding server $1$ busy waits in line for that server. Upon completion of service at server $1$, a customer either enters service with server $2$ if that server is free or else remains with server $1$ until server $2$ is free. Customers depart the system after being served by server $2$. Suppose that when you arrive there is one customer in the system and that customer is being served by server $1$. What is the expected total time you spend in the system?
\end{exercise}

\solution{
Let A be the customer initially at server $1$. You must first wait for A to finish service at server $1$, which has mean
\[
\frac1{\mu_1}.
\]
Then you enter service at server $1$, while A moves to server $2$.

Let $S_1$ be your service time at server $1$ and let $R_2$ be A's service time at server $2$. Then
\[
S_1\sim \operatorname{Exp}(\mu_1),
\qquad
R_2\sim \operatorname{Exp}(\mu_2).
\]
If you finish server $1$ before A finishes server $2$, you are blocked and must wait. The expected blocking time is
\[
E[(R_2-S_1)^+]
=
P(S_1<R_2)\frac1{\mu_2}
=
\frac{\mu_1}{\mu_1+\mu_2}\cdot \frac1{\mu_2}.
\]
Finally, your own service time at server $2$ has mean $\frac1{\mu_2}$.

Thus the expected total time is
\[
{
\frac1{\mu_1}
+\frac1{\mu_1}
+
\frac{\mu_1}{\mu_2(\mu_1+\mu_2)}
+
\frac1{\mu_2}
}
\]
or
\[
{
E[T]
=
\frac2{\mu_1}
+
\frac1{\mu_2}
+
\frac{\mu_1}{\mu_2(\mu_1+\mu_2)}.
}
\]
}

\begin{exercise}
\label{ex:counting-processes-22}
Suppose in Exercise $21$ you arrive to find two others in the system, one being served by server $1$ and one by server $2$. What is the expected time you spend in the system? Recall that if server $1$ finishes before server $2$, then server $1$'s customer will remain with him until server $2$ becomes free.
\end{exercise}

\solution{
Let $U$ be the remaining service time of the customer at server $1$, and let $V$ be the remaining service time of the customer at server $2$. Then
\[
U\sim \operatorname{Exp}(\mu_1),
\qquad
V\sim \operatorname{Exp}(\mu_2).
\]
You cannot enter server $1$ until the customer there has completed service at server $1$ and has moved to server $2$. Thus your waiting time before entering server $1$ is
\[
\max(U,V).
\]
Hence
\[
E[\max(U,V)]
=
E[U]+E[V]-E[\min(U,V)].
\]
Since
\[
E[U]=\frac1{\mu_1},
\qquad
E[V]=\frac1{\mu_2},
\qquad
E[\min(U,V)]=\frac1{\mu_1+\mu_2},
\]
we get
\[
E[\max(U,V)]
=
\frac1{\mu_1}
+
\frac1{\mu_2}
-
\frac1{\mu_1+\mu_2}.
\]

After you enter server $1$, the situation is the same as in Exercise $21$: you have service at server $1$, possible blocking before server $2$, and service at server $2$. Therefore,
\[
E[T]
=
E[\max(U,V)]
+
\frac1{\mu_1}
+
\frac{\mu_1}{\mu_2(\mu_1+\mu_2)}
+
\frac1{\mu_2}.
\]
Thus
\[
{
E[T]
=
\frac2{\mu_1}
+
\frac2{\mu_2}
-
\frac1{\mu_1+\mu_2}
+
\frac{\mu_1}{\mu_2(\mu_1+\mu_2)}.
}
\]
}

\begin{exercise}
\label{ex:counting-processes-23}
A flashlight needs two batteries to be operational. Consider such a flashlight along with a set of $n$ functional batteries---battery $1$, battery $2,\ldots,$ battery $n$. Initially, battery $1$ and battery $2$ are installed. Whenever a battery fails, it is immediately replaced by the lowest numbered functional battery that has not yet been put in use. Suppose that the lifetimes of the different batteries are independent exponential random variables each having rate $\mu$. At a random time, call it $T$, a battery will fail and our stockpile will be empty. At that moment exactly one of the batteries---which we call battery $X$---will not yet have failed.
\begin{enumerate}
\item What is $P\{X=n\}$?
\item What is $P\{X=1\}$?
\item What is $P\{X=i\}$?
\item Find $E[T]$.
\item What is the distribution of $T$?
\end{enumerate}
\end{exercise}

\solution{
At any time before $T$, two batteries are in use. Since their lifetimes are independent exponentials with the same rate $\mu$, each of the two currently installed batteries is equally likely to be the next one to fail.

Thus the process can be viewed as a sequence of pairwise competitions: the surviving battery competes with the next newly installed battery, and each has probability $\frac12$ of surviving.

\textbf{(a)} Battery $n$ is installed last. At the final failure time, battery $n$ is one of the two batteries in use, so it survives with probability
\[
{
P\{X=n\}=\frac12.
}
\]

\textbf{(b)} Battery $1$ must survive every competition against batteries $2,3,\ldots,n$. Therefore,
\[
{
P\{X=1\}=\left(\frac12\right)^{n-1}.
}
\]

\textbf{(c)} Similarly,
\[
P\{X=2\}=\left(\frac12\right)^{n-1}.
\]
For $i=3,4,\ldots,n$, battery $i$ must first beat the current surviving battery when it is installed, and then must survive against batteries $i+1,\ldots,n$. Hence
\[
{
P\{X=i\}=
\left(\frac12\right)^{n-i+1},
\qquad i=3,\ldots,n.
}
\]
Thus, in full,
\[
P\{X=i\}
=
\begin{cases}
2^{-(n-1)}, & i=1,2,\\[4pt]
2^{-(n-i+1)}, & i=3,\ldots,n.
\end{cases}
\]

\textbf{(d)} Until time $T$, there are always two batteries in use. Therefore, the time between successive failures is exponential with rate
\[
2\mu.
\]
The time $T$ occurs at the $(n-1)$st failure. Hence $T$ is the sum of $n-1$ independent exponential random variables with rate $2\mu$. Therefore,
\[
{
E[T]=\frac{n-1}{2\mu}.
}
\]

\textbf{(e)} Since $T$ is the sum of $n-1$ independent exponentials with rate $2\mu$,
\[
{
T\sim \operatorname{Gamma}(n-1,2\mu).
}
\]
Thus its density is
\[
f_T(t)
=
\frac{(2\mu)^{n-1}}{(n-2)!}t^{n-2}e^{-2\mu t},
\qquad t>0,
\]
and its distribution function is
\[
F_T(t)
=
1-e^{-2\mu t}\sum_{k=0}^{n-2}\frac{(2\mu t)^k}{k!},
\qquad t\ge 0.
\]
}

\begin{exercise}
\label{ex:counting-processes-24}
There are two servers available to process $n$ jobs. Initially, each server begins work on a job. Whenever a server completes work on a job, that job leaves the system and the server begins processing a new job, provided there are still jobs waiting to be processed. Let $T$ denote the time until all jobs have been processed. If the time that it takes server $i$ to process a job is exponentially distributed with rate $\mu_i$, $i=1,2$, find $E[T]$ and $\operatorname{Var}(T)$.
\end{exercise}

\solution{
Let
\[
\mu=\mu_1+\mu_2.
\]
While there are still waiting jobs, both servers are busy. The time between job completions is then exponential with rate $\mu_1+\mu_2$.

For the first $n-2$ completions, both servers remain busy. Hence the total time for these first $n-2$ completions is the sum of $n-2$ independent exponentials with rate $\mu$.

After that, no jobs are waiting, and there remain two jobs in service. By the memoryless property, their remaining service times are independent exponentials with rates $\mu_1$ and $\mu_2$. The final time is therefore the maximum of these two remaining service times.

Let
\[
M=\max(X_1,X_2),
\]
where
\[
X_1\sim \operatorname{Exp}(\mu_1),
\qquad
X_2\sim \operatorname{Exp}(\mu_2).
\]
Then
\[
E[M]
=
\frac1{\mu_1}
+
\frac1{\mu_2}
-
\frac1{\mu_1+\mu_2}.
\]
Thus
\[
E[T]
=
\frac{n-2}{\mu_1+\mu_2}
+
\frac1{\mu_1}
+
\frac1{\mu_2}
-
\frac1{\mu_1+\mu_2}.
\]
Therefore,
\[
{
E[T]
=
\frac1{\mu_1}
+
\frac1{\mu_2}
+
\frac{n-3}{\mu_1+\mu_2}.
}
\]

For the variance,
\[
\operatorname{Var}\left(\sum_{j=1}^{n-2} Z_j\right)
=
\frac{n-2}{(\mu_1+\mu_2)^2},
\]
where $Z_j\sim \operatorname{Exp}(\mu_1+\mu_2)$.

Also,
\[
\operatorname{Var}(M)
=
\frac1{\mu_1^2}
+
\frac1{\mu_2^2}
-
\frac{3}{(\mu_1+\mu_2)^2}.
\]
Hence
\[
\operatorname{Var}(T)
=
\frac{n-2}{(\mu_1+\mu_2)^2}
+
\frac1{\mu_1^2}
+
\frac1{\mu_2^2}
-
\frac{3}{(\mu_1+\mu_2)^2}.
\]
Thus
\[
{
\operatorname{Var}(T)
=
\frac1{\mu_1^2}
+
\frac1{\mu_2^2}
+
\frac{n-5}{(\mu_1+\mu_2)^2}.
}
\]
}

\begin{exercise}
\label{ex:counting-processes-25}
Customers can be served by any of three servers, where the service times of server $i$ are exponentially distributed with rate $\mu_i$, $i=1,2,3$. Whenever a server becomes free, the customer who has been waiting the longest begins service with that server.
\begin{enumerate}
\item If you arrive to find all three servers busy and no one waiting, find the expected time until you depart the system.
\item If you arrive to find all three servers busy and one person waiting, find the expected time until you depart the system.
\end{enumerate}
\end{exercise}

\solution{
Let
\[
\Lambda=\mu_1+\mu_2+\mu_3.
\]

\textbf{(a)} Since all three servers are busy, the time until the first server becomes free is exponential with rate $\Lambda$, so its mean is
\[
\frac1{\Lambda}.
\]
The probability that server $i$ becomes free first is
\[
\frac{\mu_i}{\Lambda}.
\]
If you enter service with server $i$, your expected service time is $\frac1{\mu_i}$. Therefore, the expected service time you receive is
\[
\sum_{i=1}^3 \frac{\mu_i}{\Lambda}\frac1{\mu_i}
=
\frac3{\Lambda}.
\]
Thus
\[
{
E[T]=\frac1{\Lambda}+\frac3{\Lambda}=\frac4{\Lambda}.
}
\]

\textbf{(b)} If one person is waiting ahead of you, then you must wait for two service completions before you enter service. Each time all three servers are busy, the time until the next completion has mean $\frac1{\Lambda}$. Hence your expected waiting time is
\[
\frac2{\Lambda}.
\]
Once you enter service, your expected service time is again
\[
\frac3{\Lambda}.
\]
Therefore,
\[
{
E[T]=\frac2{\Lambda}+\frac3{\Lambda}=\frac5{\Lambda}.
}
\]
}

\begin{exercise}
\label{ex:counting-processes-26}
Each entering customer must be served first by server $1$, then by server $2$, and finally by server $3$. The amount of time it takes to be served by server $i$ is an exponential random variable with rate $\mu_i$, $i=1,2,3$. Suppose you enter the system when it contains a single customer who is being served by server $3$.
\begin{enumerate}
\item Find the probability that server $3$ will still be busy when you move over to server $2$.
\item Find the probability that server $3$ will still be busy when you move over to server $3$.
\item Find the expected amount of time that you spend in the system.
\item Suppose that you enter the system when it contains a single customer who is being served by server $2$. Find the expected amount of time that you spend in the system.
\end{enumerate}
\end{exercise}

\solution{
Let $R_3$ be the remaining service time of the customer initially at server $3$. Then
\[
R_3\sim \operatorname{Exp}(\mu_3).
\]

\textbf{(a)} You move to server $2$ after completing service at server $1$. Let
\[
S_1\sim \operatorname{Exp}(\mu_1).
\]
Server $3$ is still busy at that time if
\[
S_1<R_3.
\]
Therefore,
\[
{
P(S_1<R_3)=\frac{\mu_1}{\mu_1+\mu_3}.
}
\]

\textbf{(b)} You move to server $3$ after completing service at servers $1$ and $2$. Let
\[
S_2\sim \operatorname{Exp}(\mu_2).
\]
Then server $3$ is still busy if
\[
S_1+S_2<R_3.
\]
Thus
\[
P(S_1+S_2<R_3)
=
E[e^{-\mu_3(S_1+S_2)}].
\]
Since $S_1$ and $S_2$ are independent,
\[
{
P(S_1+S_2<R_3)
=
\frac{\mu_1}{\mu_1+\mu_3}
\cdot
\frac{\mu_2}{\mu_2+\mu_3}.
}
\]

\textbf{(c)} Your own expected service time is
\[
\frac1{\mu_1}+\frac1{\mu_2}+\frac1{\mu_3}.
\]
The only possible waiting occurs at server $3$. The expected waiting time there is
\[
E[(R_3-S_1-S_2)^+].
\]
By the memoryless property,
\[
E[(R_3-S_1-S_2)^+]
=
P(S_1+S_2<R_3)\frac1{\mu_3}.
\]
Therefore,
\[
{
E[T]
=
\frac1{\mu_1}
+
\frac1{\mu_2}
+
\frac1{\mu_3}
+
\frac1{\mu_3}
\frac{\mu_1}{\mu_1+\mu_3}
\frac{\mu_2}{\mu_2+\mu_3}.
}
\]

\textbf{(d)} Now suppose the single customer initially in the system is at server $2$. Let $A$ be that customer's remaining service time at server $2$:
\[
A\sim \operatorname{Exp}(\mu_2).
\]
You first complete service at server $1$, taking time
\[
S_1\sim \operatorname{Exp}(\mu_1).
\]
Your expected waiting time at server $2$ is
\[
E[(A-S_1)^+]
=
P(S_1<A)\frac1{\mu_2}
=
\frac{\mu_1}{\mu_1+\mu_2}\cdot \frac1{\mu_2}.
\]

Now consider the possible waiting at server $3$. If $S_1<A$, then you and the other customer begin service at servers $2$ and $3$ at the same time. Your expected waiting at server $3$ in this case is
\[
\frac{\mu_2}{\mu_2+\mu_3}\cdot \frac1{\mu_3}.
\]
If $A<S_1$, then by the memoryless property the remaining time until you finish server $1$ after the other customer leaves server $2$ is exponential with rate $\mu_1$. Hence the expected waiting at server $3$ in this case is
\[
\frac{\mu_1}{\mu_1+\mu_3}
\frac{\mu_2}{\mu_2+\mu_3}
\frac1{\mu_3}.
\]
Therefore,
\[
E[W_3]
=
\frac{\mu_1}{\mu_1+\mu_2}
\frac{\mu_2}{\mu_2+\mu_3}
\frac1{\mu_3}
+
\frac{\mu_2}{\mu_1+\mu_2}
\frac{\mu_1}{\mu_1+\mu_3}
\frac{\mu_2}{\mu_2+\mu_3}
\frac1{\mu_3}.
\]

Thus the expected total time is
\[
{
E[T]
=
\frac1{\mu_1}
+
\frac1{\mu_2}
+
\frac1{\mu_3}
+
\frac{\mu_1}{\mu_2(\mu_1+\mu_2)}
+
E[W_3].
}
\]
}

\setcounter{exercise}{33}
\begin{exercise}
\label{ex:counting-processes-34}
Two individuals, A and B, both require kidney transplants. If she does not receive a new kidney, then A will die after an exponential time with rate $\mu_A$, and B after an exponential time with rate $\mu_B$. New kidneys arrive in accordance with a Poisson process having rate $\lambda$. It has been decided that the first kidney will go to A, or to B if B is alive and A is not at that time, and the next one to B if still living.
\begin{enumerate}
\item What is the probability that A obtains a new kidney?
\item What is the probability that B obtains a new kidney?
\end{enumerate}
\end{exercise}

\solution{
Let $K_1$ be the time until the first kidney arrives. Then
\[
K_1\sim \operatorname{Exp}(\lambda).
\]
Let $D_A$ and $D_B$ be the death times of A and B, with
\[
D_A\sim \operatorname{Exp}(\mu_A),
\qquad
D_B\sim \operatorname{Exp}(\mu_B).
\]

\textbf{(a)} A obtains a kidney if she is alive when the first kidney arrives:
\[
K_1<D_A.
\]
Therefore,
\[
{
P(\text{A obtains a kidney})
=
\frac{\lambda}{\lambda+\mu_A}.
}
\]

\textbf{(b)} B obtains a kidney in either of two ways.

First, A may die before the first kidney arrives, while B is still alive at that time. This event is
\[
D_A<K_1<D_B.
\]
Its probability is
\[
\int_0^\infty \lambda e^{-\lambda t}
(1-e^{-\mu_A t})e^{-\mu_B t}\,dt
=
\frac{\lambda}{\lambda+\mu_B}
-
\frac{\lambda}{\lambda+\mu_A+\mu_B}.
\]

Second, A may receive the first kidney, and B may still be alive when the second kidney arrives. If $K_2$ is the time between the first and second kidney arrivals, then $K_2\sim \operatorname{Exp}(\lambda)$, and this event has probability
\[
P(K_1<D_A,\ K_1+K_2<D_B).
\]
Thus
\[
P(K_1<D_A,\ K_1+K_2<D_B)
=
\int_0^\infty\int_0^\infty
\lambda^2 e^{-\lambda(x+y)}
e^{-\mu_A x}
e^{-\mu_B(x+y)}
\,dy\,dx.
\]
Hence
\[
P(K_1<D_A,\ K_1+K_2<D_B)
=
\frac{\lambda^2}{(\lambda+\mu_A+\mu_B)(\lambda+\mu_B)}.
\]
Therefore,
\[
P(\text{B obtains a kidney})
=
\frac{\lambda}{\lambda+\mu_B}
-
\frac{\lambda}{\lambda+\mu_A+\mu_B}
+
\frac{\lambda^2}{(\lambda+\mu_A+\mu_B)(\lambda+\mu_B)}.
\]
Simplifying,
\[
{
P(\text{B obtains a kidney})
=
\frac{\lambda(\lambda+\mu_A)}
{(\lambda+\mu_B)(\lambda+\mu_A+\mu_B)}.
}
\]
}

\begin{exercise}
\label{ex:counting-processes-35}
Let $\{N(t):t\ge0\}$ be a right-continuous counting process with $N(0)=0$,
independent increments, and
$N(s+t)-N(s)\sim\operatorname{Poisson}(\lambda t)$ for every $s,t\ge0$.
Show that its successive interarrival times are independent exponential
random variables with rate $\lambda$.
\end{exercise}

\solution{
The assumptions give $N(0)=0$, the process has independent increments, and
\[
N(t+s)-N(s)\sim \operatorname{Poisson}(\lambda t).
\]

Let $S_n$ be the time of the $n$th event, and let
\[
T_n=S_n-S_{n-1},
\qquad S_0=0,
\]
be the interarrival times.

For the first interarrival time,
\[
P(T_1>t)=P(N(t)=0)=e^{-\lambda t}.
\]
Therefore,
\[
T_1\sim \operatorname{Exp}(\lambda).
\]

Now condition on the past up to time $S_{n-1}$. For $t\ge 0$,
\[
P(T_n>t\mid \text{past up to }S_{n-1})
=
P(N(S_{n-1}+t)-N(S_{n-1})=0).
\]
The strong Markov property at $S_{n-1}$ gives the same fresh Poisson increment law. This property follows from independent stationary increments at deterministic times by approximating $S_{n-1}$ from above by dyadic-grid stopping times and using right continuity. Thus
\[
P(N(S_{n-1}+t)-N(S_{n-1})=0)=e^{-\lambda t}.
\]
The same stopping-time argument shows that this probability is independent of the past. Hence
\[
T_n\sim \operatorname{Exp}(\lambda),
\]
and the interarrival times are independent.

Thus the event times of the process are generated by independent exponential interarrival times with rate $\lambda$, as required.
}

\begin{exercise}
\label{ex:counting-processes-36}
Let $S(t)$ denote the price of a security at time $t$. A popular model for the process $\{S(t),t\ge 0\}$ supposes that the price remains unchanged until a shock occurs, at which time the price is multiplied by a random factor. If we let $N(t)$ denote the number of shocks by time $t$, and let $X_i$ denote the $i$th multiplicative factor, then this model supposes that
\[
S(t)=S(0)\prod_{i=1}^{N(t)}X_i,
\]
where $\prod_{i=1}^{N(t)}X_i$ is equal to $1$ when $N(t)=0$. Suppose that the $X_i$ are independent exponential random variables with rate $\mu$, that $\{N(t),t\ge 0\}$ is a Poisson process with rate $\lambda$, that $\{N(t),t\ge 0\}$ is independent of the $X_i$, and that $S(0)=s$.
\begin{enumerate}
\item Find $E[S(t)]$.
\item Find $E[S^2(t)]$.
\end{enumerate}
\end{exercise}

\solution{
Condition on $N(t)=n$. Since the $X_i$ are independent,
\[
E\left[\prod_{i=1}^n X_i\right]
=
(E[X])^n.
\]
Because $X_i\sim \operatorname{Exp}(\mu)$,
\[
E[X]=\frac1\mu.
\]
Thus
\[
E[S(t)\mid N(t)=n]
=
s\left(\frac1\mu\right)^n.
\]
Since $N(t)\sim \operatorname{Poisson}(\lambda t)$,
\[
E[S(t)]
=
sE\left[\left(\frac1\mu\right)^{N(t)}\right].
\]
Using the probability generating function of a Poisson random variable,
\[
E[a^{N(t)}]=e^{\lambda t(a-1)}.
\]
Therefore,
\[
{
E[S(t)]
=
s\exp\left\{\lambda t\left(\frac1\mu-1\right)\right\}.
}
\]

For the second moment,
\[
S^2(t)=s^2\prod_{i=1}^{N(t)}X_i^2.
\]
Since
\[
E[X^2]=\frac2{\mu^2},
\]
we obtain
\[
E[S^2(t)]
=
s^2E\left[\left(\frac2{\mu^2}\right)^{N(t)}\right].
\]
Hence
\[
{
E[S^2(t)]
=
s^2
\exp\left\{
\lambda t\left(\frac2{\mu^2}-1\right)
\right\}.
}
\]
}

\begin{exercise}
\label{ex:counting-processes-37}
A machine works for an exponentially distributed time with rate $\mu$ and then fails. A repair crew checks the machine at times distributed according to a Poisson process with rate $\lambda$; if the machine is found to have failed then it is immediately replaced. Find the expected time between replacements of machines.
\end{exercise}

\solution{
Let $F$ be the time until the machine fails. Then
\[
F\sim \operatorname{Exp}(\mu),
\]
so
\[
E[F]=\frac1\mu.
\]

After the machine fails, the replacement does not occur until the next inspection time. Since inspections occur according to a Poisson process with rate $\lambda$, the time from the failure until the next inspection is exponential with rate $\lambda$. Its mean is
\[
\frac1\lambda.
\]

Therefore, the time between replacements is the failure time plus the time from failure to the next inspection. Hence
\[
{
E[\text{time between replacements}]
=
\frac1\mu+\frac1\lambda.
}
\]
}

\begin{exercise}
\label{ex:counting-processes-38}
Let $\{M_i(t),t\ge 0\}$, $i=1,2,3$, be independent Poisson processes with respective rates $\lambda_i$, $i=1,2,3$, and set
\[
N_1(t)=M_1(t)+M_2(t),
\qquad
N_2(t)=M_2(t)+M_3(t).
\]
The stochastic process $\{(N_1(t),N_2(t)),t\ge 0\}$ is called a bivariate Poisson process.
\begin{enumerate}
\item Find $P\{N_1(t)=n,N_2(t)=m\}$.
\item Find $\operatorname{Cov}(N_1(t),N_2(t))$.
\end{enumerate}
\end{exercise}

\solution{
\textbf{(a)} Since
\[
N_1(t)=M_1(t)+M_2(t),
\qquad
N_2(t)=M_2(t)+M_3(t),
\]
we condition on
\[
M_2(t)=k.
\]
Then
\[
M_1(t)=n-k,
\qquad
M_3(t)=m-k.
\]
Thus $k$ must satisfy
\[
0\le k\le \min(n,m).
\]
Since $M_1,M_2,M_3$ are independent,
\[
P\{N_1(t)=n,N_2(t)=m\}
\]
equals
\[
\sum_{k=0}^{\min(n,m)}
P(M_1(t)=n-k)P(M_2(t)=k)P(M_3(t)=m-k).
\]
Therefore,
\[
{
P\{N_1(t)=n,N_2(t)=m\}
=
e^{-(\lambda_1+\lambda_2+\lambda_3)t}
\sum_{k=0}^{\min(n,m)}
\frac{(\lambda_1t)^{n-k}}{(n-k)!}
\frac{(\lambda_2t)^k}{k!}
\frac{(\lambda_3t)^{m-k}}{(m-k)!}.
}
\]

\textbf{(b)} We have
\[
\operatorname{Cov}(N_1(t),N_2(t))
=
\operatorname{Cov}(M_1(t)+M_2(t),M_2(t)+M_3(t)).
\]
All processes are independent except for the common term $M_2(t)$, so
\[
\operatorname{Cov}(N_1(t),N_2(t))
=
\operatorname{Var}(M_2(t)).
\]
Since
\[
M_2(t)\sim \operatorname{Poisson}(\lambda_2t),
\]
we get
\[
{
\operatorname{Cov}(N_1(t),N_2(t))=\lambda_2t.
}
\]
}

\begin{exercise}
\label{ex:counting-processes-39}
A certain scientific theory supposes that mistakes in cell division occur according to a Poisson process with rate $2.5$ per year, and that an individual dies when $196$ such mistakes have occurred. Assuming this theory, find
\begin{enumerate}
\item the mean lifetime of an individual,
\item the variance of the lifetime of an individual.
\end{enumerate}
Also approximate
\begin{enumerate}
\setcounter{enumi}{2}
\item the probability that an individual dies before age $67.2$,
\item the probability that an individual reaches age $90$,
\item the probability that an individual reaches age $100$.
\end{enumerate}
\end{exercise}

\solution{
All normal-tail calculations below are approximations. Let $T$ be the time of the $196$th event in a Poisson process with rate
\[
\lambda=2.5.
\]
Then
\[
T\sim \operatorname{Gamma}(196,2.5).
\]

\textbf{(a)} Hence
\[
{
E[T]=\frac{196}{2.5}=78.4.
}
\]

\textbf{(b)}
\[
{
\operatorname{Var}(T)=\frac{196}{(2.5)^2}=31.36.
}
\]
Thus
\[
\sigma_T=\sqrt{31.36}=5.6.
\]

For the approximations, use
\[
T\approx N(78.4,31.36).
\]

\textbf{(c)}
\[
P(T<67.2)
=
P\left(Z<\frac{67.2-78.4}{5.6}\right)
=
P(Z<-2).
\]
Thus
\[
{
P(T<67.2)\approx 0.0228.
}
\]

\textbf{(d)}
\[
P(T>90)
=
P\left(Z>\frac{90-78.4}{5.6}\right)
=
P(Z>2.0714).
\]
Hence
\[
{
P(T>90)\approx 0.0192.
}
\]

\textbf{(e)}
\[
P(T>100)
=
P\left(Z>\frac{100-78.4}{5.6}\right)
=
P(Z>3.8571).
\]
Therefore,
\[
{
P(T>100)\approx 0.00006.
}
\]
}

\begin{exercise}
\label{ex:counting-processes-40}
Show that if $\{N_i(t),t\ge 0\}$ are independent Poisson processes with rate $\lambda_i$, $i=1,2$, then $\{N(t),t\ge 0\}$ is a Poisson process with rate $\lambda_1+\lambda_2$, where
\[
N(t)=N_1(t)+N_2(t).
\]
\end{exercise}

\solution{
Clearly,
\[
N(0)=N_1(0)+N_2(0)=0.
\]

For $0\le s<t$,
\[
N(t)-N(s)
=
[N_1(t)-N_1(s)]+[N_2(t)-N_2(s)].
\]
Since $N_1$ and $N_2$ are independent Poisson processes, the two increments are independent, with
\[
N_1(t)-N_1(s)\sim \operatorname{Poisson}(\lambda_1(t-s)),
\]
and
\[
N_2(t)-N_2(s)\sim \operatorname{Poisson}(\lambda_2(t-s)).
\]
The sum of independent Poisson random variables is Poisson, so
\[
N(t)-N(s)
\sim
\operatorname{Poisson}((\lambda_1+\lambda_2)(t-s)).
\]
Also, increments over disjoint intervals are independent because the corresponding increments of $N_1$ and $N_2$ are independent.

Therefore,
\[
{
\{N(t),t\ge 0\}
}
\]
is a Poisson process with rate
\[
{\lambda_1+\lambda_2.}
\]
}

\begin{exercise}
\label{ex:counting-processes-41}
In Exercise $40$, what is the probability that the first event of the combined process is from the $N_1$ process?
\end{exercise}

\solution{
Let $T_1$ be the time of the first event of process $N_1$, and let $T_2$ be the time of the first event of process $N_2$. Then
\[
T_1\sim \operatorname{Exp}(\lambda_1),
\qquad
T_2\sim \operatorname{Exp}(\lambda_2),
\]
independently.

The first event of the combined process is from $N_1$ if
\[
T_1<T_2.
\]
Therefore,
\[
{
P(T_1<T_2)=\frac{\lambda_1}{\lambda_1+\lambda_2}.
}
\]
}

\begin{exercise}
\label{ex:counting-processes-42}
Let $\{N(t),t\ge 0\}$ be a Poisson process with rate $\lambda$. Let $S_n$ denote the time of the $n$th event. Find
\begin{enumerate}
\item $E[S_4]$,
\item $E[S_4\mid N(1)=2]$,
\item $E[N(4)-N(2)\mid N(1)=3]$.
\end{enumerate}
\end{exercise}

\solution{
\textbf{(a)} The time $S_4$ is the sum of $4$ independent exponential random variables with rate $\lambda$. Therefore,
\[
{
E[S_4]=\frac4\lambda.
}
\]

\textbf{(b)} Given that
\[
N(1)=2,
\]
there have been exactly two events by time $1$. The fourth event will be the second event after time $1$. By independent increments, the process after time $1$ starts fresh. Thus the expected additional time until two more events is
\[
\frac2\lambda.
\]
Therefore,
\[
{
E[S_4\mid N(1)=2]
=
1+\frac2\lambda.
}
\]

\textbf{(c)} The increment
\[
N(4)-N(2)
\]
is independent of $N(1)$ because it concerns the disjoint interval $(2,4]$. Also,
\[
N(4)-N(2)\sim \operatorname{Poisson}(2\lambda).
\]
Hence
\[
{
E[N(4)-N(2)\mid N(1)=3]=2\lambda.
}
\]
}

\begin{exercise}
\label{ex:counting-processes-43}
Customers arrive at a two-server service station according to a Poisson process with rate $\lambda$. Whenever a new customer arrives, any customer that is in the system immediately departs. A new arrival enters service first with server $1$ and then with server $2$. If the service times at the servers are independent exponentials with respective rates $\mu_1$ and $\mu_2$, what proportion of entering customers completes their service with server $2$?
\end{exercise}

\solution{
A customer completes service with server $2$ if both service times are completed before the next arrival occurs.

Let
\[
A\sim \operatorname{Exp}(\lambda)
\]
be the time until the next arrival. Let
\[
S_1\sim \operatorname{Exp}(\mu_1),
\qquad
S_2\sim \operatorname{Exp}(\mu_2)
\]
be the service times at servers $1$ and $2$. The customer completes service if
\[
S_1+S_2<A.
\]
Thus
\[
P(S_1+S_2<A)
=
E[e^{-\lambda(S_1+S_2)}].
\]
By independence,
\[
E[e^{-\lambda(S_1+S_2)}]
=
E[e^{-\lambda S_1}]E[e^{-\lambda S_2}].
\]
Since
\[
E[e^{-\lambda S_i}]
=
\frac{\mu_i}{\mu_i+\lambda},
\]
we get
\[
{
P(\text{customer completes service})
=
\frac{\mu_1}{\mu_1+\lambda}
\cdot
\frac{\mu_2}{\mu_2+\lambda}.
}
\]
}

\begin{exercise}
\label{ex:counting-processes-44}
Cars pass a certain street location according to a Poisson process with rate $\lambda$. A woman who wants to cross the street at that location waits until she can see that no cars will come by in the next $T$ time units.
\begin{enumerate}
\item Find the probability that her waiting time is $0$.
\item Find her expected waiting time.
\end{enumerate}
\end{exercise}

\solution{
Let $X$ be the time until the next car arrives. Since cars arrive according to a Poisson process with rate $\lambda$,
\[
X\sim \operatorname{Exp}(\lambda).
\]

\textbf{(a)} Her waiting time is $0$ if the next car will not arrive within the next $T$ time units, that is, if
\[
X>T.
\]
Thus
\[
{
P(\text{waiting time is }0)=P(X>T)=e^{-\lambda T}.
}
\]

\textbf{(b)} Let $W$ be her waiting time. Condition on the time $X$ of the first car.

If $X>T$, then
\[
W=0.
\]
If $X\le T$, then she waits until that car passes, and by the memoryless property of the Poisson process, the problem starts over. Hence
\[
W=X+W',
\]
where $W'$ has the same distribution as $W$ and is independent of $X$.

Taking expectations,
\[
E[W]
=
E[X\mathbf{1}_{\{X\le T\}}]
+
P(X\le T)E[W].
\]
Thus
\[
E[W]P(X>T)
=
E[X\mathbf{1}_{\{X\le T\}}].
\]
Now
\[
P(X>T)=e^{-\lambda T},
\]
and
\[
E[X\mathbf{1}_{\{X\le T\}}]
=
\int_0^T x\lambda e^{-\lambda x}\,dx.
\]
Evaluating the integral,
\[
\int_0^T x\lambda e^{-\lambda x}\,dx
=
\frac1\lambda(1-e^{-\lambda T})-Te^{-\lambda T}.
\]
Therefore,
\[
E[W]
=
\frac{\frac1\lambda(1-e^{-\lambda T})-Te^{-\lambda T}}
{e^{-\lambda T}}.
\]
Hence
\[
{
E[W]
=
\frac{e^{\lambda T}-1}{\lambda}-T.
}
\]
}

\setcounter{exercise}{46}

\begin{exercise}
\label{ex:counting-processes-47}
Consider a two-server parallel queueing system where customers arrive according to a Poisson process with rate $\lambda$, and where the service times are exponential with rate $\mu$. Moreover, suppose that arrivals finding both servers busy immediately depart without receiving any service, whereas those finding at least one free server immediately enter service and then depart when their service is completed.
\begin{enumerate}
\item If both servers are presently busy, find the expected time until the next customer enters the system.
\item Starting empty, find the expected time until both servers are busy.
\item Find the expected time between two successive lost customers.
\end{enumerate}
\end{exercise}

\solution{
\textbf{(a)}
If both servers are busy, the time until the first service completion is exponential with rate $2\mu$, and so has mean
\[
\frac{1}{2\mu}.
\]
After one server becomes free, the next arrival enters the system. The time until the next arrival has mean
\[
\frac1\lambda.
\]
Therefore, the expected time until the next customer enters is
\[
{
\frac{1}{2\mu}+\frac1\lambda.
}
\]

\textbf{(b)}
Let $E_i$ be the expected time until both servers are busy, starting with $i$ busy servers. Then
\[
E_2=0.
\]
Starting empty,
\[
E_0=\frac1\lambda+E_1.
\]
Starting with one busy server, the next event is either an arrival, with rate $\lambda$, or a service completion, with rate $\mu$. Hence
\[
E_1
=
\frac1{\lambda+\mu}
+
\frac{\mu}{\lambda+\mu}E_0.
\]
Substituting into $E_0=\frac1\lambda+E_1$ gives
\[
E_0
=
\frac1\lambda
+
\frac1{\lambda+\mu}
+
\frac{\mu}{\lambda+\mu}E_0.
\]
Thus
\[
\frac{\lambda}{\lambda+\mu}E_0
=
\frac1\lambda+\frac1{\lambda+\mu}.
\]
Therefore,
\[
E_0
=
\frac{2\lambda+\mu}{\lambda^2}.
\]
So the expected time until both servers are busy, starting empty, is
\[
{
\frac{2\lambda+\mu}{\lambda^2}.
}
\]

\textbf{(c)}
Let $L_i$ be the expected time until the next lost customer, starting with $i$ busy servers. A lost customer occurs only when an arrival finds both servers busy.

Starting with two busy servers,
\[
L_2
=
\frac1{\lambda+2\mu}
+
\frac{2\mu}{\lambda+2\mu}L_1.
\]
Starting with one busy server,
\[
L_1
=
\frac1{\lambda+\mu}
+
\frac{\lambda}{\lambda+\mu}L_2
+
\frac{\mu}{\lambda+\mu}L_0.
\]
Starting empty,
\[
L_0=\frac1\lambda+L_1.
\]
Using this in the equation for $L_1$,
\[
L_1=L_2+\frac1\lambda+\frac{\mu}{\lambda^2}.
\]
Substituting into the equation for $L_2$ gives
\[
L_2
=
\frac1{\lambda+2\mu}
+
\frac{2\mu}{\lambda+2\mu}
\left(
L_2+\frac1\lambda+\frac{\mu}{\lambda^2}
\right).
\]
Solving,
\[
L_2
=
\frac1\lambda+\frac{2\mu}{\lambda^2}+\frac{2\mu^2}{\lambda^3}.
\]
Thus the expected time between two successive lost customers is
\[
{
\frac1\lambda+\frac{2\mu}{\lambda^2}+\frac{2\mu^2}{\lambda^3}.
}
\]
}

\begin{exercise}
\label{ex:counting-processes-48}
Consider an $n$-server parallel queueing system where customers arrive according to a Poisson process with rate $\lambda$, where the service times are exponential random variables with rate $\mu$, and where any arrival finding all servers busy immediately departs without receiving any service. If an arrival finds all servers busy, find
\begin{enumerate}
\item the expected number of busy servers found by the next arrival,
\item the probability that the next arrival finds all servers free,
\item the probability that the next arrival finds exactly $i$ of the servers free.
\end{enumerate}
\end{exercise}

\solution{
Suppose that immediately after the lost arrival, all $n$ servers are busy. Let $A$ be the time until the next arrival. Then
\[
A\sim \operatorname{Exp}(\lambda).
\]
Each busy server independently remains busy until time $A$ with probability
\[
P(\operatorname{Exp}(\mu)>A\mid A)=e^{-\mu A}.
\]

\textbf{(a)}
Given $A=t$, the number of busy servers at the next arrival is binomial with parameters $n$ and $e^{-\mu t}$. Thus
\[
E[\text{busy servers}\mid A=t]=ne^{-\mu t}.
\]
Therefore,
\[
E[\text{busy servers}]
=
nE[e^{-\mu A}]
=
n\frac{\lambda}{\lambda+\mu}.
\]
Hence
\[
{
E[\text{busy servers found by next arrival}]
=
\frac{n\lambda}{\lambda+\mu}.
}
\]

\textbf{(b)}
The next arrival finds all servers free if all $n$ current services finish before the next arrival. Thus
\[
P(\text{all free})
=
E[(1-e^{-\mu A})^n].
\]
Therefore,
\[
P(\text{all free})
=
\int_0^\infty
\lambda e^{-\lambda t}(1-e^{-\mu t})^n\,dt.
\]
Equivalently,
\[
{
P(\text{all free})
=
\frac{\lambda}{\mu}
B\left(\frac{\lambda}{\mu},n+1\right),
}
\]
where $B(\cdot,\cdot)$ is the beta function.

\textbf{(c)}
The next arrival finds exactly $i$ free servers if exactly $i$ of the $n$ original services have completed by the arrival time. Thus
\[
P(\text{exactly }i\text{ free})
=
\binom{n}{i}
\int_0^\infty
\lambda e^{-\lambda t}
(1-e^{-\mu t})^i
e^{-\mu t(n-i)}
\,dt.
\]
Hence
\[
{
P(\text{exactly }i\text{ free})
=
\binom{n}{i}
\frac{\lambda}{\mu}
B\left(\frac{\lambda}{\mu}+n-i,i+1\right),
}
\]
for
\[
i=0,1,\ldots,n.
\]
}

\begin{exercise}
\label{ex:counting-processes-49}
Events occur according to a Poisson process with rate $\lambda$. Each time an event occurs, we must decide whether or not to stop, with our objective being to stop at the last event to occur prior to some specified time $T$, where $T>1/\lambda$. Consider the strategy that stops at the first event to occur after some fixed time $s$, $0\le s\le T$.
\begin{enumerate}
\item Using this strategy, what is the probability of winning?
\item What value of $s$ maximizes the probability of winning?
\item Show that one's probability of winning when using the preceding strategy with the value of $s$ specified in part (b) is $1/e$.
\end{enumerate}
\end{exercise}

\solution{
\textbf{(a)}
We stop at the first event after time $s$. We win if this event occurs before time $T$ and no further events occur before time $T$.

Let
\[
x=T-s.
\]
The time from $s$ until the first event after $s$ is exponential with rate $\lambda$. Thus the probability of winning is
\[
\int_0^x \lambda e^{-\lambda u}e^{-\lambda(x-u)}\,du.
\]
Therefore,
\[
P(\text{win})
=
\int_0^x \lambda e^{-\lambda x}\,du
=
\lambda x e^{-\lambda x}.
\]
Since $x=T-s$,
\[
{
P(\text{win})=\lambda(T-s)e^{-\lambda(T-s)}.
}
\]

\textbf{(b)}
We maximize
\[
\lambda x e^{-\lambda x},
\qquad 0\le x\le T.
\]
The function $x e^{-\lambda x}$ is maximized at
\[
x=\frac1\lambda.
\]
Since $T>1/\lambda$, this value is feasible. Therefore,
\[
T-s=\frac1\lambda,
\]
so
\[
{
s=T-\frac1\lambda.
}
\]

\textbf{(c)}
Substituting $T-s=\frac1\lambda$ into the winning probability gives
\[
P(\text{win})
=
\lambda\left(\frac1\lambda\right)e^{-\lambda(1/\lambda)}
=
e^{-1}.
\]
Hence
\[
{
P(\text{win})=\frac1e.
}
\]
}

\begin{exercise}
\label{ex:counting-processes-50}
The number of hours between successive train arrivals at the station is uniformly distributed on $(0,1)$. Passengers arrive according to a Poisson process with rate $7$ per hour. Suppose a train has just left the station. Let $X$ denote the number of people who get on the next train. Find
\begin{enumerate}
\item $E[X]$,
\item $\operatorname{Var}(X)$.
\end{enumerate}
\end{exercise}

\solution{
Let $U$ be the time until the next train arrives. Then
\[
U\sim \operatorname{Uniform}(0,1).
\]
Given $U=u$, the number of passengers arriving before the next train is
\[
X\mid U=u\sim \operatorname{Poisson}(7u).
\]

\textbf{(a)}
Using conditional expectation,
\[
E[X]=E[E[X\mid U]]=E[7U]=7E[U].
\]
Since $E[U]=\frac12$,
\[
{
E[X]=\frac72.
}
\]

\textbf{(b)}
Using conditional variance,
\[
\operatorname{Var}(X)
=
E[\operatorname{Var}(X\mid U)]
+
\operatorname{Var}(E[X\mid U]).
\]
Since
\[
\operatorname{Var}(X\mid U)=7U
\]
and
\[
E[X\mid U]=7U,
\]
we get
\[
\operatorname{Var}(X)
=
E[7U]+\operatorname{Var}(7U).
\]
Thus
\[
\operatorname{Var}(X)
=
7\cdot \frac12+49\cdot \frac1{12}.
\]
Therefore,
\[
{
\operatorname{Var}(X)=\frac{91}{12}.
}
\]
}

\begin{exercise}
\label{ex:counting-processes-51}
If an individual has never had a previous automobile accident, then the probability he or she has an accident in the next $h$ time units is $\beta h+o(h)$; on the other hand, if he or she has ever had a previous accident, then the probability is $\alpha h+o(h)$. Find the expected number of accidents an individual has by time $t$.
\end{exercise}

\solution{
Assume no previous accidents at time $0$ and $\beta>0$. If $\beta=0$, the individual never has a first accident and the expectation is $0$. Let $S$ be the time of the first accident. Before the first accident, accidents occur at rate $\beta$, so
\[
S\sim \operatorname{Exp}(\beta).
\]
After the first accident, accidents occur at rate $\alpha$.

Let $N(t)$ be the number of accidents by time $t$. Then
\[
E[N(t)]
=
P(S\le t)+\alpha E[(t-S)^+].
\]
Now
\[
P(S\le t)=1-e^{-\beta t}.
\]
Also,
\[
E[(t-S)^+]
=
\int_0^t P(S\le u)\,du
=
\int_0^t (1-e^{-\beta u})\,du.
\]
Therefore,
\[
E[(t-S)^+]
=
t-\frac{1-e^{-\beta t}}{\beta}.
\]
Hence
\[
{
E[N(t)]
=
1-e^{-\beta t}
+
\alpha\left(
t-\frac{1-e^{-\beta t}}{\beta}
\right).
}
\]
}

\begin{exercise}
\label{ex:counting-processes-52}
Teams $1$ and $2$ are playing a match. The teams score points according to independent Poisson processes with respective rates $\lambda_1$ and $\lambda_2$. If the match ends when one of the teams has scored $k$ more points than the other, find the probability that team $1$ wins.
\end{exercise}

\solution{
Combine the two Poisson processes. Points occur according to a Poisson process with rate
\[
\lambda_1+\lambda_2.
\]
Each point is scored by team $1$ with probability
\[
p=\frac{\lambda_1}{\lambda_1+\lambda_2},
\]
and by team $2$ with probability
\[
q=\frac{\lambda_2}{\lambda_1+\lambda_2}.
\]

Let $D$ be the difference
\[
D=(\text{team 1 score})-(\text{team 2 score}).
\]
Then $D$ is a random walk that increases by $1$ with probability $p$ and decreases by $1$ with probability $q$. The match ends when $D$ reaches $k$ or $-k$.

This is the gambler's ruin problem, starting at $0$ with absorbing barriers at $k$ and $-k$. Therefore,
\[
P(\text{team 1 wins})
=
\frac{1-(q/p)^k}{1-(q/p)^{2k}},
\qquad p\ne q.
\]
Equivalently,
\[
{
P(\text{team 1 wins})
=
\frac{\lambda_1^k}{\lambda_1^k+\lambda_2^k}.
}
\]
If $\lambda_1=\lambda_2$, this gives
\[
P(\text{team 1 wins})=\frac12.
\]
}

\begin{exercise}
\label{ex:counting-processes-53}
The water level of a certain reservoir is depleted at a constant rate of $1000$ units daily. The reservoir is refilled by randomly occurring rainfalls. Rainfalls occur according to a Poisson process with rate $0.2$ per day. The amount of water added to the reservoir by a rainfall is $5000$ units with probability $0.8$ or $8000$ units with probability $0.2$. The present water level is just slightly below $5000$ units.
\begin{enumerate}
\item What is the probability the reservoir will be empty after five days?
\item What is the probability the reservoir will be empty sometime within the next ten days?
\end{enumerate}
\end{exercise}

\solution{
Write the initial level as $1000a$, with $a<5$ and close to $5$.
Take zero as the lowest possible water level, and assume no capacity
limit affects these short paths.
\textbf{(a)} A rainfall before day $5$ leaves a positive level at day
$5$; without rainfall the initial water runs out at day $a<5$. Thus
\[
P(\text{empty at day }5)=P(N(5)=0)=e^{-1}.
\]
\textbf{(b)} Without rain the reservoir first empties at day $a$.
If a $5000$-unit rainfall arrives before then, the emptying time moves
to day $a+5<10$. A single $8000$-unit rainfall, or two rainfalls before
emptying, postpones emptying beyond day $10$ for $a$ close to $5$.
The disjoint cases of emptying within ten days are therefore:
\begin{itemize}
\item no rain before day $a$;
\item one $5000$-unit rainfall before day $a$ and no other rain before
day $a+5$.
\end{itemize}
By independent Poisson increments and independent rainfall amounts,
\[
P(\text{empty by day }10)=e^{-0.2a}+0.8(0.2a)e^{-0.2(a+5)}.
\]
Interpreting ``just slightly below $5000$'' as the limit $a\uparrow5$
gives $e^{-1}+0.8e^{-2}$. For a specified positive shortfall, the
preceding formula gives the exact value.
}

\begin{exercise}
\label{ex:counting-processes-54}
A viral linear DNA molecule of length $1$ is often known to contain a certain marked position, with the exact location of this mark being unknown. One approach to locating the marked position is to cut the molecule by agents that break it at points chosen according to a Poisson process with rate $\lambda$. Letting $m$ denote the location on the line of the marked position, if $L_1$ denotes the last Poisson event time before $m$ and $R_1$ denotes the first Poisson event time after $m$, find
\begin{enumerate}
\item $P\{L_1=0\}$,
\item $P\{L_1<x\}$, $0<x<m$,
\item $P\{R_1=1\}$,
\item $P\{R_1>x\}$, $m<x<1$.
\end{enumerate}
If the cutting procedure is utilized on $n$ identical copies of the molecule, yielding the data $L_i,R_i$, $i=1,\ldots,n$, then the marked position lies between $L$ and $R$, where
\[
L=\max_i L_i,\qquad R=\min_i R_i.
\]
Find $E[R-L]$, and show that
\[
E[R-L]\sim \frac{2}{n\lambda}.
\]
\end{exercise}

\solution{
\textbf{(a)}
Assume $0<m<1$, and use the molecule endpoints when there is no cut on a side. The event $\{L_1=0\}$ means that there are no Poisson events in $(0,m)$. Thus
\[
{
P(L_1=0)=e^{-\lambda m}.
}
\]

\textbf{(b)}
For $0<x<m$, the event $\{L_1<x\}$ means that there are no Poisson events in $(x,m)$. Hence
\[
{
P(L_1<x)=e^{-\lambda(m-x)}.
}
\]

\textbf{(c)}
The event $\{R_1=1\}$ means that there are no Poisson events in $(m,1)$. Thus
\[
{
P(R_1=1)=e^{-\lambda(1-m)}.
}
\]

\textbf{(d)}
For $m<x<1$, the event $\{R_1>x\}$ means that there are no Poisson events in $(m,x)$. Therefore,
\[
{
P(R_1>x)=e^{-\lambda(x-m)}.
}
\]

Now repeat the experiment independently on $n$ identical copies.

Let
\[
L=\max_i L_i.
\]
Then
\[
m-L
\]
is the smallest of the $n$ independent left distances from $m$ to the nearest cut on the left, each truncated at $m$. For $0<y<m$,
\[
P(m-L>y)=e^{-n\lambda y}.
\]
Therefore,
\[
E[m-L]
=
\int_0^m e^{-n\lambda y}\,dy
=
\frac{1-e^{-n\lambda m}}{n\lambda}.
\]

Similarly,
\[
R-m
\]
is the smallest of the $n$ independent right distances from $m$ to the nearest cut on the right, each truncated at $1-m$. Thus
\[
E[R-m]
=
\int_0^{1-m} e^{-n\lambda y}\,dy
=
\frac{1-e^{-n\lambda(1-m)}}{n\lambda}.
\]

Hence
\[
E[R-L]
=
E[R-m]+E[m-L].
\]
Therefore,
\[
{
E[R-L]
=
\frac{2-e^{-n\lambda m}-e^{-n\lambda(1-m)}}{n\lambda}.
}
\]
As $n\to\infty$,
\[
e^{-n\lambda m}\to 0,
\qquad
e^{-n\lambda(1-m)}\to 0,
\]
so
\[
{
E[R-L]\sim \frac{2}{n\lambda}.
}
\]
}

\begin{exercise}
\label{ex:counting-processes-55}
Consider a single server queueing system where customers arrive according to a Poisson process with rate $\lambda$, service times are exponential with rate $\mu$, and customers are served in the order of their arrival. Suppose that a customer arrives and finds $n-1$ others in the system. Let $X$ denote the number in the system at the moment that customer departs. Find the probability mass function of $X$.
\end{exercise}

\solution{
The arriving customer has $n-1$ customers ahead of him, so his departure occurs after $n$ service completions: the $n-1$ customers ahead of him plus his own service.

During this time, new arrivals may occur. By the memoryless property, each event in the combined arrival-service process is an arrival with probability
\[
\frac{\lambda}{\lambda+\mu},
\]
and a service completion with probability
\[
\frac{\mu}{\lambda+\mu}.
\]

The number $X$ of customers left in the system when the tagged customer departs equals the number of arrivals before the $n$th service completion. Hence $X$ has a negative binomial distribution.

Thus, for $k=0,1,2,\ldots$,
\[
{
P(X=k)
=
\binom{n+k-1}{k}
\left(\frac{\mu}{\lambda+\mu}\right)^n
\left(\frac{\lambda}{\lambda+\mu}\right)^k.
}
\]
}

\begin{exercise}
\label{ex:counting-processes-56}
An event independently occurs on each day with probability $p$. Let $N(n)$ denote the total number of events that occur on the first $n$ days, and let $T_r$ denote the day on which the $r$th event occurs.
\begin{enumerate}
\item What is the distribution of $N(n)$?
\item What is the distribution of $T_1$?
\item What is the distribution of $T_r$?
\item Given that $N(n)=r$, show that the set of $r$ days on which events occurred has the same distribution as a random selection without replacement of $r$ of the values $1,2,\ldots,n$.
\end{enumerate}
\end{exercise}

\solution{
\textbf{(a)}
Since each day independently has an event with probability $p$,
\[
{
N(n)\sim \operatorname{Binomial}(n,p).
}
\]
Thus
\[
P(N(n)=k)=\binom{n}{k}p^k(1-p)^{n-k}.
\]

\textbf{(b)}
The random variable $T_1$ is the day of the first event. Therefore,
\[
{
P(T_1=k)=(1-p)^{k-1}p,
\qquad k=1,2,\ldots.
}
\]
Thus $T_1$ is geometric with parameter $p$.

\textbf{(c)}
The event $\{T_r=k\}$ means that there are exactly $r-1$ events in the first $k-1$ days, and an event occurs on day $k$. Hence
\[
{
P(T_r=k)
=
\binom{k-1}{r-1}
p^r(1-p)^{k-r},
\qquad k=r,r+1,\ldots.
}
\]
Thus $T_r$ has a negative binomial distribution.

\textbf{(d)}
Given that $N(n)=r$, every specific set of $r$ event days has probability proportional to
\[
p^r(1-p)^{n-r}.
\]
This value is the same for every subset of $r$ days. Since there are
\[
\binom{n}{r}
\]
such subsets, each subset has conditional probability
\[
{
\frac1{\binom{n}{r}}.
}
\]
Therefore, given $N(n)=r$, the set of event days is distributed as a random selection without replacement of $r$ days from
\[
\{1,2,\ldots,n\}.
\]
}

\begin{exercise}
\label{ex:counting-processes-57}
Events occur according to a Poisson process with rate $\lambda=2$ per hour.
\begin{enumerate}
\item What is the probability that no event occurs between 8 P.M. and 9 P.M.?
\item Starting at noon, what is the expected time at which the fourth event occurs?
\item What is the probability that two or more events occur between 6 P.M. and 8 P.M.?
\end{enumerate}
\end{exercise}

\solution{
\textbf{(a)}
The interval from 8 P.M. to 9 P.M. has length $1$ hour. Therefore,
\[
N\sim \operatorname{Poisson}(2).
\]
Thus
\[
{
P(N=0)=e^{-2}.
}
\]

\textbf{(b)}
The time until the fourth event has mean
\[
\frac4\lambda=\frac42=2
\]
hours. Starting at noon, the expected time of the fourth event is therefore
\[
{\text{2 P.M.}}
\]

\textbf{(c)}
The interval from 6 P.M. to 8 P.M. has length $2$ hours, so the number of events has distribution
\[
N\sim \operatorname{Poisson}(4).
\]
Therefore,
\[
P(N\ge 2)=1-P(N=0)-P(N=1).
\]
Thus
\[
P(N\ge 2)
=
1-e^{-4}-4e^{-4}.
\]
Hence
\[
{
P(N\ge 2)=1-5e^{-4}.
}
\]
}

\begin{exercise}
\label{ex:counting-processes-58}
Consider the coupon collecting problem where there are $m$ distinct types of coupons, and each new coupon collected is type $j$ with probability $p_j$, where
\[
\sum_{j=1}^m p_j=1.
\]
Suppose you stop collecting when you have a complete set of at least one of each type. Show that
\[
P\{i\text{ is the last type collected}\}
=
E\left[
\prod_{j\ne i}
\left(1-U^{p_j/p_i}\right)
\right],
\]
where $U$ is a uniform random variable on $(0,1)$.
\end{exercise}

\solution{
Assume $p_j>0$ for every type, so collecting a complete set is possible. Use a Poissonization argument. Suppose coupons arrive according to a Poisson process of rate $1$, and each coupon is independently of type $j$ with probability $p_j$. Then the arrivals of type $j$ form an independent Poisson process with rate $p_j$.

Let $T_j$ be the time until the first coupon of type $j$ appears. Then
\[
T_j\sim \operatorname{Exp}(p_j),
\]
and the random variables
\[
T_1,\ldots,T_m
\]
are independent.

The event that type $i$ is the last type collected is
\[
\{T_j<T_i \text{ for all } j\ne i\}.
\]
Conditioning on $T_i=t$,
\[
P(T_j<t \text{ for all }j\ne i\mid T_i=t)
=
\prod_{j\ne i}(1-e^{-p_jt}).
\]
Hence
\[
P(i\text{ is last})
=
E\left[
\prod_{j\ne i}(1-e^{-p_jT_i})
\right].
\]

Now define
\[
U=e^{-p_iT_i}.
\]
Since $T_i\sim \operatorname{Exp}(p_i)$, the random variable $U$ is uniform on $(0,1)$. Also,
\[
e^{-p_jT_i}
=
\left(e^{-p_iT_i}\right)^{p_j/p_i}
=
U^{p_j/p_i}.
\]
Therefore,
\[
{
P\{i\text{ is the last type collected}\}
=
E\left[
\prod_{j\ne i}
\left(1-U^{p_j/p_i}\right)
\right].
}
\]
}

\begin{exercise}
\label{ex:counting-processes-59}
There are two types of claims that are made to an insurance company. Let $N_i(t)$ denote the number of type $i$ claims made by time $t$, and suppose that $\{N_1(t),t\ge 0\}$ and $\{N_2(t),t\ge 0\}$ are independent Poisson processes with rates $\lambda_1=10$ and $\lambda_2=1$. The amounts of successive type $1$ claims are independent exponential random variables with mean $\$1000$, whereas the amounts from type $2$ claims are independent exponential random variables with mean $\$5000$. A claim for $\$4000$ has just been received; what is the probability it is a type $1$ claim?
\end{exercise}

\solution{
Let $f_1$ and $f_2$ be the claim amount densities for type $1$ and type $2$ claims. Since type $1$ claims have mean $1000$,
\[
f_1(x)=\frac1{1000}e^{-x/1000}.
\]
Since type $2$ claims have mean $5000$,
\[
f_2(x)=\frac1{5000}e^{-x/5000}.
\]

Using Bayes' rule with arrival rates,
\[
P(\text{type }1\mid X=x)
=
\frac{\lambda_1 f_1(x)}
{\lambda_1 f_1(x)+\lambda_2 f_2(x)}.
\]
For $x=4000$,
\[
P(\text{type }1\mid X=4000)
=
\frac{10\cdot \frac1{1000}e^{-4}}
{10\cdot \frac1{1000}e^{-4}
+
1\cdot \frac1{5000}e^{-0.8}}.
\]
Thus
\[
P(\text{type }1\mid X=4000)
=
\frac{50e^{-4}}{50e^{-4}+e^{-0.8}}.
\]
Equivalently,
\[
{
P(\text{type }1\mid X=4000)
=
\frac{50}{50+e^{3.2}}.
}
\]
}

\begin{exercise}
\label{ex:counting-processes-60}
Customers arrive at a bank at a Poisson rate $\lambda$. Suppose two customers arrived during the first hour. What is the probability that
\begin{enumerate}
\item both arrived during the first $20$ minutes?
\item at least one arrived during the first $20$ minutes?
\end{enumerate}
\end{exercise}

\solution{
Given that two customers arrived during the first hour, their arrival times are distributed as two independent uniform random variables on the interval $(0,1)$, unordered.

The first $20$ minutes correspond to
\[
\frac13
\]
of an hour.

\textbf{(a)}
The probability that both arrivals occurred during the first $20$ minutes is
\[
\left(\frac13\right)^2.
\]
Thus
\[
{
P(\text{both in first }20\text{ minutes})=\frac19.
}
\]

\textbf{(b)}
The probability that at least one arrived during the first $20$ minutes is
\[
1-P(\text{neither arrived during the first }20\text{ minutes}).
\]
The probability that a single arrival was not in the first $20$ minutes is
\[
\frac23.
\]
Therefore,
\[
P(\text{at least one in first }20\text{ minutes})
=
1-\left(\frac23\right)^2
=
1-\frac49
=
\frac59.
\]
Hence
\[
{
P(\text{at least one in first }20\text{ minutes})=\frac59.
}
\]
}

\setcounter{exercise}{60}

\begin{exercise}
\label{ex:counting-processes-61}
A system has a random number of flaws that we will suppose is Poisson distributed with mean $c$. Each of these flaws will, independently, cause the system to fail at a random time having distribution $G$. When a system failure occurs, suppose that the flaw causing the failure is immediately located and fixed.
\begin{enumerate}
\item What is the distribution of the number of failures by time $t$?
\item What is the distribution of the number of flaws that remain in the system at time $t$?
\item Are the random variables in parts (a) and (b) dependent or independent?
\end{enumerate}
\end{exercise}

\solution{
Assume flaw failure times are independent of the initial flaw count. Let $N$ denote the initial number of flaws. Then
\[
N\sim \operatorname{Poisson}(c).
\]
For each flaw, let its failure time have distribution function $G$. By time $t$, a given flaw has caused a failure with probability
\[
G(t),
\]
and has not yet caused a failure with probability
\[
1-G(t).
\]

\textbf{(a)} By Poisson thinning, the number of flaws that have caused failures by time $t$ is Poisson with mean
\[
cG(t).
\]
Thus, if $N_1(t)$ denotes the number of failures by time $t$, then
\[
{
N_1(t)\sim \operatorname{Poisson}(cG(t)).
}
\]

\textbf{(b)} Similarly, the number of flaws that remain in the system at time $t$ is Poisson with mean
\[
c[1-G(t)].
\]
Thus, if $N_2(t)$ denotes the number of remaining flaws at time $t$, then
\[
{
N_2(t)\sim \operatorname{Poisson}(c[1-G(t)]).
}
\]

\textbf{(c)} Since the flaws are independently classified into two categories, namely those that have failed by time $t$ and those that have not, the Poisson thinning property implies that
\[
N_1(t)
\quad\text{and}\quad
N_2(t)
\]
are independent random variables.

Therefore,
\[
{\text{The two random variables are independent.}}
\]
}

\begin{exercise}
\label{ex:counting-processes-62}
Suppose that the number of typographical errors in a new text is Poisson distributed with mean $\lambda$. Two proofreaders independently read the text. Suppose that each error is independently found by proofreader $i$ with probability $p_i$, $i=1,2$. Let $X_1$ denote the number of errors that are found by proofreader $1$ but not by proofreader $2$. Let $X_2$ denote the number of errors that are found by proofreader $2$ but not by proofreader $1$. Let $X_3$ denote the number of errors that are found by both proofreaders. Finally, let $X_4$ denote the number of errors found by neither proofreader.
\begin{enumerate}
\item Describe the joint probability distribution of $X_1,X_2,X_3,X_4$.
\item Show that
\[
\frac{E[X_1]}{E[X_3]}=\frac{1-p_2}{p_2}
\qquad\text{and}\qquad
\frac{E[X_2]}{E[X_3]}=\frac{1-p_1}{p_1}.
\]
\item Suppose now that $\lambda,p_1,$ and $p_2$ are all unknown. By using $X_i$ as an estimator of $E[X_i]$, $i=1,2,3$, present estimators of $p_1,p_2,$ and $\lambda$.
\item Give an estimator of $X_4$, the number of errors not found by either proofreader.
\end{enumerate}
\end{exercise}

\solution{
Each error falls into exactly one of four categories:
\[
\begin{array}{c|c}
\text{Category} & \text{Probability} \\ \hline
\text{found by proofreader 1 but not 2} & p_1(1-p_2)\\
\text{found by proofreader 2 but not 1} & p_2(1-p_1)\\
\text{found by both proofreaders} & p_1p_2\\
\text{found by neither proofreader} & (1-p_1)(1-p_2).
\end{array}
\]

\textbf{(a)}
Since the total number of errors is Poisson with mean $\lambda$, the Poisson splitting property implies that
\[
X_1,X_2,X_3,X_4
\]
are independent Poisson random variables with means
\[
E[X_1]=\lambda p_1(1-p_2),
\]
\[
E[X_2]=\lambda p_2(1-p_1),
\]
\[
E[X_3]=\lambda p_1p_2,
\]
and
\[
E[X_4]=\lambda(1-p_1)(1-p_2).
\]
Thus the joint probability mass function is
\[
P(X_1=x_1,X_2=x_2,X_3=x_3,X_4=x_4)
=
\prod_{i=1}^4 e^{-m_i}\frac{m_i^{x_i}}{x_i!},
\]
where
\[
m_1=\lambda p_1(1-p_2),
\quad
m_2=\lambda p_2(1-p_1),
\quad
m_3=\lambda p_1p_2,
\quad
m_4=\lambda(1-p_1)(1-p_2).
\]

\textbf{(b)}
We have
\[
\frac{E[X_1]}{E[X_3]}
=
\frac{\lambda p_1(1-p_2)}{\lambda p_1p_2}
=
\frac{1-p_2}{p_2}.
\]
Similarly,
\[
\frac{E[X_2]}{E[X_3]}
=
\frac{\lambda p_2(1-p_1)}{\lambda p_1p_2}
=
\frac{1-p_1}{p_1}.
\]

\textbf{(c)}
Using $X_i$ as an estimator of $E[X_i]$, we estimate
\[
\frac{1-p_2}{p_2}
\]
by
\[
\frac{X_1}{X_3}.
\]
Solving gives
\[
{
\widehat p_2=\frac{X_3}{X_1+X_3}.
}
\]
Similarly,
\[
{
\widehat p_1=\frac{X_3}{X_2+X_3}.
}
\]

Since
\[
E[X_3]=\lambda p_1p_2,
\]
we estimate $\lambda$ by
\[
\widehat\lambda
=
\frac{X_3}{\widehat p_1\widehat p_2}.
\]
Substituting the estimators gives
\[
{
\widehat\lambda
=
\frac{(X_1+X_3)(X_2+X_3)}{X_3}.
}
\]

\textbf{(d)}
Since
\[
E[X_4]=\lambda(1-p_1)(1-p_2),
\]
and also
\[
E[X_4]
=
\frac{E[X_1]E[X_2]}{E[X_3]},
\]
a natural estimator of $X_4$ is
\[
{
\widehat X_4=\frac{X_1X_2}{X_3}.
}
\]
This estimator is meaningful when $X_3>0$.
}

\begin{exercise}
\label{ex:counting-processes-63}
Consider an infinite server queueing system in which customers arrive in accordance with a Poisson process with rate $\lambda$, and where the service distribution is exponential with rate $\mu$. Let $X(t)$ denote the number of customers in the system at time $t$. Find
\begin{enumerate}
\item $E[X(t+s)\mid X(s)=n]$,
\item $\operatorname{Var}[X(t+s)\mid X(s)=n]$.
\end{enumerate}
Also, consider an infinite server queueing system in which customers arrive according to a Poisson process with rate $\lambda$, and where the service times are all exponential random variables with rate $\mu$. If there is currently a single customer in the system, find the probability that the system becomes empty when that customer departs.
\end{exercise}

\solution{
Suppose
\[
X(s)=n.
\]
The customers present at time $s$ are the ``old'' customers. Each old customer independently remains in the system until time $s+t$ with probability
\[
p=e^{-\mu t}.
\]
Thus the number of old customers remaining at time $s+t$ is
\[
Y\sim \operatorname{Binomial}(n,e^{-\mu t}).
\]

Now consider the new arrivals during the interval $(s,s+t]$. An arrival at time $s+u$, where $0<u<t$, remains in the system until time $s+t$ with probability
\[
e^{-\mu(t-u)}.
\]
By thinning a Poisson process, the number of new customers present at time $s+t$ is Poisson with mean
\[
m
=
\int_0^t \lambda e^{-\mu(t-u)}\,du
=
\frac{\lambda}{\mu}(1-e^{-\mu t}).
\]
Call this number $Z$. Then
\[
Z\sim \operatorname{Poisson}\left(\frac{\lambda}{\mu}(1-e^{-\mu t})\right),
\]
and $Y$ and $Z$ are independent.

Therefore,
\[
X(t+s)=Y+Z.
\]

\textbf{(a)}
\[
E[X(t+s)\mid X(s)=n]
=
E[Y]+E[Z].
\]
Thus
\[
{
E[X(t+s)\mid X(s)=n]
=
ne^{-\mu t}
+
\frac{\lambda}{\mu}(1-e^{-\mu t}).
}
\]

\textbf{(b)}
Since $Y$ and $Z$ are independent,
\[
\operatorname{Var}[X(t+s)\mid X(s)=n]
=
\operatorname{Var}(Y)+\operatorname{Var}(Z).
\]
Hence
\[
{
\operatorname{Var}[X(t+s)\mid X(s)=n]
=
ne^{-\mu t}(1-e^{-\mu t})
+
\frac{\lambda}{\mu}(1-e^{-\mu t}).
}
\]

\textbf{(c)}
Suppose there is currently one customer in the system. Let $S$ be the remaining service time of that customer. Then
\[
S\sim \operatorname{Exp}(\mu).
\]
Given $S=t$, the number of new customers still present when this customer departs is Poisson with mean
\[
\frac{\lambda}{\mu}(1-e^{-\mu t}).
\]
Thus the conditional probability that none are present is
\[
\exp\left\{-\frac{\lambda}{\mu}(1-e^{-\mu t})\right\}.
\]
Therefore, the desired probability is
\[
\int_0^\infty
\mu e^{-\mu t}
\exp\left\{-\frac{\lambda}{\mu}(1-e^{-\mu t})\right\}\,dt.
\]
Let
\[
u=e^{-\mu t}.
\]
Then
\[
\mu e^{-\mu t}dt=-du.
\]
So the integral becomes
\[
\int_0^1
\exp\left\{-\frac{\lambda}{\mu}(1-u)\right\}\,du.
\]
Evaluating,
\[
{
P(\text{system empty when the customer departs})
=
\frac{\mu}{\lambda}\left(1-e^{-\lambda/\mu}\right).
}
\]
}

\begin{exercise}
\label{ex:counting-processes-64}
Suppose that people arrive at a bus stop in accordance with a Poisson process with rate $\lambda$. The bus departs at time $t$. Let $X$ denote the total amount of waiting time of all those who get on the bus at time $t$. We want to determine $\operatorname{Var}(X)$. Let $N(t)$ denote the number of arrivals by time $t$.
\begin{enumerate}
\item What is $E[X\mid N(t)]$?
\item Argue that
\[
\operatorname{Var}[X\mid N(t)]=\frac{N(t)t^2}{12}.
\]
\item What is $\operatorname{Var}(X)$?
\end{enumerate}
\end{exercise}

\solution{
Given that
\[
N(t)=n,
\]
the $n$ arrival times are distributed as $n$ independent uniform random variables on $(0,t)$, unordered.

If a person arrives at time $U$, that person's waiting time is
\[
t-U.
\]
Since $U\sim \operatorname{Uniform}(0,t)$,
\[
E[t-U]=\frac t2,
\]
and
\[
\operatorname{Var}(t-U)=\operatorname{Var}(U)=\frac{t^2}{12}.
\]

\textbf{(a)}
Given $N(t)$,
\[
X
=
\sum_{i=1}^{N(t)}(t-U_i).
\]
Therefore,
\[
{
E[X\mid N(t)]
=
\frac{N(t)t}{2}.
}
\]

\textbf{(b)}
Given $N(t)$, the waiting times are independent and each has variance $\frac{t^2}{12}$. Hence
\[
{
\operatorname{Var}[X\mid N(t)]
=
\frac{N(t)t^2}{12}.
}
\]

\textbf{(c)}
Using the conditional variance formula,
\[
\operatorname{Var}(X)
=
E[\operatorname{Var}(X\mid N(t))]
+
\operatorname{Var}(E[X\mid N(t)]).
\]
Since
\[
N(t)\sim \operatorname{Poisson}(\lambda t),
\]
we have
\[
E[N(t)]=\lambda t,
\qquad
\operatorname{Var}(N(t))=\lambda t.
\]
Therefore,
\[
E[\operatorname{Var}(X\mid N(t))]
=
E\left[\frac{N(t)t^2}{12}\right]
=
\frac{\lambda t^3}{12}.
\]
Also,
\[
\operatorname{Var}(E[X\mid N(t)])
=
\operatorname{Var}\left(\frac{N(t)t}{2}\right)
=
\frac{t^2}{4}\lambda t
=
\frac{\lambda t^3}{4}.
\]
Thus
\[
\operatorname{Var}(X)
=
\frac{\lambda t^3}{12}
+
\frac{\lambda t^3}{4}
=
\frac{\lambda t^3}{3}.
\]
Hence
\[
{
\operatorname{Var}(X)=\frac{\lambda t^3}{3}.
}
\]
}

\begin{exercise}
\label{ex:counting-processes-65}
An average of $500$ people pass the California bar exam each year. A California lawyer practices law, on average, for $30$ years. Assuming these numbers remain steady, how many lawyers would you expect California to have in $2050$?
\end{exercise}

\solution{
Under a stationary-population interpretation of the stated steady rates, this is an application of Little's law. The long-run average number of lawyers practicing equals
\[
\text{arrival rate}\times \text{average time in system}.
\]
Here, the arrival rate of new lawyers is
\[
500
\]
per year, and the average practicing lifetime is
\[
30
\]
years. Therefore, the expected number of practicing lawyers is
\[
500\cdot 30=15000.
\]
Thus,
\[
{15000}
\]
lawyers would be expected in this stationary approximation. Without a stationary assumption or an initial population and lifetime distribution, the exact expected population in a specified calendar year is not determined by these two averages alone.
}

\begin{exercise}
\label{ex:counting-processes-66}
Policyholders of a certain insurance company have accidents at times distributed according to a Poisson process with rate $\lambda$. The amount of time from when the accident occurs until a claim is made has distribution $G$.
\begin{enumerate}
\item Find the probability there are exactly $n$ incurred but as yet unreported claims at time $t$.
\item Suppose that each claim amount has distribution $F$, and that the claim amount is independent of the time that it takes to report the claim. Find the expected value of the sum of all incurred but as yet unreported claims at time $t$.
\end{enumerate}
\end{exercise}

\solution{
Assume no outstanding claims at time $0$. Reporting delays are i.i.d.
with distribution $G$ and independent of the Poisson accident process.
An accident at time $s\le t$ is still unreported at $t$ with probability
$1-G(t-s)$.

\textbf{(a)} Independent thinning gives a Poisson count of unreported
claims with mean
\[
m(t)=\lambda\int_0^t[1-G(t-s)]\,ds
=\lambda\int_0^t[1-G(u)]\,du.
\]
Consequently, for $n=0,1,2,\ldots$,
\[
P(\text{$n$ incurred but unreported claims at }t)
=e^{-m(t)}\frac{m(t)^n}{n!}.
\]
\textbf{(b)} Let $Y$ be a claim amount with distribution $F$, independent
of reporting delay and accident time, and assume $E[Y]<\infty$.
Each unreported claim has the same mean amount $E[Y]$. By conditioning
on the accident times and independent marks, the expected total amount is
\[
E[\text{total unreported claim amount at }t]
=\lambda E[Y]\int_0^t[1-G(u)]\,du.
\]
}

\end{document}
